\documentclass[11pt,a4paper]{article}
\usepackage[margin=2.5cm]{geometry}
%\usepackage[utf8]{inputenc}
\usepackage[T1]{fontenc}      % Schriftkodierung
\usepackage{fontspec}
\usepackage[ngerman]{babel}
\usepackage{lmodern}
\usepackage{microtype}
\usepackage{amsmath,amssymb}
\usepackage{booktabs,longtable,array,tabularx}
\usepackage{ragged2e}
\usepackage{enumitem}
\usepackage[table]{xcolor}
\usepackage{xurl}
\usepackage{hyperref}
\usepackage{graphicx}
\usepackage{pdflscape}
\usepackage{seqsplit}
\usepackage{fancyhdr}
\usepackage{titlesec}
\usepackage{etoolbox}
\urlstyle{same}
\setlength{\Urlmuskip}{0mu plus 1mu}
\makeatletter
\def\UrlBreaks{\do\/\do\-\do\.\do\_\do\{\do\}}
\makeatother
\hypersetup{colorlinks=true,linkcolor=blue!50!black,citecolor=blue!50!black,urlcolor=blue!50!black,pdftitle={CNNA-Roadmap v0.0605},pdfauthor={antaris; ChatGPT; Claude}}
\setlength{\headheight}{14pt}
\pagestyle{fancy}
\fancyhf{}
\fancyhead[L]{CNNA-Roadmap aus REAL-OQS v0.0605}
\fancyhead[R]{\thepage}
\titleformat{\section}{\large\bfseries\color{blue!45!black}}{\thesection}{0.6em}{}
\titleformat{\subsection}{\normalsize\bfseries\color{blue!35!black}}{\thesubsection}{0.6em}{}
\setlist[itemize]{leftmargin=1.6em,itemsep=0.25em,topsep=0.35em}
\setlist[enumerate]{leftmargin=1.8em,itemsep=0.25em,topsep=0.35em}
\newcolumntype{L}[1]{>{\RaggedRight\arraybackslash\hspace{0pt}}p{#1}}
\newcommand{\code}[1]{\ifmmode\text{\ttfamily\seqsplit{#1}}\else\begingroup\ttfamily\seqsplit{#1}\endgroup\fi}
\let\oldpath\path
\renewcommand{\path}[1]{\oldpath{#1}}
\newcommand{\goalbox}[1]{\begin{center}\fcolorbox{blue!50!black}{blue!4}{\parbox{0.95\linewidth}{#1}}\end{center}}
\newcommand{\warningbox}[1]{\begin{center}\fcolorbox{black!60}{black!2}{\parbox{0.95\linewidth}{#1}}\end{center}}
\newcommand{\doclevel}[1]{\par\small\noindent\textbf{Dokumentebene:} #1\par\normalsize\medskip}
\newcommand{\tablefontsize}{\fontsize{6.1}{6.8}\selectfont}
\setlength{\tabcolsep}{2pt}
\setlength{\LTleft}{0pt}
\setlength{\LTright}{0pt}
\setlength{\LTpre}{0.3em}
\setlength{\LTpost}{0.4em}
\setlength{\emergencystretch}{4em}
\sloppy
\AtBeginEnvironment{tabular}{\tablefontsize}
\AtBeginEnvironment{tabularx}{\tablefontsize}
\AtBeginEnvironment{longtable}{\tablefontsize}
\newcommand{\landtablecaption}[1]{\begin{center}{\tablefontsize\bfseries #1}\end{center}\vspace{0.2em}}
\newcommand{\tablecaptionline}[1]{\refstepcounter{table}\noindent{\tablefontsize\bfseries Tabelle \thetable: #1}\par\vspace{0.2em}}
\newenvironment{wideauditlandscape}{%
  \begingroup
  \setlength{\tabcolsep}{1pt}%
  \renewcommand{\arraystretch}{0.92}%
  \tablefontsize%
  \begin{landscape}%
}{%
  \end{landscape}%
  \endgroup
  \clearpage
}

\begin{document}
\selectlanguage{ngerman}
\pagenumbering{gobble}

\begin{titlepage}
\centering
{\Large \textbf{Roadmap für die Migration von REAL-OQS v0.0605 nach CNNA}}\\[0.7em]
{\large Complement-Netzarchitektur als gemeinsame Stammtheorie für mehrere QFT-Faces}\\[2em]
\begin{tabularx}{0.92\linewidth}{@{}L{0.28\linewidth}X@{}}
Dokumenttyp: & architekturgebundene Migrationsroadmap \\
Codebindung: & Archiv \path{REAL-OQS_v0.0605_full_no_comments.zip} \\
Interne Leitdokumente: & CNNA-Spezifikation v0.1; Lean-Dokumentationsregelwerk v1.8 \\
Externe Leitquellen: & Snowmass-White-Paper zur QFT-Definition; AQFT/LCQFT/FFT/CFT/AS-, Defekt-, Anomalie- und String-Mathematik-Literatur \\
Status: & keine Phasenplanung, sondern pfeilerweise Vorgehensweise mit Abhängigkeiten, Gates und Stop-Regeln \\
Datum: & 02. April 2026 \\
Editor: & Jan Seeck (antaris) \\
Ghostwriter: & ChatGPT
\end{tabularx}
\vfill
\begin{minipage}{0.92\linewidth}
\small
Diese Fassung ist \emph{regelwerksbasiert}: Sie ist explizit an ein Archiv gebunden, trennt zwischen
portierbarem Kern, neu zu schreibenden Modulen, Bright-Spezialfall und späten Seeds/Windows,
und organisiert die Migration nicht als Zeitplan, sondern als partielle Ordnung fachlicher Arbeitspakete.
Die Motivation für eine mehrgesichtige CNNA-Stammtheorie folgt sowohl aus der internen
Spezifikation als auch aus der in der Snowmass-Übersicht betonten Tatsache, dass keine einzelne
heute etablierte Axiomatik alle physikalisch relevanten QFT-Situationen abdeckt; deshalb ist das Ziel
nicht \emph{eine} einzig privilegierte Oberflächenbeschreibung, sondern eine gemeinsame tiefere
Struktur, aus der mehrere QFT-Faces sektor- und regimeabhängig emergieren.\cite[Abstract; pp.~1, 4--6, 9--12]{Dedushenko2023}\cite[pp.~1--10, 20--36]{CNNASpec}\cite[pp.~7--8]{RulebookV18}
\end{minipage}
\vfill
\end{titlepage}

\clearpage
\pagenumbering{arabic}
\tableofcontents
\clearpage

\section*{Dokument-Metadaten}
\addcontentsline{toc}{section}{Dokument-Metadaten}
\doclevel{Projekt-Ebene und Meta-Audit-Ebene}

\begin{longtable}{L{0.26\linewidth}L{0.68\linewidth}}
\toprule
Feld & Inhalt \\
\midrule
Dokumenttitel & Roadmap für die Migration von REAL-OQS v0.0605 nach CNNA \\
Codebasis & Archiv \path{REAL-OQS_v0.0605_full_no_comments.zip} \\
Prüfstatus & Archivgebundene Migrationsroadmap; aktuelles Artefakt als PDF kompilierbar. Die Roadmap selbst ist keine modulweise Vollbeweisdokumentation, sondern ein pfad- und handoffgebundenes Auditdokument. Auf dem kritischen Pfad gelten keine freien \code{sorry}/\code{admit}/Axiome als zulässige Stütze; wo private Hilfssätze, Instanz-Synthese oder \code{omit}-Kontexte für die spätere Vollformalisierung relevant werden, müssen sie in Kontextblatt, Handoff- oder Kettenblatt explizit sichtbar gemacht werden. \\
Geltungsbereich & Pfeiler A--E, aktive Migrationspfade, Gate- und Stop-Regeln, Moduldisposition, Handoffs, Kettenblätter, Audit-Anhänge; keine atomare Voll-Dokumentation jedes einzelnen Lean-Beweises des Archivs. \\
Zielpublikum & Reviewer ohne laufenden Codezugriff, Projektmitglieder, spätere Fachreviewer für Mathematik, AQFT, OQS, Geometrie und Matter-Pfade. \\
Sprachstufe & Mathematik first, Erläuterung second, Interpretation third. \\
Quellenbasis & Archivcode v0.0605; CNNA-Spezifikation v0.1; Lean-Dokumentationsregelwerk v1.8; gezielt eingesetzte Primärliteratur für AQFT/LCQFT/NCG/String/Scattering. \\
Audit-Konvention & Tragende Aussagen werden über Lean-Namen, Modulpfade, Handoff-Tabellen, Kettenblätter, Meta-Audit-Marker, Register- und Anhangsverweise rückverfolgbar gemacht. \\
Geltigkeitsbindung & Bindung an das konkrete Archiv, nicht an einen Git-Commit. Fortschreibungen müssen Delta-Änderungen an Pfaden, Gates, Handoffs, Registern und Auditpunkten kenntlich machen. \\
Reifegrad-Konvention & Projektintern: \code{Seed}, \code{Scaffold}, Kernpfad, \code{Window}. Regelwerksnahe Abbildung: \emph{vollständig bewiesen}, \emph{strukturell vollständig / Eigenschaften als Prop-Felder}, \emph{reiner Interface-Hook}. \\
Dokumentationstyp & Vollständige Projekt-/Pfad-Roadmap mit Handoff-, Ketten- und Meta-Audit-Anhängen; keine Einzeltheorem-Dokumentation im Sinne eines vollständigen Beweisbandes. \\
\bottomrule
\end{longtable}

\section{Einleitung}
\doclevel{Projekt-Ebene}

CNNA ist in dieser Roadmap kein bereits abgeschlossenes System, sondern ein \emph{geplantes Projekt}:
Gemeint ist der kontrollierte Übergang von der archivgebundenen REAL-OQS-Konfiguration v0.0605
zu einer neuen Stammarchitektur, in der die heute im Archiv dominierende bright-/boundary-first
Lesart methodisch zurücktritt und durch eine derived-only, sector-first und komplementnetzbasierte
Kernlogik ersetzt wird. Die Roadmap beschreibt daher nicht die Verteidigung eines fertigen Resultats,
sondern den formalen Rahmen eines Umbaus, in dem ontischer Kern, AQFT-Schicht, offene Dynamik,
emergente Geometrie und späte Matter-/Gauge-Faces aus einer gemeinsamen CNNA-Stammtheorie
hervorgehen sollen.

\warningbox{\textbf{Projektstatus.} CNNA wird hier ausdrücklich als \emph{geplantes} und
noch nicht vollständig realisiertes Projekt gefasst. Die Roadmap ist deshalb weder ein Abschlussbericht
noch eine bloss motivische Vision, sondern ein bindendes Zwischenartefakt zwischen vorhandenem Archiv,
interner Spezifikation, Lean-Regelwerk und den späteren Implementationsentscheidungen.}

Der Rahmen von CNNA ist durch drei methodische Entscheidungen bestimmt. Erstens gilt
\emph{derived-only}: tragende physikalische Aussagen sollen möglichst nicht durch freie Zusatzannahmen,
sondern aus formal kontrollierten Strukturen, Handoffs und Gates folgen. Zweitens gilt
\emph{sector-first}: Bright, Dark und Interface werden nicht als nachträgliche Dekoration eines bereits
fertigen boundary-Pfads behandelt, sondern als primäre Organisationsform der Architektur. Drittens gilt
\emph{stammtheoretische Mehrgesichtigkeit}: CNNA soll nicht nur einen einzigen Darstellungsstil der QFT
liefern, sondern mehrere QFT-Faces --- AQFT, lokal kovariante, modulare, streutheoretische,
renormierte und späte Matter-/Gauge-Fenster --- als regime- und sektorabhängige Ausprägungen einer
gemeinsamen tieferen Struktur tragen.

Daraus folgt zugleich eine negative Abgrenzung. CNNA ist in dieser Fassung weder als fertige Physiktheorie
noch als vollständig formalisierter Lean-Kern präsentiert. Die Roadmap beansprucht lediglich, dass der
Übergang vom Altarchiv zu einer solchen Stammtheorie \emph{architektonisch} planbar, auditierbar und
regelwerksgebunden formuliert werden kann. Gerade deshalb spielt die Unterscheidung zwischen portierbarem
Kern, neu zu schreibenden Modulen, Bright-Recovery, Seeds, Scaffolds und Windows eine zentrale Rolle:
Sie markiert, was bereits als belastbarer Trager der späteren Theorie gelten darf und was vorerst nur als
Anschluss- oder Spezialfallstruktur geführt wird.

Die Rolle dieser Roadmap ist damit doppelt. Einerseits ist sie ein Migrationsdokument, das den vorhandenen
Archivstand in operative Arbeitspakete, Stop-Regeln und Handoff-Punkte zerlegt. Andererseits ist sie ein
Projektvertrag für das geplante CNNA: Sie legt fest, in welchem Sinn die künftige Architektur als
komplementnetzbasierte Stammtheorie verstanden werden soll, welche Freiheitsgrade geschlossen werden
müssen, welche alten Boundary-Pfade nur regressionsweise weiterleben dürfen und unter welchen Bedingungen
spätere D- und E-Schichten überhaupt als derived-only anschlussfähig gelten. In diesem Sinn ist die
Einleitung nicht bloss Motivation, sondern die ausdrückliche Rahmensetzung für das Projekt CNNA als
nächste Entwicklungsstufe jenseits von REAL-OQS v0.0605.
\newpage
\section{Ziel und Geltungsbereich}
\doclevel{Projekt-Ebene}

\subsection{Ziel}
\doclevel{Projekt-Ebene}
Dieses Dokument legt eine konkrete Migrationsvorgehensweise für das neue Projekt
\emph{Complement-Netzarchitektur} (CNNA) fest. Ziel ist nicht ein kosmetischer Fork des bisherigen
boundary-first Pfads, sondern die Rekonstruktion einer gemeinsamen CNNA-Stammtheorie, aus der
mehrere etablierte QFT-Faces --- AQFT, lokal kovariante Varianten, funktoriale/glueing-orientierte
Fenster, konforme Spezialregime und späte Standardmodell-/SMEFT-Fenster --- sektorabhängig
emergieren. Das motivische Rückgrat ist die interne CNNA-Spezifikation: \emph{derived-only},
\emph{sector-first}, organisierte Komplementfamilien, eigener Interface-Träger, Parameter-Closure
anstelle freier ontischer Cutoffs und explizite Bright-Recovery nur als später Spezialfall.\cite[pp.~1--10, 20--36]{CNNASpec}

\subsection{Geltungsbereich}
\doclevel{Projekt-Ebene}
Die Aussagen sind an das Archiv
\path{REAL-OQS_v0.0605_full_no_comments.zip} gebunden.\cite[p.~7]{RulebookV18}
Die Roadmap deckt den vorhandenen Codebestand der Pfeiler A--C ab und plant darauf aufbauend die
neuen Pfeiler D und E. Sie ist \emph{keine} theorematische Voll-Dokumentation jedes vorhandenen
Lean-Beweises, sondern eine architekturgebundene Migrationsschrift mit Modul-, Herkunfts- und
Gate-Bindung.

\subsection{Was dieses Dokument nicht behauptet}
\doclevel{Projekt-Ebene}
\begin{itemize}
\item Es behauptet nicht, dass die physikalische Endtheorie bereits konstruiert ist.
\item Es behauptet nicht, dass die späten Windows (Minkowski, Conformal, FLRW, SM/SMEFT) schon zum Kern gehören.
\item Es behauptet nicht, dass einzelne Pfade, wie das Komplement-Netzwerk, AQFT, andere QFT Faces, Thermodynamik, Raumzeit, asymptotic safety, induzierte Gravitation oder gar das Standardmodell-Matching schon formal geschlossen vorliegen.
\end{itemize}

\goalbox{\textbf{Startpunkt der strikten Migration:} das aktuelle REAL-OQS-Archiv v0.0605.
\textbf{Zwischenobjekte:} portierbarer Kern aus A--C, neue Core/ToC/OQS/AQFT/Integration-Schicht,
anschliessend D- und E-Strang. \textbf{Endpunkt:} eine gemeinsame CNNA-Stammtheorie, in der der alte
bright/boundary-first Pfad nur noch als später rekonstruierbarer Spezialfall erscheint.\par
\medskip
\textbf{Leitsatz der späteren Symmetrie-/Materiepfade:} Die full-derived Route zu vN-/Typ-III-,
CPT-, Kramers- und Spin-Statistik-Aussagen läuft nicht über Bright-Recovery, sondern über die
komplementnetz-basierte modulare Darstellungsarchitektur. Das Komplementnetz ist hier kein
projektinterner Stilzug, sondern der physikalische Träger dafür, dass Lokalität, Zustand,
Darstellung, Modularität, Kausalität/Geometrie und spätere Symmetrie-/Statistiksätze aus
demselben derived-only Pfad hervorgehen. Für belastbare Charge-/Gauge-, Nuclearity-,
lokal-kovariante Dynamik- und renormierte Lokalfeld-Aussagen werden darüber hinaus explizit eine
DHR/DR-artige Sektor-/Feldrekonstruktionsschicht, eine Nuclearity-Schicht, relative
Cauchy-Evolution/dynamical locality sowie Hadamard-/mikrolokale und lokale Wick-/Zeitordnungs-
produkte benötigt.}

\section{Dokumenttyp, Quellenhierarchie und Minimal-Audit}
\doclevel{Projekt-Ebene und Meta-Audit-Ebene}

Das Regelwerk fordert für vollständige oder pfadtragende Dokumente mindestens Ziel/Geltungsbereich,
Dokumenttyp, Definitionskatalog, Projektglossar, globale Kontexte, Abhängigkeitsgraph bzw.
Strukturkarte sowie Meta-Audit-Marker.\cite[pp.~7--8]{RulebookV18}
Diese Roadmap folgt derselben Logik in einer architektonischen, nicht in einer beweisdokumentierenden,
Form.

\subsection{Quellenhierarchie}
\doclevel{Projekt-Ebene und Meta-Audit-Ebene}
\begin{enumerate}
\item primäre technische Wahrheitsquelle: der inspectierte Lean-Archivbestand v0.0605;
\item primäre Projektsemantik für das Zielsystem aber mittlerweile unterdeterminiert: \emph{Cnna Project Spec V0 1.pdf};
\item normativer Dokumentationsrahmen: \emph{lean\_doku\_regelwerk\_v1\_8.pdf};
\item externe Physik-/Math-Physik-Literatur nur dort, wo sie die Wahl der Zielarchitektur oder die Gate-Reihenfolge rational motiviert.
\end{enumerate}

\subsection{Reifegrad-Konvention und Abbildung auf das Regelwerk}
\doclevel{Projekt-Ebene und Meta-Audit-Ebene}
Die Roadmap verwendet projektintern die Klassen \code{Seed}, \code{Scaffold}, Kernpfad und
\code{Window}. Für die regelwerksnahe Auditierbarkeit werden diese Begriffe wie folgt rückgebunden:
\begin{itemize}
\item \code{Seed}/\code{Window}: in der Regel \emph{reiner Interface-Hook} oder bewusst offenes Spezialfenster;
\item \code{Scaffold}: in der Regel \emph{strukturell vollständig, Eigenschaften als Prop-Felder} oder noch nicht de-seedete Tragerstruktur;
\item Kernpfad mit geschlossenen Gates: Ziel ist \emph{vollständig bewiesen}; wo dies in der Roadmap nur projektiv geplant ist, bleibt der Status explizit als noch nicht geschlossen markiert.
\end{itemize}
Damit wird vermieden, dass projektinterne Reifegrade sprachlich stärker klingen als die regelwerksnahe
Auditklassifikation.


\subsection{Zwang zur lesbaren Notation für alle Lean-Objekte}
\doclevel{Projekt-Ebene, Pfad-Ebene und Meta-Audit-Ebene}
Das Regelwerk fordert für tragende und wiederkehrende Kernausdrücke eine explizite
Lesbarkeitsnotation; für CNNA wird diese Pflicht hier bewusst verschärft: \emph{alle Lean-Objekte
auf dem aktiven Pfad müssen entweder bereits als menschenlesbare Namen vorliegen oder von einer
explizit registrierten lesbaren Notations- bzw. Alias-Schicht begleitet werden.} Das betrifft nicht nur
wiederkehrende Formeln, sondern ebenso Struktur- und Def-Namen, projektinterne Statusmarker,
Transportabbildungen, Handoff-Objekte und Theoremgruppen, sobald ihre rohe Lean-Form die direkte
mathematische Lesbarkeit verdeckt.

Die beigefügte Notationsskizze zeigt das intendierte Muster: eine eigene Lean-Notationsdatei führt
über \code{notation} bzw. \code{scoped notation} eine rein syntaktische, semantisch neutrale
Lesbarkeitsschicht ein. Der Compiler sieht denselben Term wie zuvor; nur die Darstellung für Menschen
wird mathematisch lesbar. Genau dieses Muster wird für CNNA zum verbindlichen Designprinzip:
\begin{itemize}
\item jede tragende Objektfamilie des aktiven Pfads erhält eine dokumentierte lesbare Form;
\item die Notationsschicht lebt nach Möglichkeit \emph{in Lean selbst} (z. B. \path{CNNA/Notation.lean},
\path{CNNA/PillarA/Notation.lean}, \path{CNNA/AQFT/Notation.lean}) und nicht bloss als externe PDF-Umschrift;
\item die Notation erzeugt \emph{keinen} neuen mathematischen Gehalt, sondern expandiert auf bereits
  definierte Terme, Strukturen oder Kompositionen;
\item bei der ersten Verwendung im Text müssen Lean-Name, lesbare Form und Expansion bzw.
  semantische Bindung zusammen sichtbar sein;
\item für neue Versionen größer als v0.0605 gilt: roher Name allein ist für tragende Objekte nur dann
  zulässig, wenn er bereits ohne Zusatz voll menschenlesbar ist.
\end{itemize}

Damit verschiebt sich die Lesbarkeitsforderung von einer optionalen Dokumenthilfe zu einer echten
Architekturpflicht. Nicht nur wiederkehrende Ausdrücke wie quadratische Formen, Spuren,
Adjunktionen oder Exponentialterme, sondern \emph{alle} relevanten Lean-Objektnamen des CNNA-Kerns
müssen für Reviewer in eine mathematisch oder fachlich lesbare Schicht übersetzt werden. Die
Roadmap behandelt diese Schicht deshalb nicht als Kosmetik, sondern als Teil der Auditierbarkeit des
Codes selbst: wer einen Namen nicht lesen kann, kann auch seine Rolle in Handoff, Kette, Gate und
Beweisziel nicht verlässlich reviewen.

\warningbox{\textbf{Verbindliche Notationsregel.} Für den aktiven CNNA-Kern darf kein tragender
Lean-Objektname dauerhaft nur in seiner rohen Projektschreibweise verbleiben. Entweder der Name ist
bereits selbstsprechend, oder es muss eine registrierte lesbare Form existieren. Diese lesbare Form SOLL
über \code{notation}/\code{scoped notation} in Lean realisiert werden; eine reine PDF-Umschrift ohne
Lean-seitige Entsprechung genügt nicht als Zielzustand.}

\subsection{Warum eine Stammtheorie mit mehreren Faces?}
\doclevel{Projekt-Ebene und Meta-Audit-Ebene}
Dedushenko betont im Snowmass-White-Paper, dass keine der heute verfügbaren Axiomatik-Familien
alle physikalischen Erscheinungsformen von QFT abdeckt. Die Übersicht behandelt gerade deshalb
nebeneinander Wightman, OS, AQFT/Haag-Kastler, perturbatives AQFT, lokal kovariante QFT,
homotopisches AQFT, Functorial Field Theory und CFT/conformal nets.\cite[pp.~1, 4--12]{Dedushenko2023}
Genau daraus folgt für CNNA eine scharfe Designkonsequenz: das Ziel darf nicht sein, eine einzige
übernommene Oberflächenbeschreibung zu verabsolutieren; das Ziel muss eine tieferliegende,
regime-sensitive Stammstruktur sein, aus der mehrere Faces konsistent hervorgehen.

\subsection{Aktuelle Größenordnung des inspectierten Archivs}
\doclevel{Projekt-Ebene und Meta-Audit-Ebene}
Die Codebasis v0.0605 ist in ihrer heutigen Form noch asymmetrisch: A und B tragen den Grossteil des
mathematischen Gewichts, während C kompakt ist und D/E als neue CNNA-Struktur noch fehlen.

\tablecaptionline{Aktuelle Größenordnung des inspectierten Archivs -- Teil 1}
\begin{longtable}{L{0.3\linewidth}L{0.18\linewidth}L{0.18\linewidth}}
\toprule
Pfeiler & Lean-Dateien & LOC \\
\midrule
\endfirsthead
\toprule
Pfeiler & Lean-Dateien & LOC \\
\midrule
\endhead
\midrule
\multicolumn{3}{r}{Fortsetzung auf der nächsten Seite} \\
\midrule
\endfoot
\bottomrule
\endlastfoot
Pillar A & 127 & 11'073 \\
Pillar B & 89 & 12'412 \\
Pillar C & 18 & 648 \\
\end{longtable}

Die reine Pfeilerstatistik ist aber nicht die ganze operative Last. Im Archiv existieren zusätzlich
Root-, Meta- und Example-Dateien, die zwar nicht als fachlicher Stamm fungieren, aber Importstabilität,
Build-Gates und Regressionsfähigkeit tragen. Sie müssen deshalb in der Migration explizit
mitgeplant werden und dürfen nicht zwischen Pillar und ``blosser Infrastruktur'' verschwinden.

\tablecaptionline{Aktuelle Größenordnung des inspectierten Archivs -- Teil 2}
\begin{longtable}{L{0.46\linewidth}L{0.16\linewidth}L{0.16\linewidth}}
\toprule
Block ausserhalb der Pillar-Statistik & Lean-Dateien & LOC \\
\midrule
\endfirsthead
\toprule
Block ausserhalb der Pillar-Statistik & Lean-Dateien & LOC \\
\midrule
\endhead
\bottomrule
\endfoot
Top-Level-Infrastruktur (\code{Basic}, \code{BuildAll}, \code{PillarA/B/C}) & 5 & 8 \\
Meta & 19 & 908 \\
Examples & 14 & 1'026 \\
Summe ausserhalb A--C & 38 & 1'942 \\
\end{longtable}

Diese Asymmetrie ist inhaltlich wichtig: die Migration darf D und E nicht durch freie Platzhalter
simulieren, sondern muss sie aus dem in A--C bereits vorhandenen Material ableiten. Zugleich ist damit
klargestellt, dass \emph{auch} die ausserhalb der Pillar-Tabelle liegenden Dateien echte Migrationslast
tragen: Meta-Dateien binden Regression und Audit, Examples sichern Replay-Pfade, und die Root-Dateien
tragen Toolchain-, Import- und Build-Stabilität.

\subsection{Roadmap als lebendes Gate-Dokument: prototypisches Vorauseilen und Rückfluss von Erkenntnissen}
\doclevel{Projekt-Ebene und Meta-Audit-Ebene}
Diese Roadmap ist bewusst kein starrer Master-Phasenplan, sondern ein lebendes Gate-Dokument.
Sie wird ab jetzt ausdrücklich wie ein \emph{executable Dokument} behandelt: jede neue belastbare
Erkenntnis aus dem aktiven Lean-Codepfad erzwingt ein Update der Roadmap, und jedes belastbare
Roadmap-Update erzeugt seinerseits über Tabellen, Register, Handoffs und Kettenblätter die neue
Ausgangsbasis für Notationsschicht, Audit-Anhang und gegebenenfalls auch für die partielle Ordnung
selbst. Die Roadmap lebt damit nicht nur entlang des Vorwärtspfads, sondern \emph{nach allen Seiten}:
in Richtung Code, Dokumentation, Notation, Gate-Reife und Rückfluss auf frühere Architekturentscheide.
Gerade in Lean führen zu früh verhärtete Begriffs- und Typentscheidungen leicht in eine
\emph{Typ-Sackgasse}: eine in A scheinbar harmlose Festlegung kann später B-, C- oder D-Pfade
unverhältnismäßig teuer machen, wenn sie die benötigten Darstellungs-, Restriktions- oder
Glueing-Schichten nicht mehr trägt. Deshalb wird detaillierte Nachfolgeplanung erst dort hart,
wo der vorgelagerte Pfeiler sein Gate wirklich geschlossen hat; vorher bleibt sie begründete,
aber bewusst revidierbare Absicht.

Zugleich erlaubt die Roadmap ein \emph{prototypisches Vorauseilen}: in A dürfen früh skizzierte
Dummy- bzw. Sketch-Schnittstellen für spätere Verbraucher aus B/C/D angelegt werden, etwa für
\code{ComplementLocalNet}-nahe Restriktionen, Zustandsprojektionen, modulare Aktionen oder
Glueing-Signaturen. Solche Vorgriffe zählen ausdrücklich nicht als fertiger Kern, sondern müssen
als \code{Seed} oder \code{Scaffold} markiert bleiben. Ihr Zweck ist nicht, spätere Mathematik
vorzutäuschen, sondern früh zu testen, ob die Typ- und Exportlage von A die Last der späteren
Schichten überhaupt tragen kann.

Weil CNNA als Stammtheorie auf ungeplante Emergenz und Strukturfeedback reagieren muss, fungiert
die Roadmap zusätzlich als \emph{Logbuch von Architekturentscheidungen aus Notwendigkeit}. Wenn im
aktiven Codepfad neue Symmetrien, Selektoren oder Restriktionsnotwendigkeiten sichtbar werden,
müssen die Anforderungen an B/C/D sofort nachgeführt werden; umgekehrt darf eine spätere
Einsicht aus C, D oder E frühere A-Entscheidungen wieder öffnen, falls sich deren Schnittstellen
als zu restriktiv erweisen.

Daraus folgt eine explizite \emph{Interface-Stabilitätsregel}: Erkenntnisrückfluss nach früheren
Pfeilern ist zulässig, aber nur kontrolliert. Sobald ein späterer Pfeiler zeigt, dass ein früheres
Axiom, eine Typwahl oder ein Exportvertrag zu eng war, wird das betroffene Modul zunächst wieder
auf \code{Seed}/\code{Scaffold}-Status zurückgestuft, die Source-Maps und Imports werden angepasst,
Regression und Replay werden erneut geprüft, und erst danach darf das Gate wieder geschlossen
werden. Genau dafür existiert die Gate-Logik: sie garantiert Flexibilität \emph{vor} Closure und
Härte \emph{nach} Closure, statt einen unfertigen Unterbau durch parallele Detailplanung zu
überfahren.

\subsection{Mehrstufiges Wahrheitsprotokoll: AI, Formalisierung, Verifikation, Semantik und Physik}
\doclevel{Projekt-Ebene und Meta-Audit-Ebene}
CNNA benutzt AI nicht als fachliche Autorität, sondern als heuristische Vorschlags-,
Strukturierungs- und Übersetzungsschicht zwischen informeller Mathematik, formaler Sprache,
maschinell geprüftem Code und später physikalischer Deutung. Wahrheitsquelle auf dem kritischen
Pfad ist deshalb nicht eine plausible AI-Prosa, sondern die Kombination aus wohldefinierten
Definitionen, strikt typgeprüftem Lean-Code, Kernel-Check und einer expliziten Disziplin gegen
\code{sorry}, \code{admit} und freie Axiome im aktiven Kernpfad.

Damit verschiebt sich der methodische Engpass. Nicht das freie Formulieren eines scheinbar eleganten
Beweises ist das Primärproblem, sondern die Wahl tragfähiger Definitionen, Typ- und
Interface-Entscheidungen, sauberer Handoff-Strukturen sowie die \emph{korrekte Prosa-Semantik} der
Kommentare, Dokumente und Brückentexte. Gerade der letzte Punkt ist nicht kosmetisch: der Code kann
formell korrekt sein, während seine menschliche Beschreibung falsch, zu stark oder physikalisch
irreführend ist. Deshalb gehören Kommentar- und Dokumentations-Cleanup, semantische
Gegenprüfung und explizite Herkunftsbindung selbst zur wissenschaftlichen Absicherung. Gerade an
allen kritischen Übergangen --- Definition, Handoff, Kommentarsemantik, externe Dokumentation und
physikalische Deutung --- bleiben Reviews durch fachlich qualifizierte Expertinnen, Experten und
professionelle Reviewer unverzichtbar.

Methodisch ergibt sich daraus ein mehrstufiges Wahrheitsprotokoll:
\begin{enumerate}
\item \textbf{Vermutung / informelle Idee:} eine mathematische oder physikalische Intuition wird als
Arbeitsvermutung formuliert;
\item \textbf{AI-gestützte Strukturierung:} AI darf Begriffe sortieren, Ableitungsketten skizzieren,
Literaturpfade verdichten und Übersetzungsarbeit zwischen informeller und formaler Sprache leisten;
\item \textbf{Formalisierung:} die Aussage wird in Lean-Definitionen, Interfaces, Handoffs und saubere
Theoreme überführt;
\item \textbf{Maschinelle Verifikation:} wahr im technischen Sinn ist nur, was der Proof-Checker relativ zu
den gewählten Definitionen wirklich akzeptiert;
\item \textbf{Semantikprüfung und Fachreview:} Menschen sowie fachlich qualifizierte Expertinnen,
Experten und professionelle Reviewer prüfen, ob Kommentare, externe Dokumente, Source-Maps und
Theorembeschreibungen den kompilierten Code korrekt wiedergeben;
\item \textbf{Physikalische Interpretation:} erst danach wird durch fachlich qualifizierte Reviews
bewertet, ob der formale Satz für die intendierte Physik relevant, stark genug und nicht
fehlgedeutet ist.
\end{enumerate}

Dieses Protokoll ist bewusst strenger als die übliche Redeweise ``AI hat bewiesen, dass ....
Im CNNA-Kontext gilt stattdessen: \emph{AI darf vorschlagen; wahr und belastbar ist nur, was formal
trägt, semantisch korrekt dokumentiert ist und die physikalische Deutung übersteht.} Gerade für
derived-only-Architekturen ist diese Trennung zentral, weil ein Proof-Checker nur relativ zu den
gewählten Definitionen prüft, nicht aber automatisch deren physikalische Relevanz,
Nicht-Trivialität oder Emergenzstärke.

\warningbox{\textbf{Methodische Leitregel:} AI ist in dieser Roadmap weder mathematische noch
physikalische Letztautorität. Sie dient als Heuristik- und Übersetzungsschicht. Autoritativ sind
auf dem aktiven Pfad nur (i) die formale Sprache und der Kernel des Proof Checkers, (ii) die
expliziten Projektregeln gegen \code{sorry}/\code{admit}/freie Axiome, (iii) die semantische
Gegenprüfung von Prosa und Kommentaren und (iv) die nachgelagerte physikalische Interpretation.
Damit ist die Roadmap nicht bloss ``AI + Lean, sondern ein kontrolliertes mehrstufiges
Wahrheitsprotokoll. In allen kritischen Instanzen bleiben Reviews durch fachlich qualifizierte
Expertinnen, Experten und professionelle Reviewer unabdingbar; weder AI-Prosa noch blosses
Typchecking ersetzen diesen Schritt.}

\subsection{Operative Verdichtung des aktiven Migrationspfads}
\doclevel{Projekt-Ebene und Meta-Audit-Ebene}
Die bisherige Struktur lässt sich für den aktiven Projektvollzug auf eine knappe Kernlogik
verdichten. Erstens bleibt diese Schrift eine \emph{architekturgebundene Migrationsroadmap} für
das Archiv \path{REAL-OQS_v0.0605_full_no_comments.zip}; sie ist keine lose Wunschliste und keine
ex post erzielte Ergebnisnarration. Zweitens ist das Ziel eine \emph{CNNA-Stammtheorie} im
strict-derived-only- und sector-first-Sinn, aus der mehrere QFT-Faces emergieren; der alte
Bright-/Boundary-Pfad ist dafür später Spezialfall, Recovery und Regressionsträger, aber nicht
mehr begriffliches Zentrum. Drittens trägt die Architektur genau die Pfeilerfolge
A \(\to\) B \(\to\) C \(\to\) D \(\to\) E: ontischer Kern und Komplementgrammatik,
AQFT aus Komplementnetz, offene Systeme/Kanäle, emergente Geometrie/LocalCovariance und erst
danach späte Gauge-/Matter-/SM-Fenster.

Gerade deshalb müssen die Stop- und Gate-Bedingungen im Vollzug sichtbar bleiben. Auf dem aktiven
Pfad darf \code{ComplementSpectrumCondition} nicht hinter \code{BisognanoWichmannSeed}
verschwinden; die modulare B-Kette
\(\code{ComplementGNS} \to \code{SeparatingProperty} \to \code{TomitaTakesakiData} \to\
\code{ModularConjugation} \to \code{ComplementKMS}\) ist keine stilistische Ausdifferenzierung,
sondern die harte Vorbedingung späterer vN-/Typ-III-, CPT-, Kramers- und Spin-Statistik-Lesarten.
Ebenso ist \code{ParameterClosure} keine späte Schönheitskorrektur: ohne
\code{selectedBranching}, \code{HorizonLevel}, \code{BackreactionFixedPoint} und
\code{RegularizationClosure} bleibt derived-only unvollständig, weil ontische, dynamische oder
numerische Restparameter sonst unkontrolliert aus dem Kernpfad herausragen.

Daraus folgen für den unmittelbaren Arbeitsvollzug fünf Prioritätsregeln:
\begin{enumerate}
\item \textbf{A-Core zürst einfrieren.} \code{Approximant}, \code{RegionCore},\
\code{RegionNet}, \code{InfiniteCarrier}, \code{Selection}, dazu
\code{BranchPatch}/\code{ComplementSectorFamily}/\code{SectorSplit}, werden semantisch stabilisiert,
bevor tiefere Recovery-Pfade nachgezogen werden.
\item \textbf{GeneralizedDtN als explizite Kernstruktur bauen.} Die Mehrsektor-Schur-/DtN-Schicht
bleibt ein \code{structure}-basiertes Kernobjekt mit klaren Restriktions-, Kompositions- und
Glueing-Daten; die Verbindung zu \code{ComplementLocalNet} wird früh getestet, nicht erst am Ende.
\item \textbf{Komplement- und Interface-Netze vor Bright-Recovery.}\
\code{ComplementLocalNet} und \code{InterfaceLocalNet} entstehen vor jeder portierten
BoundaryMatrix-Comfort-Schicht. Alte \code{BoundaryMatrix*}-Dateien zählen nur als Replay-,
Regression- und Extraktionsquelle und sollten organisatorisch als \emph{bright-recovery}-Unterstrang
behandelt werden.
\item \textbf{Parameter-Closure früh aktivieren.} Der operative Closure-Pfad
\(\code{BranchingWitness} \to \code{selectedBranching} \to \code{HorizonLevel}/\code{BackreactionFixedPoint} \to \code{RegularizationClosure}\)
läuft früh mit, damit \(b\), \(L_{\max}\) und Regularisierungsreste nicht als stille
Aussenparameter stehen bleiben.
\item \textbf{Namespace- und Semantik-Sauberkeit erzwingen.} Portierte Module dürfen keine
implizite boundary-first-Semantik in Bezeichnern, Kommentaren, Source-Maps oder Imports
weitertragen; wo das Altarchiv nur Bright-Spezialfall kodiert, muss dies sichtbar gemacht und aus
dem CNNA-Kern herausgezogen werden.
\end{enumerate}

\warningbox{\textbf{Operatives Risikobild:} Die kritischen Migrationsrisiken liegen nicht mehr in
fehlenden grossen Pfeilern, sondern in stillen Altsemantiken: trunkartige BoundaryMatrix-Imports im
Kernpfad, freie Restparameter \(b, L_{\max}, \beta\) ohne Closure, sowie Namens- und
Kommentarspuren, die portierte Module weiter als boundary-first lesen lassen. Die Gegenmassnahme
ist nicht mehr Prosa, sondern frühes Einfrieren des A-Kerns, aktive Closure, harte Gate-Disziplin
und konsequente Trennung zwischen CNNA-Kern und Bright-Recovery.}


\newpage
\section{Definitionskatalog der Kernobjekte}
\doclevel{Projekt-Ebene und Pfad-Ebene}
\label{sec:definitions}
Dieser Abschnitt trägt die formale Last der im aktiven und spät aktivierten Pfad wiederkehrenden
Kernobjekte. Er ist kein erläuterndes Glossar, sondern die \emph{kanonische Objektquelle} der Roadmap.
Jede Zeile definiert genau \emph{ein} kanonisches Lean-Objekt. Register und Notations-Appendix dürfen
keine neuen Fachobjekte erzeugen, sondern müssen auf diesen Katalog zurückverweisen.

\subsection{Formaler Definitionskatalog des aktiven und spät aktivierten Pfads}
\doclevel{Projekt-Ebene und Pfad-Ebene}
Die folgende Liste führt nicht nur Name, Schreibweise und Rolle, sondern für jedes Objekt auch
Objektklasse, primäre lesbare Form, Lesbarkeitsmodus, exakte Bindung, Zielmodul der möglichen
Notations-Emission, Scope und Pfadstatus. Damit ist sie die kanonische Eingabemenge für Register,
Notations-Appendix und spätere \path{Notation.lean}-Module.

\subsection{Notationspflicht innerhalb des Definitionskatalogs}
\doclevel{Projekt-Ebene und Meta-Audit-Ebene}
Für jeden vollen Kerneintrag gelten ab jetzt die folgenden Integritätsregeln:
\begin{enumerate}
\item \textbf{Eine Zeile = ein kanonisches Objekt.} Zusammengedrückte Sammelzeilen sind für ableitungsrelevante Objekte unzulässig.
\item \textbf{Primäre lesbare Form.} Jeder Eintrag trägt genau eine kanonische Oberflächenform; weitere Aliasformen gehören ins Register.
\item \textbf{Exakte Expansion / Bindung.} Jede lesbare Form wird explizit auf den zugrunde liegenden Lean-Namen zurückgeführt.
\item \textbf{Expliziter Lesbarkeitsmodus.} Zulassungen für \emph{selfspeaking}, \emph{text alias}, \emph{symbolic notation} und \emph{symbolic notation + text alias} werden pro Objekt festgelegt.
\item \textbf{Keine technischen Alternativen in einer Zelle.} Formulierungen wie \emph{oder}, \emph{bzw.} oder \emph{später optional} sind in der ableitungsrelevanten Tabellenlogik verboten.
\item \textbf{Notations-Emission nur mit Zielmodul und Scope.} Jede exportierte Lesbarkeitsform benötigt eine eindeutige Modul- und Scope-Bindung.
\end{enumerate}
Meta-Begriffe und Reifegradmarker werden \emph{nicht} in diesen Definitionskatalog eingemischt, sondern später in eigene Registertabellen ausgelagert.

Für den Definitionskatalog gelten folgende Integritätsregeln:
\begin{enumerate}
\item Jede aktive oder spät aktivierte tragende Code-Entität erscheint genau einmal.
\item Jede Zeile besitzt genau eine Primärform; Exportformen sind strikt sekundär.
\item Jeder Eintrag mit Modus ungleich \emph{selfspeaking} oder \emph{no exported readability layer} erzeugt mindestens einen Registereintrag.
\item Jeder Eintrag mit exportierter Form erzeugt mindestens einen Appendix-Eintrag.
\item Namespace, Aktivierungsregel, Kontext, Kollisionspolitik, logische Vorgänger und Verbraucher/Nachfolger dürfen nicht ausschliesslich proseartig ausserhalb der Tabelle festgelegt werden.
\end{enumerate}

\paragraph{Tabellarische Auslagerung.}
Der vollständige tabellarische Definitionskatalog ist aus Layoutgründen in den Anhang ausgelagert; vgl. Anhang~\ref{app:defcatalog}. Im Haupttext verbleiben nur die Einordnung des Katalogs und die normativen Integritätsregeln seiner Verwendung.

\newpage
\section{Projektglossar und Verweißtruktur}

\doclevel{Projekt-Ebene}

\subsection{Kernbegriffe des Projekts}
\doclevel{Projekt-Ebene}
\tablecaptionline{Kernbegriffe des Projekts -- Teil 1}
\begin{longtable}{L{0.22\linewidth}L{0.72\linewidth}}
\toprule
Begriff & Rolle in der Roadmap \\
\midrule
\endfirsthead
\toprule
Begriff & Rolle in der Roadmap \\
\midrule
\endhead
\bottomrule
\endfoot
ToC & kanonisches ontisches Substrat; in CNNA nicht mehr nur Quelle eines Bright-Randnetzes, sondern Generator von Bright-, Dark- und Interface-Struktur. Mathematisch braucht der aktive Pfad mindestens ein verzweigungsfähiges, geankertes Cliquen-/Pfadobjekt mit Inzidenz-, Koarsenierungs-, Branching- und Source-Map-Daten, so dass \code{BranchPatch}, \code{ComplementSectorFamily} und \code{SectorSplit} nicht frei schweben. Diskrete, kontinuierliche oder graduierte ToC-Varianten sind nur dann zulässig, wenn sie genau diesen sektor-split-kompatiblen Vertrag instanziieren, statt eine neue Ontologie an der Split-Schicht vorbei einzuschmuggeln. \\
Bright & expliziter endlicher REAL-Approximant bzw. sichtbarer Patch; wichtig, aber nicht mehr privilegiertes Endprodukt. \\
Dark & organisierte Komplementfamilie inklusive Trunk, Seitenästen und eigenen UV-Tails; keine blöße Restumgebung. \\
Interface & eigener mathematischer Träger für Kopplung, Rückwirkung, Entropie, Modularität, Glueing und spätere Defekt-/Dualitätsdaten. \\
BranchPatch & erster endlicher sichtbarer Patch relativ zu einem ToC-Anker; ersetzt die alte implizite Zentrierung auf eine feste Randoberfläche. \\
ComplementSectorFamily & organisierte Familie komplementärer Sektoren; ersetzt das alte singuläre Komplement. \\
SectorSplit & expliziter Mehrsektor-Split, der Bright, Dark und Interface synchronisiert. \\
GeneralizedDtN & verallgemeinerter DtN-/Schur-Kopplungskern für mehrere Komplementsektoren; algebraisch keine blosse Zahl, sondern sektorindexierte block- bzw. operatorwertige Kopplungsstruktur mit Restriktions-, Kompositions-, Symmetrie-/Positivitäts- und Glueing-Daten. Der alte binäre DtN-Fall bleibt nur Sonderfall. Lean-seitig sollte diese Schicht primär als \code{structure} mit expliziten Feldern und nachgelagerten Instanzen/Gates für Restriktion, Lokalisierung und Verträglichkeit mit \code{ComplementLocalNet} modelliert werden, nicht als nackte Typeclass oder lose Kategoriefassade. \\
TruthProtocol & mehrstufiges Wahrheitsprotokoll für CNNA: Vermutung, AI-gestützte Strukturierung, Formalisierung, maschinelle Verifikation, Semantikprüfung und physikalische Interpretation bleiben explizit getrennt; technische Wahrheit entsteht nicht durch AI-Prosa, sondern erst nach Kernel-Check und semantischer Gegenprüfung. \\
SemanticReview & verpflichtende Kontrolle, ob Kommentare, externe Dokumente, Source-Maps und Theorembeschreibungen den kompilierten Lean-Code korrekt wiedergeben; verhindert, dass formell korrekter Code durch falsche Prosa fehlgedeutet wird. \\
ExpertReview & unabdingbare fachliche Gegenprüfung durch qualifizierte Expertinnen, Experten und professionelle Reviewer; prüft Definitionen, Relevanz, Semantik, physikalische Deutung und die Belastbarkeit des formalen Pfads über das reine Typchecking hinaus. \\
ComplementLocalNet & lokales Netz auf dem dunklen bzw. komplementären Sektor. \\
InterfaceLocalNet & eigenes lokales Netz auf dem Interface; gerade hier sitzt die spätere Glueing-, Entropie- und Defektstruktur. \\
ComplementStateNet & Zustandsfamilien auf Bright, Dark und Interface; keine freie State-Signatur auf dem kritischen Pfad. \\
ComplementGNS / ComplementKMS & Repräsentations- und Thermikschicht des komplementären Netzes; bauen auf der generischen GNS/KMS-Basis auf, dürfen aber nicht ohne Separations- und Modularitätszwischenschicht identifiziert werden. \\
GradedStarAlgebra & Z2-offene algebraische Grundschicht; macht fermionische/statistische Sprache zu einer Rückbindung an die AQFT-Basis statt zu freier später Etikettierung. \\
SeparatingProperty & Reeh--Schlieder-artige Separationsschicht; stärkt die GNS-Lage von bloss zyklisch zu zyklisch+separierend für relevante Komplementalgebren. \\
TomitaTakesakiData & extrahiert Tomita-Operator, modulare Konjugation \(J\) und modularen Operator \(\Delta\) aus der komplementnetz-basierten Darstellungslage. \\
ModularConjugation & kontrolliert die Bright/Dark/Interface-Lesart von \(J M J = M'\); eigentliche modulare Darstellungsarchitektur zwischen GNS und KMS. \\
\shortstack[l]{ComplementSpectrum\\Condition} & explizite Spektral- und Vorwärtskegelkontrolle des aktiven Komplementnetzes; Vorbedingung für BW-, Guido--Longo- und spätere Scattering-Lesarten. \\
\shortstack[l]{Bisognano\\WichmannSeed} & explizite B$\to$D-Brücke; vergleicht modulare Automorphismen mit abgeleiteten Boost-/Wedge-Strukturen und bindet Modularität an Geometrie. \\
NuclearitySeed & Phase-Space-/Energielevel-Dichteschicht jenseits blossen Split-Properties; beginnt erst nach der geschlossenen KMS-/Spektralkontrollschicht und stabilisiert spätere Typ-III-/Thermik-Lesarten. \\
DHRSectorsSeed & observablenseitige Superselection-Schicht für lokalisierte Ladungssektoren; Brücke von AQFT-Netzen zu Charge- und Statistiksprache. \\
\shortstack[l]{FieldAlgebra\\Reconstruction\\Seed} & DR-artige Rekonstruktionsstufe von Feldalgebra und kompakter Eichgruppe aus der observablen Sektorstruktur. \\
\shortstack[l]{RelativeCauchy\\Evolution\\Seed} & LCQFT-Dynamik unter Hintergrundvariation; verknüpft lokale Kovarianz mit abgeleiteter Dynamik und Stress-Tensor-Lesart. \\
HadamardStateSeed & mikrolokal kontrollierte Zustandsklasse für gekrümmte emergente Hintergründe; Vorbedingung für renormierte lokale Observablen. \\
\shortstack[l]{LocalWick\\Polynomials\\Seed} & lokale kovariante Wick- und Zeitordnungsprodukte; trägt spätere Stress-Tensor-, OPE- und Renormierungslesarten. \\
HorizonLevel & abgeleitete Horizontgröße; ersetzt ein primitives \(L_{\max}\) auf dem aktiven Pfad. \\
ParameterClosure & pfeilerübergreifender Closure-Strang für ontische Mikrostrukturparameter, dynamische Emergenzskalen und sekundäre Regularisierer; markiert die derived-only-Schwelle gegen freie Restparameter. \\
BackreactionFixedPoint & Fixpunktstruktur, in der Einfluss-, Entropie-, Skalen- und Zustandsdaten konsistent auf eine effektive IR-/Horizontgröße zurückwirken. \\
RegularizationClosure & schliesst numerische Stabilisierung entweder durch intrinsische Positivität/Koerzivitaet oder degradiert Restterme zu abgeleiteten Spektral-/Nullmodenparametern. \\
EffectiveLambda & abgeleitete kosmologische Skalargröße aus Interface-Balance, Entropie oder Fixpunktdiagnostik. \\
SpectralDimensionFlow & abgeleitete Spektral-/Dimensionsdiagnostik; im Kern erst nach Closure. \\
InfiniteCarrier / quasilokaler Limes & markiert den Übergang vom endlichen Bright-Approximanten zur unendlichen/quasilokalen AQFT-Lage; Vorbedingung für belastbare Typ-III-, Tomita--Takesaki-, BW- und Scattering-Lesarten. \\
PositiveEnergy / UnitaryDynamics & explizite Positivitäts- und Unitaritätskontrolle der effektiven Generator- und Entwicklungsschicht; dürfen im full-derived Pfad nicht nur implizit mitlaufen. \\
PrototypeInterface & bewusst vorgezogene Dummy-/Sketch-Schnittstelle für spätere Verbraucher; prüft früh, ob A-seitige Typ- und Exportentscheide B/C/D überhaupt tragen, zählt aber nie als geschlossener Kern. \\
InterfaceStability & Regel, dass spätere Einsichten frühere Export- oder Typverträge kontrolliert wieder öffnen dürfen; verlangt Rückstufung auf \code{Seed}/\code{Scaffold}, Regressionsprüfung und erneutes Gate-Closing. \\
BackPropagationOfInsights & disziplinierter Rückfluss später Architekturzwänge in frühere Pfeiler; macht die Roadmap zum Logbuch notwendiger Designentscheidungen statt zu einer starren linearen Wunschliste. \\
Seed / Scaffold & bewusst noch nicht de-seedete oder generische Modulhüllen; dürfen den kritischen Pfad nicht als fertige Kernbegriffe dominieren. \\
Window & spätes Spezialfenster (z.B. Minkowski-, Conformal-, FLRW-, SM- oder SMEFT-Fenster), das standardmäßig \emph{nicht} zum Kern de-seedet wird. \\
\end{longtable}

\subsection{Strikte Strukturkarte}
\doclevel{Pfad-Ebene}
\textbf{Symbole und Ausschlüsse vor der Kette.} Startpunkt ist der ToC-Vertrag samt \code{BranchPatch}, \code{ComplementSectorFamily}, \code{SectorSplit} und \code{GeneralizedDtN}. Historische Bright-only-Recovery-Pfade, bloss rhetorische Matter-Fenster und ungeschlossene Closure-/Window-Sprachen sind \emph{nicht} Gegenstand der folgenden strikten Kette.\par
Die alte strikte Ableitung läuft im aktiven Bright-Pfad schematisch als
\[
\text{ToC} \to (\Omega,L,\beta) \to \text{BoundaryMatrixNet} \to \text{State} \to \text{GNS} \to \text{KMS} \to \text{Modular}.
\]
Die CNNA-Stammstruktur muss demgegenüber lauten:
\begin{center}
\small
\tablecaptionline{Strikte Strukturkarte -- Teil 1}
\begin{tabular}{@{}l@{}}
$\text{ToC} \to \text{BranchPatch / ComplementSectorFamily / SectorSplit}$\\
$\to \text{GeneralizedDtN / MultiSchur}$\\
$\to \text{Complement- und Interface-Netze / InfiniteCarrier / quasilokaler Limes}$\\
$\to \text{Zustände / ComplementGNS / SeparatingProperty / TomitaTakesakiData}$\\
$\to \text{ModularConjugation / ComplementKMS / ComplementSpectrumCondition}$\\
$\to \text{Bisognano-Wichmann-Seed / Nuclearity}$\\
$\to \text{Kanäle / Backreaction}$\\
$\to \text{Parameter-Closure / Branching- und Regularization-Closure}$\\
$\to \text{BackreactionFixedPoint / HorizonLevel / EffectiveLambda}$\\
$\to \text{emergente Geometrie / LocalCovariance / RelativeCauchyEvolution}$\\
$\to \text{Hadamard-/mikrolokale Zustandskontrolle / lokale Wick-Produkte}$\\
$\to \text{DHR-Sektoren / Feldalgebra- und Gaugegruppen-Rekonstruktion}$\\
$\to \text{CPT/Spin/Statistik-Brücke / Gauge-Matter- und SM/SMEFT-Fenster}$.
\end{tabular}
\end{center}
Die starke Lesart dieser Strukturkarte ist bewusst nicht optional: vN-/Typ-III-, CPT-, Kramers- und
Spin-Statistik-Aussagen sollen im CNNA-Kern gerade \emph{nicht} aus Bright-Recovery stammen,
sondern aus der komplementnetz-basierten modularen Darstellungsarchitektur, die zwischen
\code{ComplementGNS} und \code{ComplementKMS} explizit geschlossen wird. Für isolierte toy-level-
Sätze in endlicher Quantenmechanik wäre eine solche Komplementarchitektur nicht in jedem Fall
zwingend; für \emph{full-derived emergente lokale QFT} ist sie es im angestrebten physikalischen Sinn
jedoch weitgehend. Ohne Komplement-/Interface-Struktur entstünde höchstens eine Bright-only-
Recovery auf endlichen Typ-I-Testträgern; gerade der Pfad zu modularer Konjugation, Typ-III-
Lesart, Bisognano--Wichmann, CPT, Spin-Statistik und Kramers bliebe dann unvollständig oder nur
aufgesetzt. Der alte Pfad wird also nicht gelöscht, aber architektonisch entzentriert: er ist Ziel
eines späten Recovery- und Regressionstrangs, nicht mehr die Lesart des gesamten Projekts.\cite[pp.~1--10, 30--36]{CNNASpec}\cite{GuidoLongo1995}

\newpage
\section{Pillar-Handoff-Spezifikationen}
\doclevel{Handoff-Ebene}
\label{sec:handoffs}
Die aktiven Pfeilerübergänge werden nicht nur rhetorisch, sondern als formale Handoff-Tabellen
sichtbar gemacht. Damit bleibt prüfbar, welche Daten, Eigenschaften und Nicht-Übergaben den
Paradigmenwechsel zwischen den Pfeilern tragen.

\subsection{Handoff A $\to$ B}
\doclevel{Handoff-Ebene}
\tablecaptionline{Handoff A $to$ B -- Teil 1}
\begin{longtable}{L{0.24\linewidth}L{0.68\linewidth}}
\toprule
Feld & Inhalt \\
\midrule
Quell-Pillar / -bereich & Pillar A: ToC-Vertrag, \code{BranchPatch}, \code{ComplementSectorFamily}, \code{SectorSplit}, \code{GeneralizedDtN}, \code{InfiniteCarrier} \\
Ziel-Pillar / -bereich & Pillar B: AQFT auf \code{ComplementLocalNet}, \code{InterfaceLocalNet}, \code{ComplementStateNet}, \code{ComplementGNS} \\
Übergebene Daten & sektorisierte Regionen, Komplementfamilien, Multi-Schur-/DtN-Kopplung, Glueing-/Restriktionssignaturen, Infinite-Carrier-/Limes-Lage \\
Mitgegebene Beweise / Instanzen & Symmetrie-/Positivitäts- oder Restriktionsdaten nur, soweit sie in A bereits explizit hergeleitet sind; keine freie AQFT-Semantik \\
Nicht übergebene Struktur & GNS, Separierung, Tomita-Takesaki, KMS, Spektralbedingung, Haag-Kastler-/DHR-Lesart \\
Rechtsstatus des Handoffs & definitorischer Export plus Adapter-/Repackaging-Schicht \\
Audit-Punkte & ToC-Vertrag, GeneralizedDtN-/MultiSchur-Pfad, Source-Maps von A in die B-Netze \\
\bottomrule
\end{longtable}

\subsection{Handoff B $\to$ C}
\doclevel{Handoff-Ebene}
\tablecaptionline{Handoff B $to$ C -- Teil 1}
\begin{longtable}{L{0.24\linewidth}L{0.68\linewidth}}
\toprule
Feld & Inhalt \\
\midrule
Quell-Pillar / -bereich & Pillar B: lokale Netze, Zustände, GNS, KMS, Spektral- und Modularitätsdaten \\
Ziel-Pillar / -bereich & Pillar C: Kanäle, relative Entropie, Rückwirkung, offene-System-Dynamik \\
Übergebene Daten & Algebren, Zustandslagen, relevante Generatoren, Interface-Netze, thermische und spektrale Diagnostik \\
Mitgegebene Beweise / Instanzen & Isotonie, Lokalität, Split-/Nuclearity-Vorbedingungen, positive Energie/Unitarität soweit vorhanden \\
Nicht übergebene Struktur & keine freie Interpretation der Kanäle als Physik ohne explizite Rückwirkung und Fixpunktkontrolle \\
Rechtsstatus des Handoffs & gemischter Status: Export plus neue dynamische Konstruktion \\
Audit-Punkte & ComplementStateNet, ComplementGNS/KMS, OQS-Kanalmodule, relative Entropie und Backreaction-Source-Maps \\
\bottomrule
\end{longtable}

\subsection{Handoff C $\to$ D}
\doclevel{Handoff-Ebene}
\tablecaptionline{Handoff C $to$ D -- Teil 1}
\begin{longtable}{L{0.24\linewidth}L{0.68\linewidth}}
\toprule
Feld & Inhalt \\
\midrule
Quell-Pillar / -bereich & Pillar C: Kanäle, Rückwirkung, Entropie- und Einflussdaten, Parameter-Closure \\
Ziel-Pillar / -bereich & Pillar D: emergente Geometrie, LocalCovariance, Relative Cauchy Evolution, Gravity-Routen \\
Übergebene Daten & BackreactionFixedPoint, HorizonLevel, EffectiveLambda, Glueing-/Defektdaten, ggf. Spektraldiagnostik \\
Mitgegebene Beweise / Instanzen & Closure-Gates, positive Energie/UnitaryDynamics soweit für D nötig, Interface-Kohärenz \\
Nicht übergebene Struktur & keine freie Hintergrundgeometrie; keine vorab gesetzte Gravitationstheorie \\
Rechtsstatus des Handoffs & konstruktiver Übergang mit expliziter Nicht-Übergabe von Geometriepostulaten \\
Audit-Punkte & Parameter-Closure-Kettenblatt, RelativeEntropy-/Backreaction-Module, D-Startgates \\
\bottomrule
\end{longtable}

\subsection{Handoff D $\to$ E}
\doclevel{Handoff-Ebene}
\tablecaptionline{Handoff D $to$ E -- Teil 1}
\begin{longtable}{L{0.24\linewidth}L{0.68\linewidth}}
\toprule
Feld & Inhalt \\
\midrule
Quell-Pillar / -bereich & Pillar D: InterfaceCausality, LocalCovariance, Hadamard/Wick, Jacobson-/NCG-/induced-gravity-Routen \\
Ziel-Pillar / -bereich & Pillar E: CPT/Spin/Statistik, Charge/Gauge, chiral matter, Higgs/Yukawa/Flavor, SM/SMEFT-Fenster \\
Übergebene Daten & geometrisch-dynamische Kernlage, modulare/geometrische Brückenobjekte, Hadamard-/Wick-Kontrolle, DHR/DR-nahe Voraussetzungen \\
Mitgegebene Beweise / Instanzen & D-Zulassigkeit, Kausalität, renormierte lokale Observablen, ggf. spektral-geometrische Seeds \\
Nicht übergebene Struktur & keine freie Matter-Sprache ohne DHR/DR-, CPT- und Statistikgates; keine ersatzlose Gleichsetzung von NCG mit E-Kern \\
Rechtsstatus des Handoffs & Gate-gebundene Brücke, nicht bloss rhetorischer Themenwechsel \\
Audit-Punkte & D/E-Gate-Reihenfolge, Feynman-Audit, CPT-/Statistics-/Kramers-Gates, AnomalyInflow-Ordnung \\
\bottomrule
\end{longtable}

\newpage
\section{Migrationsgesetz: was direkt portiert, was umgeschrieben und was dezentriert wird}
\doclevel{Projekt-Ebene und Pfad-Ebene}

\subsection{Direkt oder nahezu direkt portierbarer Kern}
\doclevel{Projekt-Ebene und Pfad-Ebene}
Die folgenden Module oder Modulgruppen besitzen in v0.0605 bereits eine hinreichend allgemeine
mathematische Form, um als CNNA-Unterlage zu dienen. Portierung bedeutet hier: Imports bereinigen,
Namespace umstellen, Namenspolitik anpassen, aber keinen unnötigen semantischen Bruch erzeugen.

\tablecaptionline{Direkt oder nahezu direkt portierbarer Kern -- Teil 1}
\begin{longtable}{L{0.36\linewidth}L{0.25\linewidth}L{0.29\linewidth}}
\toprule
Altmodul & Ziel in CNNA & Disposition \\
\midrule
\endfirsthead
\toprule
Altmodul & Ziel in CNNA & Disposition \\
\midrule
\endhead
\midrule
\multicolumn{3}{r}{Fortsetzung auf der nächsten Seite} \\
\midrule
\endfoot
\bottomrule
\endlastfoot
\path{REALOQS/PillarA/Core/Approximant.lean} & \path{CNNA/Core/Approximant.lean} & nahezu direkt portieren; Namespace und Importkegel bereinigen \\
\path{REALOQS/PillarA/Core/RegionCore.lean} & \path{CNNA/Core/RegionCore.lean} & nahezu direkt portieren; nur Semantik auf Sektorbild umstellen \\
\path{REALOQS/PillarA/Core/RegionNet.lean} & \path{CNNA/Core/RegionNet.lean} & als Trager fur Bright/Dark/Interface-Regionen erhalten \\
\path{REALOQS/PillarA/Core/InfiniteCarrier.lean} & \path{CNNA/Core/InfiniteCarrier.lean} & beibehalten; keine boundary-first Zusatzsemantik \\
\path{REALOQS/PillarA/Core/Selection.lean} & \path{CNNA/Core/Selection.lean} & beibehalten; Beobachterbegriff nicht frei einfuhren \\
\path{REALOQS/PillarA/OQS/DtN.lean} & \path{CNNA/OQS/DtN.lean} & als binarer Sonderfall erhalten; Ausgangspunkt fur GeneralizedDtN \\
\path{REALOQS/PillarA/OQS/DtNStabilized.lean} & \path{CNNA/OQS/DtNStabilized.lean} & als technische Hilfsschicht erhalten, aber methodisch dezentralisieren \\
\path{REALOQS/PillarB/AQFT/StarAlgebra.lean} & \path{CNNA/AQFT/StarAlgebra.lean} & direkt portieren \\
\path{REALOQS/PillarB/AQFT/State.lean} & \path{CNNA/AQFT/State.lean} & direkt portieren \\
\path{REALOQS/PillarB/AQFT/StateNet.lean} & \path{CNNA/AQFT/StateNet.lean} & direkt portieren und später sektorisieren \\
\path{REALOQS/PillarB/AQFT/LocalNet.lean} & \path{CNNA/AQFT/LocalNet.lean} & direkt portieren und später sektorisieren \\
\path{REALOQS/PillarB/AQFT/GNS.lean} & \path{CNNA/AQFT/GNS.lean} & direkt portieren; ComplementGNS darauf aufbauen \\
\path{REALOQS/PillarB/AQFT/KMS.lean} & \path{CNNA/AQFT/KMS.lean} & direkt portieren; ComplementKMS darauf aufbauen \\
\path{REALOQS/PillarB/AQFT/RelativeEntropy.lean} & \path{CNNA/AQFT/RelativeEntropy.lean} & direkt portieren \\
\path{REALOQS/PillarC/OQS/Channel.lean} & \path{CNNA/OQS/SectorChannels.lean} & semantisch verbreitern von System-Umgebung auf Sektor-Kanale \\
\path{REALOQS/PillarC/OQS/Stinespring.lean} & \path{CNNA/OQS/Stinespring-kompatible Unterlage} & direkt als mathematische Grundschicht behalten \\
\end{longtable}

\subsection{Module, die semantisch neu gelesen werden müssen}
\doclevel{Projekt-Ebene und Pfad-Ebene}
Diese Gruppe ist der eigentliche Migrationsmotor. Hier reicht kein Umbenennen; die begriffliche Achse
muss vom binären boundary-first Split auf die Mehrsektor- und Interface-Sprache gedreht werden.

\tablecaptionline{Module, die semantisch neu gelesen werden müssen -- Teil 1}
\begin{longtable}{L{0.36\linewidth}L{0.3\linewidth}L{0.28\linewidth}}
\toprule
Altmodul & Ziel in CNNA & Disposition \\
\midrule
\endfirsthead
\toprule
Altmodul & Ziel in CNNA & Disposition \\
\midrule
\endhead
\midrule
\multicolumn{3}{r}{Fortsetzung auf der nächsten Seite} \\
\midrule
\endfoot
\bottomrule
\endlastfoot
\path{REALOQS/PillarA/OQS/SysEnv.lean} & \path{CNNA/OQS/SectorSysEnv.lean} & binaren Split durch Mehrsektor-Split ersetzen \\
\path{REALOQS/PillarA/Update/TailEliminationDtN.lean} & \path{CNNA/OQS/MultiSchur.lean + CNNA/Integration/TailElimAsChannel.lean} & Tail-Elimination nicht nur eliminatorisch, sondern kanalisch lesen \\
\path{REALOQS/PillarA/Core/TailEliminationCoherence.lean} & \path{CNNA/Core/ParameterClosure.lean + CNNA/OQS/RegularizationClosure.lean} & Koherenz in Closure-/Backreaction-Sprache uberfuhren \\
\path{REALOQS/PillarC/Integration/DerivedSpacetime.lean} & \path{CNNA/PillarD/DerivedSpacetime.lean} & von C-Endpunkt zu D-Vorbereitungsobjekt umdeuten \\
\path{REALOQS/PillarC/Integration/Locality.lean} & \path{CNNA/PillarD/InterfaceCausality.lean + CNNA/PillarD/LocalCovariance.lean} & lokale Einflussrelation in lokale Kovarianz uberfuhren \\
\end{longtable}

\subsection{Was nur noch als Bright-Spezialfall dienen darf}
\doclevel{Projekt-Ebene und Pfad-Ebene}
Gerade die stärksten alten Derived-Dateien sind für Regression und Bright-Recovery wichtig. Sie dürfen
aber die neue Architektur nicht semantisch dominieren.

\tablecaptionline{Was nur noch als Bright-Spezialfall dienen darf -- Teil 1}
\begin{longtable}{L{0.39\linewidth}L{0.21\linewidth}L{0.26\linewidth}}
\toprule
Altmodul & Status & Folge \\
\midrule
\endfirsthead
\toprule
Altmodul & Status & Folge \\
\midrule
\endhead
\midrule
\multicolumn{3}{r}{Fortsetzung auf der nächsten Seite} \\
\midrule
\endfoot
\bottomrule
\endlastfoot
\path{REALOQS/PillarA/Ideal/Adapter/TreeOfCliquesAQFTHandoff.lean} & Bright-Spezialfall / Regressionsorakel & nicht Architekturzentrum, sondern Vergleichspfad \\
\path{REALOQS/PillarB/AQFT/Derived/TreeOfCliquesFromPillarA.lean} & Bright-Spezialfall / Regressionsorakel & nur replayen, nicht zentrieren \\
\path{REALOQS/PillarB/AQFT/Derived/BoundaryMatrixLocalNet.lean} & Bright-Spezialfall unter CNNA/AQFT & alte Randnetz-Semantik als abgeleitete Restriktion rekonstruieren \\
\path{REALOQS/PillarB/AQFT/Derived/BoundaryMatrixStateNet.lean} & Bright-Spezialfall unter CNNA/AQFT & nur als Restriktion auf Bright \\
\path{REALOQS/PillarB/AQFT/Derived/BoundaryMatrixGNS.lean} & Bright-Spezialfall unter CNNA/AQFT & nur als Restriktion auf Bright \\
\path{REALOQS/PillarB/AQFT/Derived/BoundaryMatrixFullKMSFromL.lean} & Bright-Spezialfall unter CNNA/AQFT & nur als Restriktion auf Bright \\
\path{REALOQS/PillarB/AQFT/Derived/BoundaryMatrixQuasiLocalClosure.lean} & Bright-Spezialfall / Regression & keine architektonische Zentralstellung \\
\path{REALOQS/PillarB/AQFT/Derived/BoundaryMatrixModular.lean} & Bright-Spezialfall / Regression & modulare Daten später aus CNNA zuruckgewinnen \\
\end{longtable}

\subsection{Späte Seeds und diagnostische Inputs für D und E}
\doclevel{Projekt-Ebene und Pfad-Ebene}
Diese Module sind wertvoll, aber nur als Input, Hook oder späte diagnostische Spur. Sie dürfen vor der
jeweiligen Closure-Schicht nicht als schon geschlossene Ontologie der neuen Pfeiler missverstanden werden.

\tablecaptionline{Späte Seeds und diagnostische Inputs für D und E -- Teil 1}
\begin{longtable}{L{0.39\linewidth}L{0.25\linewidth}L{0.28\linewidth}}
\toprule
Input aus v0.0605 & Ziel in D/AS/E & Regel \\
\midrule
\endfirsthead
\toprule
Input aus v0.0605 & Ziel in D/AS/E & Regel \\
\midrule
\endhead
\midrule
\multicolumn{3}{r}{Fortsetzung auf der nächsten Seite} \\
\midrule
\endfoot
\bottomrule
\endlastfoot
\path{REALOQS/PillarA/Kinematics/DiscreteSpacetime.lean} & \path{CNNA/PillarD/DerivedSpacetime.lean} & als Seed-Herkunft, nicht als fertige Raumzeit \\
\path{REALOQS/PillarA/Kinematics/Causality.lean} & \path{CNNA/PillarD/InterfaceCausality.lean} & nur nach expliziter Interface-Herkunft de-seeden \\
\path{REALOQS/PillarA/Geometry/DimensionHooks.lean} & \shortstack[l]{\path{CNNA/PillarD/}\\\path{SpectralDimensionFlow}\\\path{Seed.lean}} & zunächst Hook/Seed \\
\path{REALOQS/PillarA/RG/ScaleWindow.lean} & \shortstack[l]{\path{CNNA/PillarD/}\\\path{HorizonLevelSeed.lean}\\und\\AS-Strang} & \shortstack[l]{Skalenfenster nicht ontologisch,\\sondern als derived Diagnose lesen} \\
\path{REALOQS/PillarA/Physics/ScaleBreaking*.lean} & \path{CNNA/PillarD/Backreaction- und Closure-Diagnostik} & nur späte Diagnostik, kein Kernaxiom \\
\end{longtable}

\warningbox{\textbf{Harte Migrationsregel.} Kein boundary-first Dateiname, kein freies \(L_{\max}\), kein freier
State-, Dictionary- oder Observer-Begriff und kein unmarkiertes Spezialfenster darf auf dem kritischen
CNNA-Pfad als echter Kernbegriff erscheinen. Diese Regel folgt direkt aus der internen Spezifikation.\cite[pp.~2--9]{CNNASpec}}

\newpage
\section{Pillar A -- ontische Stammtheorie und Komplementarchitektur}
\doclevel{Pfad-Ebene}

\subsection{Zielkern von A}
\doclevel{Pfad-Ebene}
Pillar A muss in CNNA den ontischen Stamm liefern: ToC, Anker, sichtbarer Patch,
Komplementfamilie, Mehrsektor-Split, verallgemeinerte DtN-/Schur-Mechanik und erste
Parameter-Closure. A ist in der neuen Architektur nicht bloss Lieferant eines Randoperators, sondern
Generator der Sektorgrammatik selbst. Das Snowmass-Argument gegen eine einzig privilegierte
Axiomatik stärkt genau diesen Schritt: CNNA darf nicht auf der ersten mathematisch bequemen
Oberfläche stehenbleiben.\cite[pp.~1, 10--12]{Dedushenko2023}

\subsection{Konkrete Vorgehensweise in A}
\doclevel{Pfad-Ebene}
\paragraph{1. Tragerkern einfrieren und bereinigen.}
Portiere \path{Approximant}, \path{RegionCore}, \path{RegionNet}, \path{InfiniteCarrier},
\path{Selection}, \path{Interfaces}, \path{Determinism}, \path{BoundaryPorts},
\path{DirichletLaplacian} und \path{DirichletLaplacianBridge} in einen neuen
\path{CNNA/Core}- und \path{CNNA/ToC}-Baum. In diesem ersten Schnitt werden nur diejenigen Deklarationen
übernommen, deren mathematische Bedeutung ohne boundary-first Zusatzsprache stabil bleibt.
Alles, was implizit bereits den alten Bright-Randkanal privilegiert, wird aus dem aktiven Pfad entfernt
oder markiert.

\paragraph{2. BranchPatch und Komplementfamilie als erste Klasse einführen.}
Schreibe \path{CNNA/Core/BranchPatch.lean} und
\path{CNNA/Core/ComplementSectorFamily.lean} nicht als freie Container, sondern als direkte
Verankerung an ToC-Anker, Pfadstruktur, Zellen und Koarsenierungsdaten. Der alte Approximant ist
nur der helle Ausschnitt; der dunkle Sektor entsteht aus einer \emph{organisierten} Komplementfamilie,
nicht aus \(S\setminus R\) als namenlosem Rest. Damit verschiebt sich der Grundbegriff des Projekts von
Boundary zu Sector.

\paragraph{2a. ToC-Varianten nur über einen gemeinsamen CNNA-Vertrag zulassen.}
Die Roadmap sollte hier nicht nur verbal, sondern mathematisch explizit bleiben: der ToC ist im
CNNA-Kontext kein beliebiger Graphname, sondern ein Trager mit Ankern, Pfad-/Trunk-/Side-Branch-
Information, lokaler Inzidenz, Koarsenierung und Source-Maps in Richtung Split- und Kopplungsschicht.
Genau diese Daten müssen später \code{BranchPatch}, \code{ComplementSectorFamily},
\code{SectorSplit} und \code{GeneralizedDtN} speisen. Diskrete, kontinuierliche oder graduierte
Varianten des ToC sind daher nur als \emph{Varianten eines gemeinsamen Vertrags} zulässig: sie dürfen
Topologie, Regularität oder Zusatzsymmetrien ändern, müssen aber dieselben sektorbezogenen Witness-
und Restriktionspflichten erfüllen. Die Frage ist also nicht ``welcher ToC-Stil gefällt'', sondern ob
eine Variante die Split- und Komplementlogik funktoriell trägt.

\paragraph{2b. Prototypische B-/C-Schnittstellen früh als Seeds anlegen.}
Um Typ-Sackgassen zu vermeiden, darf A dem späteren B-/C-Pfad kontrolliert vorauseilen.
Konkret heißt das: neben den eigentlichen Kernobjekten dürfen frühe Sketch-Dateien oder
Seed-Schnittstellen für spätere Verbraucher mitlaufen, etwa für
\code{ComplementLocalNet}-Signaturen, sektorweise Restriktionen, Glueing-Daten,
komplementbasierte Zustandsprojektionen oder einfache Platzhalter für modulare Aktionen.
Solche Deklarationen sind gerade \emph{nicht} die vollständige Mathematik von B oder C; sie
müssen als \code{PrototypeInterface}, \code{Seed} oder \code{Scaffold} markiert bleiben und
werden nur daran gemessen, ob sie die A-seitigen Typen, Exporte und Indizes später tragfähig
machen. Was auf diesem Weg nur mit ad-hoc Koerzionsketten, unsauberen Universumsabkürzungen
oder boundary-first Sonderannahmen compilieren würde, ist ein Frühwarnsignal gegen den
aktuellen A-Vertrag.

\paragraph{3. Binären Split durch Mehrsektor-Split ersetzen.}
Aus \path{SysEnv}- und \path{SplitInvariants}-artigem Material wird ein neuer
\path{SectorSplit} aufgebaut. Bright, Dark und Interface müssen gemeinsam auftreten; insbesondere
muss das Interface als eigener Träger modelliert werden, nicht als implizite Grenzspur zwischen zwei
anderen Datenblöken. Genau an dieser Stelle wird später die Glueing-, Entropie- und Defektlogik
aufgesetzt.

\paragraph{4. GeneralizedDtN und MultiSchur als mathematischen Kopplungskern bauen.}
Der alte \path{DtN}-Pfad wird nicht entsorgt, sondern als binärer Grenzfall verstanden. Die
mathematische Strategie auf dem aktiven Pfad ist jedoch strenger: zürst eine wohldefinierte
Mehrsektor-Blocksicht auf Bright, Dark und Interface herstellen, dann globalen und iterierten
Schur-/DtN-Abbau gegeneinander beweisen und erst danach spezielle Stabilisierung, Positivitäts- oder
Regularisierungsfragen aufsetzen. Genau dieser Vergleich von globaler und iterierter Elimination ---
architektonisch repräsentiert durch \code{MultiSchur} und mathematisch zugespitzt im Zieltyp von
\code{iteratedSchur\_eq\_globalSchur} --- ist das nichttriviale Kerntheorem des GeneralizedDtN-
Strangs; er darf nicht hinter spätere Spezialfälle oder numerische Hilfskonstruktionen
ausweichen.

Im CNNA-Kern ist \code{GeneralizedDtN} algebraisch stärker als eine einzelne Randmatrix: gebraucht
wird mindestens eine sektorindexierte Block- bzw. Operatorkernstruktur, die Bright-, Dark- und
Interface-Komponenten zugleich koppelt, Schur-Reduktion und Restriktion entlang von Teilsektoren
zulässt und mit Symmetrie-/Positivitätszeugen auftritt, statt nur eine ad-hoc Matrix zu liefern.
Genau dieser Kern muss die spätere \code{ComplementLocalNet}-Schicht speisen: lokale Algebren dürfen
nicht losgelöst vom Kopplungskern entstehen, sondern müssen aus sektorweisen Ports, Restriktionen
und MultiSchur-Reduktionen hervorgehen. Lean-seitig spricht das klar für eine \code{structure}-
basierte Darstellung mit expliziten Feldern für Sektorindex, Kern, Restriktion, Glueing und
Kompatibilitätszeugen, während Typeclasses und kategoriale Hülle erst nachgeordnet für
Instanzen, Morphismen oder Funktorialität sinnvoll sind. Dazu gehört auch, \path{MultiSchur.lean}
als eigenes Zielmodul sichtbar zu halten, statt den global-vs.-iteriert-Vergleich implizit in
Hilfslemmata zu verstecken.

\paragraph{5. Parameter-Closure bereits in A vorbereiten.}
Freie Strukturparameter aus dem Altpfad --- vor allem \(b\), \(1<b\), Horizon-/IR-Größen,
\(L_{\max}\) und sekundäre numerische Stabilisierungsterme --- dürfen nicht unreflektiert in CNNA
übernommen werden. In A ist deshalb früh eine \path{BranchingWitness}- und
\path{ParameterClosure}-Schicht zu schreiben, die drei Ebenen sauber trennt: ontische
Mikrostrukturparameter wie \(b\), dynamische Emergenzgrößen wie \(L_{\max}\) bzw. später
\code{HorizonLevel}, sowie bloss numerische Regularisierer. Für \(b\) bedeutet das: nicht nur
\(1<b\) als Nichttrivialitätsmarker festhalten, sondern über \code{UVSpectralSelector},
\code{BranchingSelector} und \code{BackreactionSelector} einen späteren kombinierten
Selektorpfad bis zu einer abgeleiteten \code{selectedBranching}-Lesart vorbereiten. Für
\(L_{\max}\) bedeutet es, dass der Name auf dem aktiven Pfad höchstens Scaffold-Status hat und
später durch \code{HorizonLevel} bzw. einen \code{BackreactionFixedPoint} ersetzt wird. Für
Regularisierer bedeutet es schliesslich, dass sie entweder ganz verschwinden oder als abgeleitete
Spektral-/Nullmodenreste in eine \code{RegularizationClosure} degradiert werden.\cite[pp.~2, 5--9, 29--31, 44--47]{CNNASpec}



\paragraph{6. Exporte neu schreiben, nicht den alten Handoff verabsolutieren.}
Der alte \path{AQFTHandoff}-Gedanke bleibt wichtig, aber sein Output darf nicht wieder nur
\((\Omega,L,\beta)\) für einen Bright-Rand sein. Der neue Export aus A muss die Sektortriple
\((\text{Bright},\text{Dark},\text{Interface})\) samt Split-, GeneralizedDtN- und Source-Map-Daten
bereitstellen. Erst darauf dürfen B und C aufsetzen.

\subsection{Abnahmekriterien für A}
\doclevel{Pfad-Ebene}
A ist erst dann erfolgreich migriert, wenn:
\begin{itemize}
\item \code{BranchPatch}, \code{ComplementSectorFamily}, \code{SectorSplit} und \code{GeneralizedDtN} als echte Kernobjekte vorliegen;
\item der alte binäre DtN-Fall als Sonderfall erklärbar ist;
\item \(L_{\max}\) auf dem aktiven Pfad nicht mehr ontisch gebraucht wird;
\item der Export aus A sektorisch und nicht boundary-first formuliert ist;
\item prototypische B-/C-Schnittstellen auf A-Basis ohne ad-hoc Typreparaturen, boundary-first Sonderannahmen oder versteckte freie Beobachter compiliert werden können;
\item ein detaillierter B-Phasen- bzw. Unterplan erst dann als belastbar gilt, wenn das A-Gate nicht nur mathematisch, sondern auch interface-stabil gegen diese Vorgriffe ist.
\end{itemize}

Als operative \emph{Definition of Done} für den Übergang zu detaillierter B-Planung gilt damit:
\code{BranchPatch}, \code{ComplementSectorFamily}, \code{SectorSplit}, \code{GeneralizedDtN} und
sektorischer Export stehen; zugleich kompilieren die frühen \code{PrototypeInterface}-Skizzen für
Netz-, Zustands- und Glueing-Verbraucher ohne semantische Sonderhaken. Erst auf dieser Lage wird aus
B-Planung mehr als eine reversible Absichtserklärung.

\subsection{Was in A strikt zu vermeiden ist}
\doclevel{Pfad-Ebene}
\begin{itemize}
\item direkte Portierung der alten boundary-first Dateinamen als Zentrum des neuen Modulbaums;
\item Einschmuggeln eines freien Beobachters;
\item frühe numerische Stabilisierung als Ersatz für hergeleitete Struktur;
\item Bright-Recovery schon in A selbst als globales Architekturziel.
\end{itemize}

\newpage
\section{Pillar B -- AQFT aus dem Komplementnetz statt aus dem Bright-Rand}
\doclevel{Pfad-Ebene}

\subsection{Zielkern von B}
\doclevel{Pfad-Ebene}
Pillar B muss aus der in A bereitgestellten Sektorstruktur ein wirkliches Komplementnetz bauen:
\code{ComplementLocalNet}, \code{InterfaceLocalNet}, \code{ComplementStateNet},
\code{ComplementGNS}, \code{ComplementKMS}, dazu die quasilokale und Haag-Kastler-nahe
Unterlage. Snowmass liest AQFT, lokal kovariante QFT, homotopy AQFT, Functorial Field Theory und
CFT gerade als unterschiedliche QFT-Faces; CNNA muss deshalb in B die AQFT-Schicht so bauen,
dass sie später für mehrere Faces offen bleibt.\cite[pp.~4--12]{Dedushenko2023}

\subsection{Konkrete Vorgehensweise in B}
\doclevel{Pfad-Ebene}
\paragraph{1. Die generische AQFT-Basis unverzerrt portieren.}
\path{StarAlgebra}, \path{State}, \path{StateNet}, \path{LocalNet}, \path{GNS}, \path{KMS},
\path{RelativeEntropy}, \path{QuasiLocal/*} und die relevanten Haag-Kastler-Bausteine
werden zürst mathematisch sauber nach \path{CNNA/AQFT} portiert. Wichtig ist, dass diese
Grundschicht nicht stillschweigend Bright-spezifische Matrix-Semantik erbt.

\paragraph{1a. Die bestehende Gate-Schicht aus B explizit mitportieren.}
Die 20 Dateien unter \path{REALOQS/PillarB/Gates} sind keine dekorative Audit-Prosa, sondern die
formale Validierungsschicht für StarAlgebra-, Norm-, StateRestriction-, QuasiLocal- und
Haag-Kastler-Aussagen. In CNNA werden sie daher als eigene Gate-Ebene neu aufgesetzt:
\path{CNNA/Gates/AQFT/*}, \path{CNNA/Gates/HaagKastler/*},
\path{CNNA/Gates/ComplementAdditivity.lean}, \path{StateRestriction.lean},
\path{SplitProperty.lean} und funktoriale Netzgates. Besonders die Teilgates für
Isotonie, Lokalität, Additivität, starke Additivität, Kovarianz, Spektrum, Time-Slice und Vakuum
dürfen nicht in allgemeiner Prosa verschwinden, weil sie später die Abnahme des Komplementnetzes
gegen den Bright-Recovery-Pfad definieren.\cite{HaagKastler1964}\cite{Haag1996}

\paragraph{1b. Die algebraische Grundschicht von Anfang an Z2-offen anlegen.}
Die direkte Portierung von \path{StarAlgebra.lean} bleibt richtig, reicht für die spätere
Spin-Statistik-Schicht aber nicht aus. Sobald CNNA bosonische und fermionische Sektoren sauber
unterscheiden will, braucht die B-Grundschicht mindestens eine optionale Z2-Graduierung mit
Superkommutator, Paritätsmarkierung und kontrollierter Rückbindung an die ungraduierte
Sternalgebra. Deshalb wird früh zusätzlich \path{CNNA/AQFT/GradedStarAlgebra.lean} geplant.
Das ist kein später Luxus, sondern die Minimalvorkehrung dafür, dass
\code{FermionicStatisticsSeed}, \code{StatisticsGate}, \code{SpinRepresentationSeed} und später
\code{KramersGate} nicht als freie Zusatzetiketten über einer rein bosonisch gelesenen Algebra
schweben. Die mathematische Richtung ist durch die klassische Spin-Statistik-Literatur und ihre
algebraische AQFT-Fortsetzung klar vorgegeben. \code{GradedStarAlgebra} und
\code{SeparatingProperty} sind dabei bewusst \emph{logisch unabhängig}: die erste Datei betrifft
die algebraische Form der Observablenbasis, die zweite die spätere Darstellungslage des
Komplementnetzes. Weder darf daher eine künstliche Importkette
\code{GradedStarAlgebra -> SeparatingProperty} behauptet werden noch umgekehrt.
\cite{Pauli1940}\cite{GuidoLongo1995}

\paragraph{2. ComplementLocalNet vor jeder Recovery schreiben.}
Der erste echte Netzbau in B muss das dunkle Netz sein, nicht die Wiederholung des alten
BoundaryMatrixNet. Konzeptionell heißt das: Regionen aus \code{BranchPatch}/
\code{ComplementSectorFamily}/\code{SectorSplit} ableiten, Algebren aus blockweiser Operator- und
GeneralizedDtN-Struktur zuweisen, danach Isotonie und Lokalität beweisen. Erst wenn dieser Schritt
trägt, darf ein Bright-Spezialfall überhaupt als Regression dienen.

\paragraph{3. InterfaceLocalNet als eigenen mathematischen Gegenstand behandeln.}
Die größte architektonische Versuchung wäre, das Interface nur als technische Zwischenzone zu
benutzen. Das wäre strategisch falsch. Das Interface-Netz ist die Stelle, an der später
Glueing-Funktorialität, Defekte, Entropie, thermische Balance und Dualitätsdaten gesammelt werden.
Deshalb sollte \path{CNNA/AQFT/InterfaceLocalNet.lean} parallel zu
\path{ComplementLocalNet.lean} geschrieben werden, nicht als späte Nachbesserung.

\paragraph{4. Zustandsanschluss strikt aus Sektorrelationen gewinnen.}
\code{ComplementStateNet} darf kein freier \code{State}-Trager sein. Zustände müssen als Restriktionen,
Pushforwards oder Balance-Zustände auf Bright, Dark und Interface erscheinen. Erst wenn diese
Restriktionen funktionieren, lohnt sich der Schritt zu \code{ComplementGNS} und \code{ComplementKMS}.
Thermik und Modularität gehören also in B an das Ende des sektorischen Zustandsaufbaus, nicht an
dessen Anfang.

\paragraph{4a. Zwischen ComplementGNS und ComplementKMS eine echte modulare Kernschicht einfügen.}
Die im Archiv vorhandene Bright-Recovery-Datei \path{BoundaryMatrixModular.lean} ist als Regression nützlich,
aber sie extrahiert weder die Tomita-Operation noch modulare Konjugation \(J\) und modularen
Operator \(\Delta\) als eigenständige CNNA-Objekte. Für die neue Architektur reicht das nicht. Nach
\code{ComplementGNS} muss daher eine Reeh--Schlieder-artige Separationsschicht folgen, die den
GNS-Vektor nicht nur als zyklisch, sondern für die relevanten Komplementalgebren auch als
separierend kontrolliert. Erst darauf dürfen \path{CNNA/AQFT/SeparatingProperty.lean},
\path{CNNA/AQFT/TomitaTakesakiData.lean} und
\path{CNNA/AQFT/ModularConjugation.lean} aufsetzen. Der Zielpunkt ist nicht blosse Terminologie,
sondern ein Theoremstil der Form \code{modular\_conjugation\_maps\_bright\_to\_complement}, also
\(J M J = M'\) in der für CNNA passenden Bright/Dark/Interface-Lesart. Genau diese Stufe gehört in
den B-Kern, nicht in eine späte Matter-Vorzone.\cite{ReehSchlieder1961}\cite{Takesaki1970}\cite{GuidoLongo1995}

\paragraph{4b. Die modulare B->D-Brücke explizit als Bisognano--Wichmann-Seed markieren.}
Sobald Tomita--Takesaki-Daten vorhanden sind, darf CNNA den Übergang nach D nicht nur heuristisch
über ``Entropie'' oder ``Kausalität'' erzählen. Vor dem eigentlichen Bisognano--Wichmann-Vergleich
muss dabei zusätzlich eine explizite Spektralkontrolle des \emph{Komplementnetzes} vorliegen:
\path{CNNA/AQFT/ComplementSpectrumCondition.lean} sammelt Vorwärtskegel-, Positivitäts- und
Translationsdaten auf dem aktiven derived-only Pfad, und das Gate
\code{gate\_spectrum\_condition\_on\_complement\_net} stoppt jede spätere Guido--Longo-,
Bisognano--Wichmann- oder Haag--Ruelle-Lesart, solange diese Kontrolle nur aus Bright-Recovery
importiert wäre. Erst darauf darf mit \path{CNNA/AQFT/BisognanoWichmannSeed.lean} eine eigene
Brückendatei als \emph{spätes B-Modul} folgen, die modulare Automorphismen mit abgeleiteten
Boost- bzw. Wedge-Strukturen vergleicht. Der Zielstil ist
\code{modular\_automorphism\_group\_equals\_derived\_boost\_on\_wedge\_sector}. Die Stop-Regel des
Seeds selbst lautet dabei: \emph{nach} \code{ComplementKMS} und
\code{ComplementSpectrumCondition}, \emph{vor} den pfeilerübergreifenden PT/CPT-Brückenobjekten;
der spätere Abgleich mit \code{InterfaceCausality} wird \emph{nicht} als Importvoraussetzung des
Seeds selbst kodiert, sondern erst über das Gate
\code{gate\_modular\_boost\_consistent\_with\_interface\_causality} vollzogen. So wird die Relation
zwischen B-Modularität und D-Geometrie physikalisch erzwungen statt bloss architektonisch behauptet,
ohne einen B$\leftrightarrow$D-Zirkel in die Dateitopologie einzubauen.
\cite{BisognanoWichmann1975}\cite{BisognanoWichmann1976}\cite{GuidoLongo1995}\cite{Borchers1992}


\subsection{B-Kernformalismus der modularen Darstellungsarchitektur}
\doclevel{Pfad-Ebene}
Der relevante B-Kern endet deshalb nicht bei einem losen Paar
\code{ComplementGNS}/\code{ComplementKMS}. Die operative Kette ist strikt:
\begin{center}
\small
\tablecaptionline{B-Kernformalismus der modularen Darstellungsarchitektur -- Teil 1}
\begin{tabular}{@{}l@{}}
$\text{ComplementStateNet} \to \text{ComplementGNS} \to \text{SeparatingProperty}$\\
$\to \text{TomitaTakesakiData} \to \text{ModularConjugation}$\\
$\to \text{ComplementKMS} \to \text{ComplementSpectrumCondition}$\\
$\to \text{BisognanoWichmannSeed}$.
\end{tabular}
\end{center}
Die ersten fünf Pfeile gehören in den B-Kern selbst. \code{ComplementGNS} allein reicht nicht,
weil ohne Separierung weder Tomita--Takesaki noch eine kontrollierte modulare Konjugation verfügbar
sind; \code{ComplementKMS} ist dann kein Ersatz für diese Schicht, sondern ihr thermischer
Verbraucher. Erst der anschliessende \code{BisognanoWichmannSeed} schliesst den späten modularen
B-Unterpfad und macht die modulare Struktur für D als Geometrie-/Kausalitätsbrücke operational;
die Konsistenz mit \code{InterfaceCausality}/\code{LocalCovariance} wird anschliessend über das
eigene Gate \code{gate\_modular\_boost\_consistent\_with\_interface\_causality} geprüft. Die
operative Feingliederung dieser Kette als Schritte 6a--6i wird in
Abschnitt~\ref{sec:b-unterpfad} zusammengestellt. Die Nuclearity-Linie beginnt dabei \emph{nicht}
schon vor \code{ComplementKMS}, sondern streng danach: sie ist B-intern, aber eine
Nach-KMS-Schicht auf bereits geschlossener modularer Darstellungsarchitektur.

\subsection{Warum das Komplementnetzwerk für full-derived lokale QFT physikalisch zwingend wird}
\doclevel{Pfad-Ebene}
Die Roadmap muss hier eine asymmetrische, aber wichtige Präzisierung festhalten. Ein einzelnes
Kramers-Theorem in endlicher Quantenmechanik oder eine isolierte Antiunitaritätsaussage lässt sich
auch ohne voll ausgearbeitete Komplementarchitektur formulieren. CNNA zielt jedoch nicht auf lose
Einzelsätze, sondern auf eine gemeinsame Rekonstruktion von lokaler Algebra, Zustand,
Darstellung, Modularität, Kausal-/Geometriestruktur und daraus erst CPT-, Spin-Statistik- und
Kramers-Lesarten. In genau diesem physikalischen Sinn wird das Komplementnetzwerk weitgehend
zwingend.

Erstens ist das bereits durch den derived-only Kern der Spezifikation vorgegeben: Regionen,
Komplementfamilien, Interface-Kanäle, Zustandsrestriktionen und Backreaction sollen nicht als
sekundäre Dekoration auftreten, sondern als Herkunftsträger der späteren Physik.
\code{ComplementSectorFamily}, \code{SectorSplit}, \code{ComplementLocalNet},
\code{ComplementStateNet}, \code{ComplementGNS} und \code{ComplementKMS} bilden daher nicht nur
einen alternativen Pfad, sondern die Stelle, an der aus Bright/Dark/Interface überhaupt eine
physikalisch belastbare lokale Theorie emergiert.\cite[pp.~1--10, 30--36]{CNNASpec}

Zweitens hängt gerade die angestrebte vN-/Typ-III-Lesart nicht an einer isolierten Matrixalgebra,
sondern an der Verschänkung von Lokalität, Darstellung und modularer Komplementarität. Ohne
separierende Vektoren, Tomita--Takesaki-Daten und die Bright/Dark/Interface-Lesart von
\(J M J = M'\) bliebe höchstens ein Typ-I-naher Testträger übrig; die für lokale AQFT relevante
Komplementstruktur wäre dann weder algebraisch noch physikalisch geschlossen.\cite{Takesaki1970}\cite{GuidoLongo1995}

Drittens sitzen CPT, Spin-Statistik und Kramers in der hier angestrebten Lesart gerade nicht auf
einem einzelnen Hamiltonoperator, sondern auf dem Zusammenspiel von lokaler Algebra,
Darstellungsraum, Spektrum, Antiunitarität, modularer Struktur und später Kausal-/Geometrie-
Information. Ohne komplementbasierte Darstellungs- und Modularschicht könnten diese Sätze zwar
formal irgendwo nachgebildet werden; sie wären im CNNA-Sinn aber nicht mehr \emph{aus derselben
Stammtheorie} rekonstruiert.\cite{BisognanoWichmann1975}\cite{BisognanoWichmann1976}\cite{GuidoLongo1995}

Die operative Designlesart lautet deshalb scharf: Ohne Komplementnetz keine echte full-derived lokale
QFT-Physik; und ohne komplementbasierte Darstellungs- und Modularstruktur bleiben CPT,
Spin-Statistik und Kramers in CNNA entweder unvollständig oder bloss aufgesetzt. Genau deshalb
gehört die Komplementseite nicht in einen späten Appendix, sondern in den tragenden Kernpfad von B.

\paragraph{5. Vakuumfrage, Grundzustände und derived-only Ersatzbegriffe.}
Auch die Rolle des Vakuums muss im Roadmap-Text explizit geklärt werden. Im alten Bright-Recovery-
Pfad existiert ein konkreter Grund- bzw. Vakuumkandidat als Grenz- oder Niedrigtemperaturlesart des
bekannten Zustandsaufbaus. Im CNNA-Kern gibt es hingegen \emph{kein} freies Vakuumpostulat: die
aktive Architektur verlangt, dass relevante Zustandslagen aus Komplement-/Interface-Daten,
Restriktionen und Balancebedingungen hervorgehen. Die \code{SeparatingProperty} des GNS-Vektors
übernimmt hier die Rolle eines abgeleiteten Vakuumersatzes für die modulare AQFT-Lesart; auf
späteren emergenten Hintergründen in D wird diese Sonderrolle nicht durch ein global ausgezeichnetes
Vakuum, sondern durch die Hadamard-/mikrolokale Zulässigkeitsklasse ersetzt. Gerade diese
Verschiebung von postuliertem Vakuum zu abgeleiteter Zustands- und Separationslage ist für eine
strictly derived-only Lesart entscheidend.\cite{Haag1996}\cite{ReehSchlieder1961}\cite{BFV2003}

\paragraph{5a. Kontinuums- und thermodynamischen Limes explizit markieren.}
Der Bright-Approximant bleibt endlich; die spätere AQFT-Lesart von quasilokaler Algebra,
Tomita--Takesaki, Typ-III-Struktur, Bisognano--Wichmann und Haag--Ruelle ist es gerade nicht.
Darum muss auf dem aktiven Pfad zwischen endlichem Patch und unendlicher/quasilokaler
Darstellungslage ein expliziter Übergang via \code{InfiniteCarrier}, quasilokaler Vervollständigung
oder thermodynamischer Limesmarkierung auftauchen. Dieser Limes ist nicht bloss technische
Nachpolitur, sondern stoppt den Pfad, solange unendliche Kommutanten- und Modularstruktur nur aus
endlichen Bright-Matrizen gelesen würde. Offen bleiben darf, \emph{welcher} operative Schliessungsweg
am Ende gewinnt; offen bleiben darf aber nicht, \emph{dass} zwischen endlichem Approximanten und
AQFT-Lage ein eigener Schritt liegt. Die Roadmap sollte deshalb drei zulässige Alternativpfade
benennen: (i) \emph{InfiniteCarrier by construction}, falls das aktive Komplementnetz bereits als
unendliche/quasilokale Struktur aufgebaut wird und der endliche Bright-Patch nur Restriktionsfenster
bleibt; (ii) \emph{induktive/quasilokale Vervollständigung} aus einer gerichteten Familie wachsender
\code{BranchPatch}- bzw. Sektor-Ausschnitte; (iii) \emph{thermodynamischer/KMS-Limes} mit expliziten
uniformen Spektral-, Positivitäts- und Modularitätskontrollen. Nicht feststehen muss heute, welche
Route gewinnt; feststehen muss aber, dass mindestens eine dieser Routen samt Gate
\code{gate\_infinite\_carrier\_or\_quasilocal\_limit\_defined} wirklich geschlossen wird, bevor Typ-III-,
BW- oder Scattering-Lesarten als derived beansprucht werden.

\paragraph{5b. \(\beta\), Unitarität und Positivität nicht als freie Restparameter stehenlassen.}
Im alten Bright-Pfad ist \(\beta\) ein gegebener Input des Gibbs-/KMS-Aufbaus. Auf dem aktiven
CNNA-Pfad darf das nicht das Endbild bleiben. Entweder wird \(\beta\) später aus modularer
Automorphismengruppe und ihrer physikalischen Zeitnormalisierung abgeleitet --- die stärkste Lesart
entsteht genau dann, wenn Bisognano--Wichmann den modularen Fluss an emergente Geometrie bindet ---
oder \(\beta\) bleibt bis zur \code{BackreactionFixedPoint}-Schicht ein explizit markierter
Scaffold- bzw. Fixpunktparameter. Analog dürfen Unitarität der effektiven Entwicklung und positive
Energie-/Vorwärtskegelkontrolle nicht nur stillschweigend vorausgesetzt werden: \code{ComplementSpectrumCondition}
trägt die Positivitätsseite für das Komplementnetz, während C und D eine explizite Kontrolle
benötigen, dass ihre physikalisch interpretierten Flüsse und Generatoren nicht in unkontrollierte
Nichtunitarität oder negative-Energie-Heuristiken kippen.\cite{Haag1996}\cite{BisognanoWichmann1975}\cite{BisognanoWichmann1976}

\subsection{Zwingende AQFT-Modulergänzungen im Modulbaum}
\doclevel{Pfad-Ebene}
Die modulare Zwischenschicht muss im B-Modulbaum ausdrücklich und nicht nur implizit auftauchen.
Für den kritischen AQFT-Unterbaum sind daher mindestens die folgenden Dateien zwingend:

\tablecaptionline{Zwingende AQFT-Modulergänzungen im Modulbaum -- Teil 1}
\begin{longtable}{L{0.39\linewidth}L{0.53\linewidth}}
\toprule
Zieldatei & Rolle im B-Kern \\
\midrule
\endfirsthead
\toprule
Zieldatei & Rolle im B-Kern \\
\midrule
\endhead
\bottomrule
\endfoot
\path{CNNA/AQFT/GradedStarAlgebra.lean} & optionale, aber früh verankerte Z2-Graduierung der algebraischen Grundschicht; keine fermionische Semantik vor dieser Datei \\
\path{CNNA/AQFT/SeparatingProperty.lean} & Reeh--Schlieder-artige Separationsstufe zwischen \code{ComplementGNS} und modularer Maschine \\
\path{CNNA/AQFT/TomitaTakesakiData.lean} & extrahiert \(S\), \(J\) und \(\Delta\) aus der komplementnetz-basierten Darstellung \\
\path{CNNA/AQFT/ModularConjugation.lean} & kontrolliert die Bright/Dark/Interface-Lesart von \(J M J = M'\) als eigentliche modulare Darstellungsarchitektur \\
\path{CNNA/AQFT/ComplementSpectrumCondition.lean} & explizite Spektral- und Vorwärtskegelkontrolle des Komplementnetzes vor jeder BW-, Guido--Longo- oder Haag--Ruelle-Lesart \\
\path{CNNA/AQFT/BisognanoWichmannSeed.lean} & modulare B$\to$D-Brücke; verknüpft Wedge-/Boost-Strukturen mit modularen Automorphismen \\
\end{longtable}

Diese Dateigruppe ist \emph{kein} später E-Zusatz. Sie gehört in den AQFT-Unterbaum selbst und
schliesst die Lücke zwischen Zustand, Darstellung, Modularität und Geometrie. Die späteren
Symmetrie-/Statistik-Dateien wie \code{TimeReversalOnGNS}, \code{AntiunitarySeed},
\code{SpinRepresentationSeed}, \code{FermionicStatisticsSeed}, \code{CPTGate} und
\code{KramersGate} konsumieren diese B-Kernschicht, ersetzen sie aber nicht.

\paragraph{6. Die alten BoundaryMatrix-Dateien nur als Extraktionsquelle verwenden.}
Aus den alten Derived-Dateien werden genau diejenigen Teile herausgelöst, die wirklich allgemein sind:
lineare Algebra, allgemeine Netztechnik, Restriktions- oder Closure-Lemmata. Alles ToC-gebundene am
hellen Rand wird explizit in einen Bright-Recovery-Strang verschoben. Damit bleibt die mathematische
Arbeit der alten Dateien nutzbar, ohne dass deren Semantik das neue Projekt dominiert.

\paragraph{7. Quasilokale und HK-nahe Zusatzstrukturen erst nach Netzstabilität anschliessen.}
Sobald dunkles Netz, Interface-Netz und Zustandsrestriktionen tragen, wird die quasilokale
Ergänzung sowie die HK-nahe Support-/Split-/StateRestriction-Schicht angeschlossen. Für diesen
Schritt sind die Snowmass-Hinweise auf AQFT, LCQFT, dynamical locality und spätere homotopische
Verallgemeinerungen wichtig: CNNA sollte hier strukturell offen bleiben.\cite[pp.~4--12]{Dedushenko2023}\cite{HaagKastler1964}\cite{BFV2003}\cite{FewsterLang2016}\cite{BeniniSchenkelWoike2021}

\paragraph{8. Superselection-, Charge- und Feldrekonstruktionsschicht explizit vorsehen.}
Spätestens auf stabiler HK-/Quasilokal-Unterlage braucht CNNA eine observablenseitige
Superselection-Schicht. Deshalb werden \path{CNNA/AQFT/DHRSectorsSeed.lean} und
\shortstack[l]{\path{CNNA/AQFT/}\\\path{FieldAlgebraReconstruction}\\\path{Seed.lean}} ausdrücklich eingeplant: erst lokalisierte DHR-Sektoren,
dann eine kontrollierte DR-artige Rekonstruktion von Feldalgebra und kompakter Eichgruppe aus der
Sektorstruktur. Dadurch bleiben \code{ChargeConjugationSeed}, \code{GaugeSectorSeed} und spätere
Statistik-/Teilchenrede auf die Observablen- und Netzschicht rückgebunden, statt als freie Seeds zu
schweben.\cite{DHR1971I}\cite{DoplicherRoberts1990}

\paragraph{9. Split Property nicht als Endpunkt missverstehen: Nuclearity als eigene Schicht führen.}
Die bereits geplante Split-Property bleibt wichtig, reicht aber als Fernziel nicht aus. Für kontrollierte
Phase-Space-Dichte, thermische Belastbarkeit und spätere Typ-III-Lesarten wird zusätzlich eine
Nuclearity-Schicht benötigt. Deshalb werden \path{CNNA/AQFT/NuclearitySeed.lean} und
\path{CNNA/Gates/Nuclearity.lean} explizit neben Split/Support geführt und nicht als später
optionaler Kommentar behandelt.\cite{BuchholzWichmann1986}

\subsection{Abnahmekriterien für B}
\doclevel{Pfad-Ebene}
B ist erst dann tragfähig, wenn:
\begin{itemize}
\item \code{ComplementLocalNet} und \code{InterfaceLocalNet} die ersten nichttrivialen Netzobjekte sind;
\item Isotonie, Lokalität und Zustandsrestriktion für diese Netze explizit formuliert sind;
\item die Grundschicht entweder bereits \code{GradedStarAlgebra} trägt oder mindestens explizit Z2-offen für spätere fermionische Sektoren angelegt ist;
\item \code{ComplementGNS} nicht isoliert bleibt, sondern durch \code{SeparatingProperty}, \code{TomitaTakesakiData} und \code{ModularConjugation} bis zur modularen Kernschicht ergänzt ist;
\item \code{ComplementKMS} auf dieser Zustands- und Modularity-Schicht aufsetzt statt sie zu ersetzen;
\item der Übergang vom endlichen Bright-Approximanten zu \code{InfiniteCarrier}, quasilokaler Vervollständigung oder thermodynamischer Limeslage explizit markiert ist, bevor unendliche AQFT-Lesarten beansprucht werden;
\item eine explizite Spektralkontrolle des aktiven Komplementnetzes via \code{ComplementSpectrumCondition} bzw. \code{gate\_spectrum\_condition\_on\_complement\_net} vorliegt, bevor Bisognano--Wichmann-, Guido--Longo- oder Haag--Ruelle-Schritte beansprucht werden;
\item \(\beta\) auf dem aktiven Pfad nicht als stiller freier Endparameter verbleibt, sondern als modular abgeleitete oder Fixpunkt-gesteuerte Größe markiert ist;
\item Split-Property nicht das Endziel bleibt, sondern eine explizite \code{NuclearitySeed}/\code{Nuclearity}-Schicht vorgesehen ist;
\item \code{DHRSectorsSeed} und \code{FieldAlgebraReconstructionSeed} als spät-B/früh-E-Schicht für Ladung, Statistik und Eichgruppenrekonstruktion vorhanden sind;
\item ein \code{BisognanoWichmannSeed} als explizite B->D-Brücke geplant ist;
\item der alte BoundaryMatrix-Pfad als Bright-Spezialfall replayt werden kann, aber nirgends mehr als stilles Zentrum auftritt.
\end{itemize}

\subsection{String-mathematische Ideen, die B vorbereiten muss}
\doclevel{Pfad-Ebene}
Snowmass diskutiert Functorial Field Theory, conformal nets, homotopy AQFT und Defekte als reale
QFT-Faces oder Erweiterungen.\cite[pp.~9--12]{Dedushenko2023}
Für CNNA bedeutet das in B ganz konkret: Defekte, Glueing, Boundary-State-Räume und
Dualitätsdictionaries werden noch \emph{nicht} de-seedet, aber die Datenstruktur muss so gelegt werden,
dass sie später in B/E andocken können. Hier sind Atiyahs Gluing-Perspektive, Segals CFT-Begriff und
die Defekt-/Boundary-Literatur aus der String-Mathematik konzeptionell einschlägig.\cite{Atiyah1988}\cite{Segal2004}\cite{FuchsRunkelSchweigert2003}\cite{Frohlich2007}


\newpage
\section{Pillar C -- offene Quantensysteme, Kanäle und Rückwirkung}
\doclevel{Pfad-Ebene}

\subsection{Zielkern von C}
\doclevel{Pfad-Ebene}
Pillar C ist in CNNA nicht mehr nur der Ort, an dem eine bereits feststehende System-Umgebung-
Interpretation mathematisch formalisiert wird. Der neue Kern ist: sektorielle Kanäle, Stinespring-artige
Dilationslogik, Nicht-Markov-Strukturen, relative Entropie- und Flussschichten und daraus gewonnene
Backreaction. Der Schritt von A/B nach C soll also nicht \emph{Umgebung wird eliminiert}, sondern
\emph{Interface-vermittelte Wechselwirkung wird kanalisch expliziert} bedeuten.

\subsection{Konkrete Vorgehensweise in C}
\doclevel{Pfad-Ebene}
\paragraph{1. Kanalgrundschicht von der Sektorgrammatik aus neu anschliessen.}
Die mathematischen Unterlagen aus \path{REALOQS/PillarC/OQS/Channel.lean} und
\path{Stinespring.lean} bleiben brauchbar. Inhaltlich müssen sie aber von einem binären SysEnv-Schema
auf \code{SectorChannels} und \code{SectorSysEnv} umgestellt werden. Bright, Dark und Interface müssen
als mögliche Quell-, Ziel- und Zwischenlagen der Kanäle explizit auftauchen.

\paragraph{2. Tail-Elimination als Kanal neu lesen.}
\path{TailElimAsChannel.lean} ist kein Nebenthema, sondern die Scharnierdatei zwischen A und C.
Sie muss zeigen, wie ein GeneralizedDtN-/MultiSchur-Schritt in eine effektive Kanalbeschreibung
übersetzt wird. Erst dadurch wird der dunkle Sektor nicht mehr bloss \emph{wegintegriert}, sondern als
strukturierte Quelle effektiver Dynamik ernst genommen.

\paragraph{3. Endliche, semigruppenartige und nichtmarkovsche Lagen staffeln.}
\path{Finite}, \path{Lindblad}, \path{Semigroup} und \path{NonMarkov} sollten nicht gleichzeitig als
parallel konkurrierende Ontologien auftreten. Die Reihenfolge ist konkreter: zürst der allgemeinste
sektorielle Kanaltrager, dann endliche/semigruppenartige Dynamik als kontrollierte Speziallage,
anschliessend Nicht-Markov-Strukturen für genuin erinnerungsbehaftete Interface-Dynamik.

\paragraph{4. Relative Entropie und Flüsse an das Interface binden.}
CNNA braucht für D und E keine abstrakte Entropierhetorik, sondern konkrete aus Kanälen und
Zustandsrestriktionen ableitbare Flussgrößen. Deshalb sollte C früh eine
\path{RelativeEntropyFlow}-Schicht erhalten, aus der später Backreaction, EffectiveLambda,
HorizonLevel und entanglement-equilibrium-nahe Diagnosen gespeist werden.

\paragraph{5. DerivedSpacetime in C nur als Vorstufe lesen.}
Das heute vorhandene \path{DerivedSpacetime.lean} in C zeigt bereits die Richtung: aus
Einflussrelationen kann eine diskrete Raumzeitstruktur konstruiert werden. In CNNA darf dieses Objekt
aber nicht Endpunkt von C bleiben. Es wird zum Vorbereitungsobjekt für D: aus Kanalinfluss wird
emergente Kausal- und Geometriestruktur. Genau darum ist das spätere Zielmodul
\path{CNNA/PillarD/DerivedSpacetime.lean} und nicht mehr \path{CNNA/PillarC/Integration/DerivedSpacetime.lean}.

\subsection{Abnahmekriterien für C}
\doclevel{Pfad-Ebene}
\begin{itemize}
\item sektorielle Kanäle sind definiert, bevor Lindblad-/Semigroup-Fenster de-seeden;
\item Tail-Elimination ist als Kanal übersetzbar;
\item relative Entropie- und Flussdaten sind an das Interface gebunden;
\item Backreaction erscheint als abgeleitete, nicht als frei postulierte Zusatzdynamik;
\item effektive Kanal- und Entwicklungslagen behalten auf dem aktiven Pfad eine explizit prüfbare Unitaritäts-/Positivitätskontrolle; semigruppenartige Spezialfenster bleiben de-seedete Sonderfälle.
\end{itemize}

\subsection{Was in C nicht passieren darf}
\doclevel{Pfad-Ebene}
\begin{itemize}
\item der dunkle Sektor wird wieder nur als eliminierte Restumgebung gelesen;
\item \code{DerivedSpacetime} wird als fertige Ontologie in C festgeschrieben;
\item semigruppenartige Speziallagen werden als allgemeine Theorie verkauft;
\item Entropie bleibt nur metaphorisch und wird nicht an konkrete Kanal- oder Zustandsdaten gebunden.
\end{itemize}


\newpage
\section{Pillar D -- emergente Geometrie, lokale Kovarianz und Gravity-Routen}
\doclevel{Pfad-Ebene}

\subsection{Zielkern von D}
\doclevel{Pfad-Ebene}
Pillar D ist der erste genuin neue CNNA-Pfeiler. Er darf nicht nur kinematische Etiketten auf ein schon
vorhandenes Netz kleben, sondern muss den Kreis aus Kausalstruktur, Mass-/Skalenwahl,
Materie-/Zustandssektor, Backreaction, Parameter-Closure und effektiver Einstein-/Friedmann-artiger
IR-Dynamik schliessen. Die interne Spezifikation ist hier zweigleisig: primäre Jacobson-/Entanglement-
Equilibrium-Route, sekundäre Route über induzierte Gravitation; asymptotic safety bleibt ein
kompatibler, aber methodisch nachgeordneter Explikationsstrang. Gerade die Closure-Frage für \(b\),
\(L_{\max}\), \(\beta\) und sekundäre Regularisierungsterme gehört deshalb nicht an den Rand,
sondern in den D-Eingang selbst.\cite[pp.~30--31, 34--35, 44--47]{CNNASpec}

\goalbox{\textbf{Kompositionsbild der Gravitation in CNNA:} Volle dynamische Raumzeit und volle
dynamische Gravitation entstehen in CNNA nicht aus einem Einzelmechanismus, sondern aus der
Komposition mehrerer abgeleiteter Stränge: Kausal-/Einflussstruktur, Mass-/Skalenfixierung,
Zustands-/Materiesektor, Entropie-/Backreaction-Daten, sektorielle Eliminationsprozesse und ---
falls erfolgreich geschlossen --- algebraisch-spektrogeometrische Zusatzstruktur. Jacobson,
induced gravity und NCG sind deshalb nicht als monokausale Konkurrenzhypothesen, sondern als
verschiedene Beitragsrouten zur selben emergenten Gravitationslesart zu führen.}

\subsection{Konkrete Vorgehensweise in D}
\doclevel{Pfad-Ebene}
\paragraph{1. Aus Einflussrelation eine echte D-Sprache machen.}
Nimm die in C vorbereiteten Einfluss- und Kanalrelationen und überführe sie in
\path{CNNA/PillarD/DerivedSpacetime.lean}, \path{InterfaceCausality.lean} und
\path{ComplementGeometry.lean}. Bereits vorhandene Kinematik-Dateien aus A ---
\path{DiscreteSpacetime}, \path{Causality}, \path{SpatialMetric}, \path{ProperTime*} --- dürfen dabei
nur als Seed-Material erscheinen. Ihre Begriffe werden erst de-seedet, wenn sie explizit aus
Bright/Dark/Interface- und Inflünce-Daten rückgebunden sind.

\paragraph{2. Lokal-kovariante Lesart vor geometrischer Spezialwahl herstellen.}
Snowmass behandelt LCQFT als eigene, ernstzunehmende Face-Familie; zugleich zeigen BFV und
spätere Arbeiten, dass lokale Kovarianz der natürliche Rahmen für QFT auf variablen Hintergründen
ist.\cite[pp.~6--9]{Dedushenko2023}\cite{BFV2003}\cite{BrunettiFredenhagenRejzner2022}
Daher sollte D früh eine \code{LocalCovariance}-Schicht erhalten, bevor Minkowski-, Conformal- oder
FLRW-Fenster überhaupt begonnen werden. Spezialfenster bleiben standardmäßig \code{Window}s.

\paragraph{2a. Relative Cauchy Evolution und dynamical locality explizit als D-Dynamikschicht führen.}
Lokale Kovarianz allein bleibt für CNNA zu statisch, sobald Hintergrundvariation, Stress-Tensor-Lesart
und spätere Geometrie-/Backreaction-Aussagen ernst gemeint sind. Deshalb werden
\path{CNNA/PillarD/RelativeCauchyEvolutionSeed.lean} und \path{CNNA/Gates/DynamicalLocality.lean}
ausdrücklich zwischen \code{LocalCovariance} und späterer renormierter Observablen-Schicht
verankert. So wird BFVs relative Cauchy evolution als operative Dynamik unter Hintergrundvariation und
Fewster--Verchs dynamical locality als Konsistenztest zwischen kinematischer und dynamischer
Lokalität in den Pfad eingebaut.\cite{BFV2003}\cite{FewsterVerch2012}

\paragraph{2b. Parameter-Closure als eigenen Zwischenstrang zwischen C und später D-Sprache führen.}
Die Spezifikation behandelt Parameter-Closure nicht als Fussnote, sondern als eigene Schliessungsschwelle
zwischen früher Dynamik und später Geometrie. Genau das muss auch die Roadmap sichtbar machen.
Der Strang besitzt drei Ebenen: (i) ontische Mikrostrukturparameter wie \(b\), die echte Verzweigung,
Sektornichttrivialität und Mikrostruktur selektieren; (ii) dynamische Emergenzgrößen wie
\(L_{\max}\), die nur als Vorname einer späteren \code{HorizonLevel}- bzw.
\code{BackreactionFixedPoint}-Lesart verbleiben dürfen; (iii) sekundäre numerische Regularisierer,
die entweder ganz eliminiert oder als abgeleitete Restgrößen in eine \code{RegularizationClosure}
überführt werden. Operativ heißt das: erst \code{BranchingWitness} mit einer sauberen
Nichttrivialitätsbedingung wie \code{nontrivialSideBranching\_iff\_one\_lt\_b}; dann ein
Selektorpfad \code{UVSpectralSelector -> BranchingSelector -> BackreactionSelector -> selectedBranching};
danach \code{HorizonLevel}, \code{EffectiveLambda} und \code{BackreactionFixedPoint} als effektive
IR-/Horizontschicht; schliesslich \code{RegularizationClosure} als Abschluss gegen freie numerische
Restparameter. Erst wenn diese Closure-Schwelle erreicht ist, dürfen asymptotic-safety-Sprache,
späte NCG-Interpretationen oder aggressive Window-Lesarten als mehr denn heuristische Diagnostik
auftreten. Auf derselben Schiene muss auch \(\beta\) aus dem Status eines primitiven Gibbs-Inputs
herausgezogen werden: entweder über den modularen/geometrischen Fluss oder als explizit markierte
Fixpunktgröße innerhalb der Closure-Schicht.\cite[pp.~44--47]{CNNASpec}

\paragraph{3. Jacobson-/Entanglement-Route explizit am Interface aufhängen.}
Jacobsons Idee, die Einstein-Gleichung als Zustandsgleichung aus Horizon-/Entropiebilanz zu lesen,
ist für CNNA besonders passend, weil das Interface bereits der natürliche Sitz für Entropie- und
Flussdiagnostik ist.\cite{Jacobson1995}
Konkret heißt das: aus \code{RelativeEntropyFlow}, \code{Backreaction} und lokalen Einflussmassen
werden \code{DerivedStressTensorSeed}, \code{EntanglementEquilibriumSeed},
\code{HorizonThermodynamicsSeed} und erst danach \code{EinsteinLimitSeed}. Dieser Strang ist der
primäre D-Kern.

\paragraph{3a. Hadamard-/mikrolokale Zustandskontrolle und lokale Wick-Produkte vor renormierter D-Sprache schliessen.}
Sobald D über reine Kinematik hinaus auf \code{DerivedStressTensorSeed},
\code{EntanglementEquilibriumSeed} oder renormierte lokale Observablen zielt, reicht eine blosse
GNS/KMS- oder Kausalitätsschicht nicht mehr aus. Deshalb werden
\path{CNNA/PillarD/HadamardStateSeed.lean}, \shortstack[l]{\path{CNNA/PillarD/}\\\path{MicrolocalSpectrumGate.lean}} und
\path{CNNA/PillarD/LocalWickPolynomialsSeed.lean} als echte D-Schicht eingeplant. Hadamard- bzw.
mikrolokale Spektralkontrolle stabilisiert die zulässige Zustandsklasse, während lokale kovariante
Wick- und Zeitordnungsprodukte spätere Stress-Tensor-, OPE- und Renormierungslesarten tragen.\cite{Radzikowski1996}\cite{BrunettiFredenhagenKoehler1996}\cite{HollandsWald2001}\cite{HollandsWald2002}

\paragraph{3b. Eine explizite Connes-/NCG-Route als spätes D und frühe D$\leftrightarrow$E-Brücke einfügen.}
Nichtkommutative/spektrogeometrische Ans"atze gehören in CNNA logisch nicht erst nach E, sondern
setzen bereits in spätem D an: dort, wo geometrische Daten, Dirac-/Spektralstruktur und die Frage
nach einer abgeleiteten Gravitationsdynamik selbst liegen. Deshalb werden
\path{CNNA/PillarD/SpectralTripleSeed.lean}, \path{DiracOperatorSeed.lean},
\path{RealStructureSeed.lean} und \path{SpectralActionSeed.lean} als eigener D-Strang geführt.
Er beginnt erst nach stabiler D-Kausalität, D-Skalenfixierung und B-Zustands-/Darstellungsschicht;
sein E-naher Ausläufer ist nicht der Ursprung der Route, sondern höchstens eine spätere interne/
finite Geometrie für Higgs-/Yukawa-/Flavor-Lesarten. Damit wird Connes/NCG als algebraisch-
spektrogeometrische Kompositionsroute von D nach E eingeordnet, nicht als verkapptes E-Ersatzfundament.\cite{ChamseddineConnes1996}\cite{ChamseddineConnes2007}

\paragraph{3c. Induzierte Gravitation explizit als materiesektorielle Effektivroute lesen.}
Um die Architektur nicht missverständlich zu machen, wird \code{InducedGravitySeed} sprachlich
nicht als Reserveideologie gegen Jacobson geführt, sondern als materiesektorielle Effektivroute,
die neben der thermodynamischen Jacobson-Linie und der algebraisch-geometrischen NCG-Linie einen
eigenen Beitrag zur emergenten Gravitationsdynamik liefert.

\paragraph{4. Asymptotic-Safety-Kompatibilität als späte Closure-Diagnostik behandeln.}
Die CNNA-Spezifikation verlangt ausdrücklich, dass asymptotic safety den derived-only Kern nicht
früh dominiert. Das ist auch physikalisch sinnvoll: Reuter-artige Theorie-Raum-/Fixpunktlogik,
Spectral-Dimension-Diagnostik und mode counting sind stark, aber sie müssen auf abgeleiteten
Größen sitzen.\cite[pp.~30--31]{CNNASpec}\cite{Reuter1998}\cite{ReuterSaueressig2012}\cite{LauscherReuter2005}
Also: zürst \code{HorizonLevel}, \code{EffectiveLambda}, \code{BackreactionFixedPoint},
\code{RegularizationClosure} und ein explizites \code{gate\_parameter\_closure\_complete}; danach erst
\code{SpectralDimensionFlowSeed}, \code{CutoffModeCountingSeed}, \code{RunningBoundaryDataSeed},
\code{TheorySpaceSeed}. Die AS-Sprache bleibt erklärend und testend, bis Closure erreicht ist; dasselbe
gilt für späte NCG-Deutungen und aggressive Windows, sofern sie primitive Parameter voraussetzen
würden.

\paragraph{5. Windows hart von Kernmodulen trennen.}
\code{FLRWSectorWindow}, \code{ConformalWindow} und \code{MinkowskiWindow} dürfen nur dann
auf den Kern rückwirken, wenn eine explizite Source-Map und ein de-seeding-Beschluss vorliegt. Das
entspricht der internen Spezifikation und vermeidet, dass späte Regime heuristisch den ontischen
Stamm überschreiben.\cite[pp.~7--10, 31]{CNNASpec}

\subsection{Abnahmekriterien für D}
\doclevel{Pfad-Ebene}
\begin{itemize}
\item \code{DerivedSpacetime}, \code{InterfaceCausality} und \code{LocalCovariance} sind aus C-Daten rückgebunden;
\item \code{RelativeCauchyEvolutionSeed} und \code{DynamicalLocality}-Gate explizieren die D-Dynamik unter Hintergrundvariation;
\item Hadamard-/mikrolokale Zustandskontrolle sowie \code{LocalWickPolynomialsSeed} stehen vor jeder renormierten Stress-Tensor- oder Einstein-Lesart;
\item der Jacobson-Strang besitzt explizite Seeds für Stress-Tensor, Entanglement-Gleichgewicht und Einstein-Limit;
\item ein expliziter Parameter-Closure-Strang schliesst \(b\), \(L_{\max}\), \(\beta\) und sekundäre Regularisierungsterme zu \code{selectedBranching}, \code{HorizonLevel}/\code{BackreactionFixedPoint} bzw. \code{RegularizationClosure};
\item die NCG-Route ist als spätes D und frühe D$\leftrightarrow$E-Brücke mit \code{SpectralTripleSeed}, \code{DiracOperatorSeed}, \code{RealStructureSeed} und \code{SpectralActionSeed} sichtbar;
\item induzierte Gravitation ist als getrennte späte materiesektorielle Effektivroute markiert;
\item volle Gravitation wird explizit als Kompositionsresultat mehrerer Stränge gelesen und nicht monokausal behauptet;
\item AS-Kompatibilität beginnt erst nach Closure-geeigneten Größen;
\item Spezialfenster bleiben Windows.
\end{itemize}


\subsection{Minimaler D-Dateibaum und explizite D/E-Gate-Reihenfolge}
\doclevel{Pfad-Ebene}
Analog zum E-Baum sollte auch D nicht nur begrifflich, sondern als minimaler Lean-Dateibaum sichtbar
sein. Gerade dadurch wird klar, wo D als eigener Pfeiler endet und wo die pfeilerübergreifende
Brücke nach E wirklich beginnt.

\tablecaptionline{Minimaler D-Dateibaum und explizite D/E-Gate-Reihenfolge -- Teil 1}
\begin{longtable}{L{0.34\linewidth}L{0.26\linewidth}L{0.28\linewidth}}
\toprule
Zieldatei / Gate & Rolle & Stop-Regel \\
\midrule
\endfirsthead
\toprule
Zieldatei / Gate & Rolle & Stop-Regel \\
\midrule
\endhead
\midrule
\multicolumn{3}{r}{Fortsetzung auf der nächsten Seite} \\
\midrule
\endfoot
\bottomrule
\endlastfoot
\path{CNNA/PillarD/DerivedSpacetime.lean} & erster expliziter Raumzeitträger aus Einfluss-, Closure- und Backreactiondaten & nach C-Flüssen und \code{gate\_parameter\_closure\_complete}, vor \code{InterfaceCausality} \\
\path{CNNA/PillarD/InterfaceCausality.lean} & kausale/regionale Struktur auf dem Interface & nach \code{DerivedSpacetime}, vor \code{LocalCovariance}, \code{PTBridgeSeed} und späterer BW-Deutung \\
\path{CNNA/PillarD/LocalCovariance.lean} & funktoriale Lokalkovarianz des emergenten D-Layers & nach \code{InterfaceCausality}, vor \code{RelativeCauchyEvolutionSeed} und renormierter D-Sprache \\
\path{CNNA/PillarD/RelativeCauchyEvolutionSeed.lean} & Dynamik unter lokaler Hintergrundvariation & nach \code{LocalCovariance}, vor \code{DynamicalLocality} \\
\path{CNNA/Gates/DynamicalLocality.lean} & Konsistenz von kinematischer und dynamischer Lokalität & nach \code{RelativeCauchyEvolutionSeed}, vor \code{HadamardStateSeed} und \code{PTBridgeSeed} \\
\path{CNNA/PillarD/HadamardStateSeed.lean} & mikrolokal zulässige Zustandsklasse für emergente Hintergründe & nach \code{DerivedSpacetime}, vor \code{MicrolocalSpectrumGate} und \code{LocalWickPolynomialsSeed} \\
\shortstack[l]{\path{CNNA/PillarD/}\\\path{MicrolocalSpectrumGate.lean}} & prüft Hadamard-/mikrolokale Spektralkontrolle & nach \code{HadamardStateSeed}, vor \code{LocalWickPolynomialsSeed} \\
\path{CNNA/PillarD/LocalWickPolynomialsSeed.lean} & lokale Wick- und Zeitordnungsprodukte & nach \code{MicrolocalSpectrumGate}, vor \code{DerivedStressTensorSeed} \\
\path{CNNA/PillarD/DerivedStressTensorSeed.lean} & renormierte lokale Stress-Tensor-Lesart des D-Layers & nach \code{LocalWickPolynomialsSeed}, vor \code{EntanglementEquilibriumSeed} und \code{InducedGravitySeed} \\
\path{CNNA/PillarD/EntanglementEquilibriumSeed.lean} & Jacobson-/Entanglement-Vorzone & nach \code{DerivedStressTensorSeed}, vor \code{EinsteinLimitSeed} \\
\shortstack[l]{\path{CNNA/PillarD/}\\\path{EinsteinLimitSeed.lean}} & \shortstack[l]{thermodynamische/\\entanglement-getriebene\\GR-Lesart} & \shortstack[l]{nach \code{EntanglementEquilibriumSeed}} \\
\shortstack[l]{\path{CNNA/PillarD/}\\\path{SpectralTripleSeed.lean}} & spektrogeometrischer Kern der NCG-Route & \shortstack[l]{erst nach \code{InterfaceCausality},\\\code{LocalCovariance} und\\geschlossener B-Darstellungslage} \\
\path{CNNA/PillarD/DiracOperatorSeed.lean} & Dirac-Struktur des spektrogeometrischen D-Strangs & nach \code{SpectralTripleSeed}, vor \code{RealStructureSeed} \\
\path{CNNA/PillarD/RealStructureSeed.lean} & reeller/KO-Layer der NCG-Brücke & nach \code{DiracOperatorSeed}, vor \code{SpectralActionSeed} \\
\path{CNNA/PillarD/SpectralActionSeed.lean} & algebraisch-spektrogeometrische D$\leftrightarrow$E-Brücke & nach \code{RealStructureSeed}, vor \code{FiniteInternalGeometrySeed} in E \\
\path{CNNA/PillarD/InducedGravitySeed.lean} & materiesektorielle Effektivroute für Gravitation & nach \code{DerivedStressTensorSeed} und stabiler Zustands-/Materieschicht \\
\end{longtable}

Die D/E-Brücke ist dabei explizit, aber nicht monolinear. Eine minimale Gate-Reihenfolge sollte
mindestens in folgenden Lagen sichtbar bleiben:
\begin{enumerate}
\item \textbf{D-Zulässigkeit:} \code{gate\_parameter\_closure\_complete} und \code{gate\_infinite\_carrier\_or\_quasilocal\_limit\_defined} müssen stehen, bevor D als kontinuierliche/aqft-artige Lesart überhaupt belastbar wird.
\item \textbf{Geometrisch-dynamischer D-Kern:} \code{DerivedSpacetime}, \code{InterfaceCausality}, \code{LocalCovariance}, \code{RelativeCauchyEvolutionSeed}, \code{DynamicalLocality}, \code{HadamardStateSeed}, \code{MicrolocalSpectrumGate} und \code{LocalWickPolynomialsSeed} bilden die erste volle D-Schale.
\item \textbf{Symmetriebrücke B$\to$D$\to$E:} \code{BisognanoWichmannSeed}, \code{TimeReversalOnGNS}, \code{AntiunitarySeed}, \code{SpinRepresentationSeed}, \code{FermionicStatisticsSeed}, \code{StatisticsGate}, \code{SpinStatisticsTheorem}, \code{KramersGate} und \code{KramersTheorem} dürfen erst auf dieser D-zulässigen Lage aufsetzen.
\item \textbf{Charge-/CPT-Brücke:} \code{DHRSectorsSeed}, \code{FieldAlgebraReconstructionSeed}, \code{ChargeStatisticsCoherence}, \code{ChargeReconstructionTheorem}, \code{PTBridgeSeed}, \code{ChargeConjugationSeed}, \code{CPTGate} und \code{CPTTheorem} schliessen die pfeilerübergreifende Ladungs- und CPT-Linie.
\item \textbf{Volle Materieorganisation:} Erst danach dürfen \code{GaugeSectorSeed}, \code{AnomalyInflowSeed}, \code{ChiralMatterSeed}, Higgs-/Flavor- und Matching-Fenster als späte E-Schicht einsetzen. Symmetrie- und Charge-Brücke können intern teilweise interleaven, aber Matter darf nicht vor beiden geschlossenen Brücken beginnen.
\end{enumerate}

\subsection{Theoretische Rahmung für D}
\doclevel{Pfad-Ebene}
Dieser Pfeiler darf sich auf drei externe Ideenkomplexe stützen, ohne von ihnen kolonisiert zu werden:
(1) lokale Kovarianz, relative Cauchy evolution und dynamical locality als QFT-Gesicht für variable Hintergründe,\cite{BFV2003}\cite{FewsterVerch2012}\cite{FewsterLang2016}
(2) Hadamard-/mikrolokale Zustandskontrolle sowie lokale Wick-/Zeitordnungsprodukte für renormierte lokale Observablen auf gekrümmten Hintergründen,\cite{Radzikowski1996}\cite{BrunettiFredenhagenKoehler1996}\cite{HollandsWald2001}\cite{HollandsWald2002}
(3) Jacobsons thermodynamische Route und spätere entanglement-orientierte GR-Lesarten,\cite{Jacobson1995}
(4) induced gravity als materiesektorielle Effektivroute,\cite{Visser2002}
(5) Connes/NCG und Spektralaktion als algebraisch-spektrogeometrische Zusatzroute mit Brückenwirkung nach E,\cite{ChamseddineConnes1996}\cite{ChamseddineConnes2007}
(6) asymptotic safety mitsamt Theorie-Raum-, Fixpunkt- und Spektraldimensionsdiagnostik.\cite{Reuter1998}\cite{ReuterSaueressig2012}\cite{LauscherReuter2005}
CNNA übernimmt dabei die Organisationsideen, nicht deren Ontologie als schon entschiedenes Fundament.


\newpage
\section{Pillar E -- später Gauge/Matter/SM-Strang aus der Stammtheorie}
\doclevel{Pfad-Ebene}

\subsection{Zielkern von E}
\doclevel{Pfad-Ebene}
Pillar E ist in der internen Spezifikation ausdrücklich \emph{spät}. Das Standardmodell darf nicht als
fertiger Block auf die Stammtheorie geklebt werden, sondern muss in logisch getrennte Teilprobleme
zerlegt und erst danach wieder zusammengesetzt werden: emergente Eichsektoren, chirale Materie,
Anomalie-/Ladungsstruktur, Higgs/EWSB, Flavor/Yukawa, QCD-/Hadronfenster und schliesslich
SM/SMEFT-Matching.\cite[pp.~33--36]{CNNASpec}
Kramers, CPT und Spin-Statistik sind dabei keine autonomen späten E-Zusätze, sondern Verbraucher
der zuvor geschlossenen B-Kern- und B$\to$D-Brückenstruktur; E darf sie nur noch nutzen, nicht
ersetzen oder nachträglich improvisieren. Für Ladung, Eichgruppe und Teilchen/Antiteilchen-Sprache
kommt zusätzlich eine observablenseitige DHR/DR-artige Schicht hinzu: \code{DHRSectorsSeed} und
\code{FieldAlgebraReconstructionSeed} müssen vor jeder belastbaren Gauge- und Charge-Lesart den
Übergang von lokalisierter Sektorstruktur zu Feldalgebra und kompakter Eichgruppe tragen.\cite{DHR1971I}\cite{DoplicherRoberts1990}
Genau das ist mit Snowmass und mit der mathematischen String-Literatur kompatibel: QFT hat viele
Faces, und String-Mathematik ist hier nicht als neues Fundament relevant, sondern als Sprache für
Defekte, Glueing, Dualitäten, Boundary States und Inflow-Mechanismen.\cite[pp.~9--12]{Dedushenko2023}\cite{Segal2004}\cite{FuchsRunkelSchweigert2003}\cite{Frohlich2007}
Nichtkommutative/spektrogeometrische Ansätze berühren E dabei zwar stark, aber nicht als deren
eigentlichen Ursprung: Der Kern der NCG-Route sitzt in spätem D, während E nur den E-nahen
Ausläufer einer möglichen internen/finitem Geometrie für Higgs-, Yukawa- und Flavor-Strukturen
aufnimmt. Damit wird \code{FiniteInternalGeometrySeed} im Erfolgsfall als organisierende
Fortsetzung, nicht als Ersatz von \code{GaugeSectorSeed}, \code{HiggsSectorSeed} oder
\code{YukawaFlavorSeed} gelesen.\cite[pp.~35--36]{CNNASpec}\cite{ChamseddineConnes1996}\cite{ChamseddineConnes2007}


\subsection{Feynman-Audit für Matter: was schon vor E als Gate relevant ist}
\doclevel{Pfad-Ebene}
Feynmans Dirac-Ged"ankengang ist für CNNA nicht erst in einem späten Teilchenfenster relevant,
sondern bereits als fr"uhes Konsistenzschema für die Architektur von B, D und dem Übergang nach E.
Der Kernpunkt seiner Vorlesung ist dreifach: (i) positive Energie und Relativit"at erzwingen
Antiteilchen- und Paarerzeugungsstrukturen; (ii) die Rückwärtslesart von Dynamik und die
Teilchen/Antiteilchen-Beziehung sind eng mit PT/C-Umdeutung verbunden; (iii) für Spin-$\tfrac12$
hängt die Spin-Statistik-Verknüpfung an einem echten Vorzeichenmechanismus, nicht bloss an der
Existenz irgendeiner involutiven Zeitumkehr.\cite[pp.~2--11, 24--27, 33--38]{FeynmanAntiparticles}
Für CNNA folgt daraus: Matter darf nicht als thermisches oder gruppentheoretisches Label beginnen,
sondern muss auf einer zuvor explizit geschlossenen Br"uckenschicht aus Spektrum, Zeitorientierung,
Zeitumkehr, Ladung/Antiladung, Repr"asentation und Statistik aufbauen. Das ist mit der internen
Spezifikation kompatibel, die E als späten derived Strang behandelt und freie Gauge/Matter-Daten auf
dem Kernpfad ausschliesst.\cite[pp.~33--36]{CNNASpec}

\warningbox{\textbf{Feynman-basierte Designregel für CNNA:} \emph{Time reversal alone is not matter.}
Ein anti-lineares Involutionsobjekt auf Observablenebene ist nützlich, ersetzt aber weder
Teilchen/Antiteilchen-Dictionaries noch Spinor-Repr"asentationen, CAR/Fock-Strukturen,
Austauschvorzeichen oder chirale Inflow-Gates.}

Im inspectierten v0.0605-Code sind bereits mehrere Vorbausteine vorhanden. Die Module
\path{REALOQS/PillarB/AQFT/TimeReversal.lean} und
\path{REALOQS/PillarB/AQFT/Derived/BoundaryMatrixTimeReversal.lean} definieren eine anti-lineare,
sternvertr"agliche und involutive Zeitumkehr auf Observablenebene; konkret wird auf der Matrixalgebra
\(A \mapsto \mathrm{transpose}(A^{\ast})\) und damit eintragweise komplexe Konjugation realisiert.
Zusätzlich unterscheidet \path{REALOQS/PillarB/AQFT/TimeOrientation.lean} bereits zwischen einer
\code{plus}- und einer \code{minus}-Orientierung der Dynamik, und die Intertwining-S"atze in der
abgeleiteten Matrixschicht tauschen diese beiden Richtungen sauber gegeneinander aus.\cite{REALOQSV0605}
Ebenso gibt es in \path{REALOQS/PillarB/AQFT/HaagKastler/SupportFull.lean} eine formale
Vorwärtskegel-Struktur \code{forwardCone/ForwardCone}, und
\path{REALOQS/PillarB/AQFT/Derived/TreeOfCliquesExtendedFromPillarA.lean} stellt darauf ein
\code{SpectrumCondition}-Zeugnis für den strikten ToC-Pfad bereit.\cite{REALOQSV0605}

Gerade hier ist aber wissenschaftliche Disziplin wichtig. \path{REALOQS/PillarA/Core/BInterfaces/GlobalCone.lean}
definiert \code{GlobalStateCone} als kompatible Familie lokaler Zustände unter Restriktion; das ist
mathematisch sinnvoll, aber nicht dasselbe wie Feynmans Lichtkegel- bzw. spacelike-Argument.
Dass beide Begriffe das Wort ``cone tragen, darf in der Roadmap nicht zu einer semantischen
Kurzschlusslesart führen. Die derived-only-Regel der CNNA-Spezifikation verbietet genau solche
vorgezogenen physikalischen Identifikationen.\cite[pp.~4--7]{CNNASpec}\cite{REALOQSV0605}

Der heutige Matter-Vorbau in \path{REALOQS/PillarB/AQFT/Derived/EnergyFromBoundaryOp.lean}
bleibt entsprechend bewusst schwach: \code{MatterTemplate} liefert eine endliche Konfigurations- und
Gibbs-Struktur über einem aus dem Randoperator abgeleiteten Energiefunktional, aber weder Ladung,
Antiteilchen, Spinor-Geometrie noch Austauschstatistik. Als thermischer Testtr"ager ist das nützlich;
als Materiesektor im Feynman-Sinn ist es unzureichend.\cite{REALOQSV0605}

\subsection{Feynman-Audit-Matrix für die Planung von D/E}
\doclevel{Pfad-Ebene}
Die folgende Matrix legt fest, wie Feynmans Argument in die CNNA-Planung eingetragen wird. Sie ist
bewusst kein physikalischer Endbeweis, sondern eine Gate-Liste gegen vorschnelle Bedeutungszuschreibungen.

\tablecaptionline{Feynman-Audit-Matrix für die Planung von D/E -- Teil 1}
\begin{longtable}{L{0.26\linewidth}L{0.26\linewidth}L{0.38\linewidth}}
\toprule
Feynman-relevante Forderung & Status in v0.0605 & Konsequenz für die Roadmap \\
\midrule
\endfirsthead
\toprule
Feynman-relevante Forderung & Status in v0.0605 & Konsequenz für die Roadmap \\
\midrule
\endhead
\midrule
\multicolumn{3}{r}{Fortsetzung auf der nächsten Seite} \\
\midrule
\endfoot
\bottomrule
\endlastfoot
Positive Energie / Vorwärtskegel & formales Vorwärtskegel- und Spektrums-Gate vorhanden, aber nur als sehr schlanke Witness-Schicht & in D als \code{InterfaceCausality}/\code{DerivedSpacetime} vertiefen; in E nicht mit Teilchenphysik verwechseln \\
Zeitumkehr und Zeitorientierung & anti-lineare involutive Zeitumkehr auf Observablenebene sowie \code{plus/minus}-Dynamik vorhanden & als B/D-Vorbau behalten; vor E um Repr"asentations- und Zustandswirkung erweitern \\
Teilchen/Antiteilchen-Dictionary & im aktuellen Archiv nicht formalisiert; kein eigener \code{ChargeConjugation}-Layer sichtbar & vor \code{GaugeSectorSeed} und \code{ChiralMatterSeed} explizite C-/PT-/CPT-Gates einführen \\
Spin-$\tfrac12$ / $T^2=-1$-artige Struktur & im aktuellen Archiv nicht vorhanden; heutige Zeitumkehr ist involutiv auf Observablenebene & Matter darf nicht aus der bestehenden \code{TimeReversalOn}-Schicht herausinterpretiert werden \\
Austauschvorzeichen / Bose-Fermi-Statistik & keine CAR/Fock-/Austauschschicht sichtbar & eigener Statistikstrang vor vollwertiger chiraler Materie; kein stilles Importieren fermionischer Regeln \\
Paarerzeugung / Inflow / spectator-Mechanismus & physikalisch motivierend durch Feynman, aber im Code noch nicht als Materieformalismus vorhanden & in E als expliziter Inflow-/Boundary-/Defekt-Strang einplanen, nicht erst als späte Kommentierung \\
\end{longtable}

Die operative Schlussfolgerung lautet daher: Feynman ist für CNNA schon \emph{vor} E nützlich,
aber im Modus eines Architektur-Audits. Er liefert keine Abkürzung, mit der aus den vorhandenen
Modulen bereits Matter hervorgezaubert werden darf. Stattdessen liefert er harte Vorbedingungen dafür,
welche Brücken zwischen B, D und E geschlossen sein müssen, bevor von einem echten Materiesektor
überhaupt gesprochen werden kann.\cite[pp.~2--11, 33--38]{FeynmanAntiparticles}\cite[pp.~33--36]{CNNASpec}

\subsection{Konkrete Vorgehensweise in E}
\doclevel{Pfad-Ebene}
\paragraph{1. Eichsektoren aus Defekten, Glueing und Interface-Struktur emergieren lassen.}
Beginne E nicht mit einer postulierten SM-Eichgruppe. Schreibe stattdessen
\path{GaugeSectorSeed.lean} als derived Kandidat, der aus Interface-Netz, Defekten,
Boundary-State-artigen Räumen und Glueing-Daten entsteht. Hier ist der string-mathematische
Input klar: Defekte und Phasengrenzen tragen echte physikalische Information und können
Symmetrien oder Dualitäten codieren.\cite{Frohlich2007}\cite{FuchsRunkelSchweigert2003}


\paragraph{1a. Vor jedem Gauge-Sektor eine echte PT/C/CPT-Brücke erzwingen.}
Der bestehende v0.0605-Code besitzt bereits eine brauchbare Zeitumkehr- und Spektrumsprache,
aber noch keine explizite Schicht für \code{Parity}, \code{ChargeConjugation} oder deren Wirkung
auf Repr"asentationen, Zustände und spätere Feld-/Teilchenobjekte. Weil Feynmans Antiteilchenlesart
an genau dieser Rückwärts- und Umdeutungsstruktur hängt, müssen
\code{PTBridgeSeed}, \code{ChargeConjugationSeed} und später ein klarer \code{CPTGate}
vor jedem belastbaren Matter- oder Chiralitätsanspruch eingeplant werden.\cite[pp.~10--11, 38--39]{FeynmanAntiparticles}\cite{REALOQSV0605}

\paragraph{2. Chiralität als harten Engpass explizit isolieren.}
Die Nielsen--Ninomiya-Obstruktion ist für jede diskrete/lattice-nahe Stammtheorie der zentrale Test.
CNNA darf daher \code{ChiralMatterSeed} erst dann de-seeden, wenn Doubling-/Mirror-Probleme
wirklich adressiert sind. Das spricht für einen expliziten Boundary-/Interface-/Inflow-Strang und
gegen naive direkte Lattice-Identifikation.\cite{NielsenNinomiya1981}\cite{CallanHarvey1985}


\paragraph{2a. Spin-Statistik nicht als späte Fussnote, sondern als eigenes Gate behandeln.}
Feynmans Argument zeigt, dass eine relativistische Materieschicht nicht schon durch die Existenz eines
Hamiltonoperators oder einer KMS-Struktur gesichert ist. Für Spin-$\tfrac12$ muss ein echter
Vorzeichenmechanismus der Austausch- bzw. Loop-Amplituden verfügbar sein, und für Kramers-artige
Zeitumkehr darf \(T^2\) nicht als freies Vorzeichen schweben, sondern muss auf einer graduierten
Struktur im Stil \(T^2 = (-1)^F\) bzw. ihrer halbzahligen Spin-Spezialisierung ruhen. Deshalb braucht
CNNA spätestens zwischen \code{ChiralMatterSeed} und \code{SMMatchingWindow} einen expliziten
\code{SpinRepresentationSeed}, einen \code{FermionicStatisticsSeed} sowie ein
\code{StatisticsGate}, das bosonische und fermionische Sektoren nicht bloss benennt, sondern formal
auseinanderhält. Ohne diesen Schritt bleibt Matter im besten Fall thermische Klassifikation, nicht
Teilchenmaterie.\cite[pp.~33--37]{FeynmanAntiparticles}\cite{Pauli1940}\cite{GuidoLongo1995}

\paragraph{3. Anomalien und Ladungsquantisierung als echte Gates implementieren.}
\path{AnomalyCancellationGate.lean} und \path{ChargeQuantizationGate.lean} müssen vor jeder
SM-kompatiblen Redeweise bestanden sein. Inflow- und String-Literatur zeigen, dass Anomalien keine
späte kosmetische Korrektur sind, sondern Konsistenzbedingungen des ganzen Sektors.\cite{CallanHarvey1985}\cite{AlvarezGaume2022}

\paragraph{4. Higgs/EWSB als eigenen Organisationssektor behandeln.}
Aus einer emergenten Eichgruppe folgt noch kein Massenspektrum. Deshalb müssen
\code{HiggsSectorSeed} und \code{EWSBSeed} als eigene derived Baustellen geführt werden. Hier sind
nichtkommutative/spektrale Geometrie und String-inspirierte Boundary-/Defektbilder hilfreich, nicht
weil CNNA sie direkt übernimmt, sondern weil sie zeigen, dass Symmetrie, Geometrie und Higgs-Daten
algebraisch gemeinsam organisiert werden können.\cite{ChamseddineConnes1996}\cite{ChamseddineConnes2007}

\paragraph{5. Flavor/Yukawa bewusst vom Kern trennen.}
\code{YukawaFlavorSeed} und \code{GenerationSeed} gehören in einen späten Parallelstrang. Die
Familien- und Mischungsproblematik darf den Kern nicht blockieren; sie darf aber auch nicht still als
freie Matrix oder freie Familienzahl eingeführt werden. Dieser Strang beginnt erst, wenn Gauge,
Chiralität und Anomaliegates stabil sind.

\paragraph{6. IR-Matching erst ganz am Schluss.}
CNNA sollte nicht zu früh auf volle beobachtete Teilchenphänomenologie zielen. Zunächst braucht
es ein kontrolliertes IR-Matching auf einen SM- oder SMEFT-artigen Sektor. Hier ist die moderne
SMEFT-Literatur der richtige Vergleichsrahmen: erst wenn der emergente IR-Sektor sauber definiert
ist, lohnt sich die Sprache von Matching und Wilson-Koeffizienten.\cite{PDG2023SM}\cite{IsidoriWilschWyler2024}


\subsection{Minimaler Zieldateibaum für E und seine Vorzone}
\doclevel{Pfad-Ebene}
Die Spezifikation beschreibt E nicht nur begrifflich, sondern bereits als Modulraum. Die Roadmap
übernimmt das jetzt explizit, damit E nicht hinter der Detailtiefe der Spezifikation zurückfällt.
Der eigentliche Pfeiler E beginnt spät, besitzt aber eine \emph{Vorzone} zwischen D und E, die
bereits in AQFT/Gates verankert ist:

\tablecaptionline{Minimaler Zieldateibaum für E und seine Vorzone -- Teil 1}
\begin{longtable}{L{0.34\linewidth}L{0.26\linewidth}L{0.28\linewidth}}
\toprule
Zieldatei / Gate & Rolle & Stop-Regel \\
\midrule
\endfirsthead
\toprule
Zieldatei / Gate & Rolle & Stop-Regel \\
\midrule
\endhead
\midrule
\multicolumn{3}{r}{Fortsetzung auf der nächsten Seite} \\
\midrule
\endfoot
\bottomrule
\endlastfoot
\path{CNNA/AQFT/GradedStarAlgebra.lean} & Z2-graduierte Sternalgebra als algebraische Unterlage für fermionische Sektoren & vor jeder belastbaren Statistik- oder Kramers-Sprache \\
\path{CNNA/AQFT/SeparatingProperty.lean} & Reeh--Schlieder-artige Separationsschicht für relevante Komplementalgebren & nach \code{ComplementGNS}, vor \code{TomitaTakesakiData} \\
\path{CNNA/AQFT/TomitaTakesakiData.lean} & extrahiert Tomita-Operator, modulare Konjugation \(J\) und modularen Operator \(\Delta\) & nach \code{SeparatingProperty}, vor \code{ComplementKMS} \\
\path{CNNA/AQFT/ModularConjugation.lean} & kontrolliert \(J M J = M'\) in Bright/Dark/Interface-Lesart & vor PT/CPT- und Kramers-Lesart \\
\path{CNNA/AQFT/ComplementSpectrumCondition.lean} & Vorwärtskegel-/Positivitäts- und Spektralkontrolle auf dem aktiven Komplementnetz & nach \code{ComplementKMS}, vor \code{BisognanoWichmannSeed} und \code{HaagRuelleSeed} \\
\path{CNNA/AQFT/BisognanoWichmannSeed.lean} & Brücke von modularer Automorphismengruppe zu abgeleiteten Boost-/Wedge-Strukturen & nach \code{ComplementKMS} und \code{ComplementSpectrumCondition}, vor \code{PTBridgeSeed}/\code{CPTGate}; Abgleich mit \code{InterfaceCausality} erst über späteres Gate \\
\path{CNNA/AQFT/NuclearitySeed.lean} & Phase-Space-/Dichteschicht jenseits blosser Split-Property, aber erst als Nach-KMS-Schicht & nach \code{ComplementKMS}, vor später Typ-III-/Charge-Lesart \\
\path{CNNA/Gates/Nuclearity.lean} & prüft die Nuclearity-/Phase-Space-Kontrolle des Netzes & nach \code{NuclearitySeed}, vor \code{DHRSectorsSeed} \\
\path{CNNA/AQFT/DHRSectorsSeed.lean} & lokalisierte Superselection-Sektoren des Observablennetzes & nach \code{Nuclearity}, vor \code{FieldAlgebraReconstructionSeed} und \code{GaugeSectorSeed} \\
\shortstack[l]{\path{CNNA/AQFT/}\\\path{FieldAlgebraReconstruction}\\\path{Seed.lean}} & \shortstack[l]{DR-artige Rekonstruktion von\\Feldalgebra und kompakter\\Eichgruppe} & \shortstack[l]{nach \code{DHRSectorsSeed},\\vor \code{ChargeConjugationSeed} und\\belastbarer Gauge-Sprache} \\
\shortstack[l]{\path{CNNA/Gates/}\\\path{ChargeStatisticsCoherence.lean}} & \shortstack[l]{prüft Kohärenz von Ladung,\\Statistik und rekonstruiertem\\Gauge-Layer} & \shortstack[l]{nach \code{FieldAlgebraReconstructionSeed},\\vor \code{GaugeSectorSeed}} \\
\path{CNNA/AQFT/ChargeReconstructionTheorem.lean} & expliziter DHR/DR-Abschluss für Feldalgebra-, kompakte Gauge- und Ladungsrekonstruktion & nach \code{ChargeStatisticsCoherence}, vor voller Gauge-/Matter-Lesart \\
\path{CNNA/PillarD/RelativeCauchyEvolutionSeed.lean} & dynamische Reaktion auf lokale Hintergrundvariation & nach \code{LocalCovariance}, vor \code{DynamicalLocality} und renormierter D-Sprache \\
\path{CNNA/Gates/DynamicalLocality.lean} & prüft die Konsistenz von kinematischer und dynamischer Lokalität & nach \code{RelativeCauchyEvolutionSeed}, vor späterer GR-/CPT-Brücke \\
\path{CNNA/PillarD/HadamardStateSeed.lean} & Hadamard-/mikrolokale Zustandsklasse auf emergenten Hintergründen & nach \code{DerivedSpacetime}, vor renormierten lokalen Observablen \\
\shortstack[l]{\path{CNNA/PillarD/}\\\path{MicrolocalSpectrumGate.lean}} & \shortstack[l]{prüft mikrolokale\\Spektralkontrolle/\\Hadamard-Zulässigkeit} & \shortstack[l]{nach \code{HadamardStateSeed},\\vor \code{LocalWickPolynomialsSeed}} \\
\shortstack[l]{\path{CNNA/PillarD/}\\\path{LocalWickPolynomialsSeed.lean}} & \shortstack[l]{lokale kovariante Wick- und\\Zeitordnungsprodukte} & \shortstack[l]{nach \code{MicrolocalSpectrumGate},\\vor \code{DerivedStressTensorSeed}} \\
\path{CNNA/PillarD/SpectralTripleSeed.lean} & spektrogeometrischer Kern für die NCG-Route in spätem D & erst nach stabiler \code{InterfaceCausality}, \code{LocalCovariance} und B-Zustands-/Darstellungsschicht \\
\path{CNNA/PillarD/DiracOperatorSeed.lean} & abgeleitete Dirac-Struktur für die spektrogeometrische D-Linie & nach \code{SpectralTripleSeed}, vor \code{RealStructureSeed} \\
\path{CNNA/PillarD/RealStructureSeed.lean} & realer Struktur-/KO-Layer der NCG-Brücke & nach \code{DiracOperatorSeed}, vor \code{SpectralActionSeed} \\
\path{CNNA/PillarD/SpectralActionSeed.lean} & algebraisch-spektrogeometrische Wirkungsroute als D$\leftrightarrow$E-Brücke & nach \code{RealStructureSeed}, vor möglicher interner/finiter Geometrie in E \\
\path{CNNA/AQFT/TimeReversalOnGNS.lean} & hebt observablenseitige Zeitumkehr auf den GNS-Raum & nach \code{TomitaTakesakiData}, vor jeder Kramers-, PT- oder Matter-Lesart \\
\path{CNNA/AQFT/AntiunitarySeed.lean} & formalisiert Antiunitarität der implementierten Zeitumkehr & direkt nach \code{TimeReversalOnGNS} \\
\path{CNNA/AQFT/SpinRepresentationSeed.lean} & halbzahlige Spin-/double-cover-Schicht & vor \code{FermionicStatisticsSeed} \\
\path{CNNA/AQFT/FermionicStatisticsSeed.lean} & trennt fermionische von bosonischen Sektoren & vor \code{StatisticsGate} und \code{ChiralMatterSeed} \\
\path{CNNA/Gates/StatisticsGate.lean} & prüft Spin-Statistik-Konsistenz auf graduiertem algebraischen Untergrund & nach \code{FermionicStatisticsSeed}, vor \code{ChiralMatterSeed} \\
\path{CNNA/AQFT/SpinStatisticsTheorem.lean} & explizites Endtheorem für die Spin-Statistik-Schicht der rekonstruierten CNNA-AQFT & nach \code{StatisticsGate}, vor voller chiraler Matter-Lesart \\
\path{CNNA/PillarE/PTBridgeSeed.lean} & koppelt aus D stammende Paritäts-/Raumspiegelungsdaten an die B-Zeitumkehrschicht & nach \code{LocalCovariance}/\code{InterfaceCausality}, vor \code{CPTGate} \\
\path{CNNA/PillarE/ChargeConjugationSeed.lean} & Teilchen/Antiteilchen- bzw. Ladungskonjugationslayer & vor \code{CPTGate} und jedem belastbaren Gauge-Sektor \\
\path{CNNA/Gates/CPTGate.lean} & pfeilerübergreifende B+D+E-Abnahme von PT/C/CPT auf Darstellungs- und Zustandsniveau & erst nach D-Kausalität/Geometrie und vor belastbarer Matter-Sprache \\
\path{CNNA/AQFT/CPTTheorem.lean} & explizites Endtheorem für die pfeilerübergreifende CPT-Schicht & nach \code{CPTGate}, vor voller Matter-Interpretation \\
\path{CNNA/Gates/KramersGate.lean} & trennt graduiertes \(\Theta^2=(-1)^F\), Zeitumkehrkovarianz, Spektrum und Entartung & direkt vor \path{KramersTheorem.lean} \\
\path{CNNA/AQFT/KramersTheorem.lean} & explizites Endtheorem für Kramers-Entartung im graduierten Zeitumkehrsektor & nach \code{KramersGate} \\
\path{CNNA/PillarE/GaugeSectorSeed.lean} & emergente Eichsektoren aus Interface/Defekt/Glueing & erst nach PT/C/CPT-Vorzone \\
\path{CNNA/PillarE/AnomalyInflowSeed.lean} & explizite Interface-/Boundary-Inflow-Schicht für chirale Rand- und Materiesektoren & nach \code{FieldAlgebraReconstructionSeed}, \code{ChargeStatisticsCoherence} und \code{GaugeSectorSeed}, vor \code{ChiralMatterSeed} und \code{AnomalyCancellationGate} \\
\path{CNNA/Gates/AnomalyInflow.lean} & prüft, dass die chirale Anomalielast wirklich durch Inflow getragen wird & nach \code{AnomalyInflowSeed}, vor \code{ChiralMatterSeed} \\
\path{CNNA/PillarE/ChiralMatterSeed.lean} & chirale Materie unter Doubling-/Inflow-Kontrolle & nur nach \code{StatisticsGate}, \code{FieldAlgebraReconstructionSeed}, \code{AnomalyInflowSeed} und \code{AnomalyInflow} \\
\path{CNNA/PillarE/AnomalyCancellationGate.lean} & Anomalieabnahme für belastbare Gauge-/Chiralitätsrede & vor \code{ChargeQuantizationGate} und \code{SMMatchingWindow} \\
\path{CNNA/PillarE/ChargeQuantizationGate.lean} & Ladungs- und Quantisierungsgate & vor Higgs-/Flavor- und SM-Matching \\
\path{CNNA/PillarE/HiggsSectorSeed.lean} & Higgs-/symmetry-breaking-Vorzone & nur nach stabiler Gauge- und Chiralitätsschicht \\
\path{CNNA/PillarE/EWSBSeed.lean} & elektroschwache Brechung als eigener abgeleiteter Strang & nicht vor \code{HiggsSectorSeed} \\
\path{CNNA/PillarE/FiniteInternalGeometrySeed.lean} & E-naher Ausläufer der NCG-Route für interne/finite Geometrie & nach \code{SpectralActionSeed}, vor belastbarer Higgs-/Yukawa-Organisation \\
\path{CNNA/PillarE/YukawaFlavorSeed.lean} & Flavor- und Yukawa-Strang & später Parallelstrang \\
\path{CNNA/PillarE/GenerationSeed.lean} & Generationen-/Mischungsstruktur & später Parallelstrang \\
\path{CNNA/PillarE/QCDWindow.lean} & spätes IR-Fenster für starke Wechselwirkung & bleibt \code{Window} \\
\path{CNNA/PillarE/HadronWindow.lean} & hadronisches Spezialfenster & bleibt \code{Window} \\
\path{CNNA/PillarE/SMMatchingWindow.lean} & IR-Matching auf SM-artigen Sektor & erst nach allen Konsistenzgates \\
\path{CNNA/PillarE/SMEFTWindow.lean} & spätes Effektivfeldtheorie-Fenster & bleibt \code{Window} \\
\end{longtable}

\subsection{Abnahmekriterien für E}
\doclevel{Pfad-Ebene}
\begin{itemize}
\item keine freie Wahl von Eichgruppe, Hyperladungen, Higgs-Dublett, Yukawa-Matrizen oder Familienzahl auf dem Kernpfad;
\item \code{GaugeSectorSeed}, \code{AnomalyInflowSeed}, \code{ChiralMatterSeed}, \code{PTBridgeSeed}, \code{ChargeConjugationSeed}, \code{SpinRepresentationSeed}, \code{FermionicStatisticsSeed}, \code{StatisticsGate}, \code{CPTGate}, \code{KramersGate}, \code{AnomalyCancellationGate} und \code{ChargeQuantizationGate} sind explizit vorhanden;
\item \code{NuclearitySeed}/\code{Nuclearity}, \code{DHRSectorsSeed}, \code{FieldAlgebraReconstructionSeed}, \code{ChargeStatisticsCoherence} und \code{ChargeReconstructionTheorem} schliessen den Weg von lokaler Observablenstruktur zu Ladung/Gauge rekonstruierbar;
\item die B-Kernschicht liefert bereits \code{GradedStarAlgebra}, \code{SeparatingProperty}, \code{TomitaTakesakiData}, \code{ModularConjugation} und einen \code{BisognanoWichmannSeed} als belastbare Vorzone;
\item \code{CPTGate} ist ausdrücklich als pfeilerübergreifendes Gate zwischen B, D und E benannt und wird nicht vor \code{InterfaceCausality}/\code{LocalCovariance} abgeschlossen; \path{CPTTheorem.lean} bleibt als eigenes Endtheorem sichtbar;
\item \code{KramersGate} setzt graduiertes Zeitumkehrverhalten, GNS-Implementierung und Spektrumsdaten voraus und bleibt daher eine späte Vorbedingung vor \path{KramersTheorem.lean};
\item \path{SpinStatisticsTheorem.lean}, \path{CPTTheorem.lean} und \path{KramersTheorem.lean} sind als symmetrisch benannte Zieltheoreme des Symmetrie-/Statistik-Strangs explizit ausgewiesen;
\item Inflow-/Boundary-Strategien für Chiralität sind Teil des Pfads und mit \code{AnomalyInflowSeed}/\code{AnomalyInflow} nicht mehr nur verbale Motivation;
\item \code{FiniteInternalGeometrySeed} bleibt ein E-naher Ausläufer der D-seitigen NCG-Route und ersetzt weder \code{GaugeSectorSeed} noch \code{HiggsSectorSeed} oder \code{YukawaFlavorSeed};
\item die Roadmap verwechselt \code{GlobalStateCone} nicht mit physikalischem Lichtkegel und behandelt \code{TimeReversalOn} nicht als bereits hinreichende Materieschicht;
\item \code{SMMatchingWindow} und \code{SMEFTWindow} sind echte \emph{Windows}, keine voreilige Kernontologie.
\end{itemize}

\subsection{Späte Horizontschicht für Teilchen-, Scattering- und IR-Physik}
\doclevel{Pfad-Ebene}
Zwischen lokaler Algebren-/Sektorstruktur und beobachtbarer QFT liegt eine eigene
Teilchen-, Scattering- und Infrared-Schicht. Sie ist kein unmittelbarer Kernblocker mehr, soll in
dieser Roadmap aber als physikalische Horizontergänzung explizit sichtbar bleiben. In der AQFT ist dies kein blosser Komfortaufsatz. Scattering ist der primäre
experimentelle Zugang zur relativistischen QFT, und für langreichweitige Kräfte wie QED reicht
die nackte DHR-Sprache nicht bis zur realistischen Ladungsphysik; hier braucht man eine
charge-class-/IR-verfeinerte Lesart. Deshalb führt die Roadmap diese Schicht bewusst als späte,
nicht-kernblockierende Horizontzone.
\cite{BuchholzScattering2005}\cite{BuchholzRoberts2014}\cite{Dybalski2005}

\tablecaptionline{Späte Horizontschicht für Teilchen-, Scattering- und IR-Physik -- Teil 1}
\begin{longtable}{L{0.34\linewidth}L{0.26\linewidth}L{0.28\linewidth}}
\toprule
Zieldatei / Gate & Rolle & Stop-Regel \\
\midrule
\endfirsthead
\toprule
Zieldatei / Gate & Rolle & Stop-Regel \\
\midrule
\endhead
\bottomrule
\endfoot
\path{CNNA/AQFT/ClusterDecompositionSeed.lean} & optionale diagnostische Vorzone für Vakuum-/Korrelationszerfall und spätere Teilcheninterpretation & nach \code{ComplementSpectrumCondition}, vor voller \code{HaagRuelleSeed}-Lesart \\
\path{CNNA/AQFT/HaagRuelleSeed.lean} & späte Scattering-Vorzone für asymptotische Mehrteilchen- und S-Matrix-Lesart & erst nach \code{ChargeReconstructionTheorem}, geeigneter Spektralkontrolle und bevorzugt expliziter Clusterdiagnostik \\
\path{CNNA/AQFT/ParticleInterpretationWindow.lean} & Fenster für asymptotische Teilchen-, Detektor- und Streuprozessinterpretation & nach \code{HaagRuelleSeed}; bleibt \code{Window} \\
\path{CNNA/AQFT/InfraredChargeClassSeed.lean} & IR-/long-range-charge-Schicht jenseits strikt lokalisierter DHR-Sektoren & spät, nach \code{ChargeReconstructionTheorem}; besonders relevant für QED-artige Langreichweite \\
\end{longtable}

\subsection{Was die String-Mathematik hier genau liefert}
\doclevel{Pfad-Ebene}
Die String-Mathematik soll in CNNA \emph{nicht} als neues Fundamentalpostulat auftreten. Ihre
brauchbaren Ideen sind viel präziser:
\begin{itemize}
\item Glueing/Funktorialität statt monolithischer globaler Beschreibung;\cite{Atiyah1988}\cite{Segal2004}
\item Defekte, Phasengrenzen und Boundary States als Trager nichttrivialer Daten;\cite{FuchsRunkelSchweigert2003}\cite{Frohlich2007}
\item Inflow-Mechanismen für chirale/anomale Randsektoren;\cite{CallanHarvey1985}\cite{AlvarezGaume2022}
\item algebraisch-geometrische Organisation von Gauge-/Higgs-Strukturen ohne simplem Import des Endresultats.\cite{ChamseddineConnes1996}\cite{ChamseddineConnes2007}
\end{itemize}
Genau in diesem präzisen Sinn gehört die String-Mathematik in CNNA hinein.


\newpage
\section{Pfeilerübergreifende partielle Ordnung statt Phasenplan}
\doclevel{Pfad-Ebene}

Es wird explizit \emph{kein} Phasenplan vorgegeben. Trotzdem braucht die Migration eine
partielle Ordnung. Diese Ordnung ist nicht zeitlich, sondern logisch.

\warningbox{\textbf{Status von Spezifikation und Roadmap.} Die ursprüngliche CNNA-Spezifikation v0.1
bleibt motivisches und begriffliches Leitdokument, ist für die konkrete Migration aus dem Archivstand
v0.0605 inzwischen aber in mehreren Punkten \emph{unterdeterminiert}. Für Priorisierung,
Stop-Regeln, Handoff-Reihenfolge, Notationspflicht und operative Initialisierung der Pfeiler ist daher die
vorliegende Roadmap das primäre bindende Zwischenartefakt zwischen Altarchiv und Zielsystem. Wo
Spezifikation und operative Migrationslogik auseinanderlaufen, steuert die partielle Ordnung dieser
Roadmap die Priorisierung.}

Die partielle Ordnung ist deshalb nicht nur eine statische Liste von Vorbedingungen. Sie wird aus dem
aktiven Codepfad und den dazugehörigen Tabellen fortlaufend nachgeführt. Insbesondere werden immer dem aktuellen Stand des Codes entsprechend die Phasenpläne je Abschnitt der partiellen Ordnung erstellt. Wenn neue belastbare
Codeeinsichten eine Notationsschicht, einen Handoff oder einen Gate-Übergang verschieben, dann ist
dies nicht nachträgliche Dokumentkosmetik, sondern eine notwendige Aktualisierung des executable
Dokuments und ggf. der zuvor ab selbst.




\begin{enumerate}
\item[\textbf{0.}] \textbf{Setup der Notationsschicht vor jeder Fachlogik:} Vor dem ersten fachlichen Schritt
in einem Pfeiler sind \path{CNNA/Notation.lean} und die jeweils benötigten
\path{CNNA/PillarX/Notation.lean}-Module anzulegen. Diese Initialisierung gehört zur Ordner- und
Importstruktur selbst; sie darf nicht nachträglich als kosmetische Dokumentationspflege folgen.
Lesbare \code{scoped\ notation}-Schichten müssen für alle tragenden Lean-Objekte des aktiven Pfads
vorbereitet sein, bevor die erste fachliche Logik des betreffenden Pfeilers implementiert wird.
\item \textbf{A vor allem anderen:} Ohne BranchPatch, ComplementSectorFamily, SectorSplit und GeneralizedDtN darf weder ein dunkles Netz in B noch ein echter Kanal in C beginnen.
\item \textbf{Kein full-derived lokales QFT-Face ohne Komplementnetz:} Wer Typ-III-, CPT-, Spin-Statistik- oder Kramers-Aussagen auf dem aktiven Pfad will, muss die komplementbasierte Netz-, Zustands- und Modularschicht zürst schliessen; Bright-only-Recovery bleibt Regression, nicht Kernableitung.
\item \textbf{B-Algebra vor fermionischer Semantik:} Die AQFT-Grundschicht von B muss spätestens mit \code{GradedStarAlgebra} Z2-offen werden, bevor Spin-, Statistik- oder Kramers-Sprache belastbar ist.
\item \textbf{GradedStarAlgebra und SeparatingProperty nicht verwechseln:} Die Graduierung der Algebrenbasis und die Reeh--Schlieder-artige Separationslage des GNS-Vektors gehören zu verschiedenen logischen Achsen; weder darf die eine als importierte Folge der anderen behandelt werden noch umgekehrt.
\item \textbf{B-Netze vor B-Modularität/Thermik:} ComplementLocalNet/InterfaceLocalNet und Zustandsrestriktionen müssen vor \code{ComplementGNS}; \code{ComplementGNS} wiederum vor \code{SeparatingProperty}, \code{TomitaTakesakiData}, \code{ModularConjugation} und erst danach vor \code{ComplementKMS} stehen.
\item \textbf{Komplement-Spektrum vor BW/Scattering/Guido--Longo:} Eine aktive Spektralkontrolle auf dem Komplementnetz via \code{ComplementSpectrumCondition} bzw. \code{gate\_spectrum\_condition\_on\_complement\_net} muss vor \code{BisognanoWichmannSeed}, \code{HaagRuelleSeed} und jeder belastbaren Spin-Statistik-Lesart geschlossen sein.
\item \textbf{Endlicher Bright-Patch ist nicht schon die unendliche AQFT-Lage:} \code{InfiniteCarrier}, quasilokale Vervollständigung oder eine äquivalente thermodynamische Limesmarkierung müssen explizit vorliegen, bevor unendliche Kommutanten-, Typ-III-, BW- oder Scattering-Lesarten belastbar beansprucht werden.
\item \textbf{B/C vor D:} D darf erst beginnen, wenn Einfluss-, Fluss- und Backreactiondaten in C sowie die relevanten Netz-, Zustands- und modularen Daten in B vorhanden sind.
\item \textbf{Volle Gravitation nur als Kompositionsresultat lesen:} Weder Jacobson noch induced gravity noch NCG dürfen als isolierte Monokausalroute missverstanden werden; die belastbare D-Lesart entsteht erst aus dem Zusammenspiel von Kausalstruktur, Skalenfixierung, Materie-/Zustandssektor, Entropie/Backreaction und ggf. spektrogeometrischer Zusatzstruktur.
\item \textbf{Parameter-Closure vor AS-, NCG- und später Fenster-Sprache:} \code{BranchingWitness}/\code{selectedBranching}, \code{HorizonLevel}, \code{EffectiveLambda}, \code{BackreactionFixedPoint} und \code{RegularizationClosure} müssen vor \code{SpectralDimensionFlow}, mode counting, \code{TheorySpaceSeed}, späten \code{SpectralActionSeed}-Deutungen und heuristischen Spezial-\code{Window}s als primitive Steuerdaten geschlossen sein.
\item \textbf{NCG erst nach stabiler D-Kausalität und vor später E-Organisation:} \code{SpectralTripleSeed}, \code{DiracOperatorSeed}, \code{RealStructureSeed} und \code{SpectralActionSeed} dürfen erst nach \code{InterfaceCausality}, \code{LocalCovariance} sowie stabiler B-Zustands-/Darstellungsschicht einsetzen, aber vor später Higgs-/Yukawa-/Flavor-Organisation in E.
\item \textbf{D-Kausalität/Geometrie vor pfeilerübergreifender CPT-Brücke:} \code{PTBridgeSeed}, \code{ChargeConjugationSeed} und \code{CPTGate} dürfen erst nach \code{InterfaceCausality}, \code{LocalCovariance} und dem \code{BisognanoWichmannSeed} beansprucht werden.
\item \textbf{Keine belastbare Charge-/Gauge-Sprache vor DHR/DR-Rekonstruktion:} \code{GaugeSectorSeed}, \code{ChargeConjugationSeed} und spätere Teilchen-/Antiteilchen- und Eichgruppenrede dürfen erst nach \code{DHRSectorsSeed}, \code{FieldAlgebraReconstructionSeed} und \code{ChargeStatisticsCoherence} auftreten.
\item \textbf{Anomaly-Inflow erst nach rekonstruierter Ladungs-/Feldschicht:} \code{AnomalyInflowSeed} und \code{AnomalyInflow} dürfen nicht bloss nach einer verbalen Gauge-Motivation erscheinen, sondern erst nach \code{FieldAlgebraReconstructionSeed}, \code{ChargeStatisticsCoherence} und belastbarer \code{GaugeSectorSeed}-Lesart; sie bleiben zugleich Vorbedingung von \code{ChiralMatterSeed}.
\item \textbf{Keine renormierte D-Sprache vor Relative-Cauchy-/Hadamard-/Wick-Schicht:} \code{DerivedStressTensorSeed}, \code{EntanglementEquilibriumSeed}, \code{EinsteinLimitSeed} und spätere lokale Observablen verlangen vorab \code{RelativeCauchyEvolutionSeed}, \code{DynamicalLocality}, \code{HadamardStateSeed}, \code{MicrolocalSpectrumGate} und \code{LocalWickPolynomialsSeed}.
\item \textbf{Symmetrie- und Statistikbrücke nach D, aber vor voller Materie:} Zwischen D und E müssen \code{TimeReversalOnGNS}, \code{AntiunitarySeed}, \code{SpinRepresentationSeed}, \code{FermionicStatisticsSeed}, \code{StatisticsGate}, \code{CPTGate} und \code{KramersGate} explizit geschlossen werden; \code{TimeReversalOn} allein reicht dafür nicht.
\item \textbf{Prototypisches Vorauseilen nur als Seed/Scaffold:} Frühe Dummy-Schnittstellen für spätere Pfeiler sind erlaubt und erwünscht, dürfen aber niemals als bereits geschlossener Kern verkauft werden; sie prüfen Tragfähigkeit, nicht Vollständigkeit.
\item \textbf{Rückfluss von Erkenntnissen ist zulässig, aber gate-diszipliniert:} Wenn spätere Pfeiler frühere Typ- oder Exportverträge sprengen, müssen die betroffenen Module kontrolliert wieder geöffnet, reklassifiziert und neu geschlossen werden, statt die Inkonsistenz in spätere Schichten zu verschieben.
\item \textbf{Recovery zuletzt:} Bright-Recovery und Replay alter Beispiele sind Regression, nicht Entwurf des Kerns.
\end{enumerate}

\subsection{Gate-Reife statt Parallelillusion}
\doclevel{Pfad-Ebene}
Die partielle Ordnung ist zugleich eine Schutzregel gegen Scheinparallelisierung. Solange ein Modul
oder Exportvertrag nur \code{Seed}-Status hat, bleiben spätere Detailpläne bewusst weich. Erst ein
explizit geschlossenes Gate verhärtet die nächste Planungsschicht. Umgekehrt darf späterer
Architekturzwang frühere Gates wieder öffnen: wenn C oder D zeigt, dass ein A-seitiger Typvertrag
zu eng war, wird nicht stillschweigend weitergebaut, sondern kontrolliert auf
\code{Seed}/\code{Scaffold} zurückgestuft, nachgebessert und neu geschlossen. Genau diese
Rückkopplung ist kein Fehler der Roadmap, sondern ihr Flexibilitätsmechanismus. Jeder aus dieser
partiellen Ordnung später abgeleitete Meilenstein trägt daher ein eigenes \emph{Notation Review} als
Definition of Done: ein Gate gilt erst dann als belastbar, wenn die zugehörigen Kernobjekte auch über
ihre zugeordnete Notationsschicht lesbar, registriert und fachlich reviewbar sind.

\subsection{Operativer Unterpfad 6a--6f für den modularen B-Kern}\label{sec:b-unterpfad}
\doclevel{Pfad-Ebene}
Obwohl diese Roadmap keinen globalen Phasenplan verwendet, braucht der kritische Pfad im B-Kern
einen expliziten Unterpfad. Die frühe algebraische Bedingung ist bereits in \code{GradedStarAlgebra}
verankert; der eigentliche modulare Unterpfad lautet dann:

\tablecaptionline{Operativer Unterpfad 6a--6f für den modularen B-Kern -- Teil 1}
\begin{longtable}{L{0.09\linewidth}L{0.24\linewidth}L{0.57\linewidth}}
\toprule
Schritt & Kernobjekt & Funktion im kritischen Pfad \\
\midrule
\endfirsthead
\toprule
Schritt & Kernobjekt & Funktion im kritischen Pfad \\
\midrule
\endhead
\bottomrule
\endfoot
6a & \code{ComplementGNS} & erste komplementnetz-basierte Darstellungslage; noch nicht modular geschlossen \\
6b & \code{SeparatingProperty} & stärkt den GNS-Vektor für relevante Komplementalgebren von zyklisch zu zyklisch+separierend \\
6c & \code{TomitaTakesakiData} + \code{ModularConjugation} & extrahiert \(S\), \(J\), \(\Delta\) und schliesst die modulare Darstellungsarchitektur \(J M J = M'\) \\
6d & \code{ComplementKMS} & thermische Schicht auf bereits geschlossener modularer Darstellungsarchitektur \\
6e & \code{ComplementSpectrumCondition} + \code{gate\_spectrum\_condition\_on\_complement\_net} & schliesst Spektral- und Vorwärtskegelkontrolle auf dem Komplementnetz vor jeder BW-, Guido--Longo- oder Haag--Ruelle-Lesart \\
6f & \code{BisognanoWichmannSeed} & explizite B$\to$D-Brücke von modularen Automorphismen zu abgeleiteten Boost-/Wedge-Strukturen \\
\end{longtable}

Dieser Unterpfad ist nicht kosmetisch. Er verhindert genau den falschen Kurzschluss
\code{ComplementGNS -> ComplementKMS -> spaet irgendwo CPT/Kramers}, der die modulare Mittelschicht
unsichtbar machen würde.

\subsection{Spät-B-/Früh-E-Ergänzung für Nuclearity, DHR/DR und physikalische Rekonstruktion}
\doclevel{Pfad-Ebene}
Auf stabiler modularer und HK-/Quasilokal-Grundlage schliesst der Pfad noch nicht automatisch zu einer
physikalisch reicheren AQFT. Deshalb werden unmittelbar oberhalb des heutigen B-Kerns drei weitere
Stufen explizit geführt:
\tablecaptionline{Spät-B-/Früh-E-Ergänzung für Nuclearity, DHR/DR und physikalische Rekonstruktion -- Teil 1}
\begin{longtable}{L{0.09\linewidth}L{0.24\linewidth}L{0.57\linewidth}}
\toprule
Schritt & Kernobjekt & Funktion im erweiterten Pfad \\
\midrule
\endfirsthead
\toprule
Schritt & Kernobjekt & Funktion im erweiterten Pfad \\
\midrule
\endhead
\bottomrule
\endfoot
6g & \code{NuclearitySeed} + \code{Nuclearity} & führt Phase-Space- und Freiheitsgradkontrolle jenseits blossen Split-Properties ein; bleibt aber eine Nach-KMS-Schicht \\
6h & \code{DHRSectorsSeed} & lokalisiert Superselection-Sektoren des Observablennetzes als Charge-/Statistikvorzone \\
6i & \code{FieldAlgebraReconstructionSeed} + \code{ChargeStatisticsCoherence} & rekonstruiert Feldalgebra/kompakte Eichgruppe und bindet Ladung an Statistik \\
\end{longtable}
Diese Zusatzstufen sind nicht späte Dekoration, sondern die sauberste Stelle, an der CNNA den Schritt
von netz- und modularitätsbasierter Architektur zu rekonstruierbarer Charge-/Gauge-Sprache vollzieht.

\subsection{Operativer Unterpfad PC-a--PC-d für Parameter-Closure zwischen C und D}
\doclevel{Pfad-Ebene}
Die Spezifikation behandelt Parameter-Closure als eigene Schliessungsschwelle des derived-only Pfads.
Diese Roadmap führt daher zusätzlich zum modularen Unterpfad einen eigenen Closure-Unterpfad:

\tablecaptionline{Operativer Unterpfad PC-a--PC-d für Parameter-Closure zwischen C und D -- Teil 1}
\begin{longtable}{L{0.10\linewidth}L{0.26\linewidth}L{0.54\linewidth}}
\toprule
Schritt & Kernobjekt & Funktion im Closure-Pfad \\
\midrule
\endfirsthead
\toprule
Schritt & Kernobjekt & Funktion im Closure-Pfad \\
\midrule
\endhead
\bottomrule
\endfoot
PC-a & \code{BranchingWitness} + \code{nontrivialSideBranching\_iff\_one\_lt\_b} & macht echte Verzweigung und Nichttrivialität des dunklen Sektors explizit; \(b>1\) ist hier Marker, nicht Endparameter \\
PC-b & \code{UVSpectralSelector} + \code{BranchingSelector} + \code{BackreactionSelector} & führt spektrale, verzweigungs- und rückwirkungsbezogene Daten zu einer abgeleiteten \code{selectedBranching}-Lesart zusammen \\
PC-c & \code{HorizonLevel} + \code{EffectiveLambda} + \code{BackreactionFixedPoint} & ersetzt ein primitives \(L_{\max}\) durch eine effektive IR-/Horizontgröße und bindet sie an Zustand, Entropie und Interface-Balance \\
PC-d & \code{RegularizationClosure} + \code{gate\_parameter\_closure\_complete} & eliminiert sekundäre numerische Regularisierer oder degradiert sie zu abgeleiteten Restgrößen; markiert die Closure-Schwelle vor AS-/NCG-/Window-Sprache \\
\end{longtable}

Auf demselben Strang muss \(\beta\) den Status eines primitiven Gibbs-Inputs verlieren: entweder als
modular/geometrisch normalisierte Größe oder als explizit markierte Fixpunktvariable innerhalb von
\code{BackreactionFixedPoint}. Erst danach ist der Pfad gegen freie Restparameter weit genug
abgeschlossen, um genuinely derived-only genannt zu werden.\cite[pp.~44--47]{CNNASpec}


\subsection{Drei harte Zusatzregeln für Abhängigkeiten und Gates}
\doclevel{Pfad-Ebene}
\begin{enumerate}
\item \textbf{Keine fermionische Semantik vor \code{GradedStarAlgebra}.} Statistik-, Spin- und Kramers-Sprache dürfen nicht auf einer rein ungraduierten AQFT-Basis freischwebend eingeführt werden.
\item \textbf{Kein full-derived QFT-Face ohne Komplementnetz.} Wer vN-/Typ-III-, CPT-, Spin-Statistik- oder Kramers-Lesarten auf dem aktiven Pfad beansprucht, muss zuvor die komplementbasierte Kette von \code{ComplementLocalNet}/\code{InterfaceLocalNet} über \code{ComplementStateNet} bis in die modulare Darstellungsarchitektur geschlossen haben; Bright-only-Recovery zählt dafür nicht.
\item \textbf{Keine Kramers-/Spin-Statistik-/CPT-Sprache vor der modularen B-Kette.} Vor jeder belastbaren PT/C/CPT-, Spin-Statistik- oder Kramers-Lesart muss die Kette \code{ComplementGNS -> SeparatingProperty -> TomitaTakesakiData -> ModularConjugation -> ComplementKMS} geschlossen sein.
\item \textbf{Kein derived-only \(\beta\) als stiller Alt-Export.} Solange \(\beta\) nur als gegebener Gibbs-Parameter aus dem Bright-Pfad mitläuft, bleibt \code{ComplementKMS} eine Scaffold-Lesart; der aktive Pfad braucht entweder modulare/geometrische Ableitung oder explizite Fixpunktmarkierung.
\item \textbf{Kein \code{CPTGate} vor D-Kausalität/Geometrie plus \code{BisognanoWichmannSeed}.} Das \code{CPTGate} ist das erste ausdrücklich pfeilerübergreifende Gate zwischen B, D und E und darf logisch nicht vor \code{InterfaceCausality}/\code{LocalCovariance} abgeschlossen werden.
\end{enumerate}

\subsection{Neue Gates für den modularen B-Kern und die B->D-Brücke}
\doclevel{Pfad-Ebene}
Zusätzlich zu \code{gate\_complement\_GNS} und \code{gate\_complement\_KMS} sind auf dem aktiven Pfad mindestens die folgenden Gates explizit zu führen:

\tablecaptionline{Neue Gates für den modularen B-Kern und die B->D-Brücke -- Teil 1}
\begin{longtable}{L{0.36\linewidth}L{0.54\linewidth}}
\toprule
Gate & Rolle \\
\midrule
\endfirsthead
\toprule
Gate & Rolle \\
\midrule
\endhead
\bottomrule
\endfoot
\code{gate\_graded\_statistics\_defined} & prüft, dass fermionische/bosonische Sektorsprache auf einer graduierten AQFT-Basis sitzt \\
\code{gate\_complement\_separating} & stoppt den Pfad, falls die Separationsstufe zwischen \code{ComplementGNS} und Modularität fehlt \\
\code{gate\_tomita\_takesaki\_defined} & sichert, dass \(S\), \(J\) und \(\Delta\) als CNNA-Objekte vorliegen \\
\shortstack[l]{\code{gate\_modular\_conjugation\_maps\_}\\code{bright\_to\_complement}} & operationalisiert die Bright/Dark/Interface-Lesart von \(J M J = M'\) \\
\shortstack[l]{\code{gate\_spectrum\_condition\_on\_}\\code{complement\_net}} & stoppt den Pfad, solange Vorwärtskegel- und Positivitätskontrolle nur aus Bright-Recovery stammt \\
\shortstack[l]{\code{gate\_modular\_boost\_}\\\code{consistent\_with\_}\\\code{interface\_causality}} & prüft die Bisognano--Wichmann-Brücke vor jeder belastbaren CPT-Lesart \\
\end{longtable}

\subsection{Weitere Gates für physikalische Vollständigkeit jenseits des heutigen Kerns}
\doclevel{Pfad-Ebene}
Zusatzlich zur modularen B-Kette müssen für eine physikalisch breiter belastbare Rekonstruktionsroute
spätestens die folgenden Gates im Dokument benannt bleiben:

\tablecaptionline{Weitere Gates für physikalische Vollständigkeit jenseits des heutigen Kerns -- Teil 1}
\begin{longtable}{L{0.36\linewidth}L{0.54\linewidth}}
\toprule
Gate & Rolle \\
\midrule
\endfirsthead
\toprule
Gate & Rolle \\
\midrule
\endhead
\bottomrule
\endfoot
\code{gate\_nuclearity\_defined} & prüft die Phase-Space-/Energielevel-Dichte jenseits blossen Split-Properties \\
\code{gate\_dhr\_sectors\_defined} & sichert lokalisierte Superselection-Sektoren des Observablennetzes \\
\shortstack[l]{\code{gate\_field\_algebra\_}\\\code{reconstruction\_}\\\code{defined}} & stoppt den Pfad, falls Feldalgebra- und kompakte Eichgruppenrekonstruktion nur heuristisch bleiben \\
\shortstack[l]{\code{gate\_charge\_statistics\_}\\code{coherence}} & prüft die Kohärenz von Ladung, Statistik und rekonstruiertem Gauge-Layer \\
\shortstack[l]{\code{gate\_relative\_cauchy\_evolution\_}\\code{defined}} & operationalisiert die lokale Dynamik unter Hintergrundvariation \\
\code{gate\_dynamical\_locality} & vergleicht kinematische und dynamische Lokalität im LCQFT-Strang \\
\code{gate\_microlocal\_spectrum\_defined} & sichert Hadamard-/mikrolokale Zulässigkeit für D-Zustände \\
\shortstack[l]{\code{gate\_local\_wick\_polynomials\_}\\code{defined}} & verlangt lokale kovariante Wick- und Zeitordnungsprodukte vor renormierter D-Sprache \\
\shortstack[l]{\code{gate\_parameter\_closure\_}\\code{complete}} & markiert die derived-only-Schwelle gegen freie Restparameter in \(b\), \(L_{\max}\), \(\beta\) und Regularisierern \\
\shortstack[l]{\code{gate\_horizon\_level\_replaces\_}\\code{primitive\_Lmax}} & stoppt den Pfad, solange \(L_{\max}\) noch als ontischer Aktivparameter statt als abgeleitete Horizontgröße auftritt \\
\shortstack[l]{\code{gate\_infinite\_carrier\_or\_}\\code{quasilocal\_limit\_defined}} & sichert den Übergang vom endlichen Bright-Approximanten zur unendlichen/quasilokalen AQFT-Lesart \\
\shortstack[l]{\code{gate\_cluster\_decomposition\_}\\code{checked}} & optionale Spätdiagnostik für Vakuum-/Korrelationszerfall vor voller Teilchen- und Scattering-Interpretation \\
\code{gate\_prototype\_interfaces\_type\_stable} & prüft, dass frühe B-/C-Verbraucherschnittstellen auf A-Basis ohne ad-hoc Typreparaturen compiliert werden \\
\code{gate\_back\_propagation\_resolved} & markiert, dass aus späteren Pfeilern zurücklaufende Architekturzwänge sauber in frühere Source-Maps und Exporte eingearbeitet wurden \\
\end{longtable}

\goalbox{\textbf{Kurzform der Abhängigkeitsordnung:}
A-Stammtheorie $\prec$ B-Algebra/Netze + InfiniteCarrier/quasilokaler Limes $\prec$ B-Modularität/GNS/KMS + Komplement-Spektralkontrolle + Nuclearity + C-Kanäle/Backreaction $\prec$ Parameter-Closure/HorizonLevel/RegularizationClosure $\prec$ D-Geometrie/Relative-Cauchy/Hadamard/Wick + späte NCG-Brücke $\prec$ DHR/DR-Rekonstruktion + pfeilerübergreifende CPT/Spin/Statistik-Brücke $\prec$ E-Gauge/Matter/SM-Fenster,
mit einem separaten späten Recovery-Strang für den alten Bright-Pfad.}

\newpage
\section{Datei-/Kettenblätter für tragende Pfade}
\doclevel{Datei-/Ketten-Ebene}
\label{sec:chaindocs}
Die folgenden Kettenblätter heben die wichtigsten tragenden Pfade der Roadmap in das regelwerksnahe
Schema 1--6: Ziel, Definitionen, Voraussetzungen, Behauptung, Beweis der Kette, Zielabgleich.
Sie ersetzen keine spätere Voll-Dokumentation atomarer Einzeltheoreme, machen aber bereits auf
Kettenebene prüfbar, welche logischen Bögen dieses Dokument tatsächlich behauptet.

\subsection{Kettenblatt des modularen B-Kerns}
\doclevel{Datei-/Ketten-Ebene}
\textbf{1. Ziel der Kette.} Die Kette soll zeigen, wie aus dem komplementären lokalen Netz über
Zustand und Darstellung eine modulare Darstellungsarchitektur entsteht, die für Spektralkontrolle,
Bisognano--Wichmann und spätere Statistik-/CPT-Lesarten tragfähig ist.\par
\textbf{2. Definitionen der Kette.} Zentrale Objekte sind \code{ComplementLocalNet},
\code{ComplementStateNet}, \code{ComplementGNS}, \code{SeparatingProperty},
\code{TomitaTakesakiData}, \code{ModularConjugation}, \code{ComplementKMS} und
\code{ComplementSpectrumCondition}.\par
\textbf{3. Voraussetzungen der Kette.} Ein sektor-first Netzpfad aus A, kontrollierte Zustandsfamilien,
positive Energie-/Unitaritätsannahmen soweit für Spektralkontrolle nötig, keine boundary-first
Abkürzung.\par
\textbf{4. Behauptung der Kette.} Der operative Unterpfad 6a--6f stellt einen geschlossenen modularen
B-Kern bereit, dessen Endpunkt nicht bloss Thermik, sondern eine B$\to$D-tragfähige modulare Lage ist.\par
\textbf{5. Beweis der Kette.} \code{ComplementStateNet} liefert die Zustandslage, \code{ComplementGNS}
die Darstellung, \code{SeparatingProperty} stärkt diese zu zyklisch+separierend, \code{TomitaTakesakiData}
und \code{ModularConjugation} extrahieren die modulare Struktur, \code{ComplementKMS} schliesst die
Thermikschicht, \code{ComplementSpectrumCondition} bindet diese an Vorwärtskegel-/Positivitätsdaten.\par
\textbf{6. Zielabgleich der Kette.} Die Roadmap behauptet damit ausdrücklich keinen fertigen BW- oder
CPT-Beweis, wohl aber den geschlossenen Vorpfad, ohne den diese späteren Aussagen architektonisch
nicht derived-only getragen wären.

\subsection{Kettenblatt des Parameter-Closure-Strangs}
\doclevel{Datei-/Ketten-Ebene}
\textbf{1. Ziel der Kette.} Freie Restparameter sollen als offener Scaffold-Zustand sichtbar bleiben und
auf dem aktiven Pfad in einen derived-only Closure-Strang überführt werden.\par
\textbf{2. Definitionen der Kette.} \code{BranchingWitness}, \code{UVSpectralSelector},
\code{BranchingSelector}, \code{BackreactionSelector}, \code{ParameterClosure},
\code{HorizonLevel}, \code{BackreactionFixedPoint}, \code{RegularizationClosure}.\par
\textbf{3. Voraussetzungen der Kette.} sektorisierte ToC-Daten, Spektraldiagnostik, Kanal- und
Rückwirkungsdaten aus C, keine freie AS-/NCG-/Window-Sprache vor Closure.\par
\textbf{4. Behauptung der Kette.} Der Unterpfad PC-a--PC-d erzeugt eine Closure-Schwelle, nach der
ontische, dynamische und numerische Restparameter nicht mehr unkontrolliert als freie Inputs fungieren.\par
\textbf{5. Beweis der Kette.} Aus \code{BranchingWitness} und Spektraldiagnostik entstehen Selektoren,
diese speisen \code{ParameterClosure}; daraus folgen \code{HorizonLevel} und
\code{BackreactionFixedPoint}; \code{RegularizationClosure} schliesst oder degradiert die numerische
Stabilisierung zu abgeleiteten Restgrößen.\par
\textbf{6. Zielabgleich der Kette.} Das Dokument beansprucht hier keinen bereits formal geschlossenen
Fixpunktsatz, wohl aber die zwingende Position dieser Kette als derived-only Schwelle vor später
geometrischer und matter-seitiger Sprache.

\subsection{Kettenblatt der D/E-Brücke}
\doclevel{Datei-/Ketten-Ebene}
\textbf{1. Ziel der Kette.} Die D/E-Brücke soll den Übergang von geometrisch-dynamischer Zulassigkeit
zu späteren Symmetrie-, Statistik-, Charge- und Matter-Schichten explizit staffeln.\par
\textbf{2. Definitionen der Kette.} \code{InterfaceCausality}, \code{LocalCovariance},
\code{RelativeCauchyEvolutionSeed}, \code{HadamardStateSeed}, \code{LocalWickPolynomialsSeed},
\code{DHRSectorsSeed}, \code{FieldAlgebraReconstructionSeed}, \code{CPTGate},
\code{StatisticsGate}, \code{KramersGate}.\par
\textbf{3. Voraussetzungen der Kette.} geschlossener B-Kern, D-Zulassigkeit, Hadamard-/Wick-Schicht,
DHR/DR-Vorzone, keine freie Matter-Etikettierung vor CPT-/Charge-/Statistikbrücke.\par
\textbf{4. Behauptung der Kette.} Zwischen D und E liegt eine explizite Gate-Reihenfolge: geometrische
Zulassigkeit, renormierte lokale Dynamik, Symmetrie-/Statistikbrücke, rekonstruierte Charge-/Gauge-
Schicht und erst danach volle Matter-Organisation.\par
\textbf{5. Beweis der Kette.} D erzeugt erst die lokal-kovariante und renormierte Hintergrundlage;
Hadamard/Wick stabilisieren die lokale QFT-Sprache; DHR/DR rekonstruieren Charge- und Feldalgebra;
CPT-/Statistics-/Kramers-Gates binden dies an die späte Matter-Ebene.\par
\textbf{6. Zielabgleich der Kette.} Die Roadmap macht damit sichtbar, dass E nicht isoliert startet,
sondern ein kontrollierter Endverbraucher der zuvor geschlossenen B-, C- und D-Schichten ist.


\newpage
\section{Operative Regeln für De-seeding, Regression und Benennung}
\doclevel{Projekt-Ebene und Meta-Audit-Ebene}

\subsection{De-seeding-Regel}
\doclevel{Projekt-Ebene und Meta-Audit-Ebene}
Ein \code{Seed}-, \code{Scaffold}- oder \code{Window}-Modul verliert sein Suffix nur dann, wenn
(1) eine explizite Source-Map vorliegt, (2) der zentrale Definitionskern keine freien Cutoffs,
Zustände, Dictionaries oder Beobachter mehr enthält, (3) das Modul ein Derived-only-Gate auf
dem kritischen Pfad besteht und (4) für alle neu eingeführten tragenden Definitionen eine
entsprechende, dokumentierte Notation in der jeweils einschlägigen \path{Notation.lean} existiert,
die den mathematischen Standard der Spezifikation lesbar abbildet. Ohne diese Notationsbindung ist ein
De-seeding abzubrechen. Diese Regel ist direkt aus der internen Spezifikation übernommen und hier
um die operative Lesbarkeitsbedingung erweitert.\cite[pp.~6--10]{CNNASpec}

\subsection{Regression gegen den Altpfad}
\doclevel{Projekt-Ebene und Meta-Audit-Ebene}
Der alte BoundaryMatrix-/Bright-Pfad bleibt im Projekt, aber in neuer Rolle:
\begin{itemize}
\item als Regressionsorakel für Bright-Recovery,
\item als Extraktionsquelle für wirklich allgemeine algebraische oder netztheoretische Lemmata,
\item als Nachweis, dass CNNA den aktiven Altpfad später reproduzieren kann.
\end{itemize}
Nicht zulässig ist dagegen, die alte Datei- und Namensordnung schrittweise wieder zum Kern zu machen.

\subsection{Benennungsregel und Immediate Readability Policy}
\doclevel{Projekt-Ebene und Meta-Audit-Ebene}
Unsuffixed Dateinamen sind nur für genuinely derived Kernobjekte reserviert. Alles andere bleibt
bewusst als \code{Seed}, \code{Scaffold} oder \code{Window} markiert. Die Modulnamen müssen daher
selbst schon Auskunft über den epistemischen Status geben; genau das fordert die interne
CNNA-Spezifikation.\cite[pp.~6--10]{CNNASpec}

Für Neuentwicklungen jenseits des Archivstands v0.0605 gilt zusätzlich eine \emph{Immediate Readability
Policy}: roher Lean-Name allein genügt für tragende Kernobjekte nicht. Neue Definitionen müssen
entweder selbstsprechend benannt sein oder sofort über eine zugehörige Notationsschicht in
\path{CNNA/Notation.lean} bzw. \path{CNNA/PillarX/Notation.lean} lesbar gemacht werden. Ein Modul,
das nur über opake Projektnamen verständlich bleibt, gilt operativ als noch nicht reviewreif.

\subsection{Konsistente Ableitung der Notation-Dateien aus der Roadmap}
\doclevel{Projekt-Ebene und Meta-Audit-Ebene}
Die globalen und pfeilerspezifischen \texttt{Notation.lean}-Dateien sind nicht als späte manülle
Schönheitskorrektur zu behandeln, sondern relational aus drei kanonischen Tabellenlagen der Roadmap
abzuleiten:
\begin{enumerate}
\item der \textbf{Definitionskatalog} als Master-Menge aller tragenden Lean-Objekte;
\item das \textbf{Register} als normalisierte Oberflächen-, Symbol-, Status- und Meta-Lage;
\item der \textbf{Notations-Appendix} als Emissionsmatrix der tatsächlich exportierten Lesbarkeitsschicht.
\end{enumerate}
Daraus folgen vier bindende Regeln:
\begin{itemize}
\item Register und Appendix dürfen kein Fachobjekt eintragen, das nicht im Definitionskatalog steht.
\item Jede stabile sichtbare Oberflächenform muss auf genau einen kanonischen Lean-Namen oder auf einen explizit als \emph{meta only} markierten Begriff zeigen.
\item Jede exportierte Lesbarkeitsform braucht Zielmodul, Scope-/Namespace-Regel und konkrete Expansion auf den Lean-Term.
\item Jede belastbare Codeeinsicht, die Benennung, Symbolik oder Scope eines tragenden Objekts ändert, muss zürst den Definitionskatalog, danach Register und Emissionsmatrix und erst dann die realen \texttt{Notation.lean}-Module aktualisieren.
\end{itemize}
Operativ heißt das:
\begin{itemize}
\item \path{CNNA/Notation.lean} wird ausschliesslich aus globalen Kernbegriffen und global zulassigen Alias-/Symbolformen gespeist;
\item \path{CNNA/PillarA/Notation.lean} bis \path{CNNA/PillarE/Notation.lean} sowie \path{CNNA/AQFT/Notation.lean} und \path{CNNA/OQS/Notation.lean} werden aus den jeweiligen pfeiler- oder fachbereichsspezifischen Tabellenzeilen emittiert;
\item Reifegradmarker wie \code{Seed}, \code{Scaffold} und \code{Window} sowie Meta-Begriffe wie \code{TruthProtocol} und \code{ExpertReview} sind registerpflichtig, aber \emph{nicht} automatisch notationsfähig;
\item das Notationsreview ist fester Bestandteil jedes Meilensteins, der aus der partiellen Ordnung gewonnen wird.
\end{itemize}
Wo die technische Umsetzung automatisiert werden kann, ist eine generierte oder halbautomatisch
abgeglichene \path{Notation.lean}-Schicht aus den Roadmap-Tabellen ausdrücklich zu bevorzugen,
solange die resultierende Lean-Notation kollisionsfrei, \code{scoped} kontrolliert und reviewerfähig
bleibt.

\subsection{Repo-Infrastruktur, Importgraph und Rollback}
\doclevel{Projekt-Ebene und Meta-Audit-Ebene}
Vor jeder fachlichen De-seeding-Welle gelten vier nichttheorematische Pflichtbedingungen:

\begin{enumerate}
\item \textbf{pinned Toolchain:} \path{lean-toolchain} und \path{lakefile.lean} dürfen während einer
laufenden Migrationswelle kein bewegliches Ziel sein; jede harte Gate-Stufe arbeitet gegen eine
eingefrorene Toolchain;
\item \textbf{Build-Aggregatoren:} \path{CNNA/BuildAll.lean} sowie alle
\path{CNNA/PillarX/BuildAll.lean}-Dateien werden früh erzeugt und in CI bzw. lokaler
Build-Automation aufgerufen;
\item \textbf{Importgraph-Audit:} für jeden größeren Modulblock ist vor dem De-seeding ein
Importgraph mit Zyklentest Pflicht; gerade QuasiLocal-, Gates-, Meta- und Recovery-Blöcke dürfen
keine stillen Rückimporte in den Kern erzeugen;
\item \textbf{Rollback-Artefakte:} jedes harte Build-Gate erzeugt einen eingefrorenen Rückfallpunkt
(Archivsnapshot oder ähnlich strikte Versionsmarke), damit fehlgeschlagene De-seeding-Versuche nicht
unbemerkt in spätere Pfeiler einsickern.
\end{enumerate}

Diese Infrastrukturregeln sind kein Fremdthema, sondern unmittelbarer Teil der Migration. Gerade die
im Archiv bereits vorhandenen Root-, Meta- und Example-Dateien zeigen, dass Import- und
Regressionsstabilität nicht erst \emph{nach} der Mathematik beginnen.



\newpage
\section{Schlussurteil}
\doclevel{Projekt-Ebene}

Die richtige Migration von REAL-OQS v0.0605 nach CNNA ist \emph{keine} lineare Fortsetzung des alten
BoundaryMatrix-Pfads. Sie ist eine begriffliche Neuorganisation des Projekts. Der portierbare Kern aus
A--C wird beibehalten, aber semantisch gedreht:
\begin{itemize}
\item von Boundary-first zu Sector-first,
\item von einzelnem Komplement zu organisierter Komplementfamilie,
\item von eliminierter Umgebung zu kanalisch explizierter Rückwirkung,
\item von frühen Spezialfenstern zu später, streng kontrollierter Recovery,
\item von einer einzigen QFT-Oberfläche zu mehreren aus derselben Stammtheorie emergierenden Faces.
\end{itemize}

Genau dadurch werden die Snowmass-Forderung nach einer offeneren, gesichtssensitiven QFT-Begriffslandschaft,
die interne CNNA-Spezifikation und der tatsächliche Codebestand v0.0605 miteinander vertragen.\cite{Dedushenko2023}\cite{CNNASpec}\cite{RulebookV18}
Die Modulmatrix der Appendices deckt den vollständigen Archivbestand ab; die verbleibenden
offenen Punkte liegen damit nicht mehr auf der Ebene architektonischer Grundsatzfragen, sondern auf
der Ebene letzter physikalischer Präzisierungen eines im Kern bereits geschlossenen
Migrationspfads.
\newpage
\appendix

\section{Erweiterte Modulmatrix für die Migration}
\doclevel{Projekt-Ebene und Pfad-Ebene}

\subsection{Direkte oder nahezu direkte Ports}
\doclevel{Projekt-Ebene und Pfad-Ebene}
\tablecaptionline{Direkte oder nahezu direkte Ports -- Teil 1}
\begin{longtable}{L{0.36\linewidth}L{0.25\linewidth}L{0.29\linewidth}}
\toprule
Altmodul & Ziel in CNNA & Disposition \\
\midrule
\endfirsthead
\toprule
Altmodul & Ziel in CNNA & Disposition \\
\midrule
\endhead
\midrule
\multicolumn{3}{r}{Fortsetzung auf der nächsten Seite} \\
\midrule
\endfoot
\bottomrule
\endlastfoot
\path{REALOQS/PillarA/Core/Approximant.lean} & \path{CNNA/Core/Approximant.lean} & nahezu direkt portieren; Namespace und Importkegel bereinigen \\
\path{REALOQS/PillarA/Core/RegionCore.lean} & \path{CNNA/Core/RegionCore.lean} & nahezu direkt portieren; nur Semantik auf Sektorbild umstellen \\
\path{REALOQS/PillarA/Core/RegionNet.lean} & \path{CNNA/Core/RegionNet.lean} & als Trager fur Bright/Dark/Interface-Regionen erhalten \\
\path{REALOQS/PillarA/Core/InfiniteCarrier.lean} & \path{CNNA/Core/InfiniteCarrier.lean} & beibehalten; keine boundary-first Zusatzsemantik \\
\path{REALOQS/PillarA/Core/Selection.lean} & \path{CNNA/Core/Selection.lean} & beibehalten; Beobachterbegriff nicht frei einfuhren \\
\path{REALOQS/PillarA/OQS/DtN.lean} & \path{CNNA/OQS/DtN.lean} & als binarer Sonderfall erhalten; Ausgangspunkt fur GeneralizedDtN \\
\path{REALOQS/PillarA/OQS/DtNStabilized.lean} & \path{CNNA/OQS/DtNStabilized.lean} & als technische Hilfsschicht erhalten, aber methodisch dezentralisieren \\
\path{REALOQS/PillarB/AQFT/StarAlgebra.lean} & \path{CNNA/AQFT/StarAlgebra.lean} & direkt portieren \\
\path{REALOQS/PillarB/AQFT/State.lean} & \path{CNNA/AQFT/State.lean} & direkt portieren \\
\path{REALOQS/PillarB/AQFT/StateNet.lean} & \path{CNNA/AQFT/StateNet.lean} & direkt portieren und später sektorisieren \\
\path{REALOQS/PillarB/AQFT/LocalNet.lean} & \path{CNNA/AQFT/LocalNet.lean} & direkt portieren und später sektorisieren \\
\path{REALOQS/PillarB/AQFT/GNS.lean} & \path{CNNA/AQFT/GNS.lean} & direkt portieren; ComplementGNS darauf aufbauen \\
\path{REALOQS/PillarB/AQFT/KMS.lean} & \path{CNNA/AQFT/KMS.lean} & direkt portieren; ComplementKMS darauf aufbauen \\
\path{REALOQS/PillarB/AQFT/RelativeEntropy.lean} & \path{CNNA/AQFT/RelativeEntropy.lean} & direkt portieren \\
\path{REALOQS/PillarC/OQS/Channel.lean} & \path{CNNA/OQS/SectorChannels.lean} & semantisch verbreitern von System-Umgebung auf Sektor-Kanale \\
\path{REALOQS/PillarC/OQS/Stinespring.lean} & \path{CNNA/OQS/Stinespring-kompatible Unterlage} & direkt als mathematische Grundschicht behalten \\
\end{longtable}

\subsection{Semantische Umschriften}
\doclevel{Projekt-Ebene und Pfad-Ebene}
\tablecaptionline{Semantische Umschriften -- Teil 1}
\begin{longtable}{L{0.36\linewidth}L{0.3\linewidth}L{0.28\linewidth}}
\toprule
Altmodul & Ziel in CNNA & Disposition \\
\midrule
\endfirsthead
\toprule
Altmodul & Ziel in CNNA & Disposition \\
\midrule
\endhead
\midrule
\multicolumn{3}{r}{Fortsetzung auf der nächsten Seite} \\
\midrule
\endfoot
\bottomrule
\endlastfoot
\path{REALOQS/PillarA/OQS/SysEnv.lean} & \path{CNNA/OQS/SectorSysEnv.lean} & binaren Split durch Mehrsektor-Split ersetzen \\
\path{REALOQS/PillarA/Update/TailEliminationDtN.lean} & \path{CNNA/OQS/MultiSchur.lean + CNNA/Integration/TailElimAsChannel.lean} & Tail-Elimination nicht nur eliminatorisch, sondern kanalisch lesen \\
\path{REALOQS/PillarA/Core/TailEliminationCoherence.lean} & \path{CNNA/Core/ParameterClosure.lean + CNNA/OQS/RegularizationClosure.lean} & Koherenz in Closure-/Backreaction-Sprache uberfuhren \\
\path{REALOQS/PillarC/Integration/DerivedSpacetime.lean} & \path{CNNA/PillarD/DerivedSpacetime.lean} & von C-Endpunkt zu D-Vorbereitungsobjekt umdeuten \\
\path{REALOQS/PillarC/Integration/Locality.lean} & \path{CNNA/PillarD/InterfaceCausality.lean + CNNA/PillarD/LocalCovariance.lean} & lokale Einflussrelation in lokale Kovarianz uberfuhren \\
\end{longtable}

\subsection{Bright-Spezialfall und Regression}
\doclevel{Projekt-Ebene und Pfad-Ebene}
\tablecaptionline{Bright-Spezialfall und Regression -- Teil 1}
\begin{longtable}{L{0.39\linewidth}L{0.21\linewidth}L{0.26\linewidth}}
\toprule
Altmodul & Status & Folge \\
\midrule
\endfirsthead
\toprule
Altmodul & Status & Folge \\
\midrule
\endhead
\midrule
\multicolumn{3}{r}{Fortsetzung auf der nächsten Seite} \\
\midrule
\endfoot
\bottomrule
\endlastfoot
\path{REALOQS/PillarA/Ideal/Adapter/TreeOfCliquesAQFTHandoff.lean} & Bright-Spezialfall / Regressionsorakel & nicht Architekturzentrum, sondern Vergleichspfad \\
\path{REALOQS/PillarB/AQFT/Derived/TreeOfCliquesFromPillarA.lean} & Bright-Spezialfall / Regressionsorakel & nur replayen, nicht zentrieren \\
\path{REALOQS/PillarB/AQFT/Derived/BoundaryMatrixLocalNet.lean} & Bright-Spezialfall unter CNNA/AQFT & alte Randnetz-Semantik als abgeleitete Restriktion rekonstruieren \\
\path{REALOQS/PillarB/AQFT/Derived/BoundaryMatrixStateNet.lean} & Bright-Spezialfall unter CNNA/AQFT & nur als Restriktion auf Bright \\
\path{REALOQS/PillarB/AQFT/Derived/BoundaryMatrixGNS.lean} & Bright-Spezialfall unter CNNA/AQFT & nur als Restriktion auf Bright \\
\path{REALOQS/PillarB/AQFT/Derived/BoundaryMatrixFullKMSFromL.lean} & Bright-Spezialfall unter CNNA/AQFT & nur als Restriktion auf Bright \\
\path{REALOQS/PillarB/AQFT/Derived/BoundaryMatrixQuasiLocalClosure.lean} & Bright-Spezialfall / Regression & keine architektonische Zentralstellung \\
\path{REALOQS/PillarB/AQFT/Derived/BoundaryMatrixModular.lean} & Bright-Spezialfall / Regression & modulare Daten später aus CNNA zuruckgewinnen \\
\end{longtable}
Die explizite Negativregel für diesen Appendix lautet: \path{BoundaryMatrixModular.lean},
\path{BoundaryMatrixTimeReversal.lean} und verwandte BM-derived-Dateien bleiben Recovery- oder
Extraktionsquellen. Sie dürfen die neue modulare B-Kernschicht aus \code{SeparatingProperty},
\code{TomitaTakesakiData}, \code{ModularConjugation} und \code{BisognanoWichmannSeed} weder
ersetzen noch logisch abkürzen.

\newpage
\section{Vollständige Disposition des BoundaryMatrix-Derived-Blocks}
\doclevel{Projekt-Ebene und Datei-/Ketten-Ebene}
\label{app:bm-full}

Der komplette \path{BoundaryMatrix*}-Block umfasst im Archiv 28 Dateien mit zusammen 7'972 LOC. Er
wird in CNNA weder pauschal verworfen noch pauschal zum Kern erklärt. Die Roadmap trennt jetzt
explizit zwischen Bright-Recovery, Extraktionsquellen für allgemeine Mathematik und Gate-/Auditquellen.

\tablecaptionline{Vollständige Disposition des BoundaryMatrix-Derived-Blocks -- Teil 1}
\begin{longtable}{L{0.39\linewidth}L{0.09\linewidth}L{0.38\linewidth}}
\toprule
Altmodul & LOC & Disposition in CNNA \\
\midrule
\endfirsthead
\toprule
Altmodul & LOC & Disposition in CNNA \\
\midrule
\endhead
\midrule
\multicolumn{3}{r}{Fortsetzung auf der nächsten Seite} \\
\midrule
\endfoot
\bottomrule
\endlastfoot
\path{REALOQS/PillarB/AQFT/Derived/BoundaryMatrixAdditiveDynamics.lean} & 465 & Bright-Recovery plus Extraktionsquelle für spätere Komplement-Additivitätsgates; nicht Kernaxiom \\
\shortstack[l]{\path{REALOQS/PillarB/AQFT/Derived/}\\\path{BoundaryMatrixAdditiveDynamics}\\\path{Decision.lean}} & 757 & \shortstack[l]{Entscheidungs- und Auditquelle für\\additiv vs.~nichtadditiv; in\\\code{CNNA/Meta} und \code{CNNA/Gates} auswerten} \\
\path{REALOQS/PillarB/AQFT/Derived/BoundaryMatrixAdditiveProperties.lean} & 160 & Extraktionsquelle für additive Netz-/Zustandslemmata; erst nach stabilem Komplementnetz \\
\path{REALOQS/PillarB/AQFT/Derived/BoundaryMatrixAdditiveStateNet.lean} & 241 & Bright-Recovery-Unterstrang; nie vor \code{ComplementStateNet} \\
\path{REALOQS/PillarB/AQFT/Derived/BoundaryMatrixAdditiveSubnet.lean} & 259 & Bright-Recovery und Vergleichspfad für spätere Additivitätsabnahmen \\
\path{REALOQS/PillarB/AQFT/Derived/BoundaryMatrixAlgebra.lean} & 47 & algebraische Extraktionsquelle; generische Lemmata nach \code{CNNA/AQFT} ziehen \\
\path{REALOQS/PillarB/AQFT/Derived/BoundaryMatrixCStar.lean} & 163 & C*-Extraktionsquelle; nicht stillschweigend verlieren \\
\path{REALOQS/PillarB/AQFT/Derived/BoundaryMatrixDensityState.lean} & 590 & Gibbs-/Dichtezustandsquelle für Bright-Recovery und allgemeine Zustandslemmata \\
\path{REALOQS/PillarB/AQFT/Derived/BoundaryMatrixDynamicsFromL.lean} & 170 & Bright-Dynamikquelle; nur als Recovery- oder Extraktionspfad \\
\path{REALOQS/PillarB/AQFT/Derived/BoundaryMatrixFullDynamics.lean} & 534 & substanzielle Dynamikableitung; als Recovery- und Extraktionsquelle explizit behalten \\
\path{REALOQS/PillarB/AQFT/Derived/BoundaryMatrixFullExp.lean} & 305 & Exponential-/Semigroup-Hilfsschicht; generische Teile extrahieren \\
\path{REALOQS/PillarB/AQFT/Derived/BoundaryMatrixFullKMSFromL.lean} & 807 & zentrales Bright-KMS-Recovery-Orakel; nicht vor \code{ComplementKMS} \\
\path{REALOQS/PillarB/AQFT/Derived/BoundaryMatrixGNS.lean} & 173 & Bright-Recovery für GNS; kein Kernstartpunkt \\
\path{REALOQS/PillarB/AQFT/Derived/BoundaryMatrixGibbsState.lean} & 118 & thermische Zustandsquelle; generische Teile extrahieren \\
\path{REALOQS/PillarB/AQFT/Derived/BoundaryMatrixKMSFromL.lean} & 236 & thermische Extraktions- und Recovery-Quelle; späterer Bright-Vergleich \\
\path{REALOQS/PillarB/AQFT/Derived/BoundaryMatrixLocalNet.lean} & 594 & Bright-Recovery-Lokalnetz; als Restriktion rekonstruieren, nicht als CNNA-Kern lesen \\
\path{REALOQS/PillarB/AQFT/Derived/BoundaryMatrixLocality.lean} & 112 & HK-/Locality-Extraktionsquelle für Recovery und Gates \\
\path{REALOQS/PillarB/AQFT/Derived/BoundaryMatrixModular.lean} & 240 & modulare Recovery-Schicht; erst nach GNS/KMS-Vergleich \\
\path{REALOQS/PillarB/AQFT/Derived/BoundaryMatrixNontriviality.lean} & 91 & Gate-/Regressionquelle für Nichttrivialität \\
\path{REALOQS/PillarB/AQFT/Derived/BoundaryMatrixObservableBasis.lean} & 114 & Extraktionsquelle für Observable-Basis-Sprache \\
\path{REALOQS/PillarB/AQFT/Derived/BoundaryMatrixPhases.lean} & 94 & späte Phasen-/Diagnostikquelle; kein Kernmodul \\
\path{REALOQS/PillarB/AQFT/Derived/BoundaryMatrixPhysicalGibbs.lean} & 250 & thermische Extraktionsquelle mit physikalischer Lesart; nur nach Zustandskern \\
\path{REALOQS/PillarB/AQFT/Derived/BoundaryMatrixQuasiLocalClosure.lean} & 225 & Bright-Recovery für QuasiLocal-Closure; nach \code{ComplementClosure} \\
\path{REALOQS/PillarB/AQFT/Derived/BoundaryMatrixRichNetNonAdditive.lean} & 362 & explizite Nichtadditivitäts- und Auditquelle; in CNNA/Meta und Gates erhalten \\
\path{REALOQS/PillarB/AQFT/Derived/BoundaryMatrixSplitProperty.lean} & 343 & Split-Property-Quelle für Recovery und HK-Gates \\
\path{REALOQS/PillarB/AQFT/Derived/BoundaryMatrixStateNet.lean} & 187 & Bright-Recovery-StateNet; keine Kernsemantik \\
\path{REALOQS/PillarB/AQFT/Derived/BoundaryMatrixTimeReversal.lean} & 235 & Bright-Recovery plus Vorstufe für den Kramers-Strang; nicht als hinreichende Matter-Symmetrie lesen \\
\path{REALOQS/PillarB/AQFT/Derived/BoundaryMatrixToCNontriviality.lean} & 100 & ToC-spezifische Gate-/Regressionsquelle \\
\end{longtable}
\newpage
\section{Weitere Pillar-A Module aus Core, Symmetry und Kinematics}
\doclevel{Projekt-Ebene und Datei-/Ketten-Ebene}
\label{app:a-rest}

Die folgende Tabelle legt für die bisher implizit gelassenen A-Module eine explizite Disposition fest.
Sie verhindert, dass stille Importabhängigkeiten später als ``unbeabsichtigter Legacy-Schleier'' in
CNNA wieder auftauchen.

\tablecaptionline{Bisher untererfasste A-Restmodule aus Core, Symmetry und Kinematics -- Teil 1}
\begin{longtable}{L{0.39\linewidth}L{0.09\linewidth}L{0.38\linewidth}}
\toprule
Altmodul & LOC & Disposition in CNNA \\
\midrule
\endfirsthead
\toprule
Altmodul & LOC & Disposition in CNNA \\
\midrule
\endhead
\midrule
\multicolumn{3}{r}{Fortsetzung auf der nächsten Seite} \\
\midrule
\endfoot
\bottomrule
\endlastfoot
\path{REALOQS/PillarA/Core/Gates.lean} & 182 & als \code{CNNA/Gates/PillarAFoundation.lean} oder ähnliche Gate-Unterlage explizit portieren; kein stiller Root-Import \\
\path{REALOQS/PillarA/Core/SymmetryCoverageGates.lean} & 119 & Symmetrie-Auditschicht; in \code{CNNA/Gates/SymmetryCoverage.lean} überführen \\
\path{REALOQS/PillarA/Core/ResistanceMetricCanonical.lean} & 226 & D-Kinematik-/Geometrie-Seed; nicht verwerfen, aber vorerst nicht unsuffixed \\
\path{REALOQS/PillarA/Core/Policy.lean} & 116 & in \code{CNNA/ArchitecturePolicy.md} und Importregeln rückbinden; keine stille Nichtübernahme \\
\path{REALOQS/PillarA/Core/MatrixNorms.lean} & 79 & algebraische Hilfsschicht für DtN-/Operatorpfade; bei Bedarf direkt portieren \\
\path{REALOQS/PillarA/Core/SubstrateClass.lean} & 68 & ToC-/Substrat-Typeclass als explizite Kernunterlage erhalten \\
\path{REALOQS/PillarA/Core/SymmetryAction.lean} & 106 & Symmetrie-Seed für D/E; nicht vorzeitig physikalisch übersetzen \\
\path{REALOQS/PillarA/Core/DynamicsEquivariant.lean} & 39 & kleine, aber echte Äquivarianzquelle; in Symmetrie-/Dynamikgates einfliessen lassen \\
\path{REALOQS/PillarA/Core/InfiniteCarrierSymmetry.lean} & 53 & Infinite-Carrier-Symmetrie als Hilfsschicht erhalten \\
\path{REALOQS/PillarA/Core/NoetherHooks.lean} & 85 & bewusst als \code{Seed/Hook} belassen, bis echte derived Symmetriestränge existieren \\
\path{REALOQS/PillarA/Symmetry/ApproxIsometry.lean} & 117 & D-/Kinematik-Seed; nicht still in allgemeine Geometrie umetikettieren \\
\path{REALOQS/PillarA/Kinematics/DiscreteSpacetime.lean} & 55 & D-Seed; erst aus Interface/Inflünce-Daten de-seeden \\
\path{REALOQS/PillarA/Kinematics/Causality.lean} & 35 & D-Seed; nicht vor \code{InterfaceCausality} verabsolutieren \\
\path{REALOQS/PillarA/Kinematics/Worldline.lean} & 48 & D-Seed für spätere Beobachter-/Trajektorienfenster \\
\path{REALOQS/PillarA/Kinematics/ProperTime.lean} & 45 & D-Seed; kein freier Zeitbegriff auf dem Kernpfad \\
\path{REALOQS/PillarA/Kinematics/ProperTimeFromMetric.lean} & 89 & D-Seed für spätere Metrikfenster; vorerst nicht Kern \\
\path{REALOQS/PillarA/Kinematics/SpatialMetric.lean} & 23 & kleine, aber semantisch starke D-Seed-Datei \\
\path{REALOQS/PillarA/Kinematics/Clock.lean} & 79 & Beobachter-/Takt-Seed; nur spät de-seeden \\
\path{REALOQS/PillarA/Kinematics/FrameChange.lean} & 41 & Frame-Seed; nicht vor lokaler Kovarianz \\
\path{REALOQS/PillarA/Kinematics/SpacetimeRegions.lean} & 31 & D-Seed für spätere Regionen-/Window-Sprache \\
\end{longtable}
\newpage
\section{Disposition des B-Gates- und Meta-Subsystems}
\doclevel{Projekt-Ebene und Datei-/Ketten-Ebene}
\label{app:b-gates-meta}

\subsection{B-Gates als eigene Validierungsschicht}
\doclevel{Projekt-Ebene und Datei-/Ketten-Ebene}
\tablecaptionline{B-Gates als eigene Validierungsschicht -- Teil 1}
\begin{longtable}{L{0.40\linewidth}L{0.09\linewidth}L{0.37\linewidth}}
\toprule
Altmodul oder Modulgruppe & LOC & Disposition in CNNA \\
\midrule
\endfirsthead
\toprule
Altmodul oder Modulgruppe & LOC & Disposition in CNNA \\
\midrule
\endhead
\midrule
\multicolumn{3}{r}{Fortsetzung auf der nächsten Seite} \\
\midrule
\endfoot
\bottomrule
\endlastfoot
\path{REALOQS/PillarB/Gates/StarAlgebraBasic.lean} + \path{CStarNormBasic.lean} & 85 & als \code{CNNA/Gates/AQFT/StarAlgebraBasic.lean} bzw. \code{CStarNormBasic.lean} portieren \\
\path{REALOQS/PillarB/Gates/StateRestriction.lean} + \path{SplitProperty.lean} & 106 & direkt in die neue Gate-Ebene übernehmen; für Komplement- und Bright-Recovery gleichermassen relevant \\
\path{REALOQS/PillarB/Gates/QuasiLocalUP.lean} + \path{QuasiLocalStateLift.lean} & 47 & QuasiLocal-Gate-Unterlage; nicht durch blosses Umbenennen von \code{QuasiLocal/Basic} ersetzen \\
\path{REALOQS/PillarB/Gates/LocalNetFunctoriality.lean} & 28 & Vorstufe für spätere Glueing-/Functorialitätsgates \\
\path{REALOQS/PillarB/Gates/HaagKastler.lean} & 79 & Aggregator für HK-Gates; als \code{CNNA/Gates/HaagKastler.lean} beibehalten \\
\shortstack[l]{\path{REALOQS/PillarB/Gates/}\\\path{HaagKastler/\{Isotony,}\\\path{Locality,Spectrum,}\\\path{TimeSlice,Vacuum\}.lean}} & 230 & \shortstack[l]{direkte Gate-Portierung;\\zentrale Abnahmeschicht für das neue\\Komplementnetz, aber Spektrum/Vakuum müssen im\\CNNA-Kern neu als Komplementnetz-Lesart ausgefaltet werden} \\
\shortstack[l]{\path{REALOQS/PillarB/Gates/}\\\path{HaagKastler/\{Additivity,}\\\path{BasisAdditivity,}\\\path{StrongAdditivity,}\\\path{Covariance\}.lean}} & 259 & \shortstack[l]{explizite HK- und\\Additivitätsgates;\\wichtig für Vergleich rich vs.~additive net} \\
\shortstack[l]{\path{REALOQS/PillarB/Gates/}\\\path{\{ImportPolicy,Exports,}\\\path{ExampleFiniteClassical\}.lean}} & 73 & \shortstack[l]{Import-/Smoke-Test-Schicht;\\ausserhalb des Kerns, aber\\nicht unerwähnt} \\
\end{longtable}

\begingroup\small
\subsection{Meta-Dateien als Regressions- und Auditstrang}
\doclevel{Projekt-Ebene und Meta-Audit-Ebene}
\tablecaptionline{Meta-Dateien als Regressions- und Auditstrang -- Teil 1}
\begin{longtable}{L{0.40\linewidth}L{0.09\linewidth}L{0.37\linewidth}}
\toprule
Altmodul oder Modulgruppe & LOC & Disposition in CNNA \\
\midrule
\endfirsthead
\toprule
Altmodul oder Modulgruppe & LOC & Disposition in CNNA \\
\midrule
\endhead
\midrule
\multicolumn{3}{r}{Fortsetzung auf der nächsten Seite} \\
\midrule
\endfoot
\bottomrule
\endlastfoot
\path{REALOQS/Meta/PillarB_ObservableDynamicsDecision.lean} & 257 & als Legacy-Regressions- und Entscheidungsquelle für Additivität/Nichtadditivität explizit behalten \\
\path{REALOQS/Meta/PillarB_RichAudit.lean} + \path{PillarB_StrictSurface.lean} & 99 & Auditquelle für rich vs.~strict surface; nicht 1:1 Kern, aber echte Regression \\
\path{REALOQS/Meta/PillarB_ObservableAudit.lean} + \path{PillarB_ObservablePath.lean} + \path{PillarB_ObservableActivePath.lean} & 146 & dokumentieren, welche Altpfade der Bright-Recovery-Strang später replayen muss \\
\path{REALOQS/Meta/PillarB_ObservableBasisFromA.lean} + \path{PillarB_ObservableClosure\_*.lean} & 137 & Closure-/Basis-Regressionsquellen; in \code{CNNA/Meta/LegacyRegression/*} übernehmen \\
\path{REALOQS/Meta/PillarB_Closure\_*.lean} + \path{PillarB_BQ8\_Closed.lean} + \path{PillarB_BQ9\_Closed.lean} & 184 & dokumentieren, welche Closure-Zielbilder später wieder erreicht werden müssen \\
\shortstack[l]{\path{REALOQS/Meta/}\\\path{PillarB_ActivePath.lean} +\\\path{PillarB_NoConfigHarnessIn}\\\path{CriticalPath.lean}} & 83 & \shortstack[l]{direkte Vorlage für\\\code{CNNA/Meta/CriticalPath.lean}\\und \code{NoLegacyInCriticalPath.lean}} \\
\path{REALOQS/Meta/BuildAll.lean} + restliche kleine Observable-Meta-Dateien & 2 & Smoke-Test-/Aggregatorrolle; mitportieren oder im Root-Meta abbilden \\
\end{longtable}
\newpage
\section{Quell-Ziel-Abgleich für Namens- und Topologiefragen}
\doclevel{Projekt-Ebene und Meta-Audit-Ebene}
\label{app:naming-topology}

\tablecaptionline{Quell-Ziel-Abgleich für Namens- und Topologiefragen -- Teil 1}
\begin{longtable}{L{0.42\linewidth}L{0.30\linewidth}L{0.20\linewidth}}
\toprule
Quellmodul(e) & Ziel in CNNA & Bemerkung \\
\midrule
\endfirsthead
\toprule
Quellmodul(e) & Ziel in CNNA & Bemerkung \\
\midrule
\endhead
\midrule
\multicolumn{3}{r}{Fortsetzung auf der nächsten Seite} \\
\midrule
\endfoot
\bottomrule
\endlastfoot
\path{REALOQS/PillarB/AQFT/QuasiLocal/\{Carrier,Structure,StarAlgebra\}.lean} & \path{CNNA/AQFT/QuasiLocal/Basic.lean} & \code{Basic} ist neues CNNA-Aggregat, kein alter Dateiname \\
\path{REALOQS/PillarB/AQFT/QuasiLocal/\{Completion,StateLift\}.lean} & \path{CNNA/AQFT/QuasiLocal/ComplementClosure.lean} & \code{ComplementClosure} ist neue sektorisierte Lesart; bei Importdruck Hilfsunterdateien beibehalten \\
\path{REALOQS/PillarB/AQFT/HaagKastler/Support.lean} & \path{CNNA/AQFT/HaagKastler/Basic.lean} & \code{Basic} ist neues Zielminimum, nicht 1:1-Quellname \\
\path{REALOQS/PillarB/AQFT/HaagKastler/SupportFull.lean} & \path{CNNA/AQFT/HaagKastler/Complement.lean} & Voll-Support wird in die neue Komplementlesart überführt \\
\path{REALOQS/PillarB/AQFT/HaagKastler/SplitProperty.lean} & \path{CNNA/AQFT/HaagKastler/SplitProperty.lean} & 1:1-Port mit neuer Importumgebung \\
\path{REALOQS/PillarC/OQS/Finite/Kraus.lean} & \path{CNNA/OQS/Finite/Kraus.lean} & explizit in den Zielbaum aufnehmen \\
\path{REALOQS/PillarC/OQS/Finite/RectKraus.lean} & \path{CNNA/OQS/Finite/RectKraus.lean} & explizit in den Zielbaum aufnehmen \\
\path{REALOQS/PillarC/OQS/Finite.lean} & \path{CNNA/OQS/Finite.lean} & bleibt Aggregator \\
\path{REALOQS/PillarB/AQFT/Derived/LocalConfigNet.lean} + \path{StateFromConfigNet.lean} + \path{QuasiLocalStateFromConfigNet.lean} & \path{CNNA/AQFT/BrightRecovery/ConfigNet/*} & nicht verschwinden lassen; später Regression und Konfigurationsvergleich \\
\path{REALOQS/PillarB/AQFT/Derived/MatterFromPillarA.lean} + \path{EnergyFromBoundaryOp.lean} + \path{GibbsStateOnConfig.lean} & \path{CNNA/PillarE/LegacyMatterTests/*} & explizit als schwacher thermischer Matter-Vorbau klassifizieren, nicht als E-Kern \\
\path{REALOQS/PillarB/AQFT/Derived/TreeOfCliquesExtendedFromPillarA.lean} & \shortstack[l]{\path{CNNA/AQFT/BrightRecovery/}\\\path{TreeOfCliquesExtendedFrom}\\\path{PillarA.lean}} & \shortstack[l]{Regressions- und Spektrumsquelle,\\insbesondere für Vorwärtskegel-/Matter-Audits} \\
\end{longtable}
\endgroup

\subsection{Notationsmodule als Zielarchitektur}
\doclevel{Meta-Audit-Ebene}
Für die reale Codebasis bedeutet die Notationspflicht eine eigene modulare Schicht. Minimal benötigt
werden ein globales Modul \path{CNNA/Notation.lean} sowie bei Bedarf pillar- oder fachbereichsspezifische
Ergänzungen wie \path{CNNA/PillarA/Notation.lean}, \path{CNNA/OQS/Notation.lean},
\path{CNNA/AQFT/Notation.lean}, \path{CNNA/PillarD/Notation.lean} und \path{CNNA/PillarE/Notation.lean}.
Diese Module sind keine freie Stilzone, sondern direkte Emissionen der Roadmap-Tabellen. Deshalb gilt
hier zusätzlich:
\begin{itemize}
\item der Definitionskatalog ist die einzige kanonische Objektquelle;
\item das Register spaltet sich in Objekt-/Aliasregister, Symbolregister, Reifegradregister und Meta-Methodikregister;
\item der Notations-Appendix enthält die tatsächlich exportierten Alias- und Notationsentscheidungen samt Zielmodul, Scope und Emissionsmuster;
\item \code{selfspeaking}-Objekte bleiben bewusst ohne zusätzliche Notations-Emission, müssen aber im Definitionskatalog explizit als solche klassifiziert sein und können im Notations-Appendix als dokumentierte \emph{Nicht-Emission} mitgeführt werden.
\end{itemize}
Jede spätere Tabellen- oder Registeränderung der Roadmap wirkt damit unmittelbar auf diese
Notationsmodule zurück; die Notationsschicht ist also kein statischer Appendix, sondern ein direkt aus
dem executable Dokument abgeleiteter Teil der Arbeitsbasis.

\subsection{Notations- und Alias-Appendix für den aktiven CNNA-Pfad}
\doclevel{Meta-Audit-Ebene}
\label{app:notationlayer}
Die folgende Emissionsmatrix listet die im aktiven und spät aktivierten Pfad tatsächlich exportierten
Lesbarkeitsformen sowie explizit dokumentierte \code{selfspeaking}-Nicht-Emissionen. Jede Zeile ist technisch determiniert: kanonischer Lean-Name, Emissionsart,
Oberflächenform, emittiertes Muster, Zielmodul, Scope und Kollisionsregel. Nicht exportierte,
aber dennoch sichtbare Formen verbleiben ausschliesslich im Register; die einzigen Ausnahmen sind hier
explizit markierte \emph{Nicht-Emissionen} für tragende \code{selfspeaking}-Objekte. Die im Definitionskatalog
geführte Primärform bleibt die dokumentarische Leitform; weicht die exportierte Kurzform davon ab,
so ist sie als sekundäre Exportform sowohl hier als auch im Register I eigens auszuweisen.

Für jede emittierte Form sind Namespace, Aktivierungsregel, erforderlicher Kontext,
Präzedenz bzw. Bindungsstärke, Kollisionsschutz, Beispielverwendung, Status und
Review-Notiz tabellarisch festzuhalten. Keine dieser Informationen darf ausschliesslich
aus dem Zielmodul oder aus dem Scope-Feld indirekt erschlossen werden.

\begin{wideauditlandscape}
\tablecaptionline{Notations- und Alias-Appendix -- Teil 1}
\begin{tabular}{L{0.07\linewidth}L{0.06\linewidth}L{0.07\linewidth}L{0.14\linewidth}L{0.08\linewidth}L{0.07\linewidth}L{0.08\linewidth}L{0.08\linewidth}L{0.055\linewidth}L{0.07\linewidth}L{0.07\linewidth}L{0.04\linewidth}L{0.06\linewidth}}
\toprule
Lean-Objekt & Emissionsart & Oberflächenform & Lean-Emission / Expansion & Zielmodul & Namespace & Aktivierungsregel & erforderlicher Kontext & Präzedenz / Bindungsstärke & Kollisionsregel & Beispielverwendung & Status & Review-Notiz \\
\midrule
\code{TreeOfCliques} & text alias & \emph{ToC} & \code{scoped abbrev ToC := TreeOfCliques} & \path{CNNA/Notation.lean} & \code{CNNA} & \code{open scoped CNNA} & global & n/a & \emph{Tree of Cliques} bleibt Dokumentform; \emph{ToC} ist der einzige Kurzalias & \emph{ToC} & active core & Form exportiert; siehe Register I für Primär-/Sekundärtrennung \\
\code{BranchPatch} & text alias & \emph{Bright-Patch} & \code{scoped abbrev BrightPatch := BranchPatch} & \path{CNNA/PillarA/Notation.lean} & \code{CNNA.PillarA} & \code{open scoped CNNA.PillarA} & Pillar A & n/a & kein zweiter Alias \emph{Anker-Patch} & \emph{Bright-Patch} & active core & Form exportiert; siehe Register I für Primär-/Sekundärtrennung \\
\code{ComplementSectorFamily} & scoped notation & \(\mathcal C\) & \code{scoped notation "C\_comp" => ComplementSectorFamily} & \path{CNNA/PillarA/Notation.lean} & \code{CNNA.PillarA} & \code{open scoped CNNA.PillarA} & Pillar A--B & n/a & zusätzlicher Textalias \emph{Komplementsektorfamilie} nur im Register & \(\mathcal C\) & active core & Form exportiert; siehe Register I für Primär-/Sekundärtrennung \\
\code{SectorSplit} & scoped notation & \(\Sigma_{\mathrm{BDI}}\) & \code{scoped notation "SigmaBDI" => SectorSplit} & \path{CNNA/PillarA/Notation.lean} & \code{CNNA.PillarA} & \code{open scoped CNNA.PillarA} & Pillar A--B & n/a & keine konkurrierende Kurzform & \(\Sigma_{\mathrm{BDI}}\) & active core & Form exportiert; siehe Register I für Primär-/Sekundärtrennung \\
\code{GeneralizedDtN} & scoped notation & \(\Lambda_{\mathrm{gen}}\) & \code{scoped notation "LambdaGen" => GeneralizedDtN} & \path{CNNA/PillarA/Notation.lean} & \code{CNNA.PillarA} & \code{open scoped CNNA.PillarA} & Pillar A--B & n/a & Textalias \emph{generalisierter DtN-Kern} bleibt zulässig & \(\Lambda_{\mathrm{gen}}\) & active core & Form exportiert; siehe Register I für Primär-/Sekundärtrennung \\
\code{SectorChannels} & text alias & \emph{SectCh} & \code{scoped abbrev SectCh := SectorChannels} & \path{CNNA/OQS/Notation.lean} & \code{CNNA.OQS} & \code{open scoped CNNA.OQS} & Pillar C & n/a & Kurzalias für den kanalisierten OQS-Strang & \emph{SectCh} & active extension & Form exportiert; siehe Register I für Primär-/Sekundärtrennung \\
\code{SectorSysEnv} & text alias & \emph{SysEnvSplit} & \code{scoped abbrev SysEnvSplit := SectorSysEnv} & \path{CNNA/OQS/Notation.lean} & \code{CNNA.OQS} & \code{open scoped CNNA.OQS} & Pillar C & n/a & Kurzalias für den sektorischen System/Environment-Split & \emph{SysEnvSplit} & active extension & Form exportiert; siehe Register I für Primär-/Sekundärtrennung \\
\code{RelativeEntropyFlow} & text alias & \emph{RelEntFlow} & \code{scoped abbrev RelEntFlow := RelativeEntropyFlow} & \path{CNNA/OQS/Notation.lean} & \code{CNNA.OQS} & kontextgebunden; bei Export \code{open scoped CNNA.OQS} & Pillar C--D & n/a & Kurzalias für entropische Rückwirkungsflüsse & \emph{RelEntFlow} & active extension & Form exportiert; siehe Register I für Primär-/Sekundärtrennung \\
\bottomrule
\end{tabular}
\end{wideauditlandscape}

\begin{wideauditlandscape}
\tablecaptionline{Notations- und Alias-Appendix -- Teil 2}
\begin{tabular}{L{0.07\linewidth}L{0.06\linewidth}L{0.07\linewidth}L{0.14\linewidth}L{0.08\linewidth}L{0.07\linewidth}L{0.08\linewidth}L{0.08\linewidth}L{0.055\linewidth}L{0.07\linewidth}L{0.07\linewidth}L{0.04\linewidth}L{0.06\linewidth}}
\toprule
Lean-Objekt & Emissionsart & Oberflächenform & Lean-Emission / Expansion & Zielmodul & Namespace & Aktivierungsregel & erforderlicher Kontext & Präzedenz / Bindungsstärke & Kollisionsregel & Beispielverwendung & Status & Review-Notiz \\
\midrule
\code{ComplementLocalNet} & scoped notation & \(\mathcal A_{\mathrm{comp}}\) & \code{scoped notation "A\_comp" => ComplementLocalNet} & \path{CNNA/AQFT/Notation.lean} & \code{CNNA.AQFT} & \code{open scoped CNNA.AQFT} & AQFT & n/a & \emph{Komplementnetz} bleibt Registeralias & \(\mathcal A_{\mathrm{comp}}\) & active core & Form exportiert; siehe Register I für Primär-/Sekundärtrennung \\
\code{InterfaceLocalNet} & scoped notation & \(\mathcal A_{\mathrm{int}}\) & \code{scoped notation "A\_int" => InterfaceLocalNet} & \path{CNNA/AQFT/Notation.lean} & \code{CNNA.AQFT} & \code{open scoped CNNA.AQFT} & AQFT & n/a & \emph{Interfacenetz} bleibt Registeralias & \(\mathcal A_{\mathrm{int}}\) & active core & Form exportiert; siehe Register I für Primär-/Sekundärtrennung \\
\code{ComplementStateNet} & scoped notation & \(\omega_{\mathrm{comp}}\) & \code{scoped notation "omegaComp" => ComplementStateNet} & \path{CNNA/AQFT/Notation.lean} & \code{CNNA.AQFT} & \code{open scoped CNNA.AQFT} & AQFT & n/a & kein Konflikt mit \(\omega_\beta\) & \(\omega_{\mathrm{comp}}\) & active core & Form exportiert; siehe Register I für Primär-/Sekundärtrennung \\
\code{ComplementGNS} & text alias & \emph{CompGNS} & \code{scoped abbrev CompGNS := ComplementGNS} & \path{CNNA/AQFT/Notation.lean} & \code{CNNA.AQFT} & kontextgebunden; bei Export \code{open scoped CNNA.AQFT} & AQFT & n/a & die Tupelform \((\pi_\omega,\mathcal H_\omega,\Omega_\omega)\) bleibt Dokument-/Registerform & \emph{CompGNS} & active core & Form exportiert; siehe Register I für Primär-/Sekundärtrennung \\
\code{SeparatingProperty} & text alias & \emph{CycSep} & \code{scoped abbrev CycSep := SeparatingProperty} & \path{CNNA/AQFT/Notation.lean} & \code{CNNA.AQFT} & kontextgebunden; bei Export \code{open scoped CNNA.AQFT} & AQFT & n/a & lesbare Kurzform nur für den Reeh--Schlieder-artigen Kontext & \emph{CycSep} & active core & Form exportiert; siehe Register I für Primär-/Sekundärtrennung \\
\code{TomitaTakesakiData} & text alias & \emph{TTData} & \code{scoped abbrev TTData := TomitaTakesakiData} & \path{CNNA/AQFT/Notation.lean} & \code{CNNA.AQFT} & kontextgebunden; bei Export \code{open scoped CNNA.AQFT} & AQFT & n/a & die Tupelform \((S,J,\Delta)\) bleibt Registerform & \emph{TTData} & active core & Form exportiert; siehe Register I für Primär-/Sekundärtrennung \\
\code{ComplementKMS} & text alias & \emph{CompKMS} & \code{scoped abbrev CompKMS := ComplementKMS} & \path{CNNA/AQFT/Notation.lean} & \code{CNNA.AQFT} & kontextgebunden; bei Export \code{open scoped CNNA.AQFT} & AQFT & n/a & \(\omega_\beta\) nur in explizitem KMS-Scope & \emph{CompKMS} & active core & Form exportiert; siehe Register I für Primär-/Sekundärtrennung \\
\code{ComplementSpectrumCondition} & text alias & \emph{CompSpectrum} & \code{scoped abbrev CompSpectrum := ComplementSpectrumCondition} & \path{CNNA/AQFT/Notation.lean} & \code{CNNA.AQFT} & kontextgebunden; bei Export \code{open scoped CNNA.AQFT} & B--D & n/a & kein Kurzsymbol ohne Klartextkontext & \emph{CompSpectrum} & active core & Form exportiert; siehe Register I für Primär-/Sekundärtrennung \\
\bottomrule
\end{tabular}
\end{wideauditlandscape}

\begin{wideauditlandscape}
\tablecaptionline{Notations- und Alias-Appendix -- Teil 3}
\begin{tabular}{L{0.07\linewidth}L{0.06\linewidth}L{0.07\linewidth}L{0.14\linewidth}L{0.08\linewidth}L{0.07\linewidth}L{0.08\linewidth}L{0.08\linewidth}L{0.055\linewidth}L{0.07\linewidth}L{0.07\linewidth}L{0.04\linewidth}L{0.06\linewidth}}
\toprule
Lean-Objekt & Emissionsart & Oberflächenform & Lean-Emission / Expansion & Zielmodul & Namespace & Aktivierungsregel & erforderlicher Kontext & Präzedenz / Bindungsstärke & Kollisionsregel & Beispielverwendung & Status & Review-Notiz \\
\midrule
\code{ParameterClosure} & text alias & \emph{ParamClosure} & \code{scoped abbrev ParamClosure := ParameterClosure} & \path{CNNA/Notation.lean} & \code{CNNA} & \code{open scoped CNNA} & global & n/a & bindet den derived-only-Closure-Strang sichtbar & \emph{ParamClosure} & active core & Form exportiert; siehe Register I für Primär-/Sekundärtrennung \\
\code{EffectiveLambda} & scoped notation & \(\lambda_{\mathrm{eff}}\) & \code{scoped notation "lambdaEff" => EffectiveLambda} & \path{CNNA/PillarD/Notation.lean} & \code{CNNA.PillarD} & \code{open scoped CNNA.PillarD} & C--D & n/a & keine Vermischung mit nacktem \(\Lambda\) & \(\lambda_{\mathrm{eff}}\) & planned active & Form exportiert; siehe Register I für Primär-/Sekundärtrennung \\
\code{RegularizationClosure} & text alias & \emph{RegClosure} & \code{scoped abbrev RegClosure := RegularizationClosure} & \path{CNNA/PillarD/Notation.lean} & \code{CNNA.PillarD} & kontextgebunden; bei Export \code{open scoped CNNA.PillarD} & C--D & n/a & lesbare Kurzform für Closure-Reviews & \emph{RegClosure} & planned active & Form exportiert; siehe Register I für Primär-/Sekundärtrennung \\
\code{BranchingWitness} & text alias & \emph{BranchWit} & \code{scoped abbrev BranchWit := BranchingWitness} & \path{CNNA/PillarA/Notation.lean} & \code{CNNA.PillarA} & kontextgebunden; bei Export \code{open scoped CNNA.PillarA} & A--D & n/a & Witnesstyp nur in Branching-Kontexten & \emph{BranchWit} & active extension & Form exportiert; siehe Register I für Primär-/Sekundärtrennung \\
\code{selectedBranching} & text alias & \emph{SelBranch} & \code{scoped abbrev SelBranch := selectedBranching} & \path{CNNA/PillarA/Notation.lean} & \code{CNNA.PillarA} & kontextgebunden; bei Export \code{open scoped CNNA.PillarA} & A--D & n/a & keine freie Umbenennung auf \emph{chosen branching} & \emph{SelBranch} & active extension & Form exportiert; siehe Register I für Primär-/Sekundärtrennung \\
\code{UVSpectralSelector} & text alias & \emph{UVSel} & \code{scoped abbrev UVSel := UVSpectralSelector} & \path{CNNA/PillarA/Notation.lean} & \code{CNNA.PillarA} & kontextgebunden; bei Export \code{open scoped CNNA.PillarA} & A--D & n/a & Selektor-Notation bleibt reviewpflichtig & \emph{UVSel} & planned active & Form exportiert; siehe Register I für Primär-/Sekundärtrennung \\
\code{BranchingSelector} & text alias & \emph{BranchSel} & \code{scoped abbrev BranchSel := BranchingSelector} & \path{CNNA/PillarA/Notation.lean} & \code{CNNA.PillarA} & kontextgebunden; bei Export \code{open scoped CNNA.PillarA} & A--D & n/a & Selektor-Notation bleibt reviewpflichtig & \emph{BranchSel} & planned active & Form exportiert; siehe Register I für Primär-/Sekundärtrennung \\
\code{BackreactionSelector} & text alias & \emph{BackSel} & \code{scoped abbrev BackSel := BackreactionSelector} & \path{CNNA/PillarD/Notation.lean} & \code{CNNA.PillarD} & kontextgebunden; bei Export \code{open scoped CNNA.PillarD} & C--D & n/a & Selektor-Notation bleibt reviewpflichtig & \emph{BackSel} & planned active & Form exportiert; siehe Register I für Primär-/Sekundärtrennung \\
\bottomrule
\end{tabular}
\end{wideauditlandscape}

\begin{wideauditlandscape}
\tablecaptionline{Notations- und Alias-Appendix -- Teil 4}
\begin{tabular}{L{0.07\linewidth}L{0.06\linewidth}L{0.07\linewidth}L{0.14\linewidth}L{0.08\linewidth}L{0.07\linewidth}L{0.08\linewidth}L{0.08\linewidth}L{0.055\linewidth}L{0.07\linewidth}L{0.07\linewidth}L{0.04\linewidth}L{0.06\linewidth}}
\toprule
Lean-Objekt & Emissionsart & Oberflächenform & Lean-Emission / Expansion & Zielmodul & Namespace & Aktivierungsregel & erforderlicher Kontext & Präzedenz / Bindungsstärke & Kollisionsregel & Beispielverwendung & Status & Review-Notiz \\
\midrule
\code{HorizonLevel} & scoped notation & \(L_H\) & \code{scoped notation "L\_H" => HorizonLevel} & \path{CNNA/PillarD/Notation.lean} & \code{CNNA.PillarD} & \code{open scoped CNNA.PillarD} & C--D & n/a & ersetzt aktiv \(L_{\max}\); Legacyname nicht exportieren & \(L_H\) & active core & Form exportiert; siehe Register I für Primär-/Sekundärtrennung \\
\code{BackreactionFixedPoint} & text alias & \emph{BackFix} & \code{scoped abbrev BackFix := BackreactionFixedPoint} & \path{CNNA/PillarD/Notation.lean} & \code{CNNA.PillarD} & kontextgebunden; bei Export \code{open scoped CNNA.PillarD} & C--D & n/a & Kurzalias nur für die Fixpunktgleichung & \emph{BackFix} & active core & Form exportiert; siehe Register I für Primär-/Sekundärtrennung \\
\code{ComplementClosure} & text alias & \emph{CompClosure} & \code{scoped abbrev CompClosure := ComplementClosure} & \path{CNNA/AQFT/Notation.lean} & \code{CNNA.AQFT} & kontextgebunden; bei Export \code{open scoped CNNA.AQFT} & AQFT/C & n/a & keine Symbolnotation ohne klaren Schliessungskontext & \emph{CompClosure} & active extension & Form exportiert; siehe Register I für Primär-/Sekundärtrennung \\
\code{ModularConjugation} & text alias & \emph{ModConj} & \code{scoped abbrev ModConj := ModularConjugation} & \path{CNNA/AQFT/Notation.lean} & \code{CNNA.AQFT} & kontextgebunden; bei Export \code{open scoped CNNA.AQFT} & AQFT & n/a & \(J\) bleibt Symbolregistereintrag & \emph{ModConj} & active extension & Form exportiert; siehe Register I für Primär-/Sekundärtrennung \\
\code{BisognanoWichmannSeed} & text alias & \emph{BWSeed} & \code{scoped abbrev BWSeed := BisognanoWichmannSeed} & \path{CNNA/AQFT/Notation.lean} & \code{CNNA.AQFT} & kontextgebunden; bei Export \code{open scoped CNNA.AQFT} & AQFT & n/a & eindeutiger Kurzname für die BW-Brücke & \emph{BWSeed} & active extension & Form exportiert; siehe Register I für Primär-/Sekundärtrennung \\
\code{NuclearitySeed} & text alias & \emph{NucSeed} & \code{scoped abbrev NucSeed := NuclearitySeed} & \path{CNNA/AQFT/Notation.lean} & \code{CNNA.AQFT} & kontextgebunden; bei Export \code{open scoped CNNA.AQFT} & AQFT & n/a & keine Kurzform für \code{Nuclearity} selbst & \emph{NucSeed} & active extension & Form exportiert; siehe Register I für Primär-/Sekundärtrennung \\
\code{DHRSectorsSeed} & text alias & \emph{DHRSeed} & \code{scoped abbrev DHRSeed := DHRSectorsSeed} & \path{CNNA/AQFT/Notation.lean} & \code{CNNA.AQFT} & kontextgebunden; bei Export \code{open scoped CNNA.AQFT} & AQFT & n/a & kurzer Reviewname für DHR-Vorzone & \emph{DHRSeed} & active extension & Form exportiert; siehe Register I für Primär-/Sekundärtrennung \\
\code{FieldAlgebraReconstructionSeed} & text alias & \emph{FieldRecSeed} & \code{scoped abbrev FieldRecSeed := FieldAlgebraReconstructionSeed} & \path{CNNA/AQFT/Notation.lean} & \code{CNNA.AQFT} & kontextgebunden; bei Export \code{open scoped CNNA.AQFT} & AQFT & n/a & keine zweite Bedeutung für Rekonstruktionssatz & \emph{FieldRecSeed} & active extension & Form exportiert; siehe Register I für Primär-/Sekundärtrennung \\
\bottomrule
\end{tabular}
\end{wideauditlandscape}

\begin{wideauditlandscape}
\tablecaptionline{Notations- und Alias-Appendix -- Teil 5}
\begin{tabular}{L{0.07\linewidth}L{0.06\linewidth}L{0.07\linewidth}L{0.14\linewidth}L{0.08\linewidth}L{0.07\linewidth}L{0.08\linewidth}L{0.08\linewidth}L{0.055\linewidth}L{0.07\linewidth}L{0.07\linewidth}L{0.04\linewidth}L{0.06\linewidth}}
\toprule
Lean-Objekt & Emissionsart & Oberflächenform & Lean-Emission / Expansion & Zielmodul & Namespace & Aktivierungsregel & erforderlicher Kontext & Präzedenz / Bindungsstärke & Kollisionsregel & Beispielverwendung & Status & Review-Notiz \\
\midrule
\code{ChargeStatisticsCoherence} & text alias & \emph{ChargeStatCoh} & \code{scoped abbrev ChargeStatCoh := ChargeStatisticsCoherence} & \path{CNNA/AQFT/Notation.lean} & \code{CNNA.AQFT} & kontextgebunden; bei Export \code{open scoped CNNA.AQFT} & AQFT/Gates & n/a & Kurzalias nur für das Gate & \emph{ChargeStatCoh} & active extension & Form exportiert; siehe Register I für Primär-/Sekundärtrennung \\
\code{ChargeReconstructionTheorem} & text alias & \emph{ChargeRec} & \code{scoped abbrev ChargeRec := ChargeReconstructionTheorem} & \path{CNNA/AQFT/Notation.lean} & \code{CNNA.AQFT} & kontextgebunden; bei Export \code{open scoped CNNA.AQFT} & AQFT & n/a & Kurzalias für Reviewtabellen & \emph{ChargeRec} & active extension & Form exportiert; siehe Register I für Primär-/Sekundärtrennung \\
\code{DerivedSpacetime} & text alias & \emph{DerSpacetime} & \code{scoped abbrev DerSpacetime := DerivedSpacetime} & \path{CNNA/PillarD/Notation.lean} & \code{CNNA.PillarD} & kontextgebunden; bei Export \code{open scoped CNNA.PillarD} & D & n/a & keine Konkurrenz zum Prosaausdruck \emph{derived spacetime} & \emph{DerSpacetime} & active extension & Form exportiert; siehe Register I für Primär-/Sekundärtrennung \\
\code{InterfaceCausality} & text alias & \emph{IfcCaus} & \code{scoped abbrev IfcCaus := InterfaceCausality} & \path{CNNA/PillarD/Notation.lean} & \code{CNNA.PillarD} & kontextgebunden; bei Export \code{open scoped CNNA.PillarD} & D & n/a & Kurzalias nur für D-Reviews & \emph{IfcCaus} & active extension & Form exportiert; siehe Register I für Primär-/Sekundärtrennung \\
\code{LocalCovariance} & text alias & \emph{LocCov} & \code{scoped abbrev LocCov := LocalCovariance} & \path{CNNA/PillarD/Notation.lean} & \code{CNNA.PillarD} & kontextgebunden; bei Export \code{open scoped CNNA.PillarD} & D & n/a & keine Konkurrenz zu mathlib-Kontexten & \emph{LocCov} & active extension & Form exportiert; siehe Register I für Primär-/Sekundärtrennung \\
\code{HadamardStateSeed} & text alias & \emph{HadSeed} & \code{scoped abbrev HadSeed := HadamardStateSeed} & \path{CNNA/PillarD/Notation.lean} & \code{CNNA.PillarD} & kontextgebunden; bei Export \code{open scoped CNNA.PillarD} & D & n/a & Kurzalias nur für Hadamard-Strang & \emph{HadSeed} & active extension & Form exportiert; siehe Register I für Primär-/Sekundärtrennung \\
\code{MicrolocalSpectrumGate} & text alias & \emph{MicroSpec} & \code{scoped abbrev MicroSpec := MicrolocalSpectrumGate} & \path{CNNA/PillarD/Notation.lean} & \code{CNNA.PillarD} & kontextgebunden; bei Export \code{open scoped CNNA.PillarD} & D & n/a & Kurzalias nur für mikrolokales Gate & \emph{MicroSpec} & active extension & Form exportiert; siehe Register I für Primär-/Sekundärtrennung \\
\code{LocalWickPolynomialsSeed} & text alias & \emph{WickSeed} & \code{scoped abbrev WickSeed := LocalWickPolynomialsSeed} & \path{CNNA/PillarD/Notation.lean} & \code{CNNA.PillarD} & kontextgebunden; bei Export \code{open scoped CNNA.PillarD} & D & n/a & keine Symbolnotation ohne lokalen Kovarianzkontext & \emph{WickSeed} & active extension & Form exportiert; siehe Register I für Primär-/Sekundärtrennung \\
\bottomrule
\end{tabular}
\end{wideauditlandscape}

\begin{wideauditlandscape}
\tablecaptionline{Notations- und Alias-Appendix -- Teil 6}
\begin{tabular}{L{0.07\linewidth}L{0.06\linewidth}L{0.07\linewidth}L{0.14\linewidth}L{0.08\linewidth}L{0.07\linewidth}L{0.08\linewidth}L{0.08\linewidth}L{0.055\linewidth}L{0.07\linewidth}L{0.07\linewidth}L{0.04\linewidth}L{0.06\linewidth}}
\toprule
Lean-Objekt & Emissionsart & Oberflächenform & Lean-Emission / Expansion & Zielmodul & Namespace & Aktivierungsregel & erforderlicher Kontext & Präzedenz / Bindungsstärke & Kollisionsregel & Beispielverwendung & Status & Review-Notiz \\
\midrule
\code{EinsteinLimitSeed} & text alias & \emph{EinsteinSeed} & \code{scoped abbrev EinsteinSeed := EinsteinLimitSeed} & \path{CNNA/PillarD/Notation.lean} & \code{CNNA.PillarD} & kontextgebunden; bei Export \code{open scoped CNNA.PillarD} & D & n/a & Kurzalias für späten GR-Limes & \emph{EinsteinSeed} & planned active & Form exportiert; siehe Register I für Primär-/Sekundärtrennung \\
\code{TimeReversalOnGNS} & text alias & \emph{TRonGNS} & \code{scoped abbrev TRonGNS := TimeReversalOnGNS} & \path{CNNA/AQFT/Notation.lean} & \code{CNNA.AQFT} & kontextgebunden; bei Export \code{open scoped CNNA.AQFT} & AQFT & n/a & \(\Theta\) bleibt Symbolregistereintrag & \emph{TRonGNS} & active extension & Form exportiert; siehe Register I für Primär-/Sekundärtrennung \\
\code{SpinRepresentationSeed} & text alias & \emph{SpinSeed} & \code{scoped abbrev SpinSeed := SpinRepresentationSeed} & \path{CNNA/AQFT/Notation.lean} & \code{CNNA.AQFT} & kontextgebunden; bei Export \code{open scoped CNNA.AQFT} & AQFT & n/a & Kurzalias nur für den Spin-Strang & \emph{SpinSeed} & active extension & Form exportiert; siehe Register I für Primär-/Sekundärtrennung \\
\code{StatisticsGate} & text alias & \emph{StatGate} & \code{scoped abbrev StatGate := StatisticsGate} & \path{CNNA/AQFT/Notation.lean} & \code{CNNA.AQFT} & kontextgebunden; bei Export \code{open scoped CNNA.AQFT} & AQFT/Gates & n/a & Kurzalias für Statistikabnahme & \emph{StatGate} & active extension & Form exportiert; siehe Register I für Primär-/Sekundärtrennung \\
\code{SpinStatisticsTheorem} & text alias & \emph{SpinStatThm} & \code{scoped abbrev SpinStatThm := SpinStatisticsTheorem} & \path{CNNA/AQFT/Notation.lean} & \code{CNNA.AQFT} & kontextgebunden; bei Export \code{open scoped CNNA.AQFT} & AQFT & n/a & Kurzalias für Theoremtabellen & \emph{SpinStatThm} & active extension & Form exportiert; siehe Register I für Primär-/Sekundärtrennung \\
\code{PTBridgeSeed} & text alias & \emph{PTSeed} & \code{scoped abbrev PTSeed := PTBridgeSeed} & \path{CNNA/PillarE/Notation.lean} & \code{CNNA.PillarE} & kontextgebunden; bei Export \code{open scoped CNNA.PillarE} & D--E & n/a & keine Konkurrenz zu Paritäts-/Zeitumkehrsymbolen & \emph{PTSeed} & active extension & Form exportiert; siehe Register I für Primär-/Sekundärtrennung \\
\code{ChargeConjugationSeed} & text alias & \emph{CSeed} & \code{scoped abbrev CSeed := ChargeConjugationSeed} & \path{CNNA/PillarE/Notation.lean} & \code{CNNA.PillarE} & kontextgebunden; bei Export \code{open scoped CNNA.PillarE} & E & n/a & Kurzalias nur für Ladungskonjugations-Seed & \emph{CSeed} & active extension & Form exportiert; siehe Register I für Primär-/Sekundärtrennung \\
\code{CPTGate} & text alias & \emph{CPTCheck} & \code{scoped abbrev CPTCheck := CPTGate} & \path{CNNA/PillarE/Notation.lean} & \code{CNNA.PillarE} & kontextgebunden; bei Export \code{open scoped CNNA.PillarE} & B--D--E & n/a & Kurzalias für pfeilerübergreifendes Gate & \emph{CPTCheck} & active extension & Form exportiert; siehe Register I für Primär-/Sekundärtrennung \\
\bottomrule
\end{tabular}
\end{wideauditlandscape}

\begin{wideauditlandscape}
\tablecaptionline{Notations- und Alias-Appendix -- Teil 7}
\begin{tabular}{L{0.07\linewidth}L{0.06\linewidth}L{0.07\linewidth}L{0.14\linewidth}L{0.08\linewidth}L{0.07\linewidth}L{0.08\linewidth}L{0.08\linewidth}L{0.055\linewidth}L{0.07\linewidth}L{0.07\linewidth}L{0.04\linewidth}L{0.06\linewidth}}
\toprule
Lean-Objekt & Emissionsart & Oberflächenform & Lean-Emission / Expansion & Zielmodul & Namespace & Aktivierungsregel & erforderlicher Kontext & Präzedenz / Bindungsstärke & Kollisionsregel & Beispielverwendung & Status & Review-Notiz \\
\midrule
\code{CPTTheorem} & text alias & \emph{CPTThm} & \code{scoped abbrev CPTThm := CPTTheorem} & \path{CNNA/PillarE/Notation.lean} & \code{CNNA.AQFT} & kontextgebunden; bei Export \code{open scoped CNNA.AQFT} & B--D--E & n/a & Kurzalias für Theoremtabellen & \emph{CPTThm} & active extension & Form exportiert; siehe Register I für Primär-/Sekundärtrennung \\
\code{KramersGate} & text alias & \emph{KramersCheck} & \code{scoped abbrev KramersCheck := KramersGate} & \path{CNNA/PillarE/Notation.lean} & \code{CNNA.PillarE} & kontextgebunden; bei Export \code{open scoped CNNA.PillarE} & AQFT/E & n/a & \(\Theta^2=(-1)^F\) bleibt Symbolregistereintrag & \emph{KramersCheck} & active extension & Form exportiert; siehe Register I für Primär-/Sekundärtrennung \\
\code{KramersTheorem} & text alias & \emph{KramersThm} & \code{scoped abbrev KramersThm := KramersTheorem} & \path{CNNA/PillarE/Notation.lean} & \code{CNNA.AQFT} & kontextgebunden; bei Export \code{open scoped CNNA.AQFT} & AQFT/E & n/a & Kurzalias für Theoremtabellen & \emph{KramersThm} & active extension & Form exportiert; siehe Register I für Primär-/Sekundärtrennung \\
\code{GaugeSectorSeed} & text alias & \emph{GaugeSeed} & \code{scoped abbrev GaugeSeed := GaugeSectorSeed} & \path{CNNA/PillarE/Notation.lean} & \code{CNNA.PillarE} & kontextgebunden; bei Export \code{open scoped CNNA.PillarE} & E & n/a & Kurzalias für emergenten Eichstrang & \emph{GaugeSeed} & active extension & Form exportiert; siehe Register I für Primär-/Sekundärtrennung \\
\code{AnomalyInflowSeed} & text alias & \emph{InflowSeed} & \code{scoped abbrev InflowSeed := AnomalyInflowSeed} & \path{CNNA/PillarE/Notation.lean} & \code{CNNA.PillarE} & kontextgebunden; bei Export \code{open scoped CNNA.PillarE} & E & n/a & Kurzalias für Inflow-Vorzone & \emph{InflowSeed} & active extension & Form exportiert; siehe Register I für Primär-/Sekundärtrennung \\
\code{AnomalyCancellationGate} & text alias & \emph{AnomCancel} & \code{scoped abbrev AnomCancel := AnomalyCancellationGate} & \path{CNNA/PillarE/Notation.lean} & \code{CNNA.PillarE} & kontextgebunden; bei Export \code{open scoped CNNA.PillarE} & E/Gates & n/a & Kurzalias für Anomalieabnahme & \emph{AnomCancel} & active extension & Form exportiert; siehe Register I für Primär-/Sekundärtrennung \\
\code{ChargeQuantizationGate} & text alias & \emph{ChargeQuant} & \code{scoped abbrev ChargeQuant := ChargeQuantizationGate} & \path{CNNA/PillarE/Notation.lean} & \code{CNNA.PillarE} & kontextgebunden; bei Export \code{open scoped CNNA.PillarE} & E/Gates & n/a & Kurzalias für Quantisierungsgate & \emph{ChargeQuant} & active extension & Form exportiert; siehe Register I für Primär-/Sekundärtrennung \\
\code{HiggsSectorSeed} & text alias & \emph{HiggsSeed} & \code{scoped abbrev HiggsSeed := HiggsSectorSeed} & \path{CNNA/PillarE/Notation.lean} & \code{CNNA.PillarE} & kontextgebunden; bei Export \code{open scoped CNNA.PillarE} & E & n/a & Kurzalias für Higgs-Vorzone & \emph{HiggsSeed} & active extension & Form exportiert; siehe Register I für Primär-/Sekundärtrennung \\
\bottomrule
\end{tabular}
\end{wideauditlandscape}

\begin{wideauditlandscape}
\tablecaptionline{Notations- und Alias-Appendix -- Teil 8}
\begin{tabular}{L{0.07\linewidth}L{0.06\linewidth}L{0.07\linewidth}L{0.14\linewidth}L{0.08\linewidth}L{0.07\linewidth}L{0.08\linewidth}L{0.08\linewidth}L{0.055\linewidth}L{0.07\linewidth}L{0.07\linewidth}L{0.04\linewidth}L{0.06\linewidth}}
\toprule
Lean-Objekt & Emissionsart & Oberflächenform & Lean-Emission / Expansion & Zielmodul & Namespace & Aktivierungsregel & erforderlicher Kontext & Präzedenz / Bindungsstärke & Kollisionsregel & Beispielverwendung & Status & Review-Notiz \\
\midrule
\code{EWSBSeed} & text alias & \emph{EWSBSeed} & \code{scoped abbrev EWSBSeed := EWSBSeed} & \path{CNNA/PillarE/Notation.lean} & \code{CNNA.PillarE} & kontextgebunden; bei Export \code{open scoped CNNA.PillarE} & E & n/a & gleicher Kurzname wie Lean-Name, aber als gescopter Alias erfasst & \emph{EWSBSeed} & active extension & Form exportiert; siehe Register I für Primär-/Sekundärtrennung \\
\code{YukawaFlavorSeed} & text alias & \emph{YukawaSeed} & \code{scoped abbrev YukawaSeed := YukawaFlavorSeed} & \path{CNNA/PillarE/Notation.lean} & \code{CNNA.PillarE} & kontextgebunden; bei Export \code{open scoped CNNA.PillarE} & E & n/a & Kurzalias für Flavorstrang & \emph{YukawaSeed} & active extension & Form exportiert; siehe Register I für Primär-/Sekundärtrennung \\
\code{GenerationSeed} & text alias & \emph{GenSeed} & \code{scoped abbrev GenSeed := GenerationSeed} & \path{CNNA/PillarE/Notation.lean} & \code{CNNA.PillarE} & kontextgebunden; bei Export \code{open scoped CNNA.PillarE} & E & n/a & Kurzalias für Generationsstrang & \emph{GenSeed} & active extension & Form exportiert; siehe Register I für Primär-/Sekundärtrennung \\
\code{SMMatchingWindow} & text alias & \emph{SMMatch} & \code{scoped abbrev SMMatch := SMMatchingWindow} & \path{CNNA/PillarE/Notation.lean} & \code{CNNA.PillarE} & kontextgebunden; bei Export \code{open scoped CNNA.PillarE} & E & n/a & Window-Status bleibt erhalten & \emph{SMMatch} & late window & Form exportiert; siehe Register I für Primär-/Sekundärtrennung \\
\code{SMEFTWindow} & text alias & \emph{SMEFTWin} & \code{scoped abbrev SMEFTWin := SMEFTWindow} & \path{CNNA/PillarE/Notation.lean} & \code{CNNA.PillarE} & kontextgebunden; bei Export \code{open scoped CNNA.PillarE} & E & n/a & Window-Status bleibt erhalten & \emph{SMEFTWin} & late window & Form exportiert; siehe Register I für Primär-/Sekundärtrennung \\
\code{DiracOperatorSeed} & text alias & \emph{DiracSeed} & \code{scoped abbrev DiracSeed := DiracOperatorSeed} & \path{CNNA/PillarD/Notation.lean} & \code{CNNA.PillarD} & kontextgebunden; bei Export \code{open scoped CNNA.PillarD} & D & n/a & Kurzalias für die NCG-Route & \emph{DiracSeed} & active extension & Form exportiert; siehe Register I für Primär-/Sekundärtrennung \\
\bottomrule
\end{tabular}
\end{wideauditlandscape}
\newpage

\section{Variablen-, Namens- und Notationsregister}
\doclevel{Meta-Audit-Ebene}
\label{app:register}
Dieses Register ist ab jetzt explizit in vier getrennte Tabellen aufgespalten. Damit werden Objekt- und
Aliasformen, Symbole und Parameter, Reifegradmarker sowie Meta-Methodik sauber getrennt. Kein
Registerteil darf neue Fachobjekte erzeugen; jede nicht-meta Form muss auf einen kanonischen Eintrag des
Definitionskatalogs zeigen.

\subsection{Register I: Objekt- und Aliasregister}
\doclevel{Meta-Audit-Ebene}
Dieses Register führt alle stabilen sichtbaren Oberflächenformen, die über den Lean-Namen hinaus
verwendet werden. Jede Form zeigt auf genau \emph{einen} kanonischen Lean-Namen. Wird im
Notations-Appendix eine von der Primärform abweichende Kurzform exportiert, so erscheinen
Primärform und Exportform hier jeweils als getrennte Registereinträge.

Weicht die Exportform von der Primärform ab, so ist im Register I
die Primärform mit Formrang \emph{primär} und die exportierte Kurzform
mit Formrang \emph{sekundär exportiert} zu führen. Der Notations-Appendix
enthält in diesem Fall ausschliesslich die exportierte Form; die Primärform
bleibt dokumentarische Leitform, sofern sie nicht selbst emittiert wird.

\begin{wideauditlandscape}
\tablecaptionline{Register I: Objekt- und Aliasregister -- Teil 1}
\begin{tabular}{L{0.10\linewidth}L{0.06\linewidth}L{0.06\linewidth}L{0.10\linewidth}L{0.07\linewidth}L{0.08\linewidth}L{0.08\linewidth}L{0.06\linewidth}L{0.12\linewidth}L{0.14\linewidth}}
\toprule
Oberflächenform & Formtyp & Formrang & kanonischer Lean-Name & Namespace & Aktivierungsregel & Scope & Exportstatus & bindende Stelle & Expansion / Verweisregel \\
\midrule
\emph{Tree of Cliques} & text alias & primär & \code{TreeOfCliques} & \code{CNNA} & keine Exportaktivierung; Registerform & global & register only & Definitionskatalog / CNNA-Core & identischer Verweis auf \code{TreeOfCliques} \\
\emph{ToC} & text alias & sekundär dokumentarisch & \code{TreeOfCliques} & \code{CNNA} & keine Exportaktivierung; Registerform & global & register only & Register-Alias des Stammobjekts & Kurzalias nur für \code{TreeOfCliques} \\
\emph{Bright-Patch} & text alias & primär & \code{BranchPatch} & \code{CNNA.PillarA} & keine Exportaktivierung; Registerform & Pillar A & register only & Definitionskatalog / A-Core & expandiert auf \code{BranchPatch} \\
\(\mathcal C\) & symbol & primär & \code{ComplementSectorFamily} & \code{CNNA.PillarA} & keine Exportaktivierung; Registerform & Pillar A--B & register only & Definitionskatalog / Complement & kontextgebundenes Symbol für \code{ComplementSectorFamily} \\
\emph{Komplementsektorfamilie} & text alias & sekundär dokumentarisch & \code{ComplementSectorFamily} & \code{CNNA.PillarA} & keine Exportaktivierung; Registerform & Pillar A--B & register only & Definitionskatalog / Complement & identischer Verweis auf \code{ComplementSectorFamily} \\
\(\Sigma_{\mathrm{BDI}}\) & symbol & primär & \code{SectorSplit} & \code{CNNA.PillarA} & keine Exportaktivierung; Registerform & Pillar A--B & register only & Definitionskatalog / SectorSplit & kontextgebundenes Symbol für \code{SectorSplit} \\
\emph{Bright/Dark/Interface-Split} & text alias & sekundär dokumentarisch & \code{SectorSplit} & \code{CNNA.PillarA} & keine Exportaktivierung; Registerform & Pillar A--B & register only & Definitionskatalog / SectorSplit & identischer Verweis auf \code{SectorSplit} \\
\(\Lambda_{\mathrm{gen}}\) & symbol & primär & \code{GeneralizedDtN} & \code{CNNA.PillarA} & keine Exportaktivierung; Registerform & Pillar A--B & register only & Definitionskatalog / OQS & kontextgebundenes Symbol für \code{GeneralizedDtN} \\
\emph{generalisierter DtN-Kern} & text alias & sekundär dokumentarisch & \code{GeneralizedDtN} & \code{CNNA.PillarA} & keine Exportaktivierung; Registerform & Pillar A--B & register only & Definitionskatalog / OQS & identischer Verweis auf \code{GeneralizedDtN} \\
\emph{Sektorkanäle} & text alias & primär & \code{SectorChannels} & \code{CNNA.OQS} & keine Exportaktivierung; Registerform & Pillar C & register only & Definitionskatalog / primäre Form & identischer Verweis auf \code{SectorChannels} \\
\emph{SectCh} & export alias & sekundär exportiert & \code{SectorChannels} & \code{CNNA.OQS} & \code{open scoped CNNA.OQS} & Pillar C & exported & Notations-Appendix / \path{CNNA/OQS/Notation.lean} & \code{scoped abbrev SectCh := SectorChannels} \\
\emph{Sektor-SysEnv-Split} & text alias & primär & \code{SectorSysEnv} & \code{CNNA.OQS} & keine Exportaktivierung; Registerform & Pillar C & register only & Definitionskatalog / primäre Form & identischer Verweis auf \code{SectorSysEnv} \\
\emph{SysEnvSplit} & export alias & sekundär exportiert & \code{SectorSysEnv} & \code{CNNA.OQS} & \code{open scoped CNNA.OQS} & Pillar C & exported & Notations-Appendix / \path{CNNA/OQS/Notation.lean} & \code{scoped abbrev SysEnvSplit := SectorSysEnv} \\
\emph{relativer Entropiefluss} & text alias & primär & \code{RelativeEntropyFlow} & \code{CNNA.OQS} & keine Exportaktivierung; Registerform & Pillar C--D & register only & Definitionskatalog / primäre Form & identischer Verweis auf \code{RelativeEntropyFlow} \\
\bottomrule
\end{tabular}
\end{wideauditlandscape}

\begin{wideauditlandscape}
\tablecaptionline{Register I: Objekt- und Aliasregister -- Teil 2}
\begin{tabular}{L{0.10\linewidth}L{0.06\linewidth}L{0.06\linewidth}L{0.10\linewidth}L{0.07\linewidth}L{0.08\linewidth}L{0.08\linewidth}L{0.06\linewidth}L{0.12\linewidth}L{0.14\linewidth}}
\toprule
Oberflächenform & Formtyp & Formrang & kanonischer Lean-Name & Namespace & Aktivierungsregel & Scope & Exportstatus & bindende Stelle & Expansion / Verweisregel \\
\midrule
\emph{RelEntFlow} & export alias & sekundär exportiert & \code{RelativeEntropyFlow} & \code{CNNA.OQS} & kontextgebunden; bei Export \code{open scoped CNNA.OQS} & Pillar C--D & exported & Notations-Appendix / \path{CNNA/OQS/Notation.lean} & \code{scoped abbrev RelEntFlow := RelativeEntropyFlow} \\
\emph{Komplementschliessung} & text alias & primär & \code{ComplementClosure} & \code{CNNA.AQFT} & keine Exportaktivierung; Registerform & AQFT/C & register only & Definitionskatalog / Closure & identischer Verweis auf \code{ComplementClosure} \\
\emph{CompClosure} & export alias & sekundär exportiert & \code{ComplementClosure} & \code{CNNA.AQFT} & kontextgebunden; bei Export \code{open scoped CNNA.AQFT} & AQFT/C & exported & Notations-Appendix / \path{CNNA/AQFT/Notation.lean} & \code{scoped abbrev CompClosure := ComplementClosure} \\
\emph{Parameter-Closure} & text alias & primär & \code{ParameterClosure} & \code{CNNA} & keine Exportaktivierung; Registerform & global & register only & Definitionskatalog / Closure & identischer Verweis auf \code{ParameterClosure} \\
\emph{ParamClosure} & export alias & sekundär exportiert & \code{ParameterClosure} & \code{CNNA} & \code{open scoped CNNA} & global & exported & Notations-Appendix / \path{CNNA/Notation.lean} & \code{scoped abbrev ParamClosure := ParameterClosure} \\
\(\lambda_{\mathrm{eff}}\) & symbol & primär & \code{EffectiveLambda} & \code{CNNA.PillarD} & keine Exportaktivierung; Registerform & C--D & register only & Definitionskatalog / Closure & kontextgebundenes Symbol für \code{EffectiveLambda} \\
\emph{Regularisierungs-Closure} & text alias & primär & \code{RegularizationClosure} & \code{CNNA.PillarD} & keine Exportaktivierung; Registerform & C--D & register only & Definitionskatalog / Closure & identischer Verweis auf \code{RegularizationClosure} \\
\emph{RegClosure} & export alias & sekundär exportiert & \code{RegularizationClosure} & \code{CNNA.PillarD} & kontextgebunden; bei Export \code{open scoped CNNA.PillarD} & C--D & exported & Notations-Appendix / \path{CNNA/PillarD/Notation.lean} & \code{scoped abbrev RegClosure := RegularizationClosure} \\
\emph{Branching-Witness} & text alias & primär & \code{BranchingWitness} & \code{CNNA.PillarA} & keine Exportaktivierung; Registerform & A--D & register only & Definitionskatalog / Closure & identischer Verweis auf \code{BranchingWitness} \\
\emph{BranchWit} & export alias & sekundär exportiert & \code{BranchingWitness} & \code{CNNA.PillarA} & kontextgebunden; bei Export \code{open scoped CNNA.PillarA} & A--D & exported & Notations-Appendix / \path{CNNA/PillarA/Notation.lean} & \code{scoped abbrev BranchWit := BranchingWitness} \\
\emph{selected branching} & text alias & primär & \code{selectedBranching} & \code{CNNA.PillarA} & keine Exportaktivierung; Registerform & A--D & register only & Definitionskatalog / Closure & identischer Verweis auf \code{selectedBranching} \\
\emph{SelBranch} & export alias & sekundär exportiert & \code{selectedBranching} & \code{CNNA.PillarA} & kontextgebunden; bei Export \code{open scoped CNNA.PillarA} & A--D & exported & Notations-Appendix / \path{CNNA/PillarA/Notation.lean} & \code{scoped abbrev SelBranch := selectedBranching} \\
\emph{UV-Spektralselektor} & text alias & primär & \code{UVSpectralSelector} & \code{CNNA.PillarA} & keine Exportaktivierung; Registerform & A--D & register only & Definitionskatalog / Selektoren & identischer Verweis auf \code{UVSpectralSelector} \\
\emph{UVSel} & export alias & sekundär exportiert & \code{UVSpectralSelector} & \code{CNNA.PillarA} & kontextgebunden; bei Export \code{open scoped CNNA.PillarA} & A--D & exported & Notations-Appendix / \path{CNNA/PillarA/Notation.lean} & \code{scoped abbrev UVSel := UVSpectralSelector} \\
\bottomrule
\end{tabular}
\end{wideauditlandscape}

\begin{wideauditlandscape}
\tablecaptionline{Register I: Objekt- und Aliasregister -- Teil 3}
\begin{tabular}{L{0.10\linewidth}L{0.06\linewidth}L{0.06\linewidth}L{0.10\linewidth}L{0.07\linewidth}L{0.08\linewidth}L{0.08\linewidth}L{0.06\linewidth}L{0.12\linewidth}L{0.14\linewidth}}
\toprule
Oberflächenform & Formtyp & Formrang & kanonischer Lean-Name & Namespace & Aktivierungsregel & Scope & Exportstatus & bindende Stelle & Expansion / Verweisregel \\
\midrule
\emph{Branching-Selektor} & text alias & primär & \code{BranchingSelector} & \code{CNNA.PillarA} & keine Exportaktivierung; Registerform & A--D & register only & Definitionskatalog / Selektoren & identischer Verweis auf \code{BranchingSelector} \\
\emph{BranchSel} & export alias & sekundär exportiert & \code{BranchingSelector} & \code{CNNA.PillarA} & kontextgebunden; bei Export \code{open scoped CNNA.PillarA} & A--D & exported & Notations-Appendix / \path{CNNA/PillarA/Notation.lean} & \code{scoped abbrev BranchSel := BranchingSelector} \\
\emph{Backreaction-Selektor} & text alias & primär & \code{BackreactionSelector} & \code{CNNA.PillarD} & keine Exportaktivierung; Registerform & C--D & register only & Definitionskatalog / Selektoren & identischer Verweis auf \code{BackreactionSelector} \\
\emph{BackSel} & export alias & sekundär exportiert & \code{BackreactionSelector} & \code{CNNA.PillarD} & kontextgebunden; bei Export \code{open scoped CNNA.PillarD} & C--D & exported & Notations-Appendix / \path{CNNA/PillarD/Notation.lean} & \code{scoped abbrev BackSel := BackreactionSelector} \\
\(L_H\) & symbol & primär & \code{HorizonLevel} & \code{CNNA.PillarD} & keine Exportaktivierung; Registerform & C--D & register only & Definitionskatalog / Pillar D & kontextgebundenes Symbol für \code{HorizonLevel} \\
\emph{Backreaction-Fixpunkt} & text alias & primär & \code{BackreactionFixedPoint} & \code{CNNA.PillarD} & keine Exportaktivierung; Registerform & C--D & register only & Definitionskatalog / Pillar D & identischer Verweis auf \code{BackreactionFixedPoint} \\
\emph{BackFix} & export alias & sekundär exportiert & \code{BackreactionFixedPoint} & \code{CNNA.PillarD} & kontextgebunden; bei Export \code{open scoped CNNA.PillarD} & C--D & exported & Notations-Appendix / \path{CNNA/PillarD/Notation.lean} & \code{scoped abbrev BackFix := BackreactionFixedPoint} \\
\(\mathcal A_{\mathrm{comp}}\) & symbol & primär & \code{ComplementLocalNet} & \code{CNNA.AQFT} & keine Exportaktivierung; Registerform & AQFT & register only & Definitionskatalog / AQFT & kontextgebundenes Symbol für \code{ComplementLocalNet} \\
\emph{Komplementnetz} & text alias & sekundär dokumentarisch & \code{ComplementLocalNet} & \code{CNNA.AQFT} & keine Exportaktivierung; Registerform & AQFT & register only & Definitionskatalog / AQFT & identischer Verweis auf \code{ComplementLocalNet} \\
\(\mathcal A_{\mathrm{int}}\) & symbol & primär & \code{InterfaceLocalNet} & \code{CNNA.AQFT} & keine Exportaktivierung; Registerform & AQFT & register only & Definitionskatalog / AQFT & kontextgebundenes Symbol für \code{InterfaceLocalNet} \\
\emph{Interfacenetz} & text alias & sekundär dokumentarisch & \code{InterfaceLocalNet} & \code{CNNA.AQFT} & keine Exportaktivierung; Registerform & AQFT & register only & Definitionskatalog / AQFT & identischer Verweis auf \code{InterfaceLocalNet} \\
\(\omega_{\mathrm{comp}}\) & symbol & primär & \code{ComplementStateNet} & \code{CNNA.AQFT} & keine Exportaktivierung; Registerform & AQFT & register only & Definitionskatalog / AQFT & kontextgebundenes Symbol für \code{ComplementStateNet} \\
\emph{Komplementzustandsnetz} & text alias & sekundär dokumentarisch & \code{ComplementStateNet} & \code{CNNA.AQFT} & keine Exportaktivierung; Registerform & AQFT & register only & Definitionskatalog / AQFT & identischer Verweis auf \code{ComplementStateNet} \\
\emph{Komplement-GNS-Datum} & text alias & primär & \code{ComplementGNS} & \code{CNNA.AQFT} & keine Exportaktivierung; Registerform & AQFT & register only & Definitionskatalog / AQFT & identischer Verweis auf \code{ComplementGNS} \\
\bottomrule
\end{tabular}
\end{wideauditlandscape}

\begin{wideauditlandscape}
\tablecaptionline{Register I: Objekt- und Aliasregister -- Teil 4}
\begin{tabular}{L{0.10\linewidth}L{0.06\linewidth}L{0.06\linewidth}L{0.10\linewidth}L{0.07\linewidth}L{0.08\linewidth}L{0.08\linewidth}L{0.06\linewidth}L{0.12\linewidth}L{0.14\linewidth}}
\toprule
Oberflächenform & Formtyp & Formrang & kanonischer Lean-Name & Namespace & Aktivierungsregel & Scope & Exportstatus & bindende Stelle & Expansion / Verweisregel \\
\midrule
\emph{CompGNS} & export alias & sekundär exportiert & \code{ComplementGNS} & \code{CNNA.AQFT} & kontextgebunden; bei Export \code{open scoped CNNA.AQFT} & AQFT & exported & Notations-Appendix / \path{CNNA/AQFT/Notation.lean} & \code{scoped abbrev CompGNS := ComplementGNS} \\
\((\pi_\omega,\mathcal H_\omega,\Omega_\omega)\) & composite tuple form & sekundär dokumentarisch & \code{ComplementGNS} & \code{CNNA.AQFT} & keine Exportaktivierung; Registerform & AQFT & register only & GNS-Kontext & nur für \code{ComplementGNS}; kein zweites Objekt \\
\emph{zyklisch-separierende Lage} & text alias & primär & \code{SeparatingProperty} & \code{CNNA.AQFT} & keine Exportaktivierung; Registerform & AQFT & register only & Definitionskatalog / AQFT & identischer Verweis auf \code{SeparatingProperty} \\
\emph{CycSep} & export alias & sekundär exportiert & \code{SeparatingProperty} & \code{CNNA.AQFT} & kontextgebunden; bei Export \code{open scoped CNNA.AQFT} & AQFT & exported & Notations-Appendix / \path{CNNA/AQFT/Notation.lean} & \code{scoped abbrev CycSep := SeparatingProperty} \\
\emph{Tomita--Takesaki-Datum} & text alias & primär & \code{TomitaTakesakiData} & \code{CNNA.AQFT} & keine Exportaktivierung; Registerform & AQFT & register only & Definitionskatalog / AQFT & identischer Verweis auf \code{TomitaTakesakiData} \\
\emph{TTData} & export alias & sekundär exportiert & \code{TomitaTakesakiData} & \code{CNNA.AQFT} & kontextgebunden; bei Export \code{open scoped CNNA.AQFT} & AQFT & exported & Notations-Appendix / \path{CNNA/AQFT/Notation.lean} & \code{scoped abbrev TTData := TomitaTakesakiData} \\
\((S,J,\Delta)\) & composite tuple form & sekundär dokumentarisch & \code{TomitaTakesakiData} & \code{CNNA.AQFT} & keine Exportaktivierung; Registerform & AQFT & register only & modularer Kontext & nur für \code{TomitaTakesakiData}; kein zweites Objekt \\
\emph{Komplement-KMS-Zustand} & text alias & primär & \code{ComplementKMS} & \code{CNNA.AQFT} & keine Exportaktivierung; Registerform & AQFT & register only & Definitionskatalog / AQFT & identischer Verweis auf \code{ComplementKMS} \\
\emph{CompKMS} & export alias & sekundär exportiert & \code{ComplementKMS} & \code{CNNA.AQFT} & kontextgebunden; bei Export \code{open scoped CNNA.AQFT} & AQFT & exported & Notations-Appendix / \path{CNNA/AQFT/Notation.lean} & \code{scoped abbrev CompKMS := ComplementKMS} \\
\(\omega_\beta\) & symbol & sekundär dokumentarisch & \code{ComplementKMS} & \code{CNNA.AQFT} & keine Exportaktivierung; Registerform & AQFT/KMS & register only & KMS-Kontext & nur im KMS-Kontext für \code{ComplementKMS} \\
\emph{Komplement-Spektralbedingung} & text alias & primär & \code{ComplementSpectrumCondition} & \code{CNNA.AQFT} & keine Exportaktivierung; Registerform & B--D & register only & Definitionskatalog / AQFT & identischer Verweis auf \code{ComplementSpectrumCondition} \\
\emph{CompSpectrum} & export alias & sekundär exportiert & \code{ComplementSpectrumCondition} & \code{CNNA.AQFT} & kontextgebunden; bei Export \code{open scoped CNNA.AQFT} & B--D & exported & Notations-Appendix / \path{CNNA/AQFT/Notation.lean} & \code{scoped abbrev CompSpectrum := ComplementSpectrumCondition} \\
\emph{Modularkonjugation} & text alias & primär & \code{ModularConjugation} & \code{CNNA.AQFT} & keine Exportaktivierung; Registerform & AQFT & register only & Definitionskatalog / primäre Form & identischer Verweis auf \code{ModularConjugation} \\
\emph{ModConj} & export alias & sekundär exportiert & \code{ModularConjugation} & \code{CNNA.AQFT} & kontextgebunden; bei Export \code{open scoped CNNA.AQFT} & AQFT & exported & Notations-Appendix / \path{CNNA/AQFT/Notation.lean} & \code{scoped abbrev ModConj := ModularConjugation} \\
\bottomrule
\end{tabular}
\end{wideauditlandscape}

\begin{wideauditlandscape}
\tablecaptionline{Register I: Objekt- und Aliasregister -- Teil 5}
\begin{tabular}{L{0.10\linewidth}L{0.06\linewidth}L{0.06\linewidth}L{0.10\linewidth}L{0.07\linewidth}L{0.08\linewidth}L{0.08\linewidth}L{0.06\linewidth}L{0.12\linewidth}L{0.14\linewidth}}
\toprule
Oberflächenform & Formtyp & Formrang & kanonischer Lean-Name & Namespace & Aktivierungsregel & Scope & Exportstatus & bindende Stelle & Expansion / Verweisregel \\
\midrule
\emph{BW-Seed} & text alias & primär & \code{BisognanoWichmannSeed} & \code{CNNA.AQFT} & keine Exportaktivierung; Registerform & AQFT & register only & Definitionskatalog / AQFT & identischer Verweis auf \code{BisognanoWichmannSeed} \\
\emph{BWSeed} & export alias & sekundär exportiert & \code{BisognanoWichmannSeed} & \code{CNNA.AQFT} & kontextgebunden; bei Export \code{open scoped CNNA.AQFT} & AQFT & exported & Notations-Appendix / \path{CNNA/AQFT/Notation.lean} & \code{scoped abbrev BWSeed := BisognanoWichmannSeed} \\
\emph{Nuclearity-Seed} & text alias & primär & \code{NuclearitySeed} & \code{CNNA.AQFT} & keine Exportaktivierung; Registerform & AQFT & register only & Definitionskatalog / AQFT & identischer Verweis auf \code{NuclearitySeed} \\
\emph{NucSeed} & export alias & sekundär exportiert & \code{NuclearitySeed} & \code{CNNA.AQFT} & kontextgebunden; bei Export \code{open scoped CNNA.AQFT} & AQFT & exported & Notations-Appendix / \path{CNNA/AQFT/Notation.lean} & \code{scoped abbrev NucSeed := NuclearitySeed} \\
\emph{DHR-Sektoren-Seed} & text alias & primär & \code{DHRSectorsSeed} & \code{CNNA.AQFT} & keine Exportaktivierung; Registerform & AQFT & register only & Definitionskatalog / AQFT & identischer Verweis auf \code{DHRSectorsSeed} \\
\emph{DHRSeed} & export alias & sekundär exportiert & \code{DHRSectorsSeed} & \code{CNNA.AQFT} & kontextgebunden; bei Export \code{open scoped CNNA.AQFT} & AQFT & exported & Notations-Appendix / \path{CNNA/AQFT/Notation.lean} & \code{scoped abbrev DHRSeed := DHRSectorsSeed} \\
\emph{Feldalgebra-Rekonstruktions-Seed} & text alias & primär & \code{FieldAlgebraReconstructionSeed} & \code{CNNA.AQFT} & keine Exportaktivierung; Registerform & AQFT & register only & Definitionskatalog / AQFT & identischer Verweis auf \code{FieldAlgebraReconstructionSeed} \\
\emph{FieldRecSeed} & export alias & sekundär exportiert & \code{FieldAlgebraReconstructionSeed} & \code{CNNA.AQFT} & kontextgebunden; bei Export \code{open scoped CNNA.AQFT} & AQFT & exported & Notations-Appendix / \path{CNNA/AQFT/Notation.lean} & \code{scoped abbrev FieldRecSeed := FieldAlgebraReconstructionSeed} \\
\emph{Ladungs-Statistik-Kohärenz} & text alias & primär & \code{ChargeStatisticsCoherence} & \code{CNNA.AQFT} & keine Exportaktivierung; Registerform & AQFT/Gates & register only & Definitionskatalog / AQFT & identischer Verweis auf \code{ChargeStatisticsCoherence} \\
\emph{ChargeStatCoh} & export alias & sekundär exportiert & \code{ChargeStatisticsCoherence} & \code{CNNA.AQFT} & kontextgebunden; bei Export \code{open scoped CNNA.AQFT} & AQFT/Gates & exported & Notations-Appendix / \path{CNNA/AQFT/Notation.lean} & \code{scoped abbrev ChargeStatCoh := ChargeStatisticsCoherence} \\
\emph{Ladungsrekonstruktionssatz} & text alias & primär & \code{ChargeReconstructionTheorem} & \code{CNNA.AQFT} & keine Exportaktivierung; Registerform & AQFT & register only & Definitionskatalog / AQFT & identischer Verweis auf \code{ChargeReconstructionTheorem} \\
\emph{ChargeRec} & export alias & sekundär exportiert & \code{ChargeReconstructionTheorem} & \code{CNNA.AQFT} & kontextgebunden; bei Export \code{open scoped CNNA.AQFT} & AQFT & exported & Notations-Appendix / \path{CNNA/AQFT/Notation.lean} & \code{scoped abbrev ChargeRec := ChargeReconstructionTheorem} \\
\emph{derived spacetime} & text alias & primär & \code{DerivedSpacetime} & \code{CNNA.PillarD} & keine Exportaktivierung; Registerform & Pillar D & register only & Definitionskatalog / Pillar D & identischer Verweis auf \code{DerivedSpacetime} \\
\emph{DerSpacetime} & export alias & sekundär exportiert & \code{DerivedSpacetime} & \code{CNNA.PillarD} & kontextgebunden; bei Export \code{open scoped CNNA.PillarD} & D & exported & Notations-Appendix / \path{CNNA/PillarD/Notation.lean} & \code{scoped abbrev DerSpacetime := DerivedSpacetime} \\
\bottomrule
\end{tabular}
\end{wideauditlandscape}

\begin{wideauditlandscape}
\tablecaptionline{Register I: Objekt- und Aliasregister -- Teil 6}
\begin{tabular}{L{0.10\linewidth}L{0.06\linewidth}L{0.06\linewidth}L{0.10\linewidth}L{0.07\linewidth}L{0.08\linewidth}L{0.08\linewidth}L{0.06\linewidth}L{0.12\linewidth}L{0.14\linewidth}}
\toprule
Oberflächenform & Formtyp & Formrang & kanonischer Lean-Name & Namespace & Aktivierungsregel & Scope & Exportstatus & bindende Stelle & Expansion / Verweisregel \\
\midrule
\emph{Interface-Kausalität} & text alias & primär & \code{InterfaceCausality} & \code{CNNA.PillarD} & keine Exportaktivierung; Registerform & Pillar D & register only & Definitionskatalog / Pillar D & identischer Verweis auf \code{InterfaceCausality} \\
\emph{IfcCaus} & export alias & sekundär exportiert & \code{InterfaceCausality} & \code{CNNA.PillarD} & kontextgebunden; bei Export \code{open scoped CNNA.PillarD} & D & exported & Notations-Appendix / \path{CNNA/PillarD/Notation.lean} & \code{scoped abbrev IfcCaus := InterfaceCausality} \\
\emph{lokale Kovarianz} & text alias & primär & \code{LocalCovariance} & \code{CNNA.PillarD} & keine Exportaktivierung; Registerform & Pillar D & register only & Definitionskatalog / Pillar D & identischer Verweis auf \code{LocalCovariance} \\
\emph{LocCov} & export alias & sekundär exportiert & \code{LocalCovariance} & \code{CNNA.PillarD} & kontextgebunden; bei Export \code{open scoped CNNA.PillarD} & D & exported & Notations-Appendix / \path{CNNA/PillarD/Notation.lean} & \code{scoped abbrev LocCov := LocalCovariance} \\
\emph{Hadamard-State-Seed} & text alias & primär & \code{HadamardStateSeed} & \code{CNNA.PillarD} & keine Exportaktivierung; Registerform & Pillar D & register only & Definitionskatalog / Pillar D & identischer Verweis auf \code{HadamardStateSeed} \\
\emph{HadSeed} & export alias & sekundär exportiert & \code{HadamardStateSeed} & \code{CNNA.PillarD} & kontextgebunden; bei Export \code{open scoped CNNA.PillarD} & D & exported & Notations-Appendix / \path{CNNA/PillarD/Notation.lean} & \code{scoped abbrev HadSeed := HadamardStateSeed} \\
\emph{mikrolokales Spektralgate} & text alias & primär & \code{MicrolocalSpectrumGate} & \code{CNNA.PillarD} & keine Exportaktivierung; Registerform & Pillar D & register only & Definitionskatalog / Pillar D & identischer Verweis auf \code{MicrolocalSpectrumGate} \\
\emph{MicroSpec} & export alias & sekundär exportiert & \code{MicrolocalSpectrumGate} & \code{CNNA.PillarD} & kontextgebunden; bei Export \code{open scoped CNNA.PillarD} & D & exported & Notations-Appendix / \path{CNNA/PillarD/Notation.lean} & \code{scoped abbrev MicroSpec := MicrolocalSpectrumGate} \\
\emph{lokale-Wick-Polynome-Seed} & text alias & primär & \code{LocalWickPolynomialsSeed} & \code{CNNA.PillarD} & keine Exportaktivierung; Registerform & Pillar D & register only & Definitionskatalog / Pillar D & identischer Verweis auf \code{LocalWickPolynomialsSeed} \\
\emph{WickSeed} & export alias & sekundär exportiert & \code{LocalWickPolynomialsSeed} & \code{CNNA.PillarD} & kontextgebunden; bei Export \code{open scoped CNNA.PillarD} & D & exported & Notations-Appendix / \path{CNNA/PillarD/Notation.lean} & \code{scoped abbrev WickSeed := LocalWickPolynomialsSeed} \\
\emph{Einstein-Limes-Seed} & text alias & primär & \code{EinsteinLimitSeed} & \code{CNNA.PillarD} & keine Exportaktivierung; Registerform & Pillar D & register only & Definitionskatalog / Pillar D & identischer Verweis auf \code{EinsteinLimitSeed} \\
\emph{EinsteinSeed} & export alias & sekundär exportiert & \code{EinsteinLimitSeed} & \code{CNNA.PillarD} & kontextgebunden; bei Export \code{open scoped CNNA.PillarD} & D & exported & Notations-Appendix / \path{CNNA/PillarD/Notation.lean} & \code{scoped abbrev EinsteinSeed := EinsteinLimitSeed} \\
\emph{Zeitumkehr auf GNS} & text alias & primär & \code{TimeReversalOnGNS} & \code{CNNA.AQFT} & keine Exportaktivierung; Registerform & AQFT & register only & Definitionskatalog / AQFT & identischer Verweis auf \code{TimeReversalOnGNS} \\
\emph{TRonGNS} & export alias & sekundär exportiert & \code{TimeReversalOnGNS} & \code{CNNA.AQFT} & kontextgebunden; bei Export \code{open scoped CNNA.AQFT} & AQFT & exported & Notations-Appendix / \path{CNNA/AQFT/Notation.lean} & \code{scoped abbrev TRonGNS := TimeReversalOnGNS} \\
\bottomrule
\end{tabular}
\end{wideauditlandscape}

\begin{wideauditlandscape}
\tablecaptionline{Register I: Objekt- und Aliasregister -- Teil 7}
\begin{tabular}{L{0.10\linewidth}L{0.06\linewidth}L{0.06\linewidth}L{0.10\linewidth}L{0.07\linewidth}L{0.08\linewidth}L{0.08\linewidth}L{0.06\linewidth}L{0.12\linewidth}L{0.14\linewidth}}
\toprule
Oberflächenform & Formtyp & Formrang & kanonischer Lean-Name & Namespace & Aktivierungsregel & Scope & Exportstatus & bindende Stelle & Expansion / Verweisregel \\
\midrule
\emph{Spinrepräsentations-Seed} & text alias & primär & \code{SpinRepresentationSeed} & \code{CNNA.AQFT} & keine Exportaktivierung; Registerform & AQFT & register only & Definitionskatalog / AQFT & identischer Verweis auf \code{SpinRepresentationSeed} \\
\emph{SpinSeed} & export alias & sekundär exportiert & \code{SpinRepresentationSeed} & \code{CNNA.AQFT} & kontextgebunden; bei Export \code{open scoped CNNA.AQFT} & AQFT & exported & Notations-Appendix / \path{CNNA/AQFT/Notation.lean} & \code{scoped abbrev SpinSeed := SpinRepresentationSeed} \\
\emph{Statistik-Gate} & text alias & primär & \code{StatisticsGate} & \code{CNNA.AQFT} & keine Exportaktivierung; Registerform & AQFT/Gates & register only & Definitionskatalog / AQFT & identischer Verweis auf \code{StatisticsGate} \\
\emph{StatGate} & export alias & sekundär exportiert & \code{StatisticsGate} & \code{CNNA.AQFT} & kontextgebunden; bei Export \code{open scoped CNNA.AQFT} & AQFT/Gates & exported & Notations-Appendix / \path{CNNA/AQFT/Notation.lean} & \code{scoped abbrev StatGate := StatisticsGate} \\
\emph{Spin-Statistik-Theorem} & text alias & primär & \code{SpinStatisticsTheorem} & \code{CNNA.AQFT} & keine Exportaktivierung; Registerform & AQFT & register only & Definitionskatalog / AQFT & identischer Verweis auf \code{SpinStatisticsTheorem} \\
\emph{SpinStatThm} & export alias & sekundär exportiert & \code{SpinStatisticsTheorem} & \code{CNNA.AQFT} & kontextgebunden; bei Export \code{open scoped CNNA.AQFT} & AQFT & exported & Notations-Appendix / \path{CNNA/AQFT/Notation.lean} & \code{scoped abbrev SpinStatThm := SpinStatisticsTheorem} \\
\emph{PT-Bridge-Seed} & text alias & primär & \code{PTBridgeSeed} & \code{CNNA.PillarE} & keine Exportaktivierung; Registerform & D--E & register only & Definitionskatalog / Pillar E & identischer Verweis auf \code{PTBridgeSeed} \\
\emph{PTSeed} & export alias & sekundär exportiert & \code{PTBridgeSeed} & \code{CNNA.PillarE} & kontextgebunden; bei Export \code{open scoped CNNA.PillarE} & D--E & exported & Notations-Appendix / \path{CNNA/PillarE/Notation.lean} & \code{scoped abbrev PTSeed := PTBridgeSeed} \\
\emph{Ladungskonjugations-Seed} & text alias & primär & \code{ChargeConjugationSeed} & \code{CNNA.PillarE} & keine Exportaktivierung; Registerform & E & register only & Definitionskatalog / Pillar E & identischer Verweis auf \code{ChargeConjugationSeed} \\
\emph{CSeed} & export alias & sekundär exportiert & \code{ChargeConjugationSeed} & \code{CNNA.PillarE} & kontextgebunden; bei Export \code{open scoped CNNA.PillarE} & E & exported & Notations-Appendix / \path{CNNA/PillarE/Notation.lean} & \code{scoped abbrev CSeed := ChargeConjugationSeed} \\
\emph{CPT-Gate} & text alias & primär & \code{CPTGate} & \code{CNNA.PillarE} & keine Exportaktivierung; Registerform & B--D--E & register only & Definitionskatalog / Gates & identischer Verweis auf \code{CPTGate} \\
\emph{CPTCheck} & export alias & sekundär exportiert & \code{CPTGate} & \code{CNNA.PillarE} & kontextgebunden; bei Export \code{open scoped CNNA.PillarE} & B--D--E & exported & Notations-Appendix / \path{CNNA/PillarE/Notation.lean} & \code{scoped abbrev CPTCheck := CPTGate} \\
\emph{CPT-Theorem} & text alias & primär & \code{CPTTheorem} & \code{CNNA.AQFT} & keine Exportaktivierung; Registerform & B--D--E & register only & Definitionskatalog / AQFT & identischer Verweis auf \code{CPTTheorem} \\
\emph{CPTThm} & export alias & sekundär exportiert & \code{CPTTheorem} & \code{CNNA.AQFT} & kontextgebunden; bei Export \code{open scoped CNNA.AQFT} & B--D--E & exported & Notations-Appendix / \path{CNNA/PillarE/Notation.lean} & \code{scoped abbrev CPTThm := CPTTheorem} \\
\bottomrule
\end{tabular}
\end{wideauditlandscape}

\begin{wideauditlandscape}
\tablecaptionline{Register I: Objekt- und Aliasregister -- Teil 8}
\begin{tabular}{L{0.10\linewidth}L{0.06\linewidth}L{0.06\linewidth}L{0.10\linewidth}L{0.07\linewidth}L{0.08\linewidth}L{0.08\linewidth}L{0.06\linewidth}L{0.12\linewidth}L{0.14\linewidth}}
\toprule
Oberflächenform & Formtyp & Formrang & kanonischer Lean-Name & Namespace & Aktivierungsregel & Scope & Exportstatus & bindende Stelle & Expansion / Verweisregel \\
\midrule
\emph{Kramers-Gate} & text alias & primär & \code{KramersGate} & \code{CNNA.PillarE} & keine Exportaktivierung; Registerform & AQFT/E & register only & Definitionskatalog / Gates & identischer Verweis auf \code{KramersGate} \\
\emph{KramersCheck} & export alias & sekundär exportiert & \code{KramersGate} & \code{CNNA.PillarE} & kontextgebunden; bei Export \code{open scoped CNNA.PillarE} & AQFT/E & exported & Notations-Appendix / \path{CNNA/PillarE/Notation.lean} & \code{scoped abbrev KramersCheck := KramersGate} \\
\emph{Kramers-Theorem} & text alias & primär & \code{KramersTheorem} & \code{CNNA.AQFT} & keine Exportaktivierung; Registerform & AQFT/E & register only & Definitionskatalog / AQFT & identischer Verweis auf \code{KramersTheorem} \\
\emph{KramersThm} & export alias & sekundär exportiert & \code{KramersTheorem} & \code{CNNA.AQFT} & kontextgebunden; bei Export \code{open scoped CNNA.AQFT} & AQFT/E & exported & Notations-Appendix / \path{CNNA/PillarE/Notation.lean} & \code{scoped abbrev KramersThm := KramersTheorem} \\
\emph{Gauge-Sektor-Seed} & text alias & primär & \code{GaugeSectorSeed} & \code{CNNA.PillarE} & keine Exportaktivierung; Registerform & E & register only & Definitionskatalog / Pillar E & identischer Verweis auf \code{GaugeSectorSeed} \\
\emph{GaugeSeed} & export alias & sekundär exportiert & \code{GaugeSectorSeed} & \code{CNNA.PillarE} & kontextgebunden; bei Export \code{open scoped CNNA.PillarE} & E & exported & Notations-Appendix / \path{CNNA/PillarE/Notation.lean} & \code{scoped abbrev GaugeSeed := GaugeSectorSeed} \\
\emph{Anomaly-Inflow-Seed} & text alias & primär & \code{AnomalyInflowSeed} & \code{CNNA.PillarE} & keine Exportaktivierung; Registerform & E & register only & Definitionskatalog / Pillar E & identischer Verweis auf \code{AnomalyInflowSeed} \\
\emph{InflowSeed} & export alias & sekundär exportiert & \code{AnomalyInflowSeed} & \code{CNNA.PillarE} & kontextgebunden; bei Export \code{open scoped CNNA.PillarE} & E & exported & Notations-Appendix / \path{CNNA/PillarE/Notation.lean} & \code{scoped abbrev InflowSeed := AnomalyInflowSeed} \\
\emph{Anomalie-Kompensations-Gate} & text alias & primär & \code{AnomalyCancellationGate} & \code{CNNA.PillarE} & keine Exportaktivierung; Registerform & E/Gates & register only & Definitionskatalog / Gates & identischer Verweis auf \code{AnomalyCancellationGate} \\
\emph{AnomCancel} & export alias & sekundär exportiert & \code{AnomalyCancellationGate} & \code{CNNA.PillarE} & kontextgebunden; bei Export \code{open scoped CNNA.PillarE} & E/Gates & exported & Notations-Appendix / \path{CNNA/PillarE/Notation.lean} & \code{scoped abbrev AnomCancel := AnomalyCancellationGate} \\
\emph{Ladungsquantisierungs-Gate} & text alias & primär & \code{ChargeQuantizationGate} & \code{CNNA.PillarE} & keine Exportaktivierung; Registerform & E/Gates & register only & Definitionskatalog / Gates & identischer Verweis auf \code{ChargeQuantizationGate} \\
\emph{ChargeQuant} & export alias & sekundär exportiert & \code{ChargeQuantizationGate} & \code{CNNA.PillarE} & kontextgebunden; bei Export \code{open scoped CNNA.PillarE} & E/Gates & exported & Notations-Appendix / \path{CNNA/PillarE/Notation.lean} & \code{scoped abbrev ChargeQuant := ChargeQuantizationGate} \\
\emph{Higgs-Sektor-Seed} & text alias & primär & \code{HiggsSectorSeed} & \code{CNNA.PillarE} & keine Exportaktivierung; Registerform & E & register only & Definitionskatalog / Pillar E & identischer Verweis auf \code{HiggsSectorSeed} \\
\emph{HiggsSeed} & export alias & sekundär exportiert & \code{HiggsSectorSeed} & \code{CNNA.PillarE} & kontextgebunden; bei Export \code{open scoped CNNA.PillarE} & E & exported & Notations-Appendix / \path{CNNA/PillarE/Notation.lean} & \code{scoped abbrev HiggsSeed := HiggsSectorSeed} \\
\bottomrule
\end{tabular}
\end{wideauditlandscape}

\begin{wideauditlandscape}
\tablecaptionline{Register I: Objekt- und Aliasregister -- Teil 9}
\begin{tabular}{L{0.10\linewidth}L{0.06\linewidth}L{0.06\linewidth}L{0.10\linewidth}L{0.07\linewidth}L{0.08\linewidth}L{0.08\linewidth}L{0.06\linewidth}L{0.12\linewidth}L{0.14\linewidth}}
\toprule
Oberflächenform & Formtyp & Formrang & kanonischer Lean-Name & Namespace & Aktivierungsregel & Scope & Exportstatus & bindende Stelle & Expansion / Verweisregel \\
\midrule
\emph{EWSB-Seed} & text alias & primär & \code{EWSBSeed} & \code{CNNA.PillarE} & keine Exportaktivierung; Registerform & E & register only & Definitionskatalog / Pillar E & identischer Verweis auf \code{EWSBSeed} \\
\emph{EWSBSeed} & export alias & sekundär exportiert & \code{EWSBSeed} & \code{CNNA.PillarE} & kontextgebunden; bei Export \code{open scoped CNNA.PillarE} & E & exported & Notations-Appendix / \path{CNNA/PillarE/Notation.lean} & \code{scoped abbrev EWSBSeed := EWSBSeed} \\
\emph{Yukawa-Flavor-Seed} & text alias & primär & \code{YukawaFlavorSeed} & \code{CNNA.PillarE} & keine Exportaktivierung; Registerform & E & register only & Definitionskatalog / Pillar E & identischer Verweis auf \code{YukawaFlavorSeed} \\
\emph{YukawaSeed} & export alias & sekundär exportiert & \code{YukawaFlavorSeed} & \code{CNNA.PillarE} & kontextgebunden; bei Export \code{open scoped CNNA.PillarE} & E & exported & Notations-Appendix / \path{CNNA/PillarE/Notation.lean} & \code{scoped abbrev YukawaSeed := YukawaFlavorSeed} \\
\emph{Generationen-Seed} & text alias & primär & \code{GenerationSeed} & \code{CNNA.PillarE} & keine Exportaktivierung; Registerform & E & register only & Definitionskatalog / Pillar E & identischer Verweis auf \code{GenerationSeed} \\
\emph{GenSeed} & export alias & sekundär exportiert & \code{GenerationSeed} & \code{CNNA.PillarE} & kontextgebunden; bei Export \code{open scoped CNNA.PillarE} & E & exported & Notations-Appendix / \path{CNNA/PillarE/Notation.lean} & \code{scoped abbrev GenSeed := GenerationSeed} \\
\emph{SM-Matching-Window} & text alias & primär & \code{SMMatchingWindow} & \code{CNNA.PillarE} & keine Exportaktivierung; Registerform & E & register only & Definitionskatalog / Pillar E & identischer Verweis auf \code{SMMatchingWindow} \\
\emph{SMMatch} & export alias & sekundär exportiert & \code{SMMatchingWindow} & \code{CNNA.PillarE} & kontextgebunden; bei Export \code{open scoped CNNA.PillarE} & E & exported & Notations-Appendix / \path{CNNA/PillarE/Notation.lean} & \code{scoped abbrev SMMatch := SMMatchingWindow} \\
\emph{SMEFT-Window} & text alias & primär & \code{SMEFTWindow} & \code{CNNA.PillarE} & keine Exportaktivierung; Registerform & E & register only & Definitionskatalog / Pillar E & identischer Verweis auf \code{SMEFTWindow} \\
\emph{SMEFTWin} & export alias & sekundär exportiert & \code{SMEFTWindow} & \code{CNNA.PillarE} & kontextgebunden; bei Export \code{open scoped CNNA.PillarE} & E & exported & Notations-Appendix / \path{CNNA/PillarE/Notation.lean} & \code{scoped abbrev SMEFTWin := SMEFTWindow} \\
\emph{Dirac-Operator-Seed} & text alias & primär & \code{DiracOperatorSeed} & \code{CNNA.PillarD} & keine Exportaktivierung; Registerform & D & register only & Definitionskatalog / Pillar D & identischer Verweis auf \code{DiracOperatorSeed} \\
\emph{DiracSeed} & export alias & sekundär exportiert & \code{DiracOperatorSeed} & \code{CNNA.PillarD} & kontextgebunden; bei Export \code{open scoped CNNA.PillarD} & D & exported & Notations-Appendix / \path{CNNA/PillarD/Notation.lean} & \code{scoped abbrev DiracSeed := DiracOperatorSeed} \\
\bottomrule
\end{tabular}
\end{wideauditlandscape}

\subsection{Register II: Symbol- und Parameterregister}
\doclevel{Meta-Audit-Ebene}
Dieses Register enthält nur Symbole und Parameter. Besonders wichtig ist die explizite Aufspaltung der
bisher überladenen \(\Omega\)-Verwendung in einen zyklischen Vektor \(\Omega_\omega\) und einen
Carrier-/Indexkontext \(\Omega_{\mathrm{idx}}\).

\begin{wideauditlandscape}
\tablecaptionline{Register II: Symbol- und Parameterregister -- Teil 1}
\begin{tabular}{L{0.07\linewidth}L{0.09\linewidth}L{0.13\linewidth}L{0.11\linewidth}L{0.11\linewidth}L{0.11\linewidth}L{0.11\linewidth}L{0.10\linewidth}L{0.07\linewidth}}
\toprule
Symbol & Symbolklasse & kanonische Bedeutung & zulässiger Kontext & verbotener Kontext & Disambiguierungsregel & Ersatz- / Migrationsregel & bindende Stelle & Status \\
\midrule
$b$ & parameter & Branching-/Mikrostrukturparameter & Closure-/A-Kontext & niemals ohne expliziten Closure-Verweis & nur für Branching-/Mikrostrukturparameter & später durch derived Selektoren oder Closure ersetzt & Parameter-Closure-Strang & später derived oder eliminiert \\
$L_{\max}$ & legacy parameter & alter freier IR-/Horizont-Parameter & nur in Archiv- und Migrationskontexten & im aktiven Kern als Primärname verboten & nur als Legacy-Cutoffsymbol & wird langfristig durch \code{HorizonLevel} / $L_H$ ersetzt & Migration nach \code{HorizonLevel} & Legacyname; aktiv nicht normativ \\
$L_H$ & derived scale & \code{HorizonLevel} & C--D & nicht für \(L_{\max}\) verwenden & nie als freies $L_{\max}$ lesen & ersetzt aktiven Gebrauch von $L_{\max}$ & Pillar D / HorizonLevel & aktive derived Skala \\
$\beta$ & KMS parameter & inverse Temperatur bzw. KMS-Parameter & KMS-/Closure-Kontext & nicht unmarkiert für beliebige andere Parameter & nur als KMS-/Closure-Parameter & nicht unmarkiert auf andere Parameter übertragen & B-Kern und Closure-Diskussion & standardnah, aber closure-pflichtig \\
$\omega_{\mathrm{comp}}$ & state symbol & \code{ComplementStateNet} & AQFT-Zustandskontext & nicht für \code{ComplementKMS} & nur für \code{ComplementStateNet} & kein Ersatz für $\omega_\beta$ & ComplementStateNet & symbolisch erlaubt \\
$\omega_\beta$ & KMS state symbol & \code{ComplementKMS} & nur KMS-Kontext & nicht für das allgemeine Zustandsnetz & nur im KMS-Kontext & nie als Symbol für das allgemeine Zustandsnetz & ComplementKMS & symbolisch erlaubt \\
$\mathcal{C}$ & sector-family symbol & \code{ComplementSectorFamily} & nur Komplementfamilien-Kontext & nicht für beliebige Kategorien oder Konstanten & nur für die Komplementsektorfamilie des CNNA-Kerns & nicht als generisches C-Symbol ausserhalb des CNNA-Kontexts verwenden & ComplementSectorFamily & kontextgebunden \\
$\Sigma_{\mathrm{BDI}}$ & sector-split symbol & \code{SectorSplit} & nur Bright/Dark/Interface-Split-Kontext & nicht für beliebige Summen- oder Schnittoperationen & nur für den sektor-first-Split des CNNA-Kerns & keine Wiederverwendung als allgemeines Sigma-Schema & SectorSplit & kontextgebunden \\
$\Lambda_{\mathrm{gen}}$ & DtN core symbol & \code{GeneralizedDtN} & nur DtN-/Schur-/Eliminationskontext & nicht für beliebige Spektral- oder Kopplungsoperatoren & nur für den generalisierten DtN-Kern & nicht mit anderen Lambda-Parametern oder Spektren mischen & GeneralizedDtN & kontextgebunden \\
$\mathcal{A}_{\mathrm{comp}}$ & net symbol & \code{ComplementLocalNet} & nur AQFT-Netzkontext & nicht für beliebige Algebrafamilien ausserhalb des Komplementnetzes & nur für das Komplementnetz & kein Ersatz für Interface- oder Bright-Netze & ComplementLocalNet & symbolisch erlaubt \\
$\mathcal{A}_{\mathrm{int}}$ & net symbol & \code{InterfaceLocalNet} & nur Interface-/AQFT-Netzkontext & nicht für das Komplementnetz oder allgemeine Algebrafamilien & nur für das Interfacenetz & kein Ersatz für $\mathcal{A}_{\mathrm{comp}}$ & InterfaceLocalNet & symbolisch erlaubt \\
$\pi_\omega$ & representation symbol & GNS-Darstellung aus \code{ComplementGNS} & nur GNS-/Darstellungskontext & nicht für beliebige projektive Abbildungen & nur für die im GNS-Datum gebundene Darstellung & nicht als freies pi-Symbol ausserhalb des GNS-Kontexts lesen & ComplementGNS & kontextgebunden \\
$\mathcal{H}_\omega$ & Hilbert-space symbol & Hilbertraum des GNS-Datums & nur GNS-/Darstellungskontext & nicht für beliebige Hilberträume ohne GNS-Bindung & nur für den im GNS-Datum gebundenen Hilbertraum & nicht als allgemeines H-Symbol ohne Datum verwenden & ComplementGNS & kontextgebunden \\
$\Omega_\omega$ & cyclic vector & zyklischer Vektor des GNS-Datums & nur GNS-Kontext & nicht für Carrierindexmengen & nur für zyklischen Vektor im GNS-/KMS-Kontext & nie für Carrierindex verwenden & ComplementGNS & kontextgebunden \\
$\Omega_{\mathrm{idx}}$ & carrier/index symbol & Index- oder Trägermenge & nur Carrier-/Netzkontext & nicht für zyklischen Vektor & nur für Index-/Carrierkontext & ersetzt alte überladene $\Omega$-Verwendung für Carrier & Carrier-/Netzdefinitionen & Disambiguierungsersatz für überladenes \(\Omega\) \\
$S$ & Tomita operator symbol & Tomita-Operator & nur Tomita--Takesaki-Kontext & nicht für beliebige Wirkungsfunktionen, Mengen oder S-Matrizen & nur für den Tomita-Operator des modularen Datums & in Streu- oder Wirkungs-Kontexten nicht ohne Index/Klarstellung verwenden & TomitaTakesakiData & symbolisch erlaubt \\
$J$ & modular symbol & modulare Konjugation & nur Tomita--Takesaki-Kontext & nicht als beliebiges lineares Objekt & nur für modulare Konjugation & nicht als beliebiges lineares Objekt wiederverwenden & TomitaTakesakiData / ModularConjugation & symbolisch erlaubt \\
$\Delta$ & modular symbol & modularer Operator & nur Tomita--Takesaki-Kontext & nicht als Differenzenoperator o. a. & nur für modularen Operator & nicht als Differenzenoperator o.~a. lesen & TomitaTakesakiData & symbolisch erlaubt \\
$\Theta$ & time-reversal symbol & implementierte Zeitumkehr & nur Zeitumkehr-/Kramers-Kontext & nicht für beliebige Involutionen & nur für implementierte Zeitumkehr & nicht für beliebige Involutionen verwenden & TimeReversalOnGNS / KramersGate & kontextgebunden \\
$(-1)^F$ & fermion parity symbol & Fermionparität & nur gradierter Statistik-/Kramers-Kontext & nicht für allgemeine Signaturen & nur für Fermionparität & nicht für allgemeine Signaturen oder Gradierungen ausserhalb des Statistikpfads & GradedStarAlgebra / KramersGate & kontextgebunden \\
\bottomrule
\end{tabular}
\end{wideauditlandscape}

\subsection{Register III: Reifegrad- und Architekturstatusregister}
\doclevel{Meta-Audit-Ebene}
Reifegradmarker sind keine Fachobjekte. Sie dürfen deshalb nicht in Objekt- oder Symbolregister
eingemischt werden.

\begin{wideauditlandscape}
\tablecaptionline{Register III: Reifegrad- und Architekturstatusregister -- Teil 1}
\begin{tabular}{L{0.10\linewidth}L{0.17\linewidth}L{0.15\linewidth}L{0.15\linewidth}L{0.12\linewidth}L{0.13\linewidth}L{0.04\linewidth}}
\toprule
Marker & Bedeutung & zulässig für & unzulässig für & bindende Regel & Migrationspflicht & Status \\
\midrule
\code{Seed} & frühe, noch nicht abgeschlossene Vorzone & def, structure, theorem route, gate route & Endtheoreme, vollständige Closure-Objekte, reine Meta-Begriffe & Regelwerk / Reifegradkonvention & muss später de-seeded, geschlossen oder explizit als dauerhaft früh markiert werden & globaler Statusmarker \\
\code{Scaffold} & strukturell getragener, aber noch nicht voll geschlossener Träger & structure, interface carrier, transition object & blosse Prosa-Begriffe, reine Reviewerrollen & Regelwerk / Reifegradkonvention & muss entweder in aktiven Kern übergehen oder als Zwischenstufe bestehen bleiben & globaler Statusmarker \\
\code{Window} & spätes Spezial- oder Matching-Fenster ausserhalb des Kerns & phenomenology windows, IR windows, matching windows & active core, grundlegende Stammobjekte, zwingende Closure-Gates & Regelwerk / Reifegradkonvention & darf nicht stillschweigend in Kernontologie umetikettiert werden & globaler Statusmarker \\
\bottomrule
\end{tabular}
\end{wideauditlandscape}

\subsection{Register IV: Meta-Methodikregister}
\doclevel{Meta-Audit-Ebene}
Meta-Begriffe sind dokumentations- und auditrelevant, aber nicht Bestandteil der fachlichen
Notations-Emission. Deshalb werden sie getrennt registriert.

\begin{wideauditlandscape}
\tablecaptionline{Register IV: Meta-Methodikregister -- Teil 1}
\begin{tabular}{L{0.11\linewidth}L{0.17\linewidth}L{0.10\linewidth}L{0.07\linewidth}L{0.18\linewidth}L{0.10\linewidth}L{0.10\linewidth}L{0.05\linewidth}}
\toprule
Meta-Begriff & Bedeutung & Dokumenttyp & Lean-Objekt? & kanonischer Lean-Name / Ableitungsverbot & Ableitungsregel & bindende Stelle & Status \\
\midrule
\code{TruthProtocol} & mehrstufiges Wahrheits- und Falsifikationsprotokoll & Roadmap, Spec, Audit-Anhang & nein & kein Lean-Objekt; Ableitungsverbot & keine Notationsableitung in \path{Notation.lean} & Methodik-Abschnitt & meta only \\
\code{ExpertReview} & unverzichtbarer externer Fachreview jenseits der Maschinenprüfung & Roadmap, Audit, Review-Prozess & nein & kein Lean-Objekt; Ableitungsverbot & keine Notationsableitung in \path{Notation.lean} & Methodik-Abschnitt & meta only \\
\bottomrule
\end{tabular}
\end{wideauditlandscape}

\paragraph{Integritätsregeln der Lesbarkeitsschicht.}
Es gelten global:
\begin{enumerate}[label=(\arabic*)]
\item Kein Appendix-Eintrag ohne Definitionskatalog-Anker.
\item Kein Registereintrag ohne kanonischen Zielanker oder ausdrückliches Ableitungsverbot.
\item Keine exportierte Form ohne Namespace- und Aktivierungsregel.
\item Keine Symbolverwendung ohne Disambiguierungsregel.
\item Keine sekundäre Exportform ohne expliziten Formrang.
\item Keine Meta-Begriffe in objektbezogenen Notationstabellen.
\item Keine Rückdelegation wesentlicher Tabellensemantik an proseartige Erläuterungen.
\end{enumerate}

\newpage
\begingroup\small
\section{Instanz- und Kontextblatt für tragende Pfade}
\doclevel{Meta-Audit-Ebene}
\label{app:contexts}
\tablecaptionline{Instanz- und Kontextblatt für tragende Pfade -- Teil 1}
\begin{longtable}{L{0.24\linewidth}L{0.18\linewidth}L{0.18\linewidth}L{0.16\linewidth}L{0.18\linewidth}}
\toprule
Pfad / Objekt & globale Voraussetzungen & lokale Voraussetzungen & Instanz-/Kontexttyp & Dokumentrolle \\
\midrule
\endfirsthead
\toprule
Pfad / Objekt & globale Voraussetzungen & lokale Voraussetzungen & Instanz-/Kontexttyp & Dokumentrolle \\
\midrule
\endhead
\bottomrule
\endfoot
\code{GeneralizedDtN} & ToC-Vertrag, SectorSplit & Glueing-/Restriktionsdaten & strukturierter Exportkontext & \shortstack[l]{A$\to$B-\\Handoff} \\
\code{ComplementLocalNet} & Regionen, SectorSplit, DtN-Daten & Netzlokalisierung & AQFT-Kontext & Start des B-Kerns \\
\code{ComplementGNS} & ComplementStateNet & GNS-taugliche Zustandslage & Darstellungs-/Hilbertkontext & Beginn modularer Kette \\
\code{TomitaTakesakiData} & GNS-Lage & SeparatingProperty & operatoralgebraischer Kontext & B-Kern-Mitte \\
\code{ComplementSpectrumCondition} & KMS/positive Energie & Vorwärtskegelkontrolle & Spektral-/Generator-Kontext & Vorbedingung für BW/Scattering \\
\code{ParameterClosure} & Kanal-/Entropiedaten & Selektoren, Fixpunktstruktur & pfeilerübergreifender Closure-Kontext & C$\to$D-Schwelle \\
\code{HadamardStateSeed} & LocalCovariance & gekrümmter Hintergrund / Mikrolokalität & LCQFT-Kontext & renormierte D-Sprache \\
\code{FieldAlgebraReconstructionSeed} & DHR-Sektoren & Charge-/Statistik-Kohärenz & rekonstruktiver Sektorkontext & B/E-Vorzone \\
\bottomrule
\end{longtable}
\endgroup
\newpage
\section{Prüfprotokoll der Roadmap}
\doclevel{Meta-Audit-Ebene}
\label{app:protocol}
\begin{itemize}
\item Geprüfte Codebasis: archivgebunden an \path{REAL-OQS_v0.0605_full_no_comments.zip}.
\item Geprüfter Dokumenttyp: Projekt-/Pfad-Roadmap mit Handoff-, Ketten- und Audit-Anhängen; keine Einzeltheorem-Voll-Dokumentation.
\item Geprüfte Trennungen: CNNA-Kern vs. Bright-Recovery; Seed/Scaffold/Window vs. de-seedeter Kern; Dokumentation vs. physikalische Interpretation.
\item Sichtbare Auditwerkzeuge: Strukturkarte, partielle Ordnung, Gates, Handoffs, Kettenblätter, Register, Kontextblatt, Meta-Audit-Marker.
\item Umgang mit \code{sorry}/\code{admit}/freien Axiomen: im kritischen Pfad nicht zulässig; die Roadmap setzt formale Verifikation statt rhetorischer Ersatzbegründung voraus.
\item Umgang mit privaten Hilfssätzen und \code{omit}-Kontexten: müssen später in Datei-/Theorem-Dokumentationen explizit nachgetragen werden; hier nur als Auditpflicht markiert, nicht unsichtbar vorausgesetzt.
\item Delta-Fokus dieser Fassung: Ergänzung regelwerksnaher Metadaten, Ebenenmarken, Definitionskatalog, Handoff-Spezifikationen, Kettenblätter, Register und Audit-Anhänge.
\end{itemize}
\newpage
\section{Namens- und Pfadverzeichnis}
\doclevel{Meta-Audit-Ebene}
\label{app:pathindex}
\tablecaptionline{Namens- und Pfadverzeichnis -- Teil 1}
\begin{longtable}{L{0.38\linewidth}L{0.52\linewidth}}
\toprule
Name / Pfadbereich & Rolle in der Roadmap \\
\midrule
\endfirsthead
\toprule
Name / Pfadbereich & Rolle in der Roadmap \\
\midrule
\endhead
\bottomrule
\endfoot
\path{CNNA/Core/*} & ontischer und regionslogischer Kern, einschliesslich ToC-Vertrag, SectorSplit und InfiniteCarrier \\
\path{CNNA/OQS/*} & GeneralizedDtN-, MultiSchur-, Kanal- und Rückwirkungsstrang \\
\path{CNNA/AQFT/*} & lokale Netze, Zustände, GNS, modulare Schicht, DHR/DR-Vorzone \\
\path{CNNA/PillarD/*} & emergente Geometrie, LocalCovariance, Relative Cauchy Evolution, Gravity-Routen \\
\path{CNNA/PillarE/*} & späte CPT-/Statistik-/Gauge-/Matter-/SM-Schicht \\
\path{CNNA/AQFT/BrightRecovery/*} & expliziter Recovery-/Regressionstrang des alten Boundary-Pfads \\
\path{REALOQS/PillarA/*} & wichtigste Quellbasis für ontische und OQS-nahe Portierungen \\
\path{REALOQS/PillarB/*} & wichtigste Quellbasis für AQFT-, Modularitäts- und Bright-Recovery-Altpfade \\
\path{REALOQS/PillarC/*} & wichtigste Quellbasis für finite Kanäle, Kraus-/RectKraus-Strukturen und Rückwirkung \\
\bottomrule
\end{longtable}
\newpage
\section{Meta-Audit-Marker}
\doclevel{Meta-Audit-Ebene}
\tablecaptionline{Meta-Audit-Marker -- Teil 1}
\begin{longtable}{L{0.16\linewidth}L{0.76\linewidth}}
\toprule
Marker & Bedeutung \\
\midrule
\endfirsthead
\toprule
Marker & Bedeutung \\
\midrule
\endhead
\bottomrule
\endfoot
M0 & Dieses Dokument ist archivgebunden an \path{REAL-OQS_v0.0605_full_no_comments.zip}; nicht commitgebunden. \\
M1 & Kein boundary-first Dateiname darf auf dem kritischen CNNA-Pfad Architekturzentrum sein. \\
M2 & \(L_{\max}\) ist auf dem aktiven Pfad kein ontischer Grundparameter, sondern höchstens Scaffold-Name vor \code{HorizonLevel}. \\
M3 & Der alte BoundaryMatrix-Pfad ist Bright-Spezialfall und Regressionsorakel, nicht Kernsemantik. \\
M4 & D und E dürfen nur aus expliziten Source-Maps und abgeleiteten Entropie-/Spektral-/Backreactiondaten de-seeden. \\
M5 & Prototypisches Vorauseilen ist nur als \code{Seed}/\code{Scaffold} zulässig; frühe Dummy-Schnittstellen prüfen Typtragfähigkeit, nicht Kernclosure. \\
M6 & Rückfluss später Einsichten in frühere Pfeiler ist erlaubt, aber nur über kontrolliertes Re-Opening, Regressionsprüfung und erneutes Gate-Closing. \\
M7 & AI ist auf dem kritischen Pfad nur Heuristik- und Übersetzungsschicht; technische Wahrheitsquelle bleiben formale Sprache, Kernel-Check und die explizite Disziplin gegen \code{sorry}/\code{admit}/freie Axiome. \\
M8 & Formale Korrektheit des Codes garantiert noch keine korrekte Prosa-Semantik oder physikalische Relevanz; Kommentar-, Dokument- und Interpretationsprüfung sind eigener Pflichtschritt. \\
M9 & Spezialfenster (Minkowski, Conformal, FLRW, SM, SMEFT, Hadron, ChiralEFT) bleiben standardmäßig Windows. \\
M10 & Asymptotic-Safety-Sprache ist erst nach Closure diagnostisch legitim; sie ersetzt die Closure nicht. \\
M11 & String-mathematische Ideen werden als Organisations- und Glueing-/Defekt-/Inflow-Werkzeuge genutzt, nicht als neues fundamentales Postulat. \\
\end{longtable}
\newpage
\section{Tabellarischer Definitionskatalog der Kernobjekte}
\label{app:defcatalog}
\doclevel{Projekt-Ebene und Pfad-Ebene}
Dieser Anhang enthält den aus dem Haupttext ausgelagerten vollständigen tabellarischen Definitionskatalog. Er bleibt die kanonische Objektquelle der Roadmap; die normativen Regeln hierzu stehen weiterhin in Abschnitt~\ref{sec:definitions}.

\begingroup\tiny
\begin{wideauditlandscape}
\landtablecaption{Tabelle K.1-A1: Definitionskatalog der Kernobjekte -- Objektbasis, Teil 1}
\begin{tabular}{L{0.085\linewidth}L{0.043\linewidth}L{0.073\linewidth}L{0.055\linewidth}L{0.043\linewidth}L{0.125\linewidth}L{0.10\linewidth}L{0.105\linewidth}L{0.058\linewidth}L{0.095\linewidth}L{0.065\linewidth}}
\toprule
Objekt & Klasse & Primärform & Exportform & Modus & Expansion / Bindung & Herkunft & Zielmodul & Namespace & Aktivierungsregel & Kontext \\
\midrule
\code{TreeOfCliques} & structure & \emph{Tree of Cliques} & \emph{ToC} & text alias & \emph{Tree of Cliques} und \emph{ToC} verweisen identisch auf \code{TreeOfCliques}; keine konkurrierende Zweitbedeutung & CNNA/Core & \path{CNNA/Notation.lean} & \code{CNNA} & \code{open scoped CNNA} & global \\
\code{BranchPatch} & structure & \emph{Bright-Patch} & \emph{Bright-Patch} & text alias & \emph{Bright-Patch} expandiert auf \code{BranchPatch} & CNNA/Core bzw. Pillar A & \path{CNNA/PillarA/Notation.lean} & \code{CNNA.PillarA} & \code{open scoped CNNA.PillarA} & Pillar A \\
\code{ComplementSectorFamily} & structure & \(\mathcal C\) & \(\mathcal C\) & symbolic notation + text alias & \(\mathcal C\) und \emph{Komplementsektorfamilie} verweisen auf \code{ComplementSectorFamily} & CNNA/Core bzw. Pillar A/Complement & \path{CNNA/PillarA/Notation.lean} & \code{CNNA.PillarA} & \code{open scoped CNNA.PillarA} & Pillar A--B \\
\code{SectorSplit} & structure & \(\Sigma_{\mathrm{BDI}}\) & \(\Sigma_{\mathrm{BDI}}\) & symbolic notation + text alias & \(\Sigma_{\mathrm{BDI}}\) und \emph{Bright/Dark/Interface-Split} verweisen auf \code{SectorSplit} & CNNA/Core/SectorSplit & \path{CNNA/PillarA/Notation.lean} & \code{CNNA.PillarA} & \code{open scoped CNNA.PillarA} & Pillar A--B \\
\code{GeneralizedDtN} & def & \(\Lambda_{\mathrm{gen}}\) & \(\Lambda_{\mathrm{gen}}\) & symbolic notation + text alias & \(\Lambda_{\mathrm{gen}}\) und \emph{generalisierter DtN-Kern} verweisen auf \code{GeneralizedDtN} & CNNA/OQS/GeneralizedDtN & \path{CNNA/PillarA/Notation.lean} & \code{CNNA.PillarA} & \code{open scoped CNNA.PillarA} & Pillar A--B \\
\code{SectorChannels} & structure & \emph{Sektorkanäle} & \emph{SectCh} & text alias & \emph{Sektorkanäle} verweisen auf \code{SectorChannels} & CNNA/OQS & \path{CNNA/OQS/Notation.lean} & \code{CNNA.OQS} & \code{open scoped CNNA.OQS} & Pillar C \\
\code{SectorSysEnv} & structure & \emph{Sektor-SysEnv-Split} & \emph{SysEnvSplit} & text alias & \emph{Sektor-SysEnv-Split} verweist auf \code{SectorSysEnv} & CNNA/OQS & \path{CNNA/OQS/Notation.lean} & \code{CNNA.OQS} & \code{open scoped CNNA.OQS} & Pillar C \\
\code{RelativeEntropyFlow} & def & \emph{relativer Entropiefluss} & \emph{RelEntFlow} & text alias & \emph{relativer Entropiefluss} verweist auf \code{RelativeEntropyFlow} & CNNA/OQS bzw. Pillar C & \path{CNNA/OQS/Notation.lean} & \code{CNNA.OQS} & kontextgebunden; bei Export \code{open scoped CNNA.OQS} & Pillar C--D \\
\code{GlobalStateCone} & structure & \code{GlobalStateCone} & --- & selfspeaking & die lesbare Form ist identisch mit dem Lean-Namen; keine Lichtkegel-Bedeutung & pfeilerübergreifend C--E & kein eigener Aliasexport & \code{CNNA} & keine Exportaktivierung & global \\
\code{Backreaction} & def & \code{Backreaction} & --- & selfspeaking & identisch mit dem Lean-Namen & pfeilerübergreifend C--D & kein eigener Aliasexport & \code{CNNA} & keine Exportaktivierung & global \\
\code{ComplementClosure} & structure & \emph{Komplementschliessung} & \emph{CompClosure} & text alias & \emph{Komplementschliessung} verweist auf \code{ComplementClosure} & AQFT-/Closure-Pfad & \path{CNNA/AQFT/Notation.lean} & \code{CNNA.AQFT} & kontextgebunden; bei Export \code{open scoped CNNA.AQFT} & AQFT/C \\
\bottomrule
\end{tabular}
\end{wideauditlandscape}

\begin{wideauditlandscape}
\landtablecaption{Tabelle K.1-B1: Definitionskatalog der Kernobjekte -- Kollisionen, Abhängigkeiten und Rollen, Teil 1}
\begin{tabular}{L{0.11\linewidth}L{0.18\linewidth}L{0.18\linewidth}L{0.18\linewidth}L{0.07\linewidth}L{0.17\linewidth}}
\toprule
Objekt & Kollisionspolitik & logische Vorgänger & Verbraucher / Nachfolger & Status & Rolle \\
\midrule
\code{TreeOfCliques} & Exportform reserviert; keine konkurrierende Kurzform & --- & \code{BranchPatch}, \code{ComplementSectorFamily}, \code{SectorSplit} & active core & ontischer Startträger der Stammtheorie \\
\code{BranchPatch} & Exportform reserviert; keine konkurrierende Kurzform & \code{TreeOfCliques} & \code{ComplementSectorFamily}, \code{SectorSplit}, \code{GeneralizedDtN} & active core & erster expliziter sichtbarer Patch relativ zum ToC-Anker \\
\code{ComplementSectorFamily} & nur im angegebenen Namespace; keine Zweitbelegung der Symbolform & \code{TreeOfCliques}, \code{BranchPatch} & \code{SectorSplit}, \code{ComplementLocalNet}, \code{ComplementStateNet} & active core & organisierte Familie der dunklen, trunk- und seitenastbezogenen Komplemente \\
\code{SectorSplit} & nur im angegebenen Namespace; keine Zweitbelegung der Symbolform & \code{TreeOfCliques}, \code{ComplementSectorFamily} & \code{GeneralizedDtN}, \code{SectorChannels}, \code{ComplementLocalNet} & active core & synchronisierte Dreiteilung des sektor-first Pfads \\
\code{GeneralizedDtN} & nur im angegebenen Namespace; keine Zweitbelegung der Symbolform & \code{SectorSplit}, \code{ComplementSectorFamily} & \code{SectorChannels}, \code{ComplementLocalNet}, \code{InterfaceLocalNet} & active core & Kopplungs- und Eliminationskern zwischen Bright-, Dark- und Interface-Sektor \\
\code{SectorChannels} & Exportform reserviert; keine konkurrierende Kurzform & \code{GeneralizedDtN}, \code{SectorSplit} & \code{SectorSysEnv}, \code{RelativeEntropyFlow}, \code{Backreaction} & active extension & Kanalstruktur für Bright/System/Environment- bzw. Mehrsektor-Lesarten \\
\code{SectorSysEnv} & Exportform reserviert; keine konkurrierende Kurzform & \code{SectorChannels} & \code{RelativeEntropyFlow}, \code{ComplementStateNet} & active extension & explizite System/Environment-Schicht über dem sektorischen Grundbild \\
\code{RelativeEntropyFlow} & Exportform reserviert; keine konkurrierende Kurzform & \code{SectorChannels}, \code{SectorSysEnv} & \code{Backreaction}, \code{EffectiveLambda}, \code{BackreactionFixedPoint} & active extension & koppelt Zustandsänderung, Entropie und Rückwirkung \\
\code{GlobalStateCone} & kein Exportalias; Lean-Name reserviert & \code{ComplementStateNet}, \code{ComplementKMS} & \code{Backreaction}, \code{RelativeEntropyFlow} & active extension & globaler Zustandskegel, nicht physikalischer Raumzeitkegel \\
\code{Backreaction} & kein Exportalias; Lean-Name reserviert & \code{RelativeEntropyFlow}, \code{GlobalStateCone} & \code{EffectiveLambda}, \code{BackreactionFixedPoint}, \code{DerivedSpacetime} & active extension & Rückkopplung von Zustand, Entropie, Spektrum und Skale \\
\code{ComplementClosure} & Exportform reserviert; keine konkurrierende Kurzform & \code{ComplementLocalNet}, \code{ComplementStateNet} & \code{ComplementGNS}, \code{InfiniteCarrier} & active extension & Schliessungsobjekt des Komplementnetzes vor weiterem Limes- oder Quasilokalitätsübergang \\ \\
\bottomrule
\end{tabular}
\end{wideauditlandscape}

\begin{wideauditlandscape}
\landtablecaption{Tabelle K.1-A2: Definitionskatalog der Kernobjekte -- Objektbasis, Teil 2}
\begin{tabular}{L{0.085\linewidth}L{0.043\linewidth}L{0.073\linewidth}L{0.055\linewidth}L{0.043\linewidth}L{0.125\linewidth}L{0.10\linewidth}L{0.105\linewidth}L{0.058\linewidth}L{0.095\linewidth}L{0.065\linewidth}}
\toprule
Objekt & Klasse & Primärform & Exportform & Modus & Expansion / Bindung & Herkunft & Zielmodul & Namespace & Aktivierungsregel & Kontext \\
\midrule
\code{ParameterClosure} & gate & \emph{Parameter-Closure} & \emph{ParamClosure} & text alias & \emph{Parameter-Closure} verweist auf \code{ParameterClosure} & pfeilerübergreifend C--D & \path{CNNA/Notation.lean} & \code{CNNA} & \code{open scoped CNNA} & global \\
\code{EffectiveLambda} & def & \(\lambda_{\mathrm{eff}}\) & \(\lambda_{\mathrm{eff}}\) & symbolic notation + text alias & \(\lambda_{\mathrm{eff}}\) und \emph{EffectiveLambda} verweisen auf \code{EffectiveLambda} & Closure-/D-Pfad & \path{CNNA/PillarD/Notation.lean} & \code{CNNA.PillarD} & \code{open scoped CNNA.PillarD} & C--D \\
\code{RegularizationClosure} & gate & \emph{Regularisierungs-Closure} & \emph{RegClosure} & text alias & \emph{Regularisierungs-Closure} verweist auf \code{RegularizationClosure} & Closure-Pfad & \path{CNNA/PillarD/Notation.lean} & \code{CNNA.PillarD} & kontextgebunden; bei Export \code{open scoped CNNA.PillarD} & C--D \\
\code{BranchingWitness} & theorem witness & \emph{Branching-Witness} & \emph{BranchWit} & text alias & \emph{Branching-Witness} verweist auf \code{BranchingWitness} & Pillar A/D-Closure & \path{CNNA/PillarA/Notation.lean} & \code{CNNA.PillarA} & kontextgebunden; bei Export \code{open scoped CNNA.PillarA} & A--D \\
\code{selectedBranching} & def & \emph{selected branching} & \emph{SelBranch} & text alias & \emph{selected branching} verweist auf \code{selectedBranching} & Closure-/Selektorpfad & \path{CNNA/PillarA/Notation.lean} & \code{CNNA.PillarA} & kontextgebunden; bei Export \code{open scoped CNNA.PillarA} & A--D \\
\code{UVSpectralSelector} & selector & \emph{UV-Spektralselektor} & \emph{UVSel} & text alias & \emph{UV-Spektralselektor} verweist auf \code{UVSpectralSelector} & Closure-/Selektorpfad & \path{CNNA/PillarA/Notation.lean} & \code{CNNA.PillarA} & kontextgebunden; bei Export \code{open scoped CNNA.PillarA} & A--D \\
\code{BranchingSelector} & selector & \emph{Branching-Selektor} & \emph{BranchSel} & text alias & \emph{Branching-Selektor} verweist auf \code{BranchingSelector} & Closure-/Selektorpfad & \path{CNNA/PillarA/Notation.lean} & \code{CNNA.PillarA} & kontextgebunden; bei Export \code{open scoped CNNA.PillarA} & A--D \\
\code{BackreactionSelector} & selector & \emph{Backreaction-Selektor} & \emph{BackSel} & text alias & \emph{Backreaction-Selektor} verweist auf \code{BackreactionSelector} & Closure-/Selektorpfad & \path{CNNA/PillarD/Notation.lean} & \code{CNNA.PillarD} & kontextgebunden; bei Export \code{open scoped CNNA.PillarD} & C--D \\
\code{InfiniteCarrier} & structure & \code{InfiniteCarrier} & --- & selfspeaking & identisch mit dem Lean-Namen; keine freie Synonymbildung & AQFT-/Quasilokalitäts- und D-Limespfad & kein eigener Aliasexport & \code{CNNA} & keine Exportaktivierung & B--D \\
\code{HorizonLevel} & def & \(L_H\) & \(L_H\) & symbolic notation + text alias & \(L_H\) und \emph{Horizonniveau} verweisen auf \code{HorizonLevel} & CNNA/PillarD/HorizonLevel & \path{CNNA/PillarD/Notation.lean} & \code{CNNA.PillarD} & \code{open scoped CNNA.PillarD} & C--D \\
\code{BackreactionFixedPoint} & gate & \emph{Backreaction-Fixpunkt} & \emph{BackFix} & text alias & \emph{Backreaction-Fixpunkt} verweist auf \code{BackreactionFixedPoint} & CNNA/PillarD/BackreactionFixedPoint & \path{CNNA/PillarD/Notation.lean} & \code{CNNA.PillarD} & kontextgebunden; bei Export \code{open scoped CNNA.PillarD} & C--D \\
\bottomrule
\end{tabular}
\end{wideauditlandscape}

\begin{wideauditlandscape}
\landtablecaption{Tabelle K.1-B2: Definitionskatalog der Kernobjekte -- Kollisionen, Abhängigkeiten und Rollen, Teil 2}
\begin{tabular}{L{0.11\linewidth}L{0.18\linewidth}L{0.18\linewidth}L{0.18\linewidth}L{0.07\linewidth}L{0.17\linewidth}}
\toprule
Objekt & Kollisionspolitik & logische Vorgänger & Verbraucher / Nachfolger & Status & Rolle \\
\midrule
\code{ParameterClosure} & Exportform reserviert; keine konkurrierende Kurzform & \code{RelativeEntropyFlow}, \code{Backreaction} & \code{EffectiveLambda}, \code{HorizonLevel}, \code{BackreactionFixedPoint} & active core & derived-only-Schliessung freier Restparameter \\
\code{EffectiveLambda} & nur im angegebenen Namespace; keine Zweitbelegung der Symbolform & \code{Backreaction}, \code{ParameterClosure} & \code{RegularizationClosure}, \code{BackreactionFixedPoint}, \code{DerivedSpacetime} & planned active & effektive Größe für rückwirkungsstabilisierte Skalenlesarten \\
\code{RegularizationClosure} & Exportform reserviert; keine konkurrierende Kurzform & \code{EffectiveLambda} & \code{HorizonLevel}, \code{BackreactionFixedPoint}, \code{DerivedSpacetime} & planned active & Schliessung von Regularisierung und derived Skala \\
\code{BranchingWitness} & Exportform reserviert; keine konkurrierende Kurzform & \code{BranchPatch}, \code{SectorSplit} & \code{selectedBranching}, \code{UVSpectralSelector}, \code{BranchingSelector} & active extension & expliziter Zeuge nichttrivialer Verzweigung \\
\code{selectedBranching} & Exportform reserviert; keine konkurrierende Kurzform & \code{BranchingWitness} & \code{UVSpectralSelector}, \code{BranchingSelector}, \code{BackreactionSelector} & active extension & ausgewählte Verzweigung nach dokumentierter Selektorlogik \\
\code{UVSpectralSelector} & Exportform reserviert; keine konkurrierende Kurzform & \code{selectedBranching} & \code{BranchingSelector}, \code{BackreactionSelector} & planned active & wahlfreier Selektor für UV-/Spektralstufen soll später derived werden \\
\code{BranchingSelector} & Exportform reserviert; keine konkurrierende Kurzform & \code{UVSpectralSelector} & \code{selectedBranching}, \code{BackreactionSelector} & planned active & explizite, später de-freie Auswahlregel für Verzweigungsdaten \\
\code{BackreactionSelector} & Exportform reserviert; keine konkurrierende Kurzform & \code{BranchingSelector}, \code{selectedBranching} & \code{BackreactionFixedPoint}, \code{EffectiveLambda} & planned active & dokumentierte Auswahlregel für rückgekoppelte Skalenfestlegung \\
\code{InfiniteCarrier} & kein Exportalias; Lean-Name reserviert & \code{ComplementClosure} & \code{DerivedSpacetime}, \code{LocalCovariance} & active extension & explizite unendliche Tragerstruktur jenseits endlicher Approximanten \\
\code{HorizonLevel} & nur im angegebenen Namespace; keine Zweitbelegung der Symbolform & \code{RegularizationClosure}, \code{ParameterClosure} & \code{BackreactionFixedPoint}, \code{DerivedSpacetime}, \code{EinsteinLimitSeed} & active core & ersetzt rohes \(L_{\max}\) durch abgeleitete IR-/Horizontskale \\
\code{BackreactionFixedPoint} & Exportform reserviert; keine konkurrierende Kurzform & \code{HorizonLevel}, \code{Backreaction} & \code{DerivedSpacetime}, \code{InterfaceCausality}, \code{EinsteinLimitSeed} & active core & konsistente Rückkopplung von Entropie, Einfluss und Skale \\ \\
\bottomrule
\end{tabular}
\end{wideauditlandscape}

\begin{wideauditlandscape}
\landtablecaption{Tabelle K.1-A3: Definitionskatalog der Kernobjekte -- Objektbasis, Teil 3}
\begin{tabular}{L{0.085\linewidth}L{0.043\linewidth}L{0.073\linewidth}L{0.055\linewidth}L{0.043\linewidth}L{0.125\linewidth}L{0.10\linewidth}L{0.105\linewidth}L{0.058\linewidth}L{0.095\linewidth}L{0.065\linewidth}}
\toprule
Objekt & Klasse & Primärform & Exportform & Modus & Expansion / Bindung & Herkunft & Zielmodul & Namespace & Aktivierungsregel & Kontext \\
\midrule
\code{ComplementLocalNet} & structure & \(\mathcal A_{\mathrm{comp}}\) & \(\mathcal A_{\mathrm{comp}}\) & symbolic notation + text alias & \(\mathcal A_{\mathrm{comp}}\) und \emph{Komplementnetz} verweisen auf \code{ComplementLocalNet} & CNNA/AQFT/ComplementLocalNet & \path{CNNA/AQFT/Notation.lean} & \code{CNNA.AQFT} & \code{open scoped CNNA.AQFT} & AQFT \\
\code{InterfaceLocalNet} & structure & \(\mathcal A_{\mathrm{int}}\) & \(\mathcal A_{\mathrm{int}}\) & symbolic notation + text alias & \(\mathcal A_{\mathrm{int}}\) und \emph{Interfacenetz} verweisen auf \code{InterfaceLocalNet} & CNNA/AQFT/InterfaceLocalNet & \path{CNNA/AQFT/Notation.lean} & \code{CNNA.AQFT} & \code{open scoped CNNA.AQFT} & AQFT \\
\code{ComplementStateNet} & structure & \(\omega_{\mathrm{comp}}\) & \(\omega_{\mathrm{comp}}\) & symbolic notation + text alias & \(\omega_{\mathrm{comp}}\) und \emph{Komplementzustandsnetz} verweisen auf \code{ComplementStateNet} & CNNA/AQFT/ComplementStateNet & \path{CNNA/AQFT/Notation.lean} & \code{CNNA.AQFT} & \code{open scoped CNNA.AQFT} & AQFT \\
\code{ComplementGNS} & structure & \emph{Komplement-GNS-Datum} & \emph{CompGNS} & text alias & \emph{Komplement-GNS-Datum} verweist auf \code{ComplementGNS}; die Tupelform benennt dieselbe Datenlage und führt keinen zweiten Begriff ein & CNNA/AQFT/ComplementGNS & \path{CNNA/AQFT/Notation.lean} & \code{CNNA.AQFT} & kontextgebunden; bei Export \code{open scoped CNNA.AQFT} & AQFT \\
\code{SeparatingProperty} & theorem/property & \emph{zyklisch-separierende Lage} & \emph{CycSep} & text alias & \emph{zyklisch-separierende Lage} verweist auf \code{SeparatingProperty} & CNNA/AQFT/SeparatingProperty & \path{CNNA/AQFT/Notation.lean} & \code{CNNA.AQFT} & kontextgebunden; bei Export \code{open scoped CNNA.AQFT} & AQFT \\
\code{TomitaTakesakiData} & structure & \emph{Tomita--Takesaki-Datum} & \emph{TTData} & text alias & \emph{Tomita--Takesaki-Datum} verweist auf \code{TomitaTakesakiData}; die Tupelform ist nur dieselbe Datenexposition & CNNA/AQFT/TomitaTakesakiData & \path{CNNA/AQFT/Notation.lean} & \code{CNNA.AQFT} & kontextgebunden; bei Export \code{open scoped CNNA.AQFT} & AQFT \\
\code{ComplementKMS} & def/state & \emph{Komplement-KMS-Zustand} & \emph{CompKMS} & text alias & \emph{Komplement-KMS-Zustand} verweist auf \code{ComplementKMS}; \(\omega_\beta\) ist nur im KMS-Kontext erlaubt & CNNA/AQFT/ComplementKMS & \path{CNNA/AQFT/Notation.lean} & \code{CNNA.AQFT} & kontextgebunden; bei Export \code{open scoped CNNA.AQFT} & AQFT \\
\code{ComplementSpectrumCondition} & gate/property & \emph{Komplement-Spektralbedingung} & \emph{CompSpectrum} & text alias & \emph{Komplement-Spektralbedingung} verweist auf \code{ComplementSpectrumCondition} & CNNA/AQFT/ComplementSpectrumCondition & \path{CNNA/AQFT/Notation.lean} & \code{CNNA.AQFT} & kontextgebunden; bei Export \code{open scoped CNNA.AQFT} & B--D \\
\code{GradedStarAlgebra} & structure & \code{GradedStarAlgebra} & --- & selfspeaking & identisch mit dem Lean-Namen & CNNA/AQFT/GradedStarAlgebra & kein eigener Aliasexport & \code{CNNA.AQFT} & keine Exportaktivierung & AQFT \\
\code{ModularConjugation} & def/theorem layer & \emph{Modularkonjugation} & \emph{ModConj} & text alias & \emph{Modularkonjugation} verweist auf \code{ModularConjugation} & CNNA/AQFT/ModularConjugation & \path{CNNA/AQFT/Notation.lean} & \code{CNNA.AQFT} & kontextgebunden; bei Export \code{open scoped CNNA.AQFT} & AQFT \\
\code{BisognanoWichmannSeed} & seed & \emph{BW-Seed} & \emph{BWSeed} & text alias & \emph{BW-Seed} verweist auf \code{BisognanoWichmannSeed} & CNNA/AQFT/BisognanoWichmannSeed & \path{CNNA/AQFT/Notation.lean} & \code{CNNA.AQFT} & kontextgebunden; bei Export \code{open scoped CNNA.AQFT} & AQFT \\
\bottomrule
\end{tabular}
\end{wideauditlandscape}

\begin{wideauditlandscape}
\landtablecaption{Tabelle K.1-B3: Definitionskatalog der Kernobjekte -- Kollisionen, Abhängigkeiten und Rollen, Teil 3}
\begin{tabular}{L{0.11\linewidth}L{0.18\linewidth}L{0.18\linewidth}L{0.18\linewidth}L{0.07\linewidth}L{0.17\linewidth}}
\toprule
Objekt & Kollisionspolitik & logische Vorgänger & Verbraucher / Nachfolger & Status & Rolle \\
\midrule
\code{ComplementLocalNet} & nur im angegebenen Namespace; keine Zweitbelegung der Symbolform & \code{ComplementSectorFamily}, \code{GeneralizedDtN} & \code{ComplementStateNet}, \code{ComplementGNS}, \code{ComplementSpectrumCondition} & active core & lokales Observablennetz des dunklen Sektors \\
\code{InterfaceLocalNet} & nur im angegebenen Namespace; keine Zweitbelegung der Symbolform & \code{GeneralizedDtN}, \code{SectorSplit} & \code{InterfaceCausality}, \code{LocalCovariance} & active core & lokales Netz der Glueing-, Defekt- und Entropieträger \\
\code{ComplementStateNet} & nur im angegebenen Namespace; keine Zweitbelegung der Symbolform & \code{ComplementLocalNet}, \code{SectorSysEnv} & \code{GlobalStateCone}, \code{ComplementGNS}, \code{ComplementKMS} & active core & Zustandsfamilie vor GNS- und KMS-Schicht \\
\code{ComplementGNS} & Exportform reserviert; keine konkurrierende Kurzform & \code{ComplementStateNet}, \code{ComplementLocalNet} & \code{SeparatingProperty}, \code{TomitaTakesakiData}, \code{TimeReversalOnGNS} & active core & Darstellungsstart des modularen B-Kerns \\
\code{SeparatingProperty} & Exportform reserviert; keine konkurrierende Kurzform & \code{ComplementGNS} & \code{TomitaTakesakiData}, \code{ModularConjugation} & active core & Scharnier zwischen GNS und Tomita--Takesaki \\
\code{TomitaTakesakiData} & Exportform reserviert; keine konkurrierende Kurzform & \code{ComplementGNS}, \code{SeparatingProperty} & \code{ModularConjugation}, \code{ComplementKMS} & active core & modularer Operator-, Konjugations- und Tomita-Kern \\
\code{ComplementKMS} & Exportform reserviert; keine konkurrierende Kurzform & \code{ComplementStateNet}, \code{TomitaTakesakiData} & \code{ComplementSpectrumCondition}, \code{BisognanoWichmannSeed} & active core & thermische und modulare Schicht des Komplementnetzes \\
\code{ComplementSpectrumCondition} & Exportform reserviert; keine konkurrierende Kurzform & \code{ComplementKMS}, \code{ComplementLocalNet} & \code{BisognanoWichmannSeed}, \code{NuclearitySeed}, \code{CPTGate} & active core & Vorbedingung für BW-, Scattering- und Kramers-nahe Strukturen \\
\code{GradedStarAlgebra} & kein Exportalias; Lean-Name reserviert & \code{ComplementLocalNet}, \code{ComplementGNS} & \code{FermionicStatisticsSeed}, \code{KramersGate} & active extension & algebraischer Untergrund für fermionische und gradierte Sektoren \\
\code{ModularConjugation} & Exportform reserviert; keine konkurrierende Kurzform & \code{TomitaTakesakiData}, \code{SeparatingProperty} & \code{BisognanoWichmannSeed}, \code{TimeReversalOnGNS}, \code{CPTGate} & active extension & kontrolliert die konjugative Seite des modularen Kerns \\
\code{BisognanoWichmannSeed} & Exportform reserviert; keine konkurrierende Kurzform & \code{ModularConjugation}, \code{ComplementSpectrumCondition} & \code{LocalCovariance}, \code{CPTGate} & active extension & Brücke von modularer Gruppe zu Boost-/Wedge-Lesarten \\ \\
\bottomrule
\end{tabular}
\end{wideauditlandscape}

\begin{wideauditlandscape}
\landtablecaption{Tabelle K.1-A4: Definitionskatalog der Kernobjekte -- Objektbasis, Teil 4}
\begin{tabular}{L{0.085\linewidth}L{0.043\linewidth}L{0.073\linewidth}L{0.055\linewidth}L{0.043\linewidth}L{0.125\linewidth}L{0.10\linewidth}L{0.105\linewidth}L{0.058\linewidth}L{0.095\linewidth}L{0.065\linewidth}}
\toprule
Objekt & Klasse & Primärform & Exportform & Modus & Expansion / Bindung & Herkunft & Zielmodul & Namespace & Aktivierungsregel & Kontext \\
\midrule
\code{NuclearitySeed} & seed & \emph{Nuclearity-Seed} & \emph{NucSeed} & text alias & \emph{Nuclearity-Seed} verweist auf \code{NuclearitySeed} & CNNA/AQFT/NuclearitySeed & \path{CNNA/AQFT/Notation.lean} & \code{CNNA.AQFT} & kontextgebunden; bei Export \code{open scoped CNNA.AQFT} & AQFT \\
\code{Nuclearity} & gate/theorem & \code{Nuclearity} & --- & selfspeaking & identisch mit dem Lean-Namen & \path{CNNA/Gates/Nuclearity.lean} & kein eigener Aliasexport & \code{CNNA} & keine Exportaktivierung & AQFT/Gates \\
\code{DHRSectorsSeed} & seed & \emph{DHR-Sektoren-Seed} & \emph{DHRSeed} & text alias & \emph{DHR-Sektoren-Seed} verweist auf \code{DHRSectorsSeed} & CNNA/AQFT/DHRSectorsSeed & \path{CNNA/AQFT/Notation.lean} & \code{CNNA.AQFT} & kontextgebunden; bei Export \code{open scoped CNNA.AQFT} & AQFT \\
\code{FieldAlgebraReconstructionSeed} & seed & \emph{Feldalgebra-Rekonstruktions-Seed} & \emph{FieldRecSeed} & text alias & \emph{Feldalgebra-Rekonstruktions-Seed} verweist auf \code{FieldAlgebraReconstructionSeed} & CNNA/AQFT/FieldAlgebraReconstructionSeed & \path{CNNA/AQFT/Notation.lean} & \code{CNNA.AQFT} & kontextgebunden; bei Export \code{open scoped CNNA.AQFT} & AQFT \\
\code{ChargeStatisticsCoherence} & gate & \emph{Ladungs-Statistik-Kohärenz} & \emph{ChargeStatCoh} & text alias & \emph{Ladungs-Statistik-Kohärenz} verweist auf \code{ChargeStatisticsCoherence} & \path{CNNA/Gates/ChargeStatisticsCoherence.lean} & \path{CNNA/AQFT/Notation.lean} & \code{CNNA.AQFT} & kontextgebunden; bei Export \code{open scoped CNNA.AQFT} & AQFT/Gates \\
\code{ChargeReconstructionTheorem} & theorem & \emph{Ladungsrekonstruktionssatz} & \emph{ChargeRec} & text alias & \emph{Ladungsrekonstruktionssatz} verweist auf \code{ChargeReconstructionTheorem} & CNNA/AQFT/ChargeReconstructionTheorem & \path{CNNA/AQFT/Notation.lean} & \code{CNNA.AQFT} & kontextgebunden; bei Export \code{open scoped CNNA.AQFT} & AQFT \\
\code{DerivedSpacetime} & structure & \emph{derived spacetime} & \emph{DerSpacetime} & text alias & \emph{derived spacetime} verweist auf \code{DerivedSpacetime} & CNNA/PillarD/DerivedSpacetime & \path{CNNA/PillarD/Notation.lean} & \code{CNNA.PillarD} & kontextgebunden; bei Export \code{open scoped CNNA.PillarD} & Pillar D \\
\code{InterfaceCausality} & structure/gate & \emph{Interface-Kausalität} & \emph{IfcCaus} & text alias & \emph{Interface-Kausalität} verweist auf \code{InterfaceCausality} & CNNA/PillarD/InterfaceCausality & \path{CNNA/PillarD/Notation.lean} & \code{CNNA.PillarD} & kontextgebunden; bei Export \code{open scoped CNNA.PillarD} & Pillar D \\
\code{LocalCovariance} & structure/gate & \emph{lokale Kovarianz} & \emph{LocCov} & text alias & \emph{lokale Kovarianz} verweist auf \code{LocalCovariance} & CNNA/PillarD/LocalCovariance & \path{CNNA/PillarD/Notation.lean} & \code{CNNA.PillarD} & kontextgebunden; bei Export \code{open scoped CNNA.PillarD} & Pillar D \\
\code{RelativeCauchyEvolutionSeed} & seed & \code{RelativeCauchyEvolutionSeed} & --- & selfspeaking & identisch mit dem Lean-Namen; kein kurzer Symbolalias & CNNA/PillarD/RelativeCauchyEvolutionSeed & kein eigener Aliasexport & \code{CNNA.PillarD} & keine Exportaktivierung & Pillar D \\
\code{DynamicalLocality} & gate & \code{DynamicalLocality} & --- & selfspeaking & identisch mit dem Lean-Namen & \path{CNNA/Gates/DynamicalLocality.lean} & kein eigener Aliasexport & \code{CNNA} & keine Exportaktivierung & D/Gates \\
\bottomrule
\end{tabular}
\end{wideauditlandscape}

\begin{wideauditlandscape}
\landtablecaption{Tabelle K.1-B4: Definitionskatalog der Kernobjekte -- Kollisionen, Abhängigkeiten und Rollen, Teil 4}
\begin{tabular}{L{0.11\linewidth}L{0.18\linewidth}L{0.18\linewidth}L{0.18\linewidth}L{0.07\linewidth}L{0.17\linewidth}}
\toprule
Objekt & Kollisionspolitik & logische Vorgänger & Verbraucher / Nachfolger & Status & Rolle \\
\midrule
\code{NuclearitySeed} & Exportform reserviert; keine konkurrierende Kurzform & \code{ComplementSpectrumCondition} & \code{Nuclearity}, \code{HadamardStateSeed}, \code{LocalCovariance} & active extension & Phase-Space-Vorzone jenseits blossen Split-Verhaltens \\
\code{Nuclearity} & kein Exportalias; Lean-Name reserviert & \code{NuclearitySeed} & \code{DHRSectorsSeed}, \code{ChargeStatisticsCoherence} & active extension & Phase-Space- und Dichteschicht vor DHR-/Gauge-Rekonstruktion \\
\code{DHRSectorsSeed} & Exportform reserviert; keine konkurrierende Kurzform & \code{Nuclearity}, \code{ComplementSpectrumCondition} & \code{FieldAlgebraReconstructionSeed}, \code{ChargeStatisticsCoherence} & active extension & lokalisierte Superselection-Sektoren des Observablennetzes \\
\code{FieldAlgebraReconstructionSeed} & Exportform reserviert; keine konkurrierende Kurzform & \code{DHRSectorsSeed} & \code{ChargeReconstructionTheorem}, \code{CPTTheorem} & active extension & DR-artiger Weg zu Feldalgebra und kompakter Eichgruppe \\
\code{ChargeStatisticsCoherence} & Exportform reserviert; keine konkurrierende Kurzform & \code{DHRSectorsSeed}, \code{Nuclearity} & \code{ChargeReconstructionTheorem}, \code{SpinStatisticsTheorem} & active extension & prüft Kohärenz von Ladung, Statistik und rekonstruiertem Gauge-Layer \\
\code{ChargeReconstructionTheorem} & Exportform reserviert; keine konkurrierende Kurzform & \code{FieldAlgebraReconstructionSeed}, \code{ChargeStatisticsCoherence} & \code{CPTTheorem}, \code{GaugeSectorSeed} & active extension & expliziter DHR/DR-Abschluss vor voller Gauge-/Matter-Lesart \\
\code{DerivedSpacetime} & Exportform reserviert; keine konkurrierende Kurzform & \code{BackreactionFixedPoint}, \code{HorizonLevel} & \code{InterfaceCausality}, \code{LocalCovariance}, \code{EinsteinLimitSeed} & active extension & emergenter geometrischer Trager jenseits blosser Interface-Logik \\
\code{InterfaceCausality} & Exportform reserviert; keine konkurrierende Kurzform & \code{DerivedSpacetime}, \code{InterfaceLocalNet} & \code{LocalCovariance}, \code{PTBridgeSeed} & active extension & kausale Glueing- und Schnittstellenstruktur \\
\code{LocalCovariance} & Exportform reserviert; keine konkurrierende Kurzform & \code{InterfaceCausality}, \code{BisognanoWichmannSeed} & \code{RelativeCauchyEvolutionSeed}, \code{HadamardStateSeed}, \code{CPTGate} & active extension & LCQFT-nahe Kovarianzbedingung auf emergentem Hintergrund \\
\code{RelativeCauchyEvolutionSeed} & kein Exportalias; Lean-Name reserviert & \code{LocalCovariance} & \code{DynamicalLocality}, \code{DerivedStressTensorSeed} & active extension & dynamische Reaktion auf lokale Hintergrundvariation \\
\code{DynamicalLocality} & kein Exportalias; Lean-Name reserviert & \code{RelativeCauchyEvolutionSeed} & \code{HadamardStateSeed}, \code{LocalWickPolynomialsSeed} & active extension & Abgleich von kinematischer und dynamischer Lokalität \\ \\
\bottomrule
\end{tabular}
\end{wideauditlandscape}

\begin{wideauditlandscape}
\landtablecaption{Tabelle K.1-A5: Definitionskatalog der Kernobjekte -- Objektbasis, Teil 5}
\begin{tabular}{L{0.085\linewidth}L{0.043\linewidth}L{0.073\linewidth}L{0.055\linewidth}L{0.043\linewidth}L{0.125\linewidth}L{0.10\linewidth}L{0.105\linewidth}L{0.058\linewidth}L{0.095\linewidth}L{0.065\linewidth}}
\toprule
Objekt & Klasse & Primärform & Exportform & Modus & Expansion / Bindung & Herkunft & Zielmodul & Namespace & Aktivierungsregel & Kontext \\
\midrule
\code{HadamardStateSeed} & seed & \emph{Hadamard-State-Seed} & \emph{HadSeed} & text alias & \emph{Hadamard-State-Seed} verweist auf \code{HadamardStateSeed} & CNNA/PillarD/HadamardStateSeed & \path{CNNA/PillarD/Notation.lean} & \code{CNNA.PillarD} & kontextgebunden; bei Export \code{open scoped CNNA.PillarD} & Pillar D \\
\code{MicrolocalSpectrumGate} & gate & \emph{mikrolokales Spektralgate} & \emph{MicroSpec} & text alias & \emph{mikrolokales Spektralgate} verweist auf \code{MicrolocalSpectrumGate} & CNNA/PillarD/MicrolocalSpectrumGate & \path{CNNA/PillarD/Notation.lean} & \code{CNNA.PillarD} & kontextgebunden; bei Export \code{open scoped CNNA.PillarD} & Pillar D \\
\code{LocalWickPolynomialsSeed} & seed & \emph{lokale-Wick-Polynome-Seed} & \emph{WickSeed} & text alias & \emph{lokale-Wick-Polynome-Seed} verweist auf \code{LocalWickPolynomialsSeed} & CNNA/PillarD/LocalWickPolynomialsSeed & \path{CNNA/PillarD/Notation.lean} & \code{CNNA.PillarD} & kontextgebunden; bei Export \code{open scoped CNNA.PillarD} & Pillar D \\
\code{DerivedStressTensorSeed} & seed & \code{DerivedStressTensorSeed} & --- & selfspeaking & identisch mit dem Lean-Namen & CNNA/PillarD/DerivedStressTensorSeed & kein eigener Aliasexport & \code{CNNA.PillarD} & keine Exportaktivierung & Pillar D \\
\code{EntanglementEquilibriumSeed} & seed & \code{EntanglementEquilibriumSeed} & --- & selfspeaking & identisch mit dem Lean-Namen & CNNA/PillarD & kein eigener Aliasexport & \code{CNNA.PillarD} & keine Exportaktivierung & D \\
\code{EinsteinLimitSeed} & seed & \emph{Einstein-Limes-Seed} & \emph{EinsteinSeed} & text alias & \emph{Einstein-Limes-Seed} verweist auf \code{EinsteinLimitSeed} & CNNA/PillarD & \path{CNNA/PillarD/Notation.lean} & \code{CNNA.PillarD} & kontextgebunden; bei Export \code{open scoped CNNA.PillarD} & D \\
\code{InducedGravitySeed} & seed & \code{InducedGravitySeed} & --- & selfspeaking & identisch mit dem Lean-Namen & CNNA/PillarD & kein eigener Aliasexport & \code{CNNA.PillarD} & keine Exportaktivierung & D \\
\code{TimeReversalOnGNS} & def & \emph{Zeitumkehr auf GNS} & \emph{TRonGNS} & text alias & \emph{Zeitumkehr auf GNS} verweist auf \code{TimeReversalOnGNS} & CNNA/AQFT/TimeReversalOnGNS & \path{CNNA/AQFT/Notation.lean} & \code{CNNA.AQFT} & kontextgebunden; bei Export \code{open scoped CNNA.AQFT} & AQFT/E-Symmetrie \\
\code{AntiunitarySeed} & seed & \code{AntiunitarySeed} & --- & selfspeaking & identisch mit dem Lean-Namen & CNNA/AQFT/AntiunitarySeed & kein eigener Aliasexport & \code{CNNA.AQFT} & keine Exportaktivierung & AQFT/E-Symmetrie \\
\code{SpinRepresentationSeed} & seed & \emph{Spinrepräsentations-Seed} & \emph{SpinSeed} & text alias & \emph{Spinrepräsentations-Seed} verweist auf \code{SpinRepresentationSeed} & CNNA/AQFT/SpinRepresentationSeed & \path{CNNA/AQFT/Notation.lean} & \code{CNNA.AQFT} & kontextgebunden; bei Export \code{open scoped CNNA.AQFT} & AQFT/E-Symmetrie \\
\code{FermionicStatisticsSeed} & seed & \code{FermionicStatisticsSeed} & --- & selfspeaking & identisch mit dem Lean-Namen & CNNA/AQFT/FermionicStatisticsSeed & kein eigener Aliasexport & \code{CNNA.AQFT} & keine Exportaktivierung & AQFT/E-Symmetrie \\
\bottomrule
\end{tabular}
\end{wideauditlandscape}

\begin{wideauditlandscape}
\landtablecaption{Tabelle K.1-B5: Definitionskatalog der Kernobjekte -- Kollisionen, Abhängigkeiten und Rollen, Teil 5}
\begin{tabular}{L{0.11\linewidth}L{0.18\linewidth}L{0.18\linewidth}L{0.18\linewidth}L{0.07\linewidth}L{0.17\linewidth}}
\toprule
Objekt & Kollisionspolitik & logische Vorgänger & Verbraucher / Nachfolger & Status & Rolle \\
\midrule
\code{HadamardStateSeed} & Exportform reserviert; keine konkurrierende Kurzform & \code{LocalCovariance}, \code{NuclearitySeed} & \code{MicrolocalSpectrumGate}, \code{LocalWickPolynomialsSeed} & active extension & mikrolokal zulässige Zustandsklasse auf emergentem Hintergrund \\
\code{MicrolocalSpectrumGate} & Exportform reserviert; keine konkurrierende Kurzform & \code{HadamardStateSeed} & \code{LocalWickPolynomialsSeed}, \code{DerivedStressTensorSeed} & active extension & prüft Hadamard-/mikrolokale Zulässigkeit \\
\code{LocalWickPolynomialsSeed} & Exportform reserviert; keine konkurrierende Kurzform & \code{HadamardStateSeed}, \code{MicrolocalSpectrumGate} & \code{DerivedStressTensorSeed}, \code{EinsteinLimitSeed} & active extension & lokale kovariante Wick- und Zeitordnungsprodukte \\
\code{DerivedStressTensorSeed} & kein Exportalias; Lean-Name reserviert & \code{LocalWickPolynomialsSeed}, \code{RelativeCauchyEvolutionSeed} & \code{EntanglementEquilibriumSeed}, \code{EinsteinLimitSeed}, \code{InducedGravitySeed} & active extension & renormierter Energie-Impuls-Ausgang der lokalen Observablen \\
\code{EntanglementEquilibriumSeed} & kein Exportalias; Lean-Name reserviert & \code{DerivedStressTensorSeed} & \code{EinsteinLimitSeed}, \code{InducedGravitySeed} & planned active & Jacobson-nahe Route von Entropiegleichgewicht zu Gravitation \\
\code{EinsteinLimitSeed} & Exportform reserviert; keine konkurrierende Kurzform & \code{DerivedStressTensorSeed}, \code{EntanglementEquilibriumSeed} & \code{InducedGravitySeed}, \code{SpectralTripleSeed} & planned active & später Limes vom derived D-Pfad zu GR-artiger Dynamik \\
\code{InducedGravitySeed} & kein Exportalias; Lean-Name reserviert & \code{EinsteinLimitSeed} & \code{SpectralTripleSeed}, \code{GaugeSectorSeed} & planned active & alternative späte Route zu gravitativem Verhalten über Materie-/Zustandsbeitrag \\
\code{TimeReversalOnGNS} & Exportform reserviert; keine konkurrierende Kurzform & \code{ComplementGNS}, \code{ModularConjugation} & \code{AntiunitarySeed}, \code{KramersGate} & active extension & hebt die observablenseitige Zeitumkehr auf den Darstellungsraum \\
\code{AntiunitarySeed} & kein Exportalias; Lean-Name reserviert & \code{TimeReversalOnGNS} & \code{KramersGate}, \code{CPTGate} & active extension & formalisiert Antiunitarität der implementierten Zeitumkehr \\
\code{SpinRepresentationSeed} & Exportform reserviert; keine konkurrierende Kurzform & \code{GradedStarAlgebra} & \code{FermionicStatisticsSeed}, \code{SpinStatisticsTheorem} & active extension & Spin-/double-cover-Schicht vor fermionischer Statistik \\
\code{FermionicStatisticsSeed} & kein Exportalias; Lean-Name reserviert & \code{SpinRepresentationSeed}, \code{GradedStarAlgebra} & \code{StatisticsGate}, \code{KramersGate} & active extension & trennt fermionische von bosonischen Sektoren \\ \\
\bottomrule
\end{tabular}
\end{wideauditlandscape}

\begin{wideauditlandscape}
\landtablecaption{Tabelle K.1-A6: Definitionskatalog der Kernobjekte -- Objektbasis, Teil 6}
\begin{tabular}{L{0.085\linewidth}L{0.043\linewidth}L{0.073\linewidth}L{0.055\linewidth}L{0.043\linewidth}L{0.125\linewidth}L{0.10\linewidth}L{0.105\linewidth}L{0.058\linewidth}L{0.095\linewidth}L{0.065\linewidth}}
\toprule
Objekt & Klasse & Primärform & Exportform & Modus & Expansion / Bindung & Herkunft & Zielmodul & Namespace & Aktivierungsregel & Kontext \\
\midrule
\code{StatisticsGate} & gate & \emph{Statistik-Gate} & \emph{StatGate} & text alias & \emph{Statistik-Gate} verweist auf \code{StatisticsGate} & \path{CNNA/Gates/StatisticsGate.lean} & \path{CNNA/AQFT/Notation.lean} & \code{CNNA.AQFT} & kontextgebunden; bei Export \code{open scoped CNNA.AQFT} & AQFT/Gates \\
\code{SpinStatisticsTheorem} & theorem & \emph{Spin-Statistik-Theorem} & \emph{SpinStatThm} & text alias & \emph{Spin-Statistik-Theorem} verweist auf \code{SpinStatisticsTheorem} & CNNA/AQFT/SpinStatisticsTheorem & \path{CNNA/AQFT/Notation.lean} & \code{CNNA.AQFT} & kontextgebunden; bei Export \code{open scoped CNNA.AQFT} & AQFT \\
\code{PTBridgeSeed} & seed & \emph{PT-Bridge-Seed} & \emph{PTSeed} & text alias & \emph{PT-Bridge-Seed} verweist auf \code{PTBridgeSeed} & CNNA/PillarE/PTBridgeSeed & \path{CNNA/PillarE/Notation.lean} & \code{CNNA.PillarE} & kontextgebunden; bei Export \code{open scoped CNNA.PillarE} & D--E \\
\code{ChargeConjugationSeed} & seed & \emph{Ladungskonjugations-Seed} & \emph{CSeed} & text alias & \emph{Ladungskonjugations-Seed} verweist auf \code{ChargeConjugationSeed} & CNNA/PillarE/ChargeConjugationSeed & \path{CNNA/PillarE/Notation.lean} & \code{CNNA.PillarE} & kontextgebunden; bei Export \code{open scoped CNNA.PillarE} & E \\
\code{CPTGate} & gate & \emph{CPT-Gate} & \emph{CPTCheck} & text alias & \emph{CPT-Gate} verweist auf \code{CPTGate} & \path{CNNA/Gates/CPTGate.lean} & \path{CNNA/PillarE/Notation.lean} & \code{CNNA.PillarE} & kontextgebunden; bei Export \code{open scoped CNNA.PillarE} & B--D--E \\
\code{CPTTheorem} & theorem & \emph{CPT-Theorem} & \emph{CPTThm} & text alias & \emph{CPT-Theorem} verweist auf \code{CPTTheorem} & CNNA/AQFT/CPTTheorem & \path{CNNA/PillarE/Notation.lean} & \code{CNNA.AQFT} & kontextgebunden; bei Export \code{open scoped CNNA.AQFT} & B--D--E \\
\code{KramersGate} & gate & \emph{Kramers-Gate} & \emph{KramersCheck} & text alias & \emph{Kramers-Gate} verweist auf \code{KramersGate} & \path{CNNA/Gates/KramersGate.lean} & \path{CNNA/PillarE/Notation.lean} & \code{CNNA.PillarE} & kontextgebunden; bei Export \code{open scoped CNNA.PillarE} & AQFT/E-Symmetrie \\
\code{KramersTheorem} & theorem & \emph{Kramers-Theorem} & \emph{KramersThm} & text alias & \emph{Kramers-Theorem} verweist auf \code{KramersTheorem} & CNNA/AQFT/KramersTheorem & \path{CNNA/PillarE/Notation.lean} & \code{CNNA.AQFT} & kontextgebunden; bei Export \code{open scoped CNNA.AQFT} & AQFT/E-Symmetrie \\
\code{GaugeSectorSeed} & seed & \emph{Gauge-Sektor-Seed} & \emph{GaugeSeed} & text alias & \emph{Gauge-Sektor-Seed} verweist auf \code{GaugeSectorSeed} & CNNA/PillarE/GaugeSectorSeed & \path{CNNA/PillarE/Notation.lean} & \code{CNNA.PillarE} & kontextgebunden; bei Export \code{open scoped CNNA.PillarE} & E \\
\code{AnomalyInflowSeed} & seed & \emph{Anomaly-Inflow-Seed} & \emph{InflowSeed} & text alias & \emph{Anomaly-Inflow-Seed} verweist auf \code{AnomalyInflowSeed} & CNNA/PillarE/AnomalyInflowSeed & \path{CNNA/PillarE/Notation.lean} & \code{CNNA.PillarE} & kontextgebunden; bei Export \code{open scoped CNNA.PillarE} & E \\
\code{AnomalyInflow} & gate & \code{AnomalyInflow} & --- & selfspeaking & identisch mit dem Lean-Namen & \path{CNNA/Gates/AnomalyInflow.lean} & kein eigener Aliasexport & \code{CNNA} & keine Exportaktivierung & E/Gates \\
\bottomrule
\end{tabular}
\end{wideauditlandscape}

\begin{wideauditlandscape}
\landtablecaption{Tabelle K.1-B6: Definitionskatalog der Kernobjekte -- Kollisionen, Abhängigkeiten und Rollen, Teil 6}
\begin{tabular}{L{0.11\linewidth}L{0.18\linewidth}L{0.18\linewidth}L{0.18\linewidth}L{0.07\linewidth}L{0.17\linewidth}}
\toprule
Objekt & Kollisionspolitik & logische Vorgänger & Verbraucher / Nachfolger & Status & Rolle \\
\midrule
\code{StatisticsGate} & Exportform reserviert; keine konkurrierende Kurzform & \code{FermionicStatisticsSeed} & \code{SpinStatisticsTheorem}, \code{KramersTheorem} & active extension & Spin-Statistik-Konsistenztest auf gradierter Algebra \\
\code{SpinStatisticsTheorem} & Exportform reserviert; keine konkurrierende Kurzform & \code{StatisticsGate}, \code{ChargeStatisticsCoherence} & \code{CPTTheorem}, \code{KramersTheorem} & active extension & Endtheorem des Statistikstrangs \\
\code{PTBridgeSeed} & Exportform reserviert; keine konkurrierende Kurzform & \code{InterfaceCausality} & \code{ChargeConjugationSeed}, \code{CPTGate} & active extension & koppelt Raumspiegelungsdaten an die B-Zeitumkehrschicht \\
\code{ChargeConjugationSeed} & Exportform reserviert; keine konkurrierende Kurzform & \code{PTBridgeSeed} & \code{CPTGate}, \code{CPTTheorem} & active extension & Teilchen/Antiteilchen- bzw. C-Layer \\
\code{CPTGate} & Exportform reserviert; keine konkurrierende Kurzform & \code{ChargeConjugationSeed}, \code{AntiunitarySeed} & \code{CPTTheorem}, \code{KramersGate} & active extension & pfeilerübergreifende Abnahme von PT, C und CPT \\
\code{CPTTheorem} & Exportform reserviert; keine konkurrierende Kurzform & \code{CPTGate}, \code{SpinStatisticsTheorem} & \code{KramersTheorem}, \code{GaugeSectorSeed} & active extension & explizites Endtheorem des pfeilerübergreifenden CPT-Strangs \\
\code{KramersGate} & Exportform reserviert; keine konkurrierende Kurzform & \code{TimeReversalOnGNS}, \code{CPTGate} & \code{KramersTheorem}, \code{ChiralMatterSeed} & active extension & trennt \(\Theta^2=(-1)^F\), Zeitumkehrkovarianz, Spektrum und Entartung \\
\code{KramersTheorem} & Exportform reserviert; keine konkurrierende Kurzform & \code{KramersGate}, \code{StatisticsGate} & \code{GaugeSectorSeed}, \code{ChiralMatterSeed} & active extension & explizites Endtheorem des graduierten Zeitumkehrsektors \\
\code{GaugeSectorSeed} & Exportform reserviert; keine konkurrierende Kurzform & \code{CPTTheorem}, \code{KramersTheorem} & \code{AnomalyInflowSeed}, \code{ChargeQuantizationGate}, \code{HiggsSectorSeed} & active extension & emergente Eichsektoren aus Interface, Defekt und Glueing \\
\code{AnomalyInflowSeed} & Exportform reserviert; keine konkurrierende Kurzform & \code{GaugeSectorSeed} & \code{AnomalyInflow}, \code{AnomalyCancellationGate} & active extension & Interface-/Boundary-Inflow-Schicht für chirale Materie \\
\code{AnomalyInflow} & kein Exportalias; Lean-Name reserviert & \code{AnomalyInflowSeed} & \code{AnomalyCancellationGate}, \code{ChargeQuantizationGate} & active extension & prüft die Tragung der chiralen Anomalielast durch Inflow \\ \\
\bottomrule
\end{tabular}
\end{wideauditlandscape}

\begin{wideauditlandscape}
\landtablecaption{Tabelle K.1-A7: Definitionskatalog der Kernobjekte -- Objektbasis, Teil 7}
\begin{tabular}{L{0.085\linewidth}L{0.043\linewidth}L{0.073\linewidth}L{0.055\linewidth}L{0.043\linewidth}L{0.125\linewidth}L{0.10\linewidth}L{0.105\linewidth}L{0.058\linewidth}L{0.095\linewidth}L{0.065\linewidth}}
\toprule
Objekt & Klasse & Primärform & Exportform & Modus & Expansion / Bindung & Herkunft & Zielmodul & Namespace & Aktivierungsregel & Kontext \\
\midrule
\code{ChiralMatterSeed} & seed & \code{ChiralMatterSeed} & --- & selfspeaking & identisch mit dem Lean-Namen & CNNA/PillarE/ChiralMatterSeed & kein eigener Aliasexport & \code{CNNA.PillarE} & keine Exportaktivierung & E \\
\code{AnomalyCancellationGate} & gate & \emph{Anomalie-Kompensations-Gate} & \emph{AnomCancel} & text alias & \emph{Anomalie-Kompensations-Gate} verweist auf \code{AnomalyCancellationGate} & \path{CNNA/Gates/AnomalyCancellationGate.lean} & \path{CNNA/PillarE/Notation.lean} & \code{CNNA.PillarE} & kontextgebunden; bei Export \code{open scoped CNNA.PillarE} & E/Gates \\
\code{ChargeQuantizationGate} & gate & \emph{Ladungsquantisierungs-Gate} & \emph{ChargeQuant} & text alias & \emph{Ladungsquantisierungs-Gate} verweist auf \code{ChargeQuantizationGate} & \path{CNNA/Gates/ChargeQuantizationGate.lean} & \path{CNNA/PillarE/Notation.lean} & \code{CNNA.PillarE} & kontextgebunden; bei Export \code{open scoped CNNA.PillarE} & E/Gates \\
\code{FiniteInternalGeometrySeed} & seed & \code{FiniteInternalGeometrySeed} & --- & selfspeaking & identisch mit dem Lean-Namen & CNNA/PillarE/FiniteInternalGeometrySeed & kein eigener Aliasexport & \code{CNNA.PillarE} & keine Exportaktivierung & D--E \\
\code{HiggsSectorSeed} & seed & \emph{Higgs-Sektor-Seed} & \emph{HiggsSeed} & text alias & \emph{Higgs-Sektor-Seed} verweist auf \code{HiggsSectorSeed} & CNNA/PillarE/HiggsSectorSeed & \path{CNNA/PillarE/Notation.lean} & \code{CNNA.PillarE} & kontextgebunden; bei Export \code{open scoped CNNA.PillarE} & E \\
\code{EWSBSeed} & seed & \emph{EWSB-Seed} & \emph{EWSBSeed} & text alias & \emph{EWSB-Seed} verweist auf \code{EWSBSeed} & CNNA/PillarE/EWSBSeed & \path{CNNA/PillarE/Notation.lean} & \code{CNNA.PillarE} & kontextgebunden; bei Export \code{open scoped CNNA.PillarE} & E \\
\code{YukawaFlavorSeed} & seed & \emph{Yukawa-Flavor-Seed} & \emph{YukawaSeed} & text alias & \emph{Yukawa-Flavor-Seed} verweist auf \code{YukawaFlavorSeed} & CNNA/PillarE/YukawaFlavorSeed & \path{CNNA/PillarE/Notation.lean} & \code{CNNA.PillarE} & kontextgebunden; bei Export \code{open scoped CNNA.PillarE} & E \\
\code{GenerationSeed} & seed & \emph{Generationen-Seed} & \emph{GenSeed} & text alias & \emph{Generationen-Seed} verweist auf \code{GenerationSeed} & CNNA/PillarE/GenerationSeed & \path{CNNA/PillarE/Notation.lean} & \code{CNNA.PillarE} & kontextgebunden; bei Export \code{open scoped CNNA.PillarE} & E \\
\code{SMMatchingWindow} & window & \emph{SM-Matching-Window} & \emph{SMMatch} & text alias & \emph{SM-Matching-Window} verweist auf \code{SMMatchingWindow} & CNNA/PillarE/SMMatchingWindow & \path{CNNA/PillarE/Notation.lean} & \code{CNNA.PillarE} & kontextgebunden; bei Export \code{open scoped CNNA.PillarE} & E \\
\code{SMEFTWindow} & window & \emph{SMEFT-Window} & \emph{SMEFTWin} & text alias & \emph{SMEFT-Window} verweist auf \code{SMEFTWindow} & CNNA/PillarE/SMEFTWindow & \path{CNNA/PillarE/Notation.lean} & \code{CNNA.PillarE} & kontextgebunden; bei Export \code{open scoped CNNA.PillarE} & E \\
\code{SpectralTripleSeed} & seed & \code{SpectralTripleSeed} & --- & selfspeaking & identisch mit dem Lean-Namen & CNNA/PillarD/SpectralTripleSeed & kein eigener Aliasexport & \code{CNNA.PillarD} & keine Exportaktivierung & D \\
\bottomrule
\end{tabular}
\end{wideauditlandscape}

\begin{wideauditlandscape}
\landtablecaption{Tabelle K.1-B7: Definitionskatalog der Kernobjekte -- Kollisionen, Abhängigkeiten und Rollen, Teil 7}
\begin{tabular}{L{0.11\linewidth}L{0.18\linewidth}L{0.18\linewidth}L{0.18\linewidth}L{0.07\linewidth}L{0.17\linewidth}}
\toprule
Objekt & Kollisionspolitik & logische Vorgänger & Verbraucher / Nachfolger & Status & Rolle \\
\midrule
\code{ChiralMatterSeed} & kein Exportalias; Lean-Name reserviert & \code{KramersTheorem}, \code{GaugeSectorSeed} & \code{AnomalyCancellationGate}, \code{YukawaFlavorSeed} & active extension & chirale Materie unter Doubling- und Inflow-Kontrolle \\
\code{AnomalyCancellationGate} & Exportform reserviert; keine konkurrierende Kurzform & \code{AnomalyInflow}, \code{ChiralMatterSeed} & \code{ChargeQuantizationGate}, \code{SMEFTWindow} & active extension & Abnahme vor belastbarer Gauge-/Chiralitätsrede \\
\code{ChargeQuantizationGate} & Exportform reserviert; keine konkurrierende Kurzform & \code{AnomalyCancellationGate}, \code{GaugeSectorSeed} & \code{FiniteInternalGeometrySeed}, \code{SMMatchingWindow} & active extension & prüft Ladungs- und Quantisierungsstruktur vor Higgs/Flavor \\
\code{FiniteInternalGeometrySeed} & kein Exportalias; Lean-Name reserviert & \code{ChargeQuantizationGate} & \code{SpectralTripleSeed}, \code{DiracOperatorSeed} & active extension & E-nahe Fortsetzung der NCG-Route für endliche interne Geometrie \\
\code{HiggsSectorSeed} & Exportform reserviert; keine konkurrierende Kurzform & \code{GaugeSectorSeed} & \code{EWSBSeed}, \code{YukawaFlavorSeed} & active extension & symmetry-breaking-Vorzone nach stabiler Gauge-/Chiralitätsschicht \\
\code{EWSBSeed} & Exportform reserviert; keine konkurrierende Kurzform & \code{HiggsSectorSeed} & \code{YukawaFlavorSeed}, \code{GenerationSeed} & active extension & elektroschwache Brechung als eigener abgeleiteter Strang \\
\code{YukawaFlavorSeed} & Exportform reserviert; keine konkurrierende Kurzform & \code{EWSBSeed}, \code{ChiralMatterSeed} & \code{GenerationSeed}, \code{SMMatchingWindow} & active extension & Flavor- und Yukawa-Strang \\
\code{GenerationSeed} & Exportform reserviert; keine konkurrierende Kurzform & \code{YukawaFlavorSeed} & \code{SMMatchingWindow}, \code{SMEFTWindow} & active extension & später Parallelstrang für Generationen- und Mischungsstruktur \\
\code{SMMatchingWindow} & Exportform reserviert; keine konkurrierende Kurzform & \code{GenerationSeed}, \code{ChargeQuantizationGate} & \code{SMEFTWindow}, \code{DiracOperatorSeed} & late window & spätes IR-Matching auf SM-artige Sektoren \\
\code{SMEFTWindow} & Exportform reserviert; keine konkurrierende Kurzform & \code{SMMatchingWindow}, \code{AnomalyCancellationGate} & \code{TheorySpaceSeed} & late window & spätes EFT-Fenster ausserhalb des Kerns \\
\code{SpectralTripleSeed} & kein Exportalias; Lean-Name reserviert & \code{FiniteInternalGeometrySeed}, \code{EinsteinLimitSeed} & \code{DiracOperatorSeed}, \code{RealStructureSeed}, \code{SpectralActionSeed} & active extension & spektrogeometrischer Kern für die NCG-Route \\ \\
\bottomrule
\end{tabular}
\end{wideauditlandscape}

\begin{wideauditlandscape}
\landtablecaption{Tabelle K.1-A8: Definitionskatalog der Kernobjekte -- Objektbasis, Teil 8}
\begin{tabular}{L{0.085\linewidth}L{0.043\linewidth}L{0.073\linewidth}L{0.055\linewidth}L{0.043\linewidth}L{0.125\linewidth}L{0.10\linewidth}L{0.105\linewidth}L{0.058\linewidth}L{0.095\linewidth}L{0.065\linewidth}}
\toprule
Objekt & Klasse & Primärform & Exportform & Modus & Expansion / Bindung & Herkunft & Zielmodul & Namespace & Aktivierungsregel & Kontext \\
\midrule
\code{DiracOperatorSeed} & seed & \emph{Dirac-Operator-Seed} & \emph{DiracSeed} & text alias & \emph{Dirac-Operator-Seed} verweist auf \code{DiracOperatorSeed} & CNNA/PillarD/DiracOperatorSeed & \path{CNNA/PillarD/Notation.lean} & \code{CNNA.PillarD} & kontextgebunden; bei Export \code{open scoped CNNA.PillarD} & D \\
\code{RealStructureSeed} & seed & \code{RealStructureSeed} & --- & selfspeaking & identisch mit dem Lean-Namen & CNNA/PillarD/RealStructureSeed & kein eigener Aliasexport & \code{CNNA.PillarD} & keine Exportaktivierung & D \\
\code{SpectralActionSeed} & seed & \code{SpectralActionSeed} & --- & selfspeaking & identisch mit dem Lean-Namen & CNNA/PillarD/SpectralActionSeed & kein eigener Aliasexport & \code{CNNA.PillarD} & keine Exportaktivierung & D--E \\
\code{SpectralDimensionFlow} & def & \code{SpectralDimensionFlow} & --- & selfspeaking & identisch mit dem Lean-Namen & diagnostischer D-Pfad & kein eigener Aliasexport & \code{CNNA} & keine Exportaktivierung & D/Diagnostik \\
\code{SpectralDimensionFlowSeed} & seed & \code{SpectralDimensionFlowSeed} & --- & selfspeaking & identisch mit dem Lean-Namen & diagnostischer D-Pfad & kein eigener Aliasexport & \code{CNNA} & keine Exportaktivierung & D/Diagnostik \\
\code{CutoffModeCountingSeed} & seed & \code{CutoffModeCountingSeed} & --- & selfspeaking & identisch mit dem Lean-Namen & diagnostischer D-Pfad & kein eigener Aliasexport & \code{CNNA} & keine Exportaktivierung & D/Diagnostik \\
\code{RunningBoundaryDataSeed} & seed & \code{RunningBoundaryDataSeed} & --- & selfspeaking & identisch mit dem Lean-Namen & diagnostischer D-Pfad & kein eigener Aliasexport & \code{CNNA} & keine Exportaktivierung & D/Diagnostik \\
\code{TheorySpaceSeed} & seed & \code{TheorySpaceSeed} & --- & selfspeaking & identisch mit dem Lean-Namen & D/AS-Route & kein eigener Aliasexport & \code{CNNA} & keine Exportaktivierung & D/AS \\
\bottomrule
\end{tabular}
\end{wideauditlandscape}

\begin{wideauditlandscape}
\landtablecaption{Tabelle K.1-B8: Definitionskatalog der Kernobjekte -- Kollisionen, Abhängigkeiten und Rollen, Teil 8}
\begin{tabular}{L{0.11\linewidth}L{0.18\linewidth}L{0.18\linewidth}L{0.18\linewidth}L{0.07\linewidth}L{0.17\linewidth}}
\toprule
Objekt & Kollisionspolitik & logische Vorgänger & Verbraucher / Nachfolger & Status & Rolle \\
\midrule
\code{DiracOperatorSeed} & Exportform reserviert; keine konkurrierende Kurzform & \code{SpectralTripleSeed}, \code{FiniteInternalGeometrySeed} & \code{RealStructureSeed}, \code{SpectralActionSeed} & active extension & abgeleitete Dirac-Struktur der spektrogeometrischen D-Linie \\
\code{RealStructureSeed} & kein Exportalias; Lean-Name reserviert & \code{SpectralTripleSeed}, \code{DiracOperatorSeed} & \code{SpectralActionSeed}, \code{TheorySpaceSeed} & active extension & KO- und Realstruktur der NCG-Brücke \\
\code{SpectralActionSeed} & kein Exportalias; Lean-Name reserviert & \code{SpectralTripleSeed}, \code{RealStructureSeed} & \code{SpectralDimensionFlow}, \code{RunningBoundaryDataSeed} & active extension & spektrogeometrische Wirkungsroute als Brücke nach E \\
\code{SpectralDimensionFlow} & kein Exportalias; Lean-Name reserviert & \code{SpectralActionSeed} & \code{SpectralDimensionFlowSeed}, \code{TheorySpaceSeed} & planned active & laufende spektrale Dimensionsdiagnostik \\
\code{SpectralDimensionFlowSeed} & kein Exportalias; Lean-Name reserviert & \code{SpectralDimensionFlow} & \code{CutoffModeCountingSeed}, \code{RunningBoundaryDataSeed} & planned active & Vorzone für spektrale Dimensionsflüsse \\
\code{CutoffModeCountingSeed} & kein Exportalias; Lean-Name reserviert & \code{SpectralDimensionFlowSeed} & \code{RunningBoundaryDataSeed}, \code{TheorySpaceSeed} & planned active & Modenzahl- und cutoffbezogene Spektraldiagnostik \\
\code{RunningBoundaryDataSeed} & kein Exportalias; Lean-Name reserviert & \code{CutoffModeCountingSeed}, \code{SpectralActionSeed} & \code{TheorySpaceSeed}, \code{SMMatchingWindow} & planned active & laufende Randdaten entlang Skalen- oder RG-artiger Entwicklung \\
\code{TheorySpaceSeed} & kein Exportalias; Lean-Name reserviert & \code{RunningBoundaryDataSeed}, \code{SpectralDimensionFlow} & \code{SMMatchingWindow}, \code{SMEFTWindow} & planned active & asymptotic-safety-nahe Theorie- und Flussstruktur \\ \\
\bottomrule
\end{tabular}
\end{wideauditlandscape}

\endgroup
\section{Quellenangaben}
\begin{thebibliography}{99}

\bibitem{REALOQSV0605}
REAL-OQS project archive,
\emph{REAL-OQS\_v0.0605\_full\_no\_comments.zip} 2026-03-29.

\bibitem{CNNASpec}
Jan Seeck,
\emph{Complement-Netzarchitektur (CNNA) aus dem ToC}, internal project specification v0.1,
uploaded PDF, 2026.

\bibitem{RulebookV18}
Jan Seeck,
\emph{Regelwerk für codebasierte Lean-Dokumentationen. Externe Dokumentation statt Code-Kommentierung},
Revision 1.8, uploaded PDF, 2026.

\bibitem{Dedushenko2023}
Mykola Dedushenko,
\emph{Snowmass White Paper: The Quest to Define QFT},
Int. J. Mod. Phys. A \textbf{38} (2023) 2330002,
arXiv:2203.08053.

\bibitem{HaagKastler1964}
Rudolf Haag and Daniel Kastler,
\emph{An Algebraic Approach to Quantum Field Theory},
J. Math. Phys. \textbf{5} (1964), 848--861,
DOI: 10.1063/1.1704187.

\bibitem{BFV2003}
Romeo Brunetti, Klaus Fredenhagen, and Rainer Verch,
\emph{The Generally Covariant Locality Principle: A New Paradigm for Local Quantum Physics},
Commun. Math. Phys. \textbf{237} (2003), 31--68,
arXiv:math-ph/0112041.

\bibitem{FewsterLang2016}
Christopher J. Fewster and Benjamin Lang,
\emph{Dynamical Locality of the Free Maxwell Field},
Ann. Henri Poincare \textbf{17} (2016), 401--436,
arXiv:1403.7083.

\bibitem{BeniniSchenkelWoike2021}
Marco Benini, Alexander Schenkel, and Lukas Woike,
\emph{Operads for Algebraic Quantum Field Theory},
Commun. Contemp. Math. \textbf{23} (2021) 2050007,
arXiv:1709.08657.

\bibitem{Atiyah1988}
Michäl Atiyah,
\emph{Topological Quantum Field Theories},
Publ. Math. IHES \textbf{68} (1988), 175--186.

\bibitem{Segal2004}
Gräme Segal,
\emph{The Definition of Conformal Field Theory},
in \emph{Topology, Geometry and Quantum Field Theory}, London Mathematical Society Lecture Note Series,
Cambridge University Press, 2004, pp.~421--577.

\bibitem{FuchsRunkelSchweigert2003}
Jürgen Fuchs, Ingo Runkel, and Christoph Schweigert,
\emph{Boundaries, Defects and Frobenius Algebras},
arXiv:hep-th/0302200.

\bibitem{Frohlich2007}
Jürg Fröhlich, Jürgen Fuchs, Ingo Runkel, and Christoph Schweigert,
\emph{Duality and Defects in Rational Conformal Field Theory},
Nucl. Phys. B \textbf{763} (2007), 354--430,
arXiv:hep-th/0607247.

\bibitem{Jacobson1995}
Ted Jacobson,
\emph{Thermodynamics of Spacetime: The Einstein Equation of State},
Phys. Rev. Lett. \textbf{75} (1995), 1260--1263,
arXiv:gr-qc/9504004.

\bibitem{Reuter1998}
Martin Reuter,
\emph{Nonperturbative Evolution Equation for Quantum Gravity},
Phys. Rev. D \textbf{57} (1998), 971--985,
arXiv:hep-th/9605030.

\bibitem{ReuterSaueressig2012}
Martin Reuter and Frank Saueressig,
\emph{Quantum Einstein Gravity},
New J. Phys. \textbf{14} (2012) 055022,
arXiv:1202.2274.

\bibitem{LauscherReuter2005}
Oliver Lauscher and Martin Reuter,
\emph{Fractal Spacetime Structure in Asymptotically Safe Gravity},
JHEP \textbf{10} (2005) 050,
arXiv:hep-th/0508202.

\bibitem{Visser2002}
Matt Visser,
\emph{Sakharov's Induced Gravity: A Modern Perspective},
Mod. Phys. Lett. A \textbf{17} (2002), 977--992,
arXiv:gr-qc/0204062.

\bibitem{BrunettiFredenhagenRejzner2022}
Romeo Brunetti, Klaus Fredenhagen, and Kasia Rejzner,
\emph{Locally Covariant Approach to Effective Quantum Gravity},
arXiv:2212.07800.

\bibitem{ChamseddineConnes1996}
Ali H. Chamseddine and Alain Connes,
\emph{The Spectral Action Principle},
arXiv:hep-th/9606001.

\bibitem{ChamseddineConnes2007}
Ali H. Chamseddine and Alain Connes,
\emph{Why the Standard Model},
arXiv:0706.3688.

\bibitem{NielsenNinomiya1981}
Holger Bech Nielsen and Masao Ninomiya,
\emph{A No-Go Theorem for Regularizing Chiral Fermions},
Phys. Lett. B \textbf{105} (1981), 219--223,
DOI: 10.1016/0370-2693(81)91026-1.

\bibitem{CallanHarvey1985}
Curtis G. Callan Jr. and Jeffrey A. Harvey,
\emph{Anomalies and Fermion Zero Modes on Strings and Domain Walls},
Nucl. Phys. B \textbf{250} (1985), 427--436,
DOI: 10.1016/0550-3213(85)90489-4.

\bibitem{AlvarezGaume2022}
Luis Alvarez-Gaume and Migül Angel Vazquez-Mozo,
\emph{Anomalies and the Green-Schwarz Mechanism},
arXiv:2211.06467.


\bibitem{DHR1971I}
Sergio Doplicher, Rudolf Haag, and John E. Roberts,
\emph{Local Observables and Particle Statistics I},
Commun. Math. Phys. \textbf{23} (1971), 199--230,
DOI: 10.1007/BF01877742.

\bibitem{DoplicherRoberts1990}
Sergio Doplicher and John E. Roberts,
\emph{Why There is a Field Algebra with a Compact Gauge Group Describing the Superselection Structure in Particle Physics},
Commun. Math. Phys. \textbf{131} (1990), 51--107,
DOI: 10.1007/BF02097680.

\bibitem{BuchholzWichmann1986}
Detlev Buchholz and Eyvind H. Wichmann,
\emph{Causal Independence and the Energy-Level Density of States in Local Quantum Field Theory},
Commun. Math. Phys. \textbf{106} (1986), 321--344,
DOI: 10.1007/BF01454978.

\bibitem{BuchholzRoberts2014}
Detlev Buchholz and John E. Roberts,
\emph{New Light on Infrared Problems: Sectors, Statistics, Symmetries and Spectrum},
Commun. Math. Phys. \textbf{330} (2014), 935--972,
arXiv:1304.2794.

\bibitem{BuchholzScattering2005}
Detlev Buchholz,
\emph{Scattering in Relativistic Quantum Field Theory: Basic Concepts, Tools, and Results},
arXiv:math-ph/0509047.

\bibitem{Borchers1992}
H. J. Borchers,
\emph{The CPT-Theorem in Two-Dimensional Theories of Local Observables},
Commun. Math. Phys. \textbf{143} (1992), 315--332,
DOI: 10.1007/BF02099011.

\bibitem{Haag1996}
Rudolf Haag,
\emph{Local Quantum Physics: Fields, Particles, Algebras},
2nd rev. ed., Texts and Monographs in Physics, Springer, Berlin, Heidelberg, 1996,
DOI: 10.1007/978-3-642-61458-3.

\bibitem{Dybalski2005}
Wojciech Dybalski,
\emph{Haag-Ruelle Scattering Theory in Presence of Massless Particles},
Lett. Math. Phys. \textbf{72} (2005), 27--38,
arXiv:hep-th/0412226.

\bibitem{FewsterVerch2012}
Christopher J. Fewster and Rainer Verch,
\emph{Dynamical Locality and Covariance: What Makes a Physical Theory the Same in All Spacetimes?},
Ann. Henri Poincare \textbf{13} (2012), 1613--1674,
DOI: 10.1007/s00023-012-0165-0.

\bibitem{Radzikowski1996}
Marek J. Radzikowski,
\emph{Micro-Local Approach to the Hadamard Condition in Quantum Field Theory on Curved Space-Time},
Commun. Math. Phys. \textbf{179} (1996), 529--553,
DOI: 10.1007/BF02100096.

\bibitem{BrunettiFredenhagenKoehler1996}
Romeo Brunetti, Klaus Fredenhagen, and Martin Köhler,
\emph{The Microlocal Spectrum Condition and Wick Polynomials of Free Fields on Curved Spacetimes},
Commun. Math. Phys. \textbf{180} (1996), 633--652,
DOI: 10.1007/BF02099626.

\bibitem{HollandsWald2001}
Stefan Hollands and Robert M. Wald,
\emph{Local Wick Polynomials and Time Ordered Products of Quantum Fields in Curved Spacetime},
Commun. Math. Phys. \textbf{223} (2001), 289--326,
DOI: 10.1007/s002200100540.

\bibitem{HollandsWald2002}
Stefan Hollands and Robert M. Wald,
\emph{Existence of Local Covariant Time Ordered Products of Quantum Fields in Curved Spacetime},
Commun. Math. Phys. \textbf{231} (2002), 309--345,
DOI: 10.1007/s00220-002-0719-y.

\bibitem{ReehSchlieder1961}
Helmut Reeh and Siegfried Schlieder,
\emph{Bemerkungen zur Unit"aräquivalenz von Lorentzinvarianten Feldern},
Nuovo Cimento \textbf{22} (1961), 1051--1068,
DOI: 10.1007/BF02787889.

\bibitem{Takesaki1970}
Masamichi Takesaki,
\emph{Tomita's Theory of Modular Hilbert Algebras and its Applications},
Lecture Notes in Mathematics \textbf{128}, Springer, 1970,
DOI: 10.1007/BFb0065832.

\bibitem{BisognanoWichmann1975}
Joseph J. Bisognano and Eyvind H. Wichmann,
\emph{On the Duality Condition for a Hermitian Scalar Field},
J. Math. Phys. \textbf{16} (1975), 985--1007,
DOI: 10.1063/1.522605.

\bibitem{BisognanoWichmann1976}
Joseph J. Bisognano and Eyvind H. Wichmann,
\emph{On the Duality Condition for Quantum Fields},
J. Math. Phys. \textbf{17} (1976), 303--321,
DOI: 10.1063/1.522898.

\bibitem{Pauli1940}
Wolfgang Pauli,
\emph{The Connection Between Spin and Statistics},
Phys. Rev. \textbf{58} (1940), 716--722,
DOI: 10.1103/PhysRev.58.716.

\bibitem{GuidoLongo1995}
Daniele Guido and Roberto Longo,
\emph{An Algebraic Spin and Statistics Theorem},
Commun. Math. Phys. \textbf{172} (1995), 517--533,
DOI: 10.1007/BF02101806.

\bibitem{Luders1957}
Gerhart L"uders,
\emph{Proof of the TCP Theorem},
Ann. Phys. \textbf{2} (1957), 1--15,
DOI: 10.1016/0003-4916(57)90032-5.

\bibitem{FeynmanAntiparticles}
Richard P. Feynman,
\emph{The Reason for Antiparticles},
in \emph{Paul Dirac: The Man and His Work},
Cambridge University Press, published online 2014,
DOI: 10.1017/CBO9781107590076.002.

\bibitem{PDG2023SM}
Particle Data Group,
\emph{Review of Particle Physics: The Standard Model and Electroweak Model and Constraints on New Physics},
PDG review article, 2023 edition.

\bibitem{IsidoriWilschWyler2024}
Gino Isidori, Felix Wilsch, and Daniel Wyler,
\emph{The Standard Model Effective Field Theory at Work},
Rev. Mod. Phys. \textbf{96} (2024) 015006,
arXiv:2303.16922.

\end{thebibliography}

\end{document}